% !TeX root = ../Book5.tex
% 纲要标题：# 第二编　对称群、模表示与箭图
% 纲要位置：工作记录/计划纲要.md，第 3899 行。

\BookPart{对称群、模表示与箭图}{b5:part:02}
\BookPartGuide
有限群的一般表示论提供了分解、特征标与诱导的方法；要把它们用于具体分类，还需要利用对象本身的结构。本编先研究对称群，借分拆和表把表示写成可以计算的模；随后转向模特征和箭图，考察完全可约性失效以后，怎样描述不可分解项及其扩张。

\ProseParagraphBreak
第6、7章组成“组合构造与张量运算”单元，回答怎样显式构造表示并计算其维数与乘法。\chapref{b5:ch:06}先构造 Specht 模，在整数系数上证明标准基，再用不变配对与支配序完成特征零分类。限制滤过、分支与钩长公式随后给出维数计算。\chapref{b5:ch:07}借助对称函数计算这些表示：Frobenius 特征映射把诱导乘积转成函数乘法，Schur--Weyl 对偶把同一分拆与张量幂上的一般线性群作用联系起来。

\ProseParagraphBreak
第8—10章组成“非半单现象与不可分解分类”单元，研究完全可约性失效后，合成因子之外还须记录什么。\chapref{b5:ch:08}从循环 \(p\) 群的不可分解模出发，介绍投射覆盖、Cartan 矩阵与块，再用离散赋值环中的格说明普通表示如何约化。Brauer 特征标保留合成因子的整数重数，但不能区别全部非分裂扩张；本章的具体矩阵与格例子会说明这种限制。这里的 Brauer 特征标与\chapref{b5:ch:05}的特征零 Brauer 诱导定理分别定义，不能混用。

\ProseParagraphBreak
箭图部分从\chapref{b5:ch:09}开始。前四节从模、张量积和有限维线性代数出发，建立投射表示与一阶扩张；\secref{b5:sec:0905}再联系第二卷第15—16章的根、幂等元提升与 Morita 理论\footnote{森田紀一（假名：もりた きいち，罗马音：Kiiti Morita，1915-1995），日本数学家，建立以模范畴等价为核心的森田理论。}。路径乘积统一按映射复合的方向书写，全部采用左模。\chapref{b5:ch:10}在有限无环箭图及代数闭基域的条件下证明 Gabriel 定理；ADE 图、整数根和反射都在该章建立。

\ProseParagraphBreak
阅读与计算可以相互穿插：在对称群中检查低阶标准表和钩长；在箭图中先把一支箭按秩分解，再检查几支箭的复合秩和扩张空间。Etingof 等人的《Introduction to Representation Theory》\cite[Sections 5.12--5.23; Chapter 6]{EtingofEtAl2011RepresentationTheory}可用于对照这两个部分；模表示的约化与 Brauer 特征标另见本编第8章所列资料。各章习题包括具体计算及条件失效的例子，解答附在题后。

% !TeX root = ../../Book5.tex
% 纲要标题：## 第6章　对称群的表示
% 纲要位置：工作记录/计划纲要.md，第 3901 行。
% > 本章按分拆与 Young 图、Specht 模、Garnir 关系与标准基、不可约分类、分支与钩长公式、维数计算顺读。分类证明不依赖后置的钩长公式或 Frobenius 特征映射。
\chapter{对称群的表示}
\label{b5:ch:06}
\BookChapterNumber{b5:ch:06}

\begin{chapterabstract}
本章用 Young 表把对称群的表示问题变成可以逐步计算的组合问题。先由行等价类与列反对称化构造 Specht 模，再用 Garnir 关系及三角性建立整系数标准基。随后通过不变配对、子模定理和支配序完成特征零下的不可约分类。最后从标准表中最大数字的位置构造限制滤过，证明分支规则与钩长公式，并计算低阶表示的矩阵、维数和限制分解。
\end{chapterabstract}

\begin{learninggoals}
\begin{itemize}
\item 区分 Young 表、行表类与列交错元，能在置换模和群代数左理想之间转换。
\item 能实际使用 Garnir 关系，并说明标准化为何终止、标准基为何在整数及任意系数环上成立。
\item 分别证明 Specht 模的不可约性、不同形状的非同构性及分类的完备性，指出特征条件的作用。
\item 能构造删格滤过、应用限制与诱导分支，并用钩长公式核验维数。
\end{itemize}
预备知识为有限群的置换表示、Maschke 定理、Schur 引理、不可约复特征标与共轭类计数，以及诱导与限制的 Frobenius 互反性。组合构造在交换含幺环上进行；不可约分类和直和分支明确限于特征零。
\end{learninggoals}

对称群已经在第一编的小群算例中出现，但逐个制作特征标表不能处理任意 \(n\)。本章的构造适用于所有 \(n\)：分拆决定图形，填数表给出向量，置换数字产生群作用。先证明这些向量的标准基定理，再讨论不可约分类，最后推导计数公式。

% !TeX root = ../../Book5.tex
% 纲要标题：### 6.1 分拆与 Young 图
% 纲要位置：工作记录/计划纲要.md，第 3905 行。
% * 分拆、Young 图与标准表
% * 共轭类与分拆
% * Young 子群
\section{分拆与 Young 图}
\label{b5:sec:0601}

对称群的共轭类由轮换长度决定。令人感兴趣的是，同一组整数资料也能标记它的不可约表示。不过，共轭类与表示并不是把名称一一配上便有了联系：我们还要从每一个分拆构造出一个向量空间和群作用，再证明这些空间不可约且没有遗漏。

\subsection{分拆、Young 图与标准表}
\label{b5:subsec:0601:01}

\begin{definition}\label{b5:def:partition-tableau}
设 \(n\) 为非负整数。有限的非增正整数列 \(\lambda=(\lambda_1,\ldots,\lambda_r)\) 满足 \(\sum_i\lambda_i=n\) 时，称为 \(n\) 的一个\term[zhengshu-fenchai]{分拆}[partition]，记为 \(\lambda\vdash n\)；\(n=0\) 时允许空列。集合
\[
[\lambda]=\bigl\{(i,j):1\leq i\leq r,\ 1\leq j\leq\lambda_i\bigr\}
\]
称为其 \term[Young-tu]{Young 图}[Young diagram]。第一个坐标由上向下表示行，第二个坐标由左向右表示列。从 \([\lambda]\) 到 \(\{1,\ldots,n\}\) 的双射称为形状为 \(\lambda\) 的 \term[Young-biao]{Young 表}[Young tableau]。若每行从左向右、每列从上向下都严格递增，则称其为\term[biaozhun-Young-biao]{标准 Young 表}[standard Young tableau]。
\end{definition}
\symidx{\(\lambda\vdash n\)：整数分拆}

只要求列递增时称为列标准表；只要求行递增时称为行标准表。这两个条件不可混同。例如形状 \((3,2)\) 的表
\[
t=\begin{array}{ccc}1&2&4\\3&5&\end{array},
\qquad
u=\begin{array}{ccc}1&4&2\\3&5&\end{array}
\]
都列标准，只有 \(t\) 标准。图的格子仍固定在原位，置换作用改变的是格内的数字。

\begin{definition}\label{b5:def:conjugate-dominance}
设 \(n\) 为非负整数，\(\lambda,\mu\) 为 \(n\) 的分拆，并在各序列的末尾补零。令 \(\lambda'_j=\#\{i:\lambda_i\geq j\}\)，则 \(\lambda'=(\lambda'_1,\lambda'_2,\ldots)\) 去掉末尾的零后称为 \(\lambda\) 的\term[gonge-fenchai]{共轭分拆}[conjugate partition]。若对每个正整数 \(a\) 都有
\[
\lambda_1+\cdots+\lambda_a\geq\mu_1+\cdots+\mu_a,
\]
便称 \(\lambda\) 在\term[zhipei-xu]{支配序}[dominance order]下大于或等于 \(\mu\)，记为 \(\lambda\trianglerighteq\mu\)。
\end{definition}

转置 Young 图便得到共轭分拆，所以 \((\lambda')'=\lambda\)。支配序是偏序：反对称性由各个部分和相减得到。它一般不是全序，例如 \((4,1,1)\) 与 \((3,3)\) 的第一部分和及前两部分和呈相反的大小关系。

\begin{lemma}\label{b5:lem:standard-remove-largest}
设 \(n\ge1\) 为整数，\(\lambda=(\lambda_1,\ldots,\lambda_r)\vdash n\)，约定 \(\lambda_{r+1}=0\)，\(t\) 为形状 \(\lambda\) 的标准表。数字 \(n\) 所在的格子 \((i,\lambda_i)\) 满足 \(\lambda_i>\lambda_{i+1}\)。删去这个格子，得到一个大小为 \(n-1\) 的标准表。反之，任何满足该条件的格子都能由这个过程出现。
\end{lemma}
\begin{proof}
最大数字右边、下边都不能还有格子，因此它在行末，而且下一行更短。删去后仍保持行列递增。反过来，给删格后的任意标准表补上数字 \(n\)，新格子的左边、上边数字均小于 \(n\)，右边、下边没有格子，所以仍为标准表。
\end{proof}

满足引理条件的格子称为\term[keyi-fangge]{可移方格}[removable box]，删去它以后仍是 Young 图。对交换含幺环上的 Specht 模，可移方格将标记限制到 \(S_{n-1}\) 时的一个滤过的相继商；在特征零域上，这个滤过分裂为直和。

\subsection{共轭类与分拆}
\label{b5:subsec:0601:02}

\begin{proposition}\label{b5:prop:symmetric-conjugacy-partitions}
对每个整数 \(n\geq1\)，\(S_n\) 的共轭类与 \(n\) 的分拆一一对应：把置换写成不交轮换之积，并把轮换长度按非增次序排列。若分拆 \(\lambda\vdash n\) 中数字 \(j\) 出现 \(m_j\) 次，则对应的中心化子阶与共轭类大小分别为
\[
z_\lambda=\prod_{j\geq1}j^{m_j}m_j!,
\qquad
\frac{n!}{z_\lambda}.
\]
\end{proposition}
\begin{proof}
共轭把轮换 \((a_1\,\cdots\,a_j)\) 变成 \(\bigl(g(a_1)\,\cdots\,g(a_j)\bigr)\)，所以保持各个轮换的长度。反过来，长度相同的两组轮换可以逐个配对，再逐位置定义一个置换 \(g\)，这个 \(g\) 就把第一组共轭到第二组。

与给定置换对易的元素必须把一个轮换轨道送到同长度轨道。长度为 \(j\) 的 \(m_j\) 个轨道可以有 \(m_j!\) 种排列，每个轨道的像又由一个起点确定，有 \(j\) 种选择；各个长度之间相互独立，得到中心化子的阶。轨道与稳定子群公式给出共轭类大小。
\end{proof}

例如 \(S_4\) 的五个分拆为 \((4),(3,1),(2,2),(2,1,1),(1,1,1,1)\)，相应共轭类大小为 \(6,8,3,6,1\)。它们的和为 \(24\)。这里用分拆描述的是轮换类型；后面用分拆标记 Specht 模，是另一项需要证明的构造。

\subsection{Young 子群}
\label{b5:subsec:0601:03}

\begin{definition}\label{b5:def:row-column-groups}
设 \(n\) 为非负整数，\(\lambda=(\lambda_1,\ldots,\lambda_r)\vdash n\)，\(t\) 为形状 \(\lambda\) 的 Young 表。使每一行的数字集合分别保持不变的置换组成子群 \(R_t\leq S_n\)，使每一列的数字集合分别保持不变的置换组成子群 \(C_t\leq S_n\)，分别称为 \(t\) 的行群与列群。按行依次填入 \(1,\ldots,n\) 所得表的行群记为
\[
S_\lambda=S_{\lambda_1}\times\cdots\times S_{\lambda_r}\leq S_n,
\]
称为相应的 \term[Young-ziqun]{Young 子群}[Young subgroup]；这里各个因子置换互不相交的连续数字段。
\end{definition}

行群的阶为 \(\prod_i\lambda_i!\)，列群的阶为 \(\prod_j\lambda'_j!\)。若置换同时保持每行和每列，便固定每一个行列交点，故
\begin{equation}\label{b5:eq:row-column-intersection}
R_t\cap C_t=\{1\}.
\end{equation}
对任意 \(g\in S_n\)，有 \(R_{gt}=gR_tg^{-1}\) 和 \(C_{gt}=gC_tg^{-1}\)。因此表的具体填数会改变子群在 \(S_n\) 中的位置，却不会改变它们的共轭类型。

% !TeX root = ../../Book5.tex
% 纲要标题：### 6.2 Specht 模
% 纲要位置：工作记录/计划纲要.md，第 3911 行。
% * 行表类置换模、行表类（tabloid）及行表类的列交错元（polytabloid）；区分等价类与列交错化所得线性组合，不使用“多重表”作为后者的译名
% * 明确 Young 表、按行置换取等价类的行表类与列交错化所得列交错元的区别，统一表模和 Specht 模的符号
% * 表类置换模、行表类（按行置换的表类）与列交错元（列反对称化的表类线性组合）；原“多重表”不再作为列交错元的译名，multitableau 另作术语辨析
% * Young 对称化子的特征零构造
% * Specht 模的基本性质
% * 明确 Young 表、按行置换取等价类的行表类与列反对称化所得列交错元的区别，统一表模和 Specht 模的符号
% * 标准列交错元基与 Garnir 关系的证明集中到第6.3节；Young 对称化子构造与该模型的关系同时交代
\section{Specht 模}
\label{b5:sec:0602}

若只要求同一行内的数字不分先后，Young 表就给出一种集合划分，因而产生一个置换表示。但这个置换表示通常很大，也并不不可约。Specht 模在其中选取列方向交替变号的线性组合，使行对称与列反对称同时参与构造。

\subsection{行表类置换模与列交错元}
\label{b5:subsec:0602:01}

\begin{definition}\label{b5:def:specht-module}
设 \(R\) 为交换含幺环，\(n\) 为非负整数，\(\lambda\vdash n\)。若两个形状为 \(\lambda\) 的 Young 表的每一行数字集合分别相同，则称这两个表行等价；表 \(t\) 的等价类记为 \(\{t\}\)，称为一个\term[hang-biaolei]{行表类}[tabloid]。令 \(M_R^\lambda\) 为以全部这些表类为基的自由 \(R\)-模，规定左 \(S_n\) 作用
\[
g\{t\}=\{gt\}\qquad(g\in S_n).
\]
记 \(C_t\) 为表 \(t\) 的列群，在群代数 \(R[S_n]\) 中令
\[
b_t=\sum_{c\in C_t}\operatorname{sgn}(c)c,\qquad
e_t=b_t\{t\}.
\]
线性组合 \(e_t\) 称为这个\term[hang-biaolei-lie-jiaocuo-yuan]{行表类的列交错元}[polytabloid]，以下在对象明确时简称列交错元。\termidx[lie-jiaocuo-yuan]{列交错元}由所有 \(e_t\) 生成的 \(R\)-子模
\[
S_R^\lambda=\sum_t R e_t\subseteq M_R^\lambda
\]
称为形状 \(\lambda\) 的 \term[Specht-mo]{Specht 模}[Specht module]。
\end{definition}
\symidx{\(M_R^\lambda\)：表类置换模}
\symidx{\(S_R^\lambda\)：Specht 模}

因为 \(g b_t g^{-1}=b_{gt}\)，所以 \(g e_t=e_{gt}\)；上述子模确实被 \(S_n\) 保持。任何两个同形状表由某个置换联系，因此一个 \(e_t\) 就在 \(R[S_n]\) 作用下生成整个 \(S_R^\lambda\)。\(M_R^\lambda\) 的基数为 \(n!/\prod_i\lambda_i!\)，并且
\[
M_R^\lambda\cong R[S_n]\otimes_{R[S_\lambda]}R,
\]
其中右边的 \(R\) 采用平凡作用；同构把 \(g\otimes1\) 送到按行填数表的表类 \(g\{t\}\)。

\begin{example}\label{b5:exm:specht-21}
对 \(\lambda=(2,1)\)，用 \(v_i\) 表示第二行只有数字 \(i\) 的行表类。表
\[
t=\begin{array}{cc}1&2\\3&\end{array}
\]
的列群为 \(\bigl\{1,(13)\bigr\}\)，所以 \(e_t=v_3-v_1\)。遍取表后得到所有 \(v_i-v_j\)，从而
\[
S_R^{(2,1)}
=\{a_1v_1+a_2v_2+a_3v_3:a_1+a_2+a_3=0\}.
\]
它以 \(v_2-v_1,v_3-v_1\) 为自由基，在任意 \(R\) 上秩均为 \(2\)。
\end{example}

\subsection{特征零下的 Young 对称化子}
\label{b5:subsec:0602:02}

\begin{definition}\label{b5:def:young-symmetrizer}
设 \(k\) 为特征零域，\(n\) 为非负整数，\(\lambda\vdash n\)，\(t\) 为形状 \(\lambda\) 的 Young 表，\(R_t,C_t\) 分别为其行群、列群。在 \(k[S_n]\) 中令
\[
a_t=\sum_{r\in R_t}r,\qquad
b_t=\sum_{c\in C_t}\operatorname{sgn}(c)c,\qquad c_t=b_ta_t.
\]
本书把 \(c_t\) 称为 \term[Young-duichenghuazi]{Young 对称化子}[Young symmetrizer]，并固定先作行求和、再作列反对称化的乘法次序。
\end{definition}

单独的 \(a_t/|R_t|\) 与 \(b_t/|C_t|\) 都是幂等元；它们的乘积一般还不是幂等元。不同文献也会采用 \(a_tb_t\)，或预先除去两个子群的阶。阅读公式时须同时核对乘法次序与归一化。本书采用的表类模型与 \(b_ta_t\) 直接对应：

\begin{proposition}\label{b5:prop:young-ideal-model}
设 \(R\) 为交换含幺环，\(n\) 为非负整数，\(\lambda\vdash n\)，\(t\) 为形状 \(\lambda\) 的 Young 表，\(R_t,C_t\) 分别为其行群、列群。令 \(a_t=\sum_{r\in R_t}r\)、\(b_t=\sum_{c\in C_t}\operatorname{sgn}(c)c\)，记 \(M_R^\lambda\) 为行表类置换模、\(S_R^\lambda\) 为 Specht 模。存在左 \(R[S_n]\)-模同构
\[
M_R^\lambda\longrightarrow R[S_n]a_t,\qquad
\{gt\}\longmapsto ga_t,
\]
它把 \(S_R^\lambda\) 同构地送到 \(R[S_n]b_ta_t\)。
\end{proposition}
\begin{proof}
\(\{g_1t\}=\{g_2t\}\) 当且仅当 \(g_2^{-1}g_1\in R_t\)，这时 \(g_1a_t=g_2a_t\)。不同左陪集上的元素和具有互不相交的群元素支集，因此选取陪集代表后，所有 \(ga_t\) 在任意 \(R\) 上线性无关，并生成左理想。映射于是良定义且为同构。它与左乘相容，并把 \(e_t=b_t\{t\}\) 送到 \(b_ta_t\)，故也识别两者生成的子模。
\end{proof}

这里的同构甚至没有用到特征零。特征零的作用将在不可约性和幂等元归一化时出现；不能把含分母的归一化照搬到任意特征。

\subsection{Specht 模的基本性质}
\label{b5:subsec:0602:03}

\begin{proposition}\label{b5:prop:specht-column-sign}
设 \(R\) 为交换含幺环，\(n\) 为非负整数，\(\lambda\vdash n\)，\(t\) 为形状 \(\lambda\) 的 Young 表，\(C_t\) 为其列群，\(e_t\in M_R^\lambda\) 为对应的列交错元。对 \(c\in C_t\)，有
\[
e_{ct}=\operatorname{sgn}(c)e_t.
\]
此外，\(e_t\) 中基向量 \(\{t\}\) 的系数为 \(1\)，故当 \(R\ne0\) 时 \(e_t\ne0\)。形状 \((n)\) 与 \((1^n)\) 的 Specht 模分别是平凡表示与符号表示。
\end{proposition}
\begin{proof}
等式 \(cb_t=\operatorname{sgn}(c)b_t\) 由列群求和中重排下标得到，再用 \(e_{ct}=ce_t\)。若 \(c\in C_t\) 满足 \(\{ct\}=\{t\}\)，则 \(c\in R_t\cap C_t\)，由\eqref{b5:eq:row-column-intersection}知 \(c=1\)，所以指定系数恰为 \(1\)。

单行时只有一个行表类，列群平凡；单列时行群平凡，反对称和在每个置换下乘以其符号，且由一个非零向量生成。两种情形分别给出所述一维表示。
\end{proof}

把每一列内部的数字排成递增次序，最多改变列交错元的符号。因此研究全部 \(e_t\) 时可以先限制到列标准表；行内的乱序则需要下一节的 Garnir 关系\footnote{加尼尔（Henri Georges Garnir，1921-1985），比利时数学家，研究代数与分析，以对称群表示中的关系式著称。}，不能用相同的符号法直接消去。

\subsection{Young 表与表类模型的符号}
\label{b5:subsec:0602:04}

这三种对象处于不同层次：\(t\) 是每个格子都有指定数字的表；\(\{t\}\) 忘掉各行内部次序，是置换模的一个基向量；\(e_t\) 则是多个这样的基向量带符号相加所得的向量。一般不能把 \(e_t\) 当成另一个等价类。

\ProseParagraphBreak

例如\cref{b5:exm:specht-21}中的 \(e_t=v_3-v_1\) 同时涉及第二行为 \(3\) 和第二行为 \(1\) 的两个表类。英文 multitableau 指由若干表组成的另一类组合对象，不是这里的列交错元；本章保留后者原文，避免仅从词形推断构造。

现在已经知道 Specht 模由有限多个显式向量生成，却还不知道它们之间的关系。例如 \((2,1)\) 形状共有六个表，但模的秩只有二。下一节要证明，恰好选取所有标准表所对应的 \(e_t\)，就得到一组基，而且这个结论首先在整数环上成立。

\ProseParagraphBreak

标准基定理需要分别证明生成性和线性无关性。Garnir 关系把非标准表改写成次序严格下降的表，由有限性得到生成性；行中小数字的分布则给出一个对角元为 \(1\) 的三角矩阵，用来证明线性无关。仅仅计算标准表的数目，还不能得到这两项结论。

% !TeX root = ../../Book5.tex
% 纲要标题：### 6.3 Garnir 关系与标准基
% 纲要位置：工作记录/计划纲要.md，第 3919 行。
% * 建立 Garnir 关系和标准化过程，证明标准列交错元的生成性及线性无关性
% * 固定表的次序和标准化规则，分别说明重写为何终止、标准列交错元为何生成，以及任意交换含幺系数环上的基变换。
% * 以首项和三角性组织线性无关性证明，区分整系数标准基与基变换后的域上模型；不要仅凭标准表个数推断线性无关。
\section{Garnir 关系与标准基}
\label{b5:sec:0603}

本节先在整数系数下处理关系。这样既能避免用除法掩盖特征限制，也能同时得到任意域和任意交换含幺环上的标准基。整个过程只移动表中数字，不改变 Young 图的形状。

\subsection{Garnir 关系与标准列交错元基}
\label{b5:subsec:0603:01}

\begin{lemma}[Garnir 关系]\label{b5:lem:garnir}
设 \(n\) 为非负整数，\(\lambda\vdash n\)，\(\lambda'\) 为共轭分拆，\(t\) 为形状 \(\lambda\) 的 Young 表。在相邻的第 \(j,j+1\) 列中，分别选数字子集 \(A,B\)，并假定
\[
|A|+|B|>\lambda'_j.
\]
令 \(H=S_A\times S_B\leq S_{A\cup B}\)，选取左陪集 \(S_{A\cup B}/H\) 的代表元集 \(D\)，其中包含单位元。记 \(e_u\) 为表 \(u\) 的列交错元，则在整数行表类置换模 \(M_{\mathbb Z}^\lambda\) 中有
\begin{equation}\label{b5:eq:garnir}
\sum_{\sigma\in D}\operatorname{sgn}(\sigma)e_{\sigma t}=0.
\end{equation}
同一等式经系数映射后，在任意交换含幺环上仍成立。它称为 \term[Garnir-guanxi]{Garnir 关系}[Garnir relation]。
\end{lemma}
\begin{proof}
记 \(\gamma=\sum_{\pi\in S_{A\cup B}}\operatorname{sgn}(\pi)\pi\)。对任意 \(c\in C_t\)，数字集合 \(A\cup B\) 在 \(ct\) 中仍全部位于这两列，因而位于前 \(\lambda'_j\) 行。由所假定的不等式，其中必有两个数字处于同一行；交换这两个数字的换位 \(\tau\in S_{A\cup B}\) 固定 \(\{ct\}\)。在 \(\gamma\{ct\}\) 的求和中把 \(\pi\) 与 \(\pi\tau\) 配成一对，两项符号相反、向量相同，故它们相消。因此 \(\gamma e_t=0\)。

另一方面，按左陪集分组得
\[
\gamma=\sum_{\sigma\in D}\operatorname{sgn}(\sigma)\sigma
\sum_{h\in H}\operatorname{sgn}(h)h.
\]
\(H\) 包含在 \(C_t\) 中，故内层和作用于 \(e_t\) 时得到 \(|A|!\,|B|!\,e_t\)。于是
\[
|A|!\,|B|!\sum_{\sigma\in D}\operatorname{sgn}(\sigma)e_{\sigma t}=0.
\]
自由交换群 \(M_{\mathbb Z}^\lambda\) 无挠，故可消去这个非零整数，得到\eqref{b5:eq:garnir}。这是整数系数向量恒等式，随后可映到任意 \(R\)；这一顺序并未在 \(R\) 中除以阶乘。
\end{proof}

\begin{theorem}[标准基]\label{b5:thm:specht-standard-basis}
设 \(R\) 为交换含幺环，\(n\) 为非负整数，\(\lambda\vdash n\)。所有形状为 \(\lambda\) 的标准表 \(t\) 所给的列交错元 \(e_t\)，构成 Specht 模 \(S_R^\lambda\) 的一组 \(R\)-基。因而 \(S_{\mathbb Z}^\lambda\) 是有限秩自由交换群，且自然映射
\[
R\otimes_{\mathbb Z}S_{\mathbb Z}^\lambda\longrightarrow S_R^\lambda
\]
为同构。
\end{theorem}
\begin{proof}
先证明生成性。利用\cref{b5:prop:specht-column-sign}，可先把每列排成递增次序。对列标准表 \(t\)，把第一列从上到下的数字、第二列从上到下的数字依次接起来，所得长度为 \(n\) 的列字记为 \(w(t)\)，按通常字典序比较。可能的列字只有有限多个。

若 \(t\) 还不是标准表，则某一行有相邻逆序
\(t(i,j)>t(i,j+1)\)。取
\[
\begin{aligned}
A&=\bigl\{t(i,j),t(i+1,j),\ldots,t(\lambda'_j,j)\bigr\},\\
B&=\bigl\{t(1,j+1),\ldots,t(i,j+1)\bigr\}.
\end{aligned}
\]
由于两列递增，\(A\) 中每个数字都大于 \(B\) 中每个数字，且
\(|A|+|B|=\lambda'_j+1\)。在 Garnir 关系中单独移出单位元项，得到
\begin{equation}\label{b5:eq:garnir-rewrite}
e_t=-\sum_{\substack{\sigma\in D\\\sigma\ne1}}
\operatorname{sgn}(\sigma)e_{\sigma t}.
\end{equation}
一个非单位陪集满足 \(\sigma A\ne A\)，所以新的第 \(j\) 列从 \(A\) 中移出若干数字，并换入同样多个来自 \(B\) 的较小数字；前 \(j-1\) 列不变。将新表的每列重新排序后，列字严格小于 \(w(t)\)。列排序只改变列交错元的符号。因此重复\eqref{b5:eq:garnir-rewrite}时，每一条重写链的列字都严格下降；有限性保证过程终止。终止项不能还有相邻行逆序，所以全部为标准表，生成性得证。

再证明线性无关。对任意行表类 \(\{u\}\)，令
\[
N_u(a,m)=\#\{\text{前 }a\text{ 行中不大于 }m\text{ 的数字}\}.
\]
这只依赖表类。若 \(t\) 标准，则对每个 \(c\in C_t\) 都有
\begin{equation}\label{b5:eq:tabloid-triangular}
N_{ct}(a,m)\leq N_t(a,m)\qquad\text{对所有 }a,m.
\end{equation}
事实上，第 \(j\) 列中不大于 \(m\) 的数字在 \(t\) 中占据最上面的若干位置，若其个数为 \(q_j\)，则前 \(a\) 行贡献 \(\min(a,q_j)\)；列置换以后的贡献不会更大。对各列相加即得不等式。

在标准表集合上，用所有这些不等式定义偏序。若双向成立，则每一行中小于等于每个 \(m\) 的数字个数都相同，因此各行的数字集合相同；标准表各行递增，故两个表相同。把这个有限偏序延伸为全序。在 \(e_t\) 的表类展开中，若标准表 \(u\) 的表类 \(\{u\}\) 出现，则\eqref{b5:eq:tabloid-triangular}说明 \(u\) 不大于 \(t\)；而 \(\{t\}\) 的系数由\cref{b5:prop:specht-column-sign}知为 \(1\)。于是只取标准表所对应的那些表类坐标，得到一个对角元全为 \(1\) 的三角矩阵。它在任意交换环上可逆，所以这些 \(e_t\) 线性无关。

最后，生成性所用系数全为整数；自然映射把整数标准基变成 \(R\) 上刚证的标准基，故为同构。
\end{proof}

\subsection{表的次序、标准化终止与生成性}
\label{b5:subsec:0603:02}

上述证明给出实际的标准化算法。输入为一个填数表，或若干列交错元的有限 \(R\)-线性组合；输出为标准列交错元的线性组合。每次先排序各列并记下符号，再找到一处相邻行逆序，取证明中规定的两段数字，按\eqref{b5:eq:garnir-rewrite}展开，最后合并相同标准表的系数。这个算法只用加减法，不要求能在 \(R\) 中求逆。

\begin{example}\label{b5:exm:garnir-21}
对列标准而非标准的表
\[
t=\begin{array}{cc}2&1\\3&\end{array}
\]
取 \(A=\{2,3\}\)、\(B=\{1\}\)，左陪集代表可取 \(1,(12),(13)\)。Garnir 关系为
\(e_t-e_{(12)t}-e_{(13)t}=0\)。令
\[
u=\begin{array}{cc}1&2\\3&\end{array},
\qquad
v=\begin{array}{cc}1&3\\2&\end{array}.
\]
则 \((12)t=u\)，而 \((13)t\) 在第一列交换上下数字后变成 \(v\)，所以
\[
e_t=e_u-e_v.
\]
在\cref{b5:exm:specht-21}的记号下，这正是
\(v_3-v_2=(v_3-v_1)-(v_2-v_1)\)。
\end{example}

两段数字都含多个元素时，陪集选择和列排序的符号也可以逐项跟踪。取形状 \((3,2,1)\) 的列标准表
\[
t=\begin{array}{ccc}1&2&5\\4&3&\\6&&\end{array},
\qquad A=\{4,6\},\quad B=\{2,3\}.
\]
第二行在前两列逆序，而 \(|A|+|B|=4>3=\lambda'_1\)。左陪集 \(S_{\{2,3,4,6\}}/(S_A\times S_B)\) 由二元子集 \(\sigma A\) 唯一编号；六个代表可取
\[
D=\{1,(24),(26),(34),(36),(24)(36)\}.
\]
它们的 \(\sigma A\) 依次为 \(\{4,6\},\{2,6\},\{2,4\},\{3,6\},\{3,4\},\{2,3\}\)，恰好穷尽六种二元子集。对 \(i<j\)、\(\{i,j\}\subset\{2,3,4,6\}\)，令 \(u_{ij}\) 的第一列为递增的 \(1,i,j\)，第二列为该集合剩下的两个数字并递增排列，第三列只有 \(5\)。

\ProseParagraphBreak
例如，代表 \((24)\) 给出
\[
(24)t=\begin{array}{ccc}1&4&5\\2&3&\\6&&\end{array},
\qquad
u_{26}=\begin{array}{ccc}1&3&5\\2&4&\\6&&\end{array}.
\]
置换 \((24)\) 的符号为负；将第二列的 \(4,3\) 排序又交换一次，所以 \(e_{(24)t}=-e_{u_{26}}\)。该项在 Garnir 和中为 \(+e_{u_{26}}\)，移到右边后才成为 \(-e_{u_{26}}\)。对其余代表同样记下列排序符号，\eqref{b5:eq:garnir}成为
\[
e_t+e_{u_{26}}-e_{u_{24}}-e_{u_{36}}
+e_{u_{34}}+e_{u_{23}}=0,
\]
因而
\[
e_t=-e_{u_{26}}+e_{u_{24}}+e_{u_{36}}
-e_{u_{34}}-e_{u_{23}}.
\]
这五个 \(u_{ij}\) 的行列都递增，因此这一次重写已经把 \(e_t\) 表成标准列交错元的组合。

\ProseParagraphBreak
算法终止并不自动保证计算很快。每次 Garnir 展开可能产生多个新项，本章只证明有限步骤内能完成重写，不给出随 \(n\) 增长的复杂度界。

\subsection{首项三角性与整系数线性无关性}
\label{b5:subsec:0603:03}

在独立性证明中，起作用的不仅是“存在一个首项”，还包括首项系数是单位 \(1\)。若这个系数只是某个非零整数，约化到正特征后可能消失；这里的三角矩阵在所有系数环上仍可逆。因此标准基不因特征改变而丢失，但由它给出的表示却可能变得可约。

\ProseParagraphBreak
三角形并不意味着不同标准列交错元的表类坐标互不重叠。对形状 \((3,2)\) 的标准表 \(t=123|45\)，列群由 \((14),(25)\) 生成，故
\[
e_t=\{123|45\}-\{234|15\}
-\{135|24\}+\{345|12\}.
\]
其中 \(u=135|24\) 也是标准表，它的行表类在 \(e_t\) 中的系数为 \(-1\)。它在证明中的偏序下小于 \(t\)：所有 \(N_u(a,m)\le N_t(a,m)\)，而 \(N_u(1,2)=1<2=N_t(1,2)\)。因此标准表类坐标矩阵会有这样的非零非对角项；线性无关性依靠的是这些项都位于同一三角形内，且对角系数为一。

\ProseParagraphBreak

记形状 \(\lambda\) 的标准表数为 \(f^\lambda\)。标准基定理现在才允许我们写
\[
\operatorname{rank}_{\mathbb Z}S_{\mathbb Z}^\lambda=f^\lambda,
\qquad
\dim_kS_k^\lambda=f^\lambda
\]
对任意域 \(k\) 成立。后面计算 \(f^\lambda\) 的公式用于求维数；标准基的存在已经证明，不可约分类将另作论证。

% !TeX root = ../../Book5.tex
% 纲要标题：### 6.4 不可约分类
% 纲要位置：工作记录/计划纲要.md，第 3925 行。
% * 特征零下的不可约性
% * 不同分拆给出不同不可约表示
% * 分类完备性
% * 调用第6.3节已建立的 Garnir 关系及标准基，在本节依次完成不可约性、非同构性与完备性；不使用尚未证明的钩长公式或 Frobenius 特征映射
% * 用 Specht 模上的不变配对与子模定理证明特征零不可约性，明确正特征下哪一步失效
% * 用支配序和 Young 子群分析不同分拆的非同构性；不把参数个数相等当作非同构性的证明
% * 在不可约性和非同构性建立后，结合共轭类个数与半单群代数完成分类完备性；处理空分拆，并用整数标准基把有理数上的作用同态扩张到任意特征零域
% * 分类之后分别讨论 Young 对称化子的归一化与幂等元、正特征中的配对退化与可约性、Young 置换模分解的支配三角性及维数平方和的核验，使标题区分新构造、反例、推论与事后检验
\section{不可约分类}
\label{b5:sec:0604}

标准基解决的是向量空间的大小，尚未说明群作用能否保持一个更小的非零子空间。本节在表类基上引入不变配对；列反对称化对这个配对有一种特别简单的秩一作用，从而控制 Specht 模的所有子模。

\subsection{特征零下的不可约性}
\label{b5:subsec:0604:01}

在 \(M_R^\lambda\) 上把不同行表类规定为正交，并令每个基向量与自身的配对为 \(1\)，得到对称双线性型
\[
B\left(\sum_u a_u\{u\},\sum_u d_u\{u\}\right)=\sum_u a_ud_u.
\]
群作用置换这组基，所以 \(B(gv,gw)=B(v,w)\)。这里采用双线性型；在复数域上不把它与共轭线性的 Hermitian 内积混为一谈。

\begin{lemma}\label{b5:lem:specht-rank-one}
设 \(R\) 为交换含幺环，\(n\) 为非负整数，\(\lambda\vdash n\)，\(t\) 为形状 \(\lambda\) 的 Young 表，\(C_t\) 为其列群。记 \(M_R^\lambda\) 为行表类置换模，\(b_t=\sum_{c\in C_t}\operatorname{sgn}(c)c\)、\(e_t=b_t\{t\}\)，\(B\) 为以行表类基为标准正交基的对称双线性型。对任意 \(v\in M_R^\lambda\)，有
\begin{equation}\label{b5:eq:specht-rank-one}
b_tv=B(v,e_t)e_t.
\end{equation}
\end{lemma}
\begin{proof}
先令 \(v=\{u\}\)。若 \(t\) 的某一列有两个数字位于 \(u\) 的同一行，交换它们的换位既属于 \(C_t\) 又固定 \(\{u\}\)；把列群求和配成符号相反的二项，即得 \(b_t\{u\}=0\)。

若不存在这样的两个数字，则 \(u\) 的第一行共有 \(\lambda_1\) 个数字，它们恰从 \(t\) 的每列各取一个。可用列置换把这些数字送到 \(t\) 的第一行。删去已匹配的第一行，对剩余各列与第二行重复：第二行必须恰从尚非空的每列各取一个。逐行进行，得到 \(c\in C_t\) 使 \(\{u\}=\{ct\}\)。于是 \(b_t\{u\}=\operatorname{sgn}(c)e_t\)。这证明 \(b_tv\) 总是 \(e_t\) 的标量倍数。

因为 \(c\) 与 \(c^{-1}\) 符号相同，\(b_t\) 关于 \(B\) 自伴随。故 \(b_tv\) 中 \(\{t\}\) 的系数为
\(B\bigl(b_tv,\{t\}\bigr)=B\bigl(v,b_t\{t\}\bigr)=B(v,e_t)\)。而 \(e_t\) 中 \(\{t\}\) 的系数为 \(1\)，所求标量即为此数。
\end{proof}

\begin{theorem}[子模定理]\label{b5:thm:specht-submodule}
设 \(k\) 为任意域，\(n\) 为非负整数，\(\lambda\vdash n\)，\(M_k^\lambda\) 为行表类置换模，\(S_k^\lambda\subseteq M_k^\lambda\) 为 Specht 模，\(U\) 为 \(M_k^\lambda\) 的左 \(k[S_n]\)-子模。则至少有下列一项成立：
\[
S_k^\lambda\subseteq U,\qquad
U\subseteq(S_k^\lambda)^\perp,
\]
其中正交补取在 \(M_k^\lambda\) 中，\(B\) 为以行表类基为标准正交基的对称双线性型。
\end{theorem}
\begin{proof}
若第二项不成立，因为全部 \(e_t\) 生成 \(S_k^\lambda\)，存在 \(u\in U\) 与某个 \(t\) 使 \(B(u,e_t)\ne0\)。由\eqref{b5:eq:specht-rank-one}，
\(b_tu=B(u,e_t)e_t\in U\)。域中非零标量可逆，故 \(e_t\in U\)；这个向量的 \(S_n\)-轨道生成整个 Specht 模，所以第一项成立。
\end{proof}

\begin{proposition}\label{b5:prop:specht-char-zero-simple}
设 \(k\) 为特征零域，\(n\) 为非负整数，\(\lambda\vdash n\)，\(M_k^\lambda\) 为行表类置换模，\(S_k^\lambda\) 为其中的 Specht 模。以 \(M_k^\lambda\) 的行表类基为标准正交基的对称双线性型 \(B\)，在 \(S_k^\lambda\) 上非退化，且 \(S_k^\lambda\) 作为左 \(k[S_n]\)-模绝对不可约。
\end{proposition}
\begin{proof}
在整数标准基中写出 \(B\) 的 Gram 矩阵 \(G_\lambda\)，其条目为整数。把它看作实矩阵时，由\cref{b5:thm:specht-standard-basis}，这些基向量在实表类空间中线性无关，而表类配对为正定型。因此 \(\det G_\lambda\) 是非零整数。它在任意特征零域中的像仍非零，故限制型非退化。

若 \(0\ne U\subseteq S_k^\lambda\) 是子模，子模定理或者给出 \(S_k^\lambda\subseteq U\)，或者给出
\(U\subseteq S_k^\lambda\cap(S_k^\lambda)^\perp=0\)。后者不可能，故 \(U=S_k^\lambda\)。对 \(k\) 的任何扩域，标准基与同一个非零 Gram 行列式仍成立，所以上述不可约性在扩域后保持，即为绝对不可约。
\end{proof}

\subsection{不同分拆与非同构表示}
\label{b5:subsec:0604:02}

\begin{proposition}\label{b5:prop:specht-dominance}
设 \(k\) 为特征零域，\(n\) 为非负整数，\(\lambda,\mu\vdash n\)，\(S_k^\lambda,S_k^\mu\) 为相应 Specht 左 \(k[S_n]\)-模，\(M_k^\lambda,M_k^\mu\) 为相应行表类置换模。若
\(\operatorname{Hom}_{k[S_n]}(S_k^\lambda,M_k^\mu)\ne0\)，则 \(\lambda\trianglerighteq\mu\)。此外，
\[
\dim_k\operatorname{Hom}_{k[S_n]}(S_k^\lambda,M_k^\lambda)=1.
\]
特别地，\(S_k^\lambda\cong S_k^\mu\) 当且仅当 \(\lambda=\mu\)。
\end{proposition}
\begin{proof}
固定一个 \(\lambda\)-表 \(t\)，并令 \(\varphi:S_k^\lambda\rightarrow M_k^\mu\) 非零。因为 \(e_t\) 生成整个模，\(\varphi(e_t)\ne0\)。列群对 \(e_t\) 按符号作用，因此
\[
b_t\cdot\varphi(e_t)=\varphi(b_te_t)=|C_t|\cdot\varphi(e_t)\ne0.
\]
所以至少有一个 \(\mu\)-行表类 \(\{u\}\) 满足 \(b_t\{u\}\ne0\)。与秩一引理中的配对消去相同，\(u\) 的任何一行不能含有 \(t\) 同一列的两个数字。前 \(a\) 行从第 \(j\) 列至多取 \(\min(a,\lambda'_j)\) 个数字，故
\[
\mu_1+\cdots+\mu_a
\leq\sum_j\min(a,\lambda'_j)
=\lambda_1+\cdots+\lambda_a.
\]
这就是支配条件。

当 \(\mu=\lambda\) 时，\cref{b5:lem:specht-rank-one}又说明 \(b_t\cdot\varphi(e_t)\in ke_t\)。因 \(|C_t|\) 可逆，\(\varphi(e_t)\) 为 \(e_t\) 的标量倍数；由生成性，\(\varphi\) 是包含映射的标量倍数。包含映射确实非零，故 Hom 空间恰为一维。

最后，若两个 Specht 模同构，把同构与 \(S_k^\mu\subseteq M_k^\mu\) 复合，得到 \(\lambda\trianglerighteq\mu\)；用逆同构得到 \(\mu\trianglerighteq\lambda\)。支配序的反对称性给出 \(\lambda=\mu\)。反方向显然成立。
\end{proof}

\subsection{分类完备性}
\label{b5:subsec:0604:03}

\begin{theorem}[对称群的不可约分类]\label{b5:thm:specht-classification}
设 \(k\) 为特征零域，\(n\geq0\) 为整数。全部 Specht 模 \(S_k^\lambda\)，其中 \(\lambda\vdash n\)，是互不同构的绝对不可约左 \(k[S_n]\)-模，并穷尽所有不可约有限维左 \(k[S_n]\)-模。
\end{theorem}
\begin{proof}
当 \(n=0\) 时，\(S_0\) 是平凡群，唯一的空分拆对应 Specht 模 \(k\)，结论成立。以下设 \(n\geq1\)。不可约性和非同构性已经分别由\cref{b5:prop:specht-char-zero-simple}及\cref{b5:prop:specht-dominance}证明。先取 \(k=\mathbb C\)。\cref{b5:prop:symmetric-conjugacy-partitions}说明共轭类数等于分拆数，而\cref{b5:thm:irreducible-character-basis}说明不可约复表示的类型数等于共轭类数。我们已经为每个分拆构造出互不同构的不可约表示，因此没有其他复不可约类型。

为把完备性推广到任意特征零域，不能直接使用只对复表示证明的类函数论。改用整数标准基写出所有表示矩阵，得到一个在有理数上定义的代数同态
\[
\Phi_{\mathbb Q}:\mathbb Q[S_n]\longrightarrow
\prod_{\lambda\vdash n}
\operatorname{End}_{\mathbb Q}(S_{\mathbb Q}^\lambda).
\]
对任意特征零域 \(k\)，标准基定理给出
\[
k\otimes_{\mathbb Q}S_{\mathbb Q}^\lambda
\xrightarrow{\ \sim\ }S_k^\lambda,
\qquad a\otimes e_t\longmapsto ae_t.
\]
两边的群作用矩阵都由同一组整数矩阵得到，故此同构等变。在这些基下，\(\operatorname{End}_{\mathbb Q}(S_{\mathbb Q}^\lambda)\) 的矩阵单位基也扩张为 \(\operatorname{End}_k(S_k^\lambda)\) 的基。因此 \(k\) 上的作用同态 \(\Phi_k\) 正是 \(k\otimes_{\mathbb Q}\Phi_{\mathbb Q}\)，其矩阵由 \(\Phi_{\mathbb Q}\) 的矩阵逐项扩张系数得到。

扩张到复数后，此映射单射：若群代数元素在所有不可约上作用为零，Maschke 定理使它在正则表示上也作用为零；作用于单位元即知该元素为零。由\cref{b5:thm:row-column-orthogonality}在单位共轭类处得到的维数平方和，目标维数为
\(\sum_\lambda(f^\lambda)^2=n!\)，所以这个单射还是同构。

因此 \(\Phi_{\mathbb Q}\) 在所选基中是方阵，其有理行列式非零。在任意特征零域 \(k\) 中，这个行列式的像仍非零，故
\begin{equation}\label{b5:eq:symmetric-split-algebra}
k[S_n]\cong\prod_{\lambda\vdash n}M_{f^\lambda}(k).
\end{equation}
一个不可约模只能由这组正交中心幂等元中的一个非零作用，所以只属于一个矩阵块。设 \(V\) 是矩阵代数 \(M_d(k)\) 的不可约模。取 \(0\ne v\in V\)，由 \(1=\sum_iE_{ii}\) 可选 \(i\) 使 \(E_{ii}v\ne0\)。记 \(e_1,\ldots,e_d\) 为 \(k^d\) 的标准列向量基，定义
\[
\psi:k^d\longrightarrow V,
\qquad e_j\longmapsto E_{ji}v.
\]
矩阵单位关系 \(E_{ab}E_{ji}=\delta_{bj}E_{ai}\) 说明 \(\psi\) 是模同态，而 \(\psi(e_i)=E_{ii}v\ne0\)。列向量模 \(k^d\) 简单，所以这个非零同态的核为零；其像是不可约模 \(V\) 的非零子模，故等于 \(V\)。于是 \(\psi\) 为同构。这证明\eqref{b5:eq:symmetric-split-algebra}所列的 Specht 模也穷尽 \(k\) 上全部类型。
\end{proof}

这里先证明“每个都是不可约”及“彼此不同”，最后才计数。没有前两项，参数个数与共轭类个数相等并不能证明分类。

\subsection{Young 对称化子的归一化与幂等元}
\label{b5:subsec:0604:04}

标准基使 Gram 行列式、基域扩张与群作用矩阵都有整系数的共同来源。它也使 Young 对称化子可以精确归一化：

\begin{proposition}\label{b5:prop:young-quasi-idempotent}
设 \(k\) 为特征零域，\(n\) 为非负整数，\(\lambda\vdash n\)，\(t\) 为形状 \(\lambda\) 的 Young 表，\(R_t,C_t\) 分别为其行群、列群。令 \(a_t=\sum_{r\in R_t}r\)、\(b_t=\sum_{c\in C_t}\operatorname{sgn}(c)c\)、\(c_t=b_ta_t\in k[S_n]\)，并记 \(f^\lambda\) 为形状 \(\lambda\) 的标准表数。则
\[
c_t^2=\frac{n!}{f^\lambda}c_t,
\]
且 \(\dfrac{f^\lambda}{n!}c_t\) 是一个幂等元，它生成的左理想同构于 Specht 模 \(S_k^\lambda\)。
\end{proposition}
\begin{proof}
由\cref{b5:prop:young-ideal-model}，计算 \(c_t^2\) 等价于计算
\(b_ta_te_t\)。秩一公式给出
\(b_ta_te_t=B(a_te_t,e_t)e_t\)。记这个整数系数为 \(\alpha\)，便有 \(c_t^2=\alpha c_t\)。

在实表类空间中，\(a_t\) 自伴随且 \(a_t^2=|R_t|a_t\)，所以
\[
\alpha=\frac{B(a_te_t,a_te_t)}{|R_t|}>0.
\]
严格不等号的理由是 \(a_te_t\) 中 \(\{t\}\) 的系数为 \(|R_t|\)：每个行置换都固定 \(\{t\}\)，而 \(e_t\) 在该处系数为 \(1\)。因此 \(\alpha\ne0\)。

右乘 \(c_t/\alpha\) 是 \(k[S_n]\) 上的投影，像是 \(k[S_n]c_t\)，其维数为 \(f^\lambda\)。由 \(R_t\cap C_t=\{1\}\)，\(c_t\) 中单位元系数为 \(1\)，所以在群元素基下，右乘 \(c_t/\alpha\) 的每个对角元都是 \(1/\alpha\)，迹为 \(n!/\alpha\)。投影的迹等于像的维数，得到 \(\alpha=n!/f^\lambda\)。
\end{proof}

这项构造可与 Etingof 等人的 Young 对称化子证明对照\cite[Sections 5.12--5.13, pp. 110--113]{EtingofEtAl2011RepresentationTheory}。该书采用预先归一化的行、列求和及相反乘法次序，因而相应的倍数也随约定改变；本节的归一化是在当前模型内算出的。

\subsection{正特征中的配对退化与可约性}
\label{b5:subsec:0604:05}

子模定理本身允许任意特征；不可约性证明需要的额外事实是 Specht 模上配对非退化。非零整数 Gram 行列式在正特征中可能变成零，这正是不能照搬证明的位置。

\begin{example}\label{b5:exm:specht-mod-three}
设 \(k\) 为特征 \(3\) 的域。由\cref{b5:exm:specht-21}，
\(S_k^{(2,1)}\) 是三点置换模中的坐标和为零子空间。向量
\(v_1+v_2+v_3\) 非零，因 \(3=0\) 而属于这个二维子空间，并生成一维平凡子模，所以 \(S_k^{(2,1)}\) 可约。在基 \(v_2-v_1,v_3-v_1\) 中，
\[
G_{(2,1)}=
\begin{pmatrix}2&1\\1&2\end{pmatrix},
\qquad
\det G_{(2,1)}=3.
\]
其根恰好是上述平凡子模。标准基仍有两个向量，退化的是配对与不可约性。
\end{example}

\subsection{Young 置换模分解的支配三角性}
\label{b5:subsec:0604:06}

由于 \(M_k^\mu=\operatorname{Ind}_{S_\mu}^{S_n}k\)，Frobenius 互反性把它与 Young 子群不动向量联系起来。在特征零下，完全可约性和 Specht 模的绝对不可约性又说明，一个不可约模在 \(M_k^\mu\) 中的重数等于
\(\dim_k\operatorname{Hom}_{k[S_n]}(S_k^\lambda,M_k^\mu)\)。因此\cref{b5:prop:specht-dominance}给出三角性：
\[
M_k^\mu\cong S_k^\mu\oplus
\bigoplus_{\substack{\lambda\trianglerighteq\mu\\\lambda\ne\mu}}
(S_k^\lambda)^{\oplus m_{\lambda\mu}},
\qquad m_{\lambda\mu}\in\mathbb Z_{\geq0}.
\]
这里尚未计算所有重数，但对角重数已经精确等于 \(1\)。后面建立 Frobenius 特征映射时，正是这项有方向的三角性帮助识别 Schur 函数。

\subsection{维数平方和的核验}
\label{b5:subsec:0604:07}

分类还给出一个事后检验维数的方法：
\[
n!=\sum_{\lambda\vdash n}(f^\lambda)^2.
\]
左边来自群的阶，右边来自各个标准表集合的大小。这个恒等式在分类完成时已经成立；钩长公式将使每一项都能迅速算出。它适合检查是否漏了分拆或算错维数，却不能在没有不可约性和非同构性证明时单独充当分类论证。

% !TeX root = ../../Book5.tex
% 纲要标题：### 6.5 分支规则与钩长公式的证明
% 纲要位置：工作记录/计划纲要.md，第 3935 行。
% * 从标准基构造限制到 \(S_{n-1}\) 的自然滤过；在特征零半单性下得到去掉可移方格的分支分解
% * 由分支递推或等价的组合论证计算标准 Young 表的个数，再证明乘积形式 \(f^\lambda=n!/\prod_{u\in\lambda}h(u)\)
% * 核对钩形、两行分拆和低阶分拆的维数，保持递推、乘积公式与不可约表示维数的一致性
% * 诱导分支通过 Frobenius 互反获得；与第7章对称函数的联系只作后续再解释，不反向承担本节证明
\section{分支规则与钩长公式的证明}
\label{b5:sec:0605}

把 \(S_n\) 中固定数字 \(n\) 的子群识别为 \(S_{n-1}\)。在标准表中，\(n\) 只能占据可移方格；我们将证明，删去这个方格不仅递推地计数标准表，也准确描述 Specht 模的限制。先建立模的滤过，再在特征零下利用完全可约性分裂它。

\subsection{限制滤过与可移方格的分支}
\label{b5:subsec:0605:01}

\begin{theorem}[限制分支]\label{b5:thm:specht-branching}
设 \(R\) 为交换含幺环，\(n\geq1\) 为整数，\(\lambda\vdash n\)，\(S_R^\lambda\) 为 Specht 左 \(R[S_n]\)-模，\(S_{n-1}\le S_n\) 为固定数字 \(n\) 的子群。把 \(\lambda\) 的可移方格按行号由小到大记为
\(x_1,\ldots,x_q\)，令 \(\mu^{(j)}=\lambda\setminus x_j\)。则
\(\operatorname{Res}_{S_{n-1}}^{S_n}S_R^\lambda\) 有自然滤过
\[
0=F_0\subseteq F_1\subseteq\cdots\subseteq F_q=S_R^\lambda,
\qquad F_j/F_{j-1}\cong S_R^{\mu^{(j)}}.
\]
其中 \(F_j\) 由数字 \(n\) 位于 \(x_1,\ldots,x_j\) 的全部标准 \(e_t\) 生成。特别地，若 \(R=k\) 为特征零域，则
\begin{equation}\label{b5:eq:specht-restriction}
\operatorname{Res}_{S_{n-1}}^{S_n}S_k^\lambda
\cong\bigoplus_{x\text{ 为可移方格}}S_k^{\lambda\setminus x}.
\end{equation}
\end{theorem}
\begin{proof}
设 \(x_j\) 在第 \(r_j\) 行。在表类置换模中，令 \(L_r\) 由数字 \(n\) 位于前 \(r\) 行的全部行表类生成。由于 \(S_{n-1}\) 固定 \(n\)，每个 \(L_r\) 都是该子群下的不变子模。

若标准表 \(t\) 中的 \(n\) 位于可移方格 \((r,\lambda_r)\)，则该格所在列的高度恰为 \(r\)。列置换只能把 \(n\) 移到这一列的前 \(r\) 行，所以 \(e_t\in L_r\)。由此已有 \(F_j\subseteq S_R^\lambda\cap L_{r_j}\)。

对一个可移行 \(r\)，定义 \(R\)-线性映射
\[
d_r:M_R^\lambda\longrightarrow M_R^{\lambda\setminus(r,\lambda_r)}:
\]
若一个行表类的第 \(r\) 行含 \(n\)，就删掉这个数字；否则送到零。删除以行的数字集合定义，与行内次序无关，且与 \(S_{n-1}\) 作用相容。对于标准表 \(t\)，若 \(n\) 在第 \(r\) 行，则只有固定 \(n\) 的列置换在 \(d_r(e_t)\) 中留下贡献；这些置换恰是删格后的列群，所以
\begin{equation}\label{b5:eq:specht-delete}
d_r(e_t)=e_{t\setminus n}.
\end{equation}
若 \(n\) 所在可移行小于 \(r\)，则 \(d_r(e_t)=0\)。

现证明上面的包含是等号。将
\(v\in S_R^\lambda\cap L_{r_j}\) 按标准基展开。假如有系数非零的标准表把 \(n\) 放在大于 \(r_j\) 的行，取其中最大的行号 \(s\)。由于 \(v\in L_{r_j}\)，有 \(d_s(v)=0\)；展开中行号小于 \(s\) 的项均被 \(d_s\) 消去，行号等于 \(s\) 的项由\eqref{b5:eq:specht-delete}变成形状 \(\lambda\setminus(s,\lambda_s)\) 的互不相同标准列交错元。它们由\cref{b5:thm:specht-standard-basis}线性无关，故这些系数都为零，矛盾。因此
\(F_j=S_R^\lambda\cap L_{r_j}\)，尤其 \(F_j\) 是子模。

限制 \(d_{r_j}\) 到 \(F_j\)。上述公式说明它的核恰为 \(F_{j-1}\)：其余标准基向量的像是互不相同的标准基向量。而删格后的每个标准表补上 \(n\) 都成为原形状的标准表，故映射满射。这证明滤过商的同构。最后，在特征零域上，\((n-1)!\) 可逆，\cref{b5:thm:maschke}使每个短正合列分裂，遂得\eqref{b5:eq:specht-restriction}。
\end{proof}

滤过由固定数字所在行的条件确定，不需要选择不变补空间。特征零下的直和同构则一般没有这样的唯一性；在模特征中，滤过仍存在，却不能据此断言各层必能直和分开。

\subsection{标准表递推与钩长公式}
\label{b5:subsec:0605:02}

由\cref{b5:lem:standard-remove-largest}，标准表数满足
\begin{equation}\label{b5:eq:standard-tableau-recursion}
f^\varnothing=1,\qquad
f^\lambda=\sum_{x\text{ 为可移方格}}f^{\lambda\setminus x}.
\end{equation}
下面从这个递推推出乘积公式，不借用后文的对称函数或 Frobenius 特征映射。

\begin{lemma}\label{b5:lem:tableau-vandermonde}
设 \(n\) 为非负整数，\(\lambda\vdash n\)，\(f^\lambda\) 为形状 \(\lambda\) 的标准表数。取不小于 \(\lambda\) 长度的正整数 \(r\)，将 \(\lambda\) 末尾补零至 \(r\) 项，并令
\(l_i=\lambda_i+r-i\)。则
\begin{equation}\label{b5:eq:tableau-vandermonde}
f^\lambda
=n!\frac{\prod_{1\leq i<j\leq r}(l_i-l_j)}
{\prod_{i=1}^r l_i!}.
\end{equation}
\end{lemma}
\begin{proof}
记等式右边为 \(D(\lambda)\)。当 \(n=0\) 时，\(l_i=r-i\)，Vandermonde 乘积\footnote{范德蒙德（Alexandre-Théophile Vandermonde，1735-1796），法国数学家，研究行列式与对称函数。}恰为 \(\prod_i(r-i)!\)，所以 \(D(\varnothing)=1\)。

对 \(n>0\)，若第 \(i\) 行末格可移，则删格把 \(l_i\) 换成 \(l_i-1\)，直接相除得到
\[
\frac{D(\lambda\setminus x_i)}{D(\lambda)}
=\frac{l_i}{n}
\prod_{j\ne i}\frac{l_i-l_j-1}{l_i-l_j}.
\]
当该行不可移时，相应表达式为零：若 \(i<r\) 且 \(\lambda_i=\lambda_{i+1}\)，则 \(l_i-l_{i+1}=1\)，乘积中有零因子；若 \(i=r\) 且 \(\lambda_r=0\)，则 \(l_r=0\)。所以要证明 \(D\) 满足\eqref{b5:eq:standard-tableau-recursion}，只须证明
\begin{equation}\label{b5:eq:hook-rational-identity}
\sum_{i=1}^r l_i
\prod_{j\ne i}\frac{l_i-l_j-1}{l_i-l_j}
=\sum_i l_i-\frac{r(r-1)}2=n.
\end{equation}

令 \(P(z)=\prod_i(z-l_i)\)。不同的 \(l_i\) 互异，因而有理函数 \(z\cdot P(z-1)/P(z)\) 的部分分式只可能含这些点对应的一阶分母；在 \(l_i\) 处的留数为
\[
\frac{l_i\cdot P(l_i-1)}{P'(l_i)}
=-l_i\prod_{j\ne i}\frac{l_i-l_j-1}{l_i-l_j}.
\]
这里“留数”只指部分分式中 \(1/(z-l_i)\) 的系数，不需要复分析。另一方面，在形式变量 \(z^{-1}\) 中展开，
\[
\begin{aligned}
\frac{P(z-1)}{P(z)}
&=\prod_i\left(1-\frac1{z-l_i}\right)\\
&=1-\frac rz+
\frac{r(r-1)/2-\sum_i l_i}{z^2}+O(z^{-3}).
\end{aligned}
\]
这里 \(O(z^{-3})\) 只表示关于 \(z^{-1}\) 的次数至少为 \(3\) 的形式幂级数。乘以 \(z\) 后，\(z^{-1}\) 的系数等于全部部分分式系数之和。代入刚算的系数即得\eqref{b5:eq:hook-rational-identity}。因此 \(D\) 与 \(f\) 有同一个初值和同一个降低格子数的递推，对 \(n\) 归纳得到所求公式。
\end{proof}

\begin{definition}\label{b5:def:hook}
设 \(\lambda\) 为分拆，\(\lambda'\) 为共轭分拆，\([\lambda]\) 为 Young 图，\(x=(i,j)\in[\lambda]\)。由 \(x\)、它右方同行的格子及下方同列的格子组成的集合称为该格的钩，其格子数
\[
h_\lambda(i,j)=\lambda_i-j+\lambda'_j-i+1
\]
称为\term[gouchang]{钩长}[hook length]。
\end{definition}

\begin{theorem}[钩长公式]\label{b5:thm:hook-length}
设 \(n\) 为非负整数，\(\lambda\vdash n\)，\(k\) 为任意域，\(h_\lambda(x)\) 为 Young 图 \([\lambda]\) 中方格 \(x\) 的钩长。形状 \(\lambda\) 的标准表数 \(f^\lambda\)，以及 Specht 模 \(S_k^\lambda\) 的 \(k\)-维数，等于
\begin{equation}\label{b5:eq:hook-length}
f^\lambda=\frac{n!}{\prod_{x\in[\lambda]}h_\lambda(x)}.
\end{equation}
这个等式称为\term[gouchang-gongshi]{钩长公式}[hook-length formula]。
\end{theorem}
\begin{proof}
标准表数与维数的相等已经由标准基定理给出。只须把\eqref{b5:eq:tableau-vandermonde}改写为钩长乘积。仍取 \(l_i=\lambda_i+r-i\)。固定第 \(i\) 行，对该行的第 \(j\) 格令
\[
q_j=j-1+r-\lambda'_j,\qquad
h_\lambda(i,j)=l_i-q_j.
\]
这些 \(q_j\) 严格递增，且属于 \(\{0,\ldots,l_i-1\}\)。它们不等于任何 \(l_a\)：若 \(l_a=q_j\)，则
\(\lambda_a-j+1=a-\lambda'_j\)；当 \(a\leq\lambda'_j\) 时左边为正、右边非正，当 \(a>\lambda'_j\) 时左边非正、右边为正，均矛盾。

对 \(a>i\)，有 \(0\leq l_a<l_i\)，且这些数互异。在长度为 \(l_i\) 的整数区间中，已有 \(\lambda_i\) 个 \(q_j\) 及 \(r-i\) 个 \(l_a\)，总数恰为 \(l_i\)，所以它们正好将该区间分完。因此
\[
\prod_{j=1}^{\lambda_i}h_\lambda(i,j)
=\frac{l_i!}{\prod_{a>i}(l_i-l_a)}.
\]
对 \(i\) 连乘，得到
\[
\prod_{x\in[\lambda]}h_\lambda(x)
=\frac{\prod_i l_i!}{\prod_{i<a}(l_i-l_a)}.
\]
代入\eqref{b5:eq:tableau-vandermonde}即得\eqref{b5:eq:hook-length}。
\end{proof}

Etingof 等人在\cite[Section 5.17, pp. 118--119]{EtingofEtAl2011RepresentationTheory}中也把 Vandermonde 形式改写成钩长乘积，其前一步使用 Frobenius 特征标公式；本节的前一步由标准表递推给出。

\subsection{钩形、两行与低阶分拆的维数}
\label{b5:subsec:0605:03}

钩形分拆 \(\lambda=(n-r,1^r)\)，其中 \(0\leq r<n\)，左上角的钩长是 \(n\)，其余第一行的钩长为 \(n-r-1,\ldots,1\)，第一列其余钩长为 \(r,\ldots,1\)。因此
\begin{equation}\label{b5:eq:hook-shape-dimension}
f^{(n-r,1^r)}
=\frac{n!}{n(n-r-1)!r!}
=\binom{n-1}{r}.
\end{equation}
也可直接计数：左上角必为 \(1\)，只需从 \(2,\ldots,n\) 选 \(r\) 个数字放到第一列，行列内均按递增次序排定。

\ProseParagraphBreak

对于两行分拆 \((n-r,r)\)，其中 \(0\leq r\leq n/2\)，取两个参数
\(l_1=n-r+1,l_2=r\)，由\eqref{b5:eq:tableau-vandermonde}得到
\begin{equation}\label{b5:eq:two-row-dimension}
f^{(n-r,r)}
=\frac{n!(n-2r+1)}{(n-r+1)!r!}
=\binom nr-\binom n{r-1},
\end{equation}
约定 \(\binom n{-1}=0\)。例如 \(f^{(3,2)}=10-5=5\)，
\(f^{(3,3)}=20-15=5\)。这两项分别计算 \(S_5\) 与 \(S_6\) 的表示维数，不能用它们说明同一个群的等维表示未必同构。

若要比较同一个群的表示，可取特征零域 \(k\) 上的两个 \(S_5\) 模
\(S_k^{(3,2)}\) 与 \(S_k^{(2,2,1)}\)。两个图形互为转置，对应格子的钩长相同，钩长乘积均为 \(24\)，所以两者的维数均为 \(5!/24=5\)。但两个分拆不同，由\cref{b5:prop:specht-dominance}，这两个 \(k[S_5]\) 模并不同构。因此，即使群、基域和维数都相同，也不能只凭维数判定表示同构。

\subsection{Frobenius 互反与诱导分支}
\label{b5:subsec:0605:04}

\begin{corollary}[诱导分支]\label{b5:cor:specht-induction-branching}
设 \(k\) 为特征零域，\(n\ge1\) 为整数，\(\mu\vdash n-1\)，\(S_k^\mu\) 为 Specht 左 \(k[S_{n-1}]\)-模，\(S_{n-1}\le S_n\) 为固定数字 \(n\) 的子群。则
\[
\operatorname{Ind}_{S_{n-1}}^{S_n}S_k^\mu
\cong\bigoplus_{\substack{\lambda\vdash n\\
\lambda\setminus\mu\text{ 为一个方格}}}S_k^\lambda,
\]
各项重数都为 \(1\)。
\end{corollary}
\begin{proof}
由分类与 Maschke 定理，诱导模可以分成 Specht 模的直和。绝对不可约性说明每个 Specht 模的自同态环为 \(k\)，所以 \(S_k^\lambda\) 的重数等于
\[
\dim_k\operatorname{Hom}_{k[S_n]}
\left(\operatorname{Ind}_{S_{n-1}}^{S_n}S_k^\mu,S_k^\lambda\right).
\]
\cref{b5:thm:frobenius-reciprocity}把它化为
\(\dim_k\operatorname{Hom}_{k[S_{n-1}]}
(S_k^\mu,\operatorname{Res}S_k^\lambda)\)。
由\cref{b5:thm:specht-branching}和不同形状的非同构性，这个数恰在删去 \(\lambda\) 的一个可移方格能得到 \(\mu\) 时为 \(1\)，否则为 \(0\)。
\end{proof}

% !TeX root = ../../Book5.tex
% 纲要标题：### 6.6 钩长公式与分支规则的应用
% 纲要位置：工作记录/计划纲要.md，第 3942 行。
% * 钩长公式的应用
% * 限制、诱导与迭代分支
% * 标准表与具体维数计算
% * 应用第6.5节已经证明的钩长公式与分支规则，系统计算低阶、钩形和两行分拆的表示维数，并核验分支两侧的维数恒等式
% ---
\section{钩长公式与分支规则的应用}
\label{b5:sec:0606}

公式的用途不止是得到一个整数。钩长乘积、标准表递推和表示的分支分解分别从图形、组合和群作用三方面描述同一维数；交叉检查它们，可以发现遗漏的可移格，也能帮助我们选取便于计算的模型。

\subsection{钩长公式的应用}
\label{b5:subsec:0606:01}

考虑 \(\lambda=(3,2,1)\)。按各格原位置排列，钩长为
\[
\begin{array}{ccc}
5&3&1\\
3&1&\\
1&&
\end{array}.
\]
因此钩长乘积为 \(45\)，在特征零下，对应的 \(S_6\)-不可约表示维数为
\(6!/45=16\)。钩长计算中，每个格子本身只数一次；不能把同行格数与同列格数直接相加而忘记减去重复的格子。

\ProseParagraphBreak

转置图形保持所有钩长的多重集合，所以
\(f^\lambda=f^{\lambda'}\)。例如 \((4,1,1)\) 与 \((3,1,1,1)\) 的维数都为 \(10\)。这一维数相等可以直接从公式读出；它本身尚未说明两种表示之间存在何种张量关系。

\subsection{限制、诱导与迭代分支}
\label{b5:subsec:0606:02}

\((3,2,1)\) 有三个可移格，故特征零下
\[
\operatorname{Res}_{S_5}^{S_6}S^{(3,2,1)}
\cong S^{(2,2,1)}\oplus S^{(3,1,1)}\oplus S^{(3,2)}.
\]
三项维数为 \(5,6,5\)，合计 \(16\)。反过来，从 \((3,2)\) 添格得到
\[
\operatorname{Ind}_{S_5}^{S_6}S^{(3,2)}
\cong S^{(4,2)}\oplus S^{(3,3)}\oplus S^{(3,2,1)},
\]
维数为 \(9+5+16=30=6\cdot5\)，与诱导模维数相符。

\begin{corollary}\label{b5:cor:iterated-specht-branching}
设 \(k\) 为特征零域，\(m,n\) 为满足 \(0\leq m\leq n\) 的整数，\(\lambda\vdash n\)、\(\mu\vdash m\)，\(S_k^\lambda,S_k^\mu\) 为相应 Specht 模，\(S_m\le S_n\) 为固定 \(m+1,\ldots,n\) 的子群。\(S_k^\mu\) 在
\(\operatorname{Res}_{S_m}^{S_n}S_k^\lambda\) 中的重数等于从 \(\mu\) 出发每次添一格到达 \(\lambda\) 的图形链条数。若 \([\mu]\subseteq[\lambda]\)，这个数也等于在差图 \([\lambda]\setminus[\mu]\) 中填入 \(m+1,\ldots,n\)，使行列均递增的表的个数；若没有该包含关系，重数为零。
\end{corollary}
\begin{proof}
连续应用\cref{b5:thm:specht-branching}，每次删格所给成分重数均为 \(1\)，所以最终重数恰为经过 \(\mu\) 的删格链条数。反向读取链条，在从大小 \(a-1\) 到大小 \(a\) 新添的格子写 \(a\)。Young 图只能向右或向下扩张，所以每个新数字都大于它左边、上边已经填入的数字，得到所述表。

反之，给定这种填法，依次保留 \([\mu]\) 及数字不超过 \(a\) 的格子。行列递增保证所得格子集向左、向上封闭，所以它仍是 Young 图；这恢复了原来的添格链条。两种操作互逆。
\end{proof}

差图称为\term[xie-Young-tu]{斜 Young 图}[skew Young diagram]。这里仅用标准填数计算迭代分支重数，并未预先使用下一章的斜 Schur 函数理论。

\subsection{标准表与具体维数计算}
\label{b5:subsec:0606:03}

标准基还能直接给出生成元的矩阵。对 \(S^{(2,1)}\)，取
\[
u=e_{\begin{smallmatrix}1&2\\3&\end{smallmatrix}}=v_3-v_1,
\qquad
v=e_{\begin{smallmatrix}1&3\\2&\end{smallmatrix}}=v_2-v_1.
\]
置换 \(s_1=(12)\)、\(s_2=(23)\) 满足
\[
s_1u=u-v,\quad s_1v=-v,\qquad
s_2u=v,\quad s_2v=u.
\]
因此在有序基 \(u,v\) 中，
\[
\rho(s_1)=
\begin{pmatrix}1&0\\-1&-1\end{pmatrix},
\qquad
\rho(s_2)=
\begin{pmatrix}0&1\\1&0\end{pmatrix}.
\]
这两个整矩阵直接满足
\(\rho(s_i)^2=I\) 与
\(\rho(s_1)\rho(s_2)\rho(s_1)=
\rho(s_2)\rho(s_1)\rho(s_2)\)。其迹在单位元、换位、三轮换上分别为 \(2,0,-1\)，与第一编的 \(S_3\) 特征标表一致。即使换到特征 \(3\)，这些矩阵仍定义一个表示，只是它会出现\cref{b5:exm:specht-mod-three}中的不变直线。

\ProseParagraphBreak

再看特征零下的 \(S^{(2,2)}\)。它有两个标准表
\[
t_1=\begin{array}{cc}1&2\\3&4\end{array},
\qquad
t_2=\begin{array}{cc}1&3\\2&4\end{array}.
\]
令 \(x,y,z\) 分别表示无序配对
\(12|34,13|24,14|23\)，得到一个三点置换模。把 \((2,2)\)-行表类的两行视为无序的两对，给出等变映射
\(M^{(2,2)}\rightarrow kx\oplus ky\oplus kz\)。直接展开四项列反对称和可得
\[
e_{t_1}\longmapsto2(x-z),\qquad
e_{t_2}\longmapsto2(y-z).
\]
因 \(2\) 可逆，这两个像构成坐标和为零平面的基，因此上述映射限制为同构
\[
S^{(2,2)}\cong W,
\qquad W=\{ax+by+cz:a+b+c=0\}.
\]
在当前特征零假设下，\(3\) 也可逆，故三点置换模分解为
\[
kx\oplus ky\oplus kz=k(x+y+z)\oplus W.
\]
其中 \(k(x+y+z)\) 是平凡子模，\(W\) 是它的不变补空间。这个模型还表明 Klein 四元子群在 \(W\) 上平凡作用。

\subsection{低阶分拆与分支维数恒等式}
\label{b5:subsec:0606:04}

特征零下，\(S_4\) 的五个不可约维数按分拆次序为
\[
\begin{array}{c|ccccc}
\lambda &(4)&(3,1)&(2,2)&(2,1,1)&(1^4)\\ \hline
f^\lambda&1&3&2&3&1
\end{array},
\qquad
1^2+3^2+2^2+3^2+1^2=24.
\]
例如 \((2,2)\) 虽有两行，但只有第二行末格可移，故它限制到 \(S_3\) 后仍为一个 \(S^{(2,1)}\)；\((3,1)\) 有两个可移格，限制则为
\(S^{(2,1)}\oplus S^{(3)}\)。

\ProseParagraphBreak

对 \(S_5\)，按
\[
(5),(4,1),(3,2),(3,1,1),(2,2,1),(2,1,1,1),(1^5)
\]
排列，维数为 \(1,4,5,6,5,4,1\)，平方和为 \(120\)。做这种核验时，应当先分别列出全部分拆及各自的钩长，再求和；如果只凭平方和反推一张维数表，就失去了检验各个构造的作用。


\begin{chaptersummary}
Specht 模把行方向的对称和列方向的反对称放进同一个表示。Garnir 关系按严格下降的列字完成标准化，表类展开的三角性则保证标准列交错元线性无关。因此标准表给出一组整系数基，改变系数环时不需要重新寻找基。

\ProseParagraphBreak

在特征零下，表类配对限制到 Specht 模仍非退化，子模定理遂给出不可约性；支配序排除不同分拆间的同构，最后用共轭类数完成分类。限制到 \(S_{n-1}\) 时，最大数字所在的可移格组织出一个滤过，完全可约性把它分裂成删格直和。钩长公式从标准表递推得到，是计算维数的结果，而不是前面分类证明的假设。
\end{chaptersummary}

\BookExercises

\begin{exercise}\label{b5:ex:06-diagram-groups}
设 \(\lambda=(4,2,1)\)。求共轭分拆、全部可移方格、行群与列群的阶，以及 \(M_k^\lambda\) 的维数。
\end{exercise}
\begin{solution}
共轭分拆为 \((3,2,1,1)\)，可移方格是 \((1,4),(2,2),(3,1)\)。行群阶为 \(4!2!1!=48\)，列群阶为 \(3!2!1!1!=12\)。行等价类与行群左陪集一一对应，所以维数为 \(7!/48=105\)。这些数均来自集合基，与域的特征无关。
\end{solution}

\begin{exercise}\label{b5:ex:06-cycle-type}
在 \(S_8\) 中，求轮换类型为 \((3,2,2,1)\) 的共轭类大小。求这个类型的置换的阶与符号。
\end{exercise}
\begin{solution}
中心化子阶为 \(3\cdot2^2\cdot2!=24\)，故类大小为 \(8!/24=1680\)。不交轮换之积的阶为长度的最小公倍数，因而为 \(6\)。三轮换为偶置换，两个二轮换的符号相乘也为正，所以总体符号为 \(1\)。
\end{solution}

\begin{exercise}\label{b5:ex:06-same-tabloid}
在形状 \((2,1)\) 中，取 \(t=\begin{smallmatrix}1&2\\3&\end{smallmatrix}\) 与
\(u=\begin{smallmatrix}2&1\\3&\end{smallmatrix}\)。证明它们的行表类相同，而在任意非零交换环上，对应列交错元不同。
\end{exercise}
\begin{solution}
二者第一行都是集合 \(\{1,2\}\)，第二行都是 \(\{3\}\)，所以表类相同。用第二行数字标记基向量，有
\(e_t=v_3-v_1\)、\(e_u=v_3-v_2\)，差为 \(v_2-v_1\)。这是两个不同自由基向量之差，在非零环上不为零。因此列反对称化所需的资料超过行等价类本身。
\end{solution}

\begin{exercise}\label{b5:ex:06-standard-hook-module}
设 \(R\) 为交换含幺环，\(n\geq3\)。证明 \(S_R^{(n-1,1)}\) 等于 \(n\) 点置换模中的坐标和为零子模，并写出一组自由基。
\end{exercise}
\begin{solution}
用 \(v_i\) 标记第二行只含 \(i\) 的行表类。每个表只有第一列含两格，若其中上、下数字分别为 \(j,i\)，则 \(e_t=v_i-v_j\)。因此 Specht 模由全部差向量生成，包含于坐标和为零子模。反过来，若 \(\sum_i a_i=0\)，则
\[
\sum_{i=1}^n a_iv_i=\sum_{i=1}^{n-1}a_i(v_i-v_n).
\]
这些 \(n-1\) 个差向量独立，因为各个 \(v_i\) 的系数直接读出 \(a_i\)。故它们给出所求基。
\end{solution}

\begin{exercise}\label{b5:ex:06-garnir-worked}
把形状 \((3,1)\) 的表
\(t=\begin{smallmatrix}2&1&3\\4&&\end{smallmatrix}\)
对应的 \(e_t\) 用标准列交错元表示，并指出一次 Garnir 重写中用到的数字集合。
\end{exercise}
\begin{solution}
取 \(A=\{2,4\}\)、\(B=\{1\}\)，代表元为 \(1,(12),(14)\)，故
\(e_t=e_{(12)t}+e_{(14)t}\)。第一项的表为
\(\begin{smallmatrix}1&2&3\\4&&\end{smallmatrix}\)。第二项的第一列为上 \(2\)、下 \(1\)，排序后变为
\(\begin{smallmatrix}1&4&3\\2&&\end{smallmatrix}\)，同时出现负号。再交换只含一格的第二、第三列中的数字，不改变第一列及表类，所以该列交错元等于
\(e_{\begin{smallmatrix}1&3&4\\2&&\end{smallmatrix}}\)。因此
\[
e_t=e_{\begin{smallmatrix}1&2&3\\4&&\end{smallmatrix}}
-e_{\begin{smallmatrix}1&3&4\\2&&\end{smallmatrix}}.
\]
也可用第二行数字基核验为 \(v_4-v_2=(v_4-v_1)-(v_2-v_1)\)。
\end{solution}

\begin{exercise}\label{b5:ex:06-standard-22-expansion}
将形状 \((2,2)\) 的两个标准列交错元展开为表类基，并从展开式直接证明它们在 \(\mathbb Z/4\mathbb Z\) 上仍独立。
\end{exercise}
\begin{solution}
用 \(\{ab|cd\}\) 记录两行数字集合。对
\(t_1=\begin{smallmatrix}1&2\\3&4\end{smallmatrix}\)、
\(t_2=\begin{smallmatrix}1&3\\2&4\end{smallmatrix}\)，有
\[
\begin{aligned}
e_{t_1}&=\{12|34\}-\{23|14\}-\{14|23\}+\{34|12\},\\
e_{t_2}&=\{13|24\}-\{23|14\}-\{14|23\}+\{24|13\}.
\end{aligned}
\]
第一式在 \(\{12|34\}\) 处的系数为 \(1\)，第二式在该处为 \(0\)；在 \(\{13|24\}\) 处情形相反。因此任何关系 \(ae_{t_1}+be_{t_2}=0\) 都迫使 \(a=b=0\)，即使系数环有零因子也一样。
\end{solution}

\begin{exercise}\label{b5:ex:06-standard-action-integral}
证明任意 \(S_n\) 元素在整系数标准基上的矩阵属于
\(\operatorname{GL}_{f^\lambda}(\mathbb Z)\)。由此证明特征零 Specht 模的特征标取整数值。
\end{exercise}
\begin{solution}
群元素把 \(e_t\) 送到 \(e_{gt}\)，Garnir 标准化只用整数系数，所以作用矩阵的条目为整数。逆置换也有整矩阵，两个矩阵乘积为单位阵，故原矩阵在整数环上可逆，其行列式为 \(1\) 或 \(-1\)。特征标是该整矩阵的迹，因此为整数；把矩阵送到任意特征零域不会改变这个整数来源。
\end{solution}

\begin{exercise}\label{b5:ex:06-young-mod-three}
取形状 \((2,1)\) 的任意表 \(t\)。证明 Young 对称化子在整数群代数中满足 \(c_t^2=3c_t\)，但在 \(\mathbb F_3[S_3]\) 中是非零的平方零元。这与标准基定理是否矛盾？
\end{exercise}
\begin{solution}
特征零下，标准表数为 \(2\)，故\cref{b5:prop:young-quasi-idempotent}给出
\(c_t^2=(3!/2)c_t=3c_t\)。两边都是整系数群代数元素，等式在 \(\mathbb Q[S_3]\) 成立即说明每个整数系数都相同，故也在 \(\mathbb Z[S_3]\) 成立。约化模 \(3\) 后平方为零，但单位元的系数仍为 \(1\)，所以元素非零。标准基定理只断言 Specht 模自由及其基的形式，并未断言对称化子在任意特征都可归一化成幂等元，故无矛盾。
\end{solution}

\begin{exercise}\label{b5:ex:06-dominance-incomparable}
在特征零域上，证明
\[
\begin{aligned}
\operatorname{Hom}_{S_6}\bigl(S^{(4,1,1)},M^{(3,3)}\bigr)&=0,\\
\operatorname{Hom}_{S_6}\bigl(S^{(3,3)},M^{(4,1,1)}\bigr)&=0.
\end{aligned}
\]
再说明支配条件成立时，为何不能仅凭本章的条件就宣称相应 Hom 必非零。
\end{exercise}
\begin{solution}
\((4,1,1)\) 的前两项和为 \(5<6\)，所以它不支配 \((3,3)\)；\((3,3)\) 的第一项 \(3<4\)，所以也不支配 \((4,1,1)\)。两项消失均由\cref{b5:prop:specht-dominance}的逆否命题得到。该命题对不同形状只证明了非零 Hom 的必要条件，没有证明充分性；把必要条件倒过来使用还需要另一个证明。
\end{solution}

\begin{exercise}\label{b5:ex:06-permutation-split}
设 \(k\) 为域，\(n\geq3\)，且 \(n\ne0\) 于 \(k\)。证明 \(n\) 点置换模分解为平凡直线与 \(S_k^{(n-1,1)}\) 的直和。若 \(\operatorname{char}k\mid n\)，说明这个特定平凡直线为什么不能作为补空间。
\end{exercise}
\begin{solution}
令 \(\mathbf v=v_1+\cdots+v_n\)。向量 \(w=\sum_i a_iv_i\) 唯一分解为
\[
w=\frac{\sum_i a_i}{n}\mathbf v+
\left(w-\frac{\sum_i a_i}{n}\mathbf v\right).
\]
第二项的坐标和为零，第一项属于固定直线。二者交为零，因为 \(c\mathbf v\) 的坐标和为 \(nc\)。当特征整除 \(n\) 时，\(\mathbf v\) 本身坐标和为零，固定直线反而包含在 Specht 模中，所以不能成为它的补空间。
\end{solution}

\begin{exercise}\label{b5:ex:06-hook-gram}
设 \(k\) 为域，\(n\geq3\)。在 \(S_k^{(n-1,1)}\) 的基
\(v_1-v_n,\ldots,v_{n-1}-v_n\) 中计算表类配对的 Gram 行列式。证明特征不整除 \(n\) 时这个 Specht 模不可约，特征整除 \(n\) 时可约。
\end{exercise}
\begin{solution}
对角项为 \(2\)，非对角项为 \(1\)，所以矩阵为 \(I_{n-1}+J_{n-1}\)，其中 \(J\) 各项均为 \(1\)。在整数上，其特征值为 \(n\) 一次、\(1\) 共 \(n-2\) 次，故行列式为 \(n\)；这个整数多项式恒等式可约化到任意域。若 \(n\ne0\)，配对非退化，子模定理与正文的论证同样给出不可约性，并不需要特征零。若 \(n=0\)，上题的非零平凡直线包含在该模中；因模维数 \(n-1\geq2\)，它是真子模，故可约。
\end{solution}

\begin{exercise}\label{b5:ex:06-gram-22}
计算 \(S_{\mathbb Z}^{(2,2)}\) 在两个标准列交错元基下的 Gram 矩阵。求出在哪些素数特征下此配对退化，并求特征 \(3\) 时的根。
\end{exercise}
\begin{solution}
每个标准向量有四项系数为 \(\pm1\)，故各自的平方范数为 \(4\)。两者共有的两项系数均为 \(-1\)，所以配对为 \(2\)。Gram 矩阵与行列式为
\[
\begin{pmatrix}4&2\\2&4\end{pmatrix},
\qquad 12.
\]
故且仅在特征 \(2\) 或 \(3\) 退化。特征 \(3\) 下矩阵为
\(\begin{psmallmatrix}1&-1\\-1&1\end{psmallmatrix}\)，根由 \(e_{t_1}+e_{t_2}\) 生成。特征 \(2\) 时矩阵为零；配对退化本身不能反向推出所有相应 Specht 模必可约。
\end{solution}

\begin{exercise}\label{b5:ex:06-self-dual}
设 \(k\) 为特征零域。证明每个 \(S_k^\lambda\) 同构于其对偶，并推出 \(\chi^\lambda(g)=\chi^\lambda(g^{-1})\)。
\end{exercise}
\begin{solution}
映射 \(v\mapsto B(v,-)\) 从 \(S_k^\lambda\) 到其线性对偶，由非退化性为同构。对偶作用定义为
\((g\varphi)(w)=\varphi(g^{-1}w)\)，而不变性给出
\(B(gv,w)=B(v,g^{-1}w)\)，所以映射等变。对偶矩阵为原来逆矩阵的转置，其迹等于 \(\chi^\lambda(g^{-1})\)；同构表示特征标相等，得到结论。
\end{solution}

\begin{exercise}\label{b5:ex:06-s4-algebra}
设 \(k\) 为特征零域。写出 \(k[S_4]\) 的矩阵代数分解，并求其中心维数及交换子所生成的线性子空间的维数。
\end{exercise}
\begin{solution}
维数列表给出
\[
k[S_4]\cong k\times M_3(k)\times M_2(k)\times M_3(k)\times k.
\]
各矩阵块中心均为标量，故中心维数为 \(5\)。\(M_d(k)\) 中交换子都迹为零；反过来，
\(E_{ij}=[E_{ii},E_{ij}]\) 对 \(i\ne j\) 成立，
\(E_{ii}-E_{jj}=[E_{ij},E_{ji}]\)，这些矩阵生成全部迹零子空间。因此交换子线性生成空间的维数为
\(0+8+3+8+0=19\)。
\end{solution}

\begin{exercise}\label{b5:ex:06-restriction-42}
在特征零域上分解
\(\operatorname{Res}_{S_5}^{S_6}S^{(4,2)}\)，并独立核验两边维数。
\end{exercise}
\begin{solution}
可移格为第一行末格与第二行末格，故限制为
\(S^{(3,2)}\oplus S^{(4,1)}\)。两行公式给出
\(f^{(4,2)}=\binom62-\binom61=9\)，而
\(f^{(3,2)}=5\)、\(f^{(4,1)}=4\)，所以 \(9=5+4\)。
\end{solution}

\begin{exercise}\label{b5:ex:06-induction-211}
在特征零域上分解
\(\operatorname{Ind}_{S_4}^{S_5}S^{(2,1,1)}\)，并核验其维数。
\end{exercise}
\begin{solution}
可以添格的位置分别产生 \((3,1,1)\)、\((2,2,1)\) 和 \((2,1,1,1)\)；不能在第三行末添格而不同时增长第二行。诱导分支给出这三个 Specht 模各一次。其维数分别为 \(6,5,4\)，和为 \(15\)；原模是维数 \(\binom32=3\) 的钩形表示，诱导指数为 \(5\)，故总维数同样为 \(15\)。
\end{solution}

\begin{exercise}\label{b5:ex:06-nonsplit-modular-branch}
在特征 \(2\) 的域上，把 \(S^{(2,1)}\) 限制到固定数字 \(3\) 的 \(S_2\)。证明分支滤过的两个商均为平凡一维表示，但该限制不是两个平凡表示的直和。
\end{exercise}
\begin{solution}
在正文的有序基 \(u=v_3-v_1,v=v_2-v_1\) 中，生成元 \((12)\) 的矩阵约化为
\(\begin{psmallmatrix}1&0\\1&1\end{psmallmatrix}\)。子空间 \(kv\) 固定，商空间也固定；它正是最高可移行给出的滤过层，符号表示在特征 \(2\) 等于平凡表示。但是生成元并非单位算子，两个平凡表示的直和中每个群元素都必须为单位算子，因此分裂不成立。
\end{solution}

\begin{exercise}\label{b5:ex:06-rectangle}
计算 \(f^{(3,3,3)}\)，并说明为什么
\(f^{(3,3,3)}=f^{(3,3,2)}\)。后一个等式是否说明两个不同阶对称群上的表示同构？
\end{exercise}
\begin{solution}
九个钩长按行是 \((5,4,3),(4,3,2),(3,2,1)\)，乘积为 \(8640\)，所以 \(f^{(3,3,3)}=9!/8640=42\)。矩形只有右下角一格可移，故标准表递推或限制分支给出 \(f^{(3,3,2)}=42\)。两个模原本分别属于 \(S_9\) 与 \(S_8\) 的表示范畴；正确结论是前者限制到 \(S_8\) 后同构于后者，不能省去限制操作。
\end{solution}

\begin{exercise}\label{b5:ex:06-421-dimension}
用钩长公式求 \(f^{(4,2,1)}\)，再逐个计算其删格分支的维数，核验递推。
\end{exercise}
\begin{solution}
钩长按行是 \((6,4,2,1),(3,1),(1)\)，乘积为 \(144\)，所以维数为 \(7!/144=35\)。三个删格形状为 \((3,2,1),(4,1,1),(4,2)\)，维数分别为 \(16,\binom52=10,9\)，因此 \(35=16+10+9\)。钩形 \((4,1,1)\) 的总大小为 \(6\)，代入公式时应使用它自身的总大小。
\end{solution}

\begin{exercise}\label{b5:ex:06-iterated-branch}
在特征零域上，把 \(S^{(3,2,1)}\) 从 \(S_6\) 限制到 \(S_4\)，求所有不可约重数。
\end{exercise}
\begin{solution}
第一步产生 \((2,2,1),(3,1,1),(3,2)\)。第二步中，
\((2,2,1)\) 给出 \((2,1,1)\) 与 \((2,2)\)，
\((3,1,1)\) 给出 \((2,1,1)\) 与 \((3,1)\)，
\((3,2)\) 给出 \((2,2)\) 与 \((3,1)\)。因此
\[
\operatorname{Res}_{S_4}^{S_6}S^{(3,2,1)}
\cong \bigl(S^{(3,1)}\bigr)^{\oplus2}
\oplus\bigl(S^{(2,2)}\bigr)^{\oplus2}
\oplus\bigl(S^{(2,1,1)}\bigr)^{\oplus2}.
\]
维数为 \(2\cdot3+2\cdot2+2\cdot3=16\)。一步限制无重数，不意味着多步限制也无重数。
\end{solution}

\begin{exercise}\label{b5:ex:06-one-dimensional}
设 \(k\) 为特征零域，\(n\geq2\)。不使用钩长公式，证明 \(S_n\) 只有平凡表示和符号表示两个一维表示，并说明哪些 Specht 模为一维。
\end{exercise}
\begin{solution}
一维表示取值于交换群 \(k^\times\)，所以共轭的换位有相同的像 \(a\)。换位平方为单位，故 \(a^2=1\)，在特征零域中只能 \(a=1\) 或 \(a=-1\)。换位生成 \(S_n\)，所以表示由该值唯一决定，分别为平凡或符号表示。这两个表示在一个换位上的取值不同，故不同构。已经构造出的 \(S^{(n)}\) 与 \(S^{(1^n)}\) 分别给出它们；分类的非同构性排除其他一维 Specht 模。
\end{solution}

\begin{exercise}\label{b5:ex:06-22-character}
用三种无序配对的置换模型，计算 \(S_4\) 的 \(S^{(2,2)}\) 特征标在下列轮换类型上的值：
\[
1^4,\qquad(2,1,1),\qquad(2,2),\qquad(3,1),\qquad(4).
\]
\end{exercise}
\begin{solution}
配对置换模分成平凡直线与二维坐标和为零空间，所以目标特征标等于固定配对数减 \(1\)。单位元固定三个配对；换位 \((12)\) 只固定 \(12|34\)；双换位 \((12)(34)\) 固定三个；三轮换循环排列三种配对而没有固定点；四轮换 \((1234)\) 只固定 \(13|24\)。因此所求值依次为
\[
2,\quad0,\quad2,\quad-1,\quad0.
\]
\end{solution}

\begin{exercise}\label{b5:ex:06-s3-char-two}
设 \(k\) 为特征 \(2\) 的域。证明 \(S_k^{(2,1)}\) 仍绝对不可约，并与特征 \(3\) 的情形比较。
\end{exercise}
\begin{solution}
整数 Gram 行列式为 \(3\)，在特征 \(2\) 中为 \(1\)，故配对非退化。子模定理适用于任意域，因此该二维模不可约；任意扩域后行列式仍为 \(1\)，所以绝对不可约。特征 \(3\) 时行列式为零，并出现由 \(v_1+v_2+v_3\) 生成的非零真子模。因此“特征整除群阶”只说明不能统一使用 Maschke 定理，并不说明每一个具体 Specht 模都必可约。
\end{solution}

\begin{exercise}\label{b5:ex:06-add-box-dimensions}
设 \(\mu\vdash n-1\)。证明纯组合恒等式
\[
\sum_{\lambda:\,\lambda\setminus\mu\text{ 为一格}} f^\lambda=n f^\mu.
\]
说明为什么只用删格递推的形式，不能不加论证地把这个式子看作同一个递推。
\end{exercise}
\begin{solution}
在复数域上，诱导模 \(\operatorname{Ind}_{S_{n-1}}^{S_n}S^\mu\) 的维数等于指数 \(n\) 乘 \(f^\mu\)。由诱导分支，它是式中各个 \(S^\lambda\) 各一次的直和，比较维数得到等式。各个维数已经等于标准表数，所以这是关于整数的组合恒等式。删格递推固定大图形、对其所有小图形求和；本式固定小图形、对所有大图形求和，求和方向不同。Frobenius 互反性正是本章把两个方向联系起来的额外论证。
\end{solution}

% !TeX root = ../../Book5.tex
% 纲要标题：## 第7章　对称函数与 Schur–Weyl 对偶
% 纲要位置：工作记录/计划纲要.md，第 3951 行。

\chapter{对称函数与 Schur–Weyl 对偶}
\label{b5:ch:07}
\BookChapterNumber{b5:ch:07}

\begin{chapterabstract}
对称函数把对称群的特征标计算变成基变换与多项式乘法。本章建立 Schur 函数的行列式公式、Hall 内积和 Frobenius 特征映射，证明 Specht 特征标恰好对应 Schur 函数。随后从张量幂上的两个交换作用出发，用极化和半单矩阵块证明 Schur--Weyl 对偶，并确定哪些分拆在给定维数中出现。Schur 函子将对称幂与外幂统一起来，Pieri 规则及低阶 Littlewood--Richardson 计算给出可实际操作的张量分解。
\end{chapterabstract}

\begin{learninggoals}
\begin{itemize}
\item 在初等、完全、幂和与 Schur 基之间作计算，区分稳定对称函数与固定变量数的多项式。
\item 证明特征映射的等距性及诱导乘法公式，并据此读取 Specht 特征标。
\item 用极化计算张量幂中心化子，理解双中心化子结论所需的半单性。
\item 按分拆行数判断 Schur 函子是否为零，并用 Pieri 规则计算张量积。
\end{itemize}
\end{learninggoals}

本章使用上一章已证的 Specht 分类及 Young 置换模的支配三角性，也使用第一编的特征标正交关系与诱导表示。对称函数先在有理系数上讨论，再按需要复化；Schur--Weyl 对偶及 Schur 函子则在任意特征零域上成立。

\ProseParagraphBreak
本章不要求 Lie 代数的最高权理论。证明从有限群与有限维线性代数出发，最后得到一般线性群的一族不可约表示；更一般的分类在后文讨论。特征标公式、双中心化子及 Schur 多项式的联系可对照 Etingof 等人的《Introduction to Representation Theory》\cite[Sections 5.14--5.15, pp. 114--118; Sections 5.18--5.19, pp. 119--122; Sections 5.21--5.22, pp. 128--130]{EtingofEtAl2011RepresentationTheory}。

% !TeX root = ../../Book5.tex
% 纲要标题：### 7.1 对称函数
% 纲要位置：工作记录/计划纲要.md，第 3953 行。
% * 初等对称函数、完全对称函数与幂和对称函数
% * Schur 函数及行列式公式
% * 次数分次与内积
\section{对称函数}
\label{b5:sec:0701}

一个置换的共轭类由各长度的轮换数决定，对称多项式的基也常以分拆编号。要联系这两种参数，须先定义与变量个数变化相容的运算。本节中的无穷多个变量只是一种逐次稳定的记法，不涉及无穷级数的解析收敛。

\subsection{初等、完全与幂和对称函数}
\label{b5:subsec:0701:01}

\begin{definition}\label{b5:def:symmetric-function}
设 \(n\ge0\) 为整数。对每个整数 \(N\ge0\) 给定一个 \(N\) 元、\(n\) 次齐次有理系数对称多项式 \(f_N\)。若对所有 \(N\ge0\) 有
\[
f_{N+1}(x_1,\ldots,x_N,0)=f_N(x_1,\ldots,x_N).
\]
则称相容族 \((f_N)_{N\ge0}\) 为 \(n\) 次\term[duicheng-hanshu]{对称函数}[symmetric function]。这些相容族构成有理向量空间 \(\Lambda_{\mathbb Q}^n\)。令
\[
\Lambda_{\mathbb Q}=\bigoplus_{n\ge0}\Lambda_{\mathbb Q}^n.
\]
逐个变量数相乘给出这个空间上的分次代数结构。将系数扩充为复数，记为 \(\Lambda_{\mathbb C}\)。
\end{definition}

对分拆 \(\lambda\vdash n\)，令 \(m_\lambda(x_1,\ldots,x_N)\) 为指数多重集等于 \(\lambda\)（补零）的不同单项式之和；若 \(\ell(\lambda)>N\)，规定它为零。对称多项式在同一变量置换轨道上的单项式系数相同，所以 \(m_\lambda\) 构成 \(\Lambda_{\mathbb Q}^n\) 的一组基。特别地，固定次数 \(n\) 后，取 \(N\ge n\) 已能检测任意恒等式；无须同时处理无限多个非零项。

\begin{definition}\label{b5:def:ehp-symmetric}
在有理系数对称函数代数 \(\Lambda_{\mathbb Q}\) 中，对整数 \(r\ge1\)，由每个有限变量数 \(N\) 下的多项式
\[
e_r=\sum_{i_1<\cdots<i_r}x_{i_1}\cdots x_{i_r},\qquad
h_r=\sum_{i_1\le\cdots\le i_r}x_{i_1}\cdots x_{i_r},\qquad
p_r=\sum_i x_i^r
\]
组成的相容族，依次称为\term[chudeng-duicheng-hanshu]{初等对称函数}[elementary symmetric function]、\term[wanquan-duicheng-hanshu]{完全对称函数}[complete symmetric function]和\term[mihe-duicheng-hanshu]{幂和对称函数}[power-sum symmetric function]；各求和指标均在 \(\{1,\ldots,N\}\) 中取值。规定 \(e_0=h_0=1\)，而 \(e_r=h_r=0\) 对 \(r<0\) 成立。对分拆 \(\lambda\)，用 \(e_\lambda,h_\lambda,p_\lambda\) 分别表示按其各部分相乘。
\end{definition}

例如 \(h_2=m_{(2)}+m_{(1,1)}\)，\(e_2=m_{(1,1)}\)，\(p_2=m_{(2)}\)。完全对称函数允许指标重复，初等对称函数不允许；这个差别后来对应对称幂与外幂。

\begin{proposition}\label{b5:prop:symmetric-generators}
设 \(\Lambda_{\mathbb Q}\) 为有理系数对称函数代数，\(e_r,h_r,p_r\) 分别为初等、完全、幂和对称函数，\(t\) 为形式变量。在 \(\Lambda_{\mathbb Q}[\![t]\!]\) 中有
\begin{align}
E(t)&\coloneqq\sum_{r\ge0}e_rt^r=\prod_i(1+x_it),&
H(t)&\coloneqq\sum_{r\ge0}h_rt^r=\prod_i(1-x_it)^{-1},
\label{b5:eq:eh-generating}\\
E(-t)H(t)&=1,&
H(t)&=\exp\left(\sum_{r\ge1}\frac{p_r}{r}t^r\right).
\label{b5:eq:hp-generating}
\end{align}
乘积按每个有限变量数下的相容多项式族理解。对每个非负整数 \(n\)，三族 \(e_\lambda,h_\lambda,p_\lambda\)，其中 \(\lambda\vdash n\)，分别为 \(\Lambda_{\mathbb Q}^n\) 的基。
\end{proposition}
\begin{proof}
在有限变量中展开每个因子即可得到前两式；乘积相消得到第三式。对 \(H\) 取形式对数，用 \(-\log(1-z)=\sum_{r\ge1}z^r/r\) 得最后一式。每个系数只需有限次运算，故这些形式恒等式也在稳定族中成立。

先说明 \(e_\lambda\) 的基性质。在 \(N\) 个变量中取字典序 \(x_1>\cdots>x_N\)。任意非零对称多项式的首项指数可写成 \(a_1\ge\cdots\ge a_N\ge0\)。乘积
\[
e_1^{a_1-a_2}e_2^{a_2-a_3}\cdots e_N^{a_N}
\]
的首项恰为 \(x_1^{a_1}\cdots x_N^{a_N}\)，系数为一。依次减去相应倍数，首项严格下降；固定总次数只有有限个单项式，故过程终止，得到生成性。不同的上述乘积有不同首项，故线性无关。取 \(N\ge n\)，得到稳定的次数 \(n\) 结论。

由 \(E(-t)H(t)=1\)，\(h_r=(-1)^{r-1}e_r\) 加上只含 \(e_1,\ldots,e_{r-1}\) 的多项式；这是可逆三角变换。因此 \(h_r\) 也自由生成。对最后一式求导并乘以 \(t\)，得到
\begin{equation}\label{b5:eq:newton-h}
r h_r=\sum_{i=1}^r p_i h_{r-i}.
\end{equation}
其中 \(p_r=rh_r\) 加上较低指标的多项式；由于系数域为 \(\mathbb Q\)，还可反向除以 \(r\)。所以 \(p_r\) 同样自由生成，三族的次数 \(n\) 单项式均为基。
\end{proof}

有理系数在幂和基的结论中不可省略。例如 \(h_2=(p_1^2+p_2)/2\)，一般不能把完全对称函数写成幂和的整系数多项式。

\subsection{Schur 函数及行列式公式}
\label{b5:subsec:0701:02}

令 \(\delta_N=(N-1,N-2,\ldots,0)\)。对严格递减的非负整数列 \(\alpha\)，记
\[
a_\alpha(x)=\det(x_i^{\alpha_j})_{1\le i,j\le N},
\qquad
\Delta_N(x)=a_{\delta_N}(x)=\prod_{i<j}(x_i-x_j).
\]
交错多项式在 \(x_i=x_j\) 时为零，因而被每个 \(x_i-x_j\) 整除。这些一次因子两两互素，所以被它们的乘积整除。

\begin{definition}\label{b5:def:schur-function}
设 \(\lambda\) 为分拆，\(\ell(\lambda)\) 为其长度，\(N\ge\ell(\lambda)\) 为整数，并把 \(\lambda\) 补零到 \(N\) 项。在 \(\mathbb Q[x_1,\ldots,x_N]\) 中，对称多项式
\[
s_\lambda(x_1,\ldots,x_N)
=\frac{\det(x_i^{\lambda_j+N-j})_{1\le i,j\le N}}
{\prod_{1\le i<j\le N}(x_i-x_j)}
\]
称为相应的\term[Schur-duoxiangshi]{Schur 多项式}[Schur polynomial]。当 \(N<\ell(\lambda)\) 时规定其为零。它们形成的相容族称为\term[Schur-hanshu]{Schur 函数}[Schur function]，仍记 \(s_\lambda\)。
\end{definition}

分子交错，所以上述商确为多项式；交换两个变量时分子、分母同时变号，故商对称。若最后一个变量为零且 \(\lambda_N=0\)，展开分子、分母的最后一行，剩余行列式各提出 \(\prod_{i<N}x_i\)，恰得到 \(N-1\) 元版本。若 \(\lambda_N>0\)，分子被 \(\prod_i x_i^{\lambda_N}\) 整除，商也如此，代入 \(x_N=0\) 为零。这就证明了定义中的相容性。

\begin{lemma}[Cauchy 恒等式]\label{b5:lem:schur-cauchy}
设 \(x=(x_i)\)、\(y=(y_j)\) 为两组独立形式变量，\(s_\lambda\) 为有理系数 Schur 函数，\(p_\mu\) 为幂和对称函数的分拆乘积。按每个有限变量数下的形式幂级数、逐次理解无穷和，有
\begin{equation}\label{b5:eq:schur-cauchy}
\begin{aligned}
\prod_{i,j}(1-x_i y_j)^{-1}
&=\sum_\lambda s_\lambda(x)s_\lambda(y)
=\sum_\mu\frac{p_\mu(x)p_\mu(y)}{z_\mu},
\\
z_\mu&=\prod_{r\ge1}r^{m_r(\mu)}\cdot m_r(\mu)!.
\end{aligned}
\end{equation}
两和均遍历全部分拆；这里 \(m_r(\mu)\) 为分拆 \(\mu\) 中部分 \(r\) 的重数。
\end{lemma}
\begin{proof}
先取两组各 \(N\) 个变量。Cauchy 行列式恒等式为
\[
\det\left(\frac1{1-x_i y_j}\right)_{i,j=1}^N
=\frac{\Delta_N(x)\Delta_N(y)}{\prod_{i,j}(1-x_i y_j)}.
\]
确实，通分后的分子分别对两组变量交错；其关于每个 \(x_i\)、每个 \(y_j\) 的次数至多 \(N-1\)。除去两组 Vandermonde 因子后只能剩常数。该常数为一：在每个条目中展开 \(\sum_{a\ge0}x_i^a y_j^a\)，行列式最低的非零总次数来自互异指数 \(0,\ldots,N-1\)，其贡献就是 \(\Delta_N(x)\Delta_N(y)\)。

再对这个级数行列式使用行列式的多线性展开。重复指数的项为零，互异指数按递减顺序排列后得到
\[
\det\left(\sum_{a\ge0}x_i^a y_j^a\right)_{i,j=1}^N
=\sum_{\alpha_1>\cdots>\alpha_N\ge0}a_\alpha(x)a_\alpha(y).
\]
每个 \(\alpha\) 唯一写为 \(\lambda+\delta_N\)。除去 Vandermonde 因子，得到第一条展开。不同变量数可先补变量，再令多余变量为零。

另一方面，对乘积取形式对数，得 \(\sum_{r\ge1}p_r(x)p_r(y)/r\)。指数化后，第 \(r\) 项选取 \(m_r\) 次的系数为 \(1/(r^{m_r}m_r!)\)。逐个 \(r\) 相乘，得到最后一条展开。
\end{proof}

\begin{theorem}[Jacobi--Trudi 公式]\label{b5:thm:jacobi-trudi}
设 \(\lambda\) 为分拆，\(N\ge\ell(\lambda)\) 为整数，把 \(\lambda\) 补零至 \(N\) 项，\(s_\lambda\) 为 Schur 函数，\(h_r\) 为完全对称函数。在有理系数对称函数代数 \(\Lambda_{\mathbb Q}\) 中有
\begin{equation}\label{b5:eq:jacobi-trudi}
s_\lambda=\det(h_{\lambda_i-i+j})_{1\le i,j\le N}.
\end{equation}
其中 \(h_0=1\)，负指标的 \(h_r\) 为零。对每个非负整数 \(n\)，\(s_\lambda\)（\(\lambda\vdash n\)）构成 \(\Lambda_{\mathbb Q}^n\) 的基。
\end{theorem}
\begin{proof}
将\eqref{b5:eq:schur-cauchy}乘以 \(\Delta_N(y)\)，并取 \(y_1^{\lambda_1+N-1}\cdots y_N^{\lambda_N}\) 的系数。右端中 \(a_{\mu+\delta_N}(y)\) 只有在 \(\mu=\lambda\) 时贡献一，因为其指数列严格递减；所得系数是 \(s_\lambda(x)\)。左端可写成
\[
\det(y_i^{N-j})_{i,j=1}^N\prod_{i=1}^N
\left(\sum_{r\ge0}h_r(x)y_i^r\right).
\]
展开行列式后，该系数恰为 \(\sum_\sigma\operatorname{sgn}(\sigma)\prod_i h_{\lambda_i-i+\sigma(i)}(x)\)，即所需行列式。对任意有限变量数均成立，故得到稳定恒等式。

最后，在 \(N\ge n\) 个变量中，乘以 \(\Delta_N\) 将 \(n\) 次对称多项式一一对应到次数 \(n+N(N-1)/2\) 的交错多项式。后者按单项式的变量置换轨道分组，每个非零轨道恰由一个严格递减指数列代表，因此以 \(a_{\lambda+\delta_N}\) 为基。除以 \(\Delta_N\)，得到 \(s_\lambda\) 的基性质。
\end{proof}

例如
\[
s_{(r)}=h_r,\qquad
s_{(2,1)}=h_2h_1-h_3,\qquad
s_{(2,2)}=h_2^2-h_1h_3.
\]
又 \(s_{(1^r)}=e_r\)：从\eqref{b5:eq:jacobi-trudi}沿首行展开，所得行列式 \(D_r\) 满足
\[
D_r=h_1D_{r-1}-h_2D_{r-2}+\cdots+(-1)^{r-1}h_r,\qquad D_0=1.
\]
这正是 \(E(-t)H(t)=1\) 对 \(e_r\) 给出的递推及初值。所以单行和单列分别给出前面两种对称函数。

\subsection{次数分次与内积}
\label{b5:subsec:0701:03}

\begin{definition}\label{b5:def:hall-inner-product}
设 \(\Lambda_{\mathbb Q}\) 为有理系数对称函数代数，\(p_\lambda\) 为幂和基，\(m_r(\lambda)\) 为分拆 \(\lambda\) 中部分 \(r\) 的重数，\(z_\lambda=\prod_{r\ge1}r^{m_r(\lambda)}\cdot m_r(\lambda)!\)。在 \(\Lambda_{\mathbb Q}\) 上，以
\[
\langle p_\lambda,p_\mu\rangle=\delta_{\lambda\mu}z_\lambda
\]
确定一个对称双线性配对，称为\term[Hall-peidui]{Hall 配对}[Hall pairing]。在 \(\Lambda_{\mathbb C}\) 上，沿同一正交基改用第一变量共轭线性、第二变量线性的正定内积，称为相应的 Hall 内积。不同次数的齐次分量相互正交。
\end{definition}

本章涉及复特征标时始终使用后一种内积；有理系数上的计算在两种约定下相同。比如 \(\langle h_2,h_2\rangle=1\)，而 \(\langle p_1^2,p_1^2\rangle=2\)。幂和基便于读取循环型，Schur 基则会对应不可约表示。

\begin{proposition}\label{b5:prop:schur-orthonormal}
在有理系数对称函数代数 \(\Lambda_{\mathbb Q}\) 上，对任意分拆 \(\lambda,\mu\) 对应的 Schur 函数 \(s_\lambda,s_\mu\)，Hall 配对满足
\[
\langle s_\lambda,s_\mu\rangle=\delta_{\lambda\mu}.
\]
复化后，Schur 函数同样构成各齐次分量的一组标准正交基。
\end{proposition}
\begin{proof}
固定次数 \(n\)。以 \(p_\lambda\) 为基时，配对矩阵为对角矩阵 \(D=\operatorname{diag}(z_\lambda)\)；因而张量 \(\sum_\lambda p_\lambda\otimes p_\lambda/z_\lambda\) 的系数矩阵是逆矩阵 \(D^{-1}\)。一般地，若一个非退化配对在基 \(v_i\) 下的矩阵为 \(G\)，同一张量就写成 \(\sum_{i,j}(G^{-1})_{ij}v_i\otimes v_j\)：这是换基矩阵两次相乘的直接结果。

\cref{b5:lem:schur-cauchy}证明这个张量在 Schur 基下为 \(\sum_\lambda s_\lambda\otimes s_\lambda\)。其系数矩阵是单位矩阵，所以配对矩阵也是单位矩阵。所有换基系数为有理数，复化并在第一变量取共轭后结论不变。
\end{proof}

% !TeX root = ../../Book5.tex
% 纲要标题：### 7.2 Frobenius 特征映射
% 纲要位置：工作记录/计划纲要.md，第 3959 行。
% * 对称群表示环与对称函数
% * 诱导乘积与对称函数乘法
% * 特征标计算的组合化
\section{Frobenius 特征映射}
\label{b5:sec:0702}

幂和函数适合记录轮换长度：一个 \(r\) 轮换对应一个 \(p_r\)，互不相交的轮换对应乘积。然而同一循环型的共轭类大小不同，直接把特征标值作为系数会破坏内积。中心化子阶 \(z_\mu\) 正好给出所需的归一化。

\subsection{对称群表示环与对称函数}
\label{b5:subsec:0702:01}

\begin{definition}\label{b5:def:frobenius-characteristic}
设 \(n\) 为非负整数，\(f\) 为 \(S_n\) 上的复类函数，\(f(\mu)\) 表示它在循环型 \(\mu\vdash n\) 上的值，\(z_\mu\) 为该循环型的中心化子阶，\(p_\mu\in\Lambda_{\mathbb C}^n\) 为幂和对称函数的分拆乘积。定义
\[
\operatorname{ch}_n(f)
=\sum_{\mu\vdash n}\frac{f(\mu)}{z_\mu}p_\mu.
\]
从复类函数空间 \(\operatorname{CF}(S_n)\) 到复对称函数的 \(n\) 次齐次分量 \(\Lambda_{\mathbb C}^n\) 的这个线性映射称为\term[Frobenius-tezheng-yingshe]{Frobenius 特征映射}[Frobenius characteristic map]。对有限维复 \(S_n\)-表示 \(U\)，也记 \(\operatorname{ch}(U)=\operatorname{ch}_n(\chi_U)\)。
\end{definition}

这里的 \(z_\mu\) 与\cref{b5:lem:schur-cauchy}一致。事实上，与一个循环型为 \(\mu\) 的置换交换的置换，可以交换等长轮换，再在每个轮换内部独立转动。因此其中心化子阶为 \(\prod_r r^{m_r(\mu)}\cdot m_r(\mu)!\)，共轭类大小为 \(n!/z_\mu\)。

\begin{proposition}\label{b5:prop:frobenius-isometry}
对任意非负整数 \(n\)，Frobenius 特征映射 \(\operatorname{ch}_n:\operatorname{CF}(S_n)\to\Lambda_{\mathbb C}^n\) 是复向量空间同构。对任意复类函数 \(f,g\in\operatorname{CF}(S_n)\)，以第一变量共轭线性的 Hall 内积配对它们的像，有
\[
\bigl\langle\operatorname{ch}_n(f),\operatorname{ch}_n(g)\bigr\rangle
=\frac1{n!}\sum_{\sigma\in S_n}\overline{f(\sigma)}g(\sigma)
=\sum_{\mu\vdash n}\frac{\overline{f(\mu)}g(\mu)}{z_\mu}.
\]
这里 \(f(\mu),g(\mu)\) 为循环型 \(\mu\) 上的值，\(z_\mu\) 为相应中心化子阶。并且平凡表示与符号表示的像分别为完全对称函数 \(h_n\)、初等对称函数 \(e_n\)。
\end{proposition}
\begin{proof}
各循环型的示性函数构成类函数基，像是非零倍数的 \(p_\mu\)，所以映射为同构。Hall 内积对幂和基的对角公式立即给出所写等式；共轭类大小解释了两种求和形式。

平凡表示的像为 \(\sum_{\mu\vdash n}p_\mu/z_\mu\)。把所有 \(n\) 汇总，其生成函数为
\[
\prod_{r\ge1}\left(\sum_{a\ge0}\frac{p_r^a t^{ra}}{r^a a!}\right)
=\exp\left(\sum_{r\ge1}\frac{p_r}{r}t^r\right)=H(t).
\]
符号在循环型 \(\mu\) 上的值为 \(\prod_r(-1)^{(r-1)m_r(\mu)}\)。相应地把上式中 \(p_r\) 换成 \((-1)^{r-1}p_r\)，得到 \(E(t)\)，所以次数 \(n\) 项为 \(e_n\)。
\end{proof}

普通表示环 \(R_{\mathbb C}(S_n)\) 的乘法来自同一个群上的张量积。本节还需要另一种乘法，它会改变对称群的次数；这两种运算不能混用。

\subsection{诱导乘积与对称函数乘法}
\label{b5:subsec:0702:02}

\begin{definition}\label{b5:def:induction-product}
设 \(R_{\mathbb C}(S_n)\) 为对称群 \(S_n\) 的复表示环，\(\mathcal R=\bigoplus_{n\ge0}R_{\mathbb C}(S_n)\)。给定非负整数 \(a,b\)、有限维复 \(S_a\)-表示 \(U\) 和 \(S_b\)-表示 \(W\)，记 \([U],[W]\) 为其在表示环中的类，\(U\boxtimes W\) 为 \(S_a\times S_b\) 在 \(U\otimes_{\mathbb C}W\) 上按各因子分别作用的表示。把 \(S_a\times S_b\) 按前 \(a\) 个与后 \(b\) 个字母的置换嵌入 \(S_{a+b}\)，定义
\[
[U]*[W]
=\bigl[\operatorname{Ind}_{S_a\times S_b}^{S_{a+b}}(U\boxtimes W)\bigr],
\]
并对虚表示双线性延拓。这称为\term[youdao-chengji]{诱导乘积}[induction product]。次数零的平凡表示是单位。
\end{definition}

\begin{proposition}\label{b5:prop:frobenius-induction-product}
设 \(a,b\) 为非负整数，\(f,g\) 分别为 \(S_a,S_b\) 的复类函数，\((f\boxtimes g)(\sigma,\tau)=f(\sigma)g(\tau)\)，\(S_a\times S_b\le S_{a+b}\) 按两个连续字母块嵌入。对 Frobenius 特征映射，有
\[
\operatorname{ch}_{a+b}
\bigl(\operatorname{Ind}_{S_a\times S_b}^{S_{a+b}}(f\boxtimes g)\bigr)
=\operatorname{ch}_a(f)\operatorname{ch}_b(g).
\]
因而诱导乘积使 \(\mathcal R=\bigoplus_{n\ge0}R_{\mathbb C}(S_n)\) 成为分次交换环，Frobenius 特征映射保持该乘法。
\end{proposition}
\begin{proof}
固定循环型 \(\nu\vdash a+b\) 的置换 \(\sigma\)。陪集空间 \(S_{a+b}/(S_a\times S_b)\) 可以看作全部 \(a\) 元子集。诱导特征标的求和只有在 \(\sigma\) 保持该子集时才有贡献；这样的子集必为若干完整轮换之并。若它及其补集的循环型分别为 \(\alpha,\beta\)，则
\[
m_r(\alpha)+m_r(\beta)=m_r(\nu),\qquad
|\alpha|=a,\quad|\beta|=b,
\]
而选择这些轮换的方法数为
\[
\prod_r\binom{m_r(\nu)}{m_r(\alpha)}
=\frac{z_\nu}{z_\alpha z_\beta}.
\]
每种选择贡献 \(f(\alpha)g(\beta)\)。所以诱导类函数在 \(\nu\) 处的值除以 \(z_\nu\)，等于
\[
\sum_{\alpha\cup\beta=\nu}
\frac{f(\alpha)g(\beta)}{z_\alpha z_\beta}.
\]
这正是右端幂和乘积中 \(p_\nu\) 的系数，证明了等式。类函数上的特征映射单射，而右端乘法结合、交换且有次数零的单位，因此同样性质在虚表示上成立。
\end{proof}

设 \(\mu=(\mu_1,\ldots,\mu_r)\vdash n\)，Young 置换模为
\[
M^\mu=\operatorname{Ind}_{S_{\mu_1}\times\cdots\times S_{\mu_r}}^{S_n}1.
\]
由刚证明的乘法公式及 \(\operatorname{ch}(1_{S_a})=h_a\)，有
\begin{equation}\label{b5:eq:young-permutation-characteristic}
\operatorname{ch}(M^\mu)=h_\mu.
\end{equation}
同样的记法可用于任意非负整数组合：重排各部分只改变块的排列，诱导表示同构；零部分可以删去。

\begin{theorem}\label{b5:thm:frobenius-characteristic}
设 \(n\) 为非负整数，\(\lambda\vdash n\)，\(S^\lambda\) 为对应的复 Specht 左 \(\mathbb C[S_n]\)-模，\(s_\lambda\) 为 Schur 函数，\(\operatorname{ch}_n\) 为 Frobenius 特征映射，记 \(\operatorname{ch}(S^\lambda)=\operatorname{ch}_n(\chi_{S^\lambda})\)。则
\begin{equation}\label{b5:eq:specht-schur}
\operatorname{ch}(S^\lambda)=s_\lambda.
\end{equation}
因此直和映射 \(\operatorname{ch}=\bigoplus_{n\ge0}\operatorname{ch}_n\) 将带诱导乘积的 \(\mathcal R=\bigoplus_{n\ge0}R_{\mathbb C}(S_n)\) 同构到 \(\Lambda_{\mathbb Q}\) 中以 \(s_\lambda\) 为整数基的子环，并将不可约表示的类对应到 Schur 基。
\end{theorem}
\begin{proof}
取 \(N\ge n\)，将 \(\lambda\) 补零到 \(N\) 项。由 Jacobi--Trudi 公式和\eqref{b5:eq:young-permutation-characteristic}，\(s_\lambda\) 是下列虚特征标的像：
\[
\theta_\lambda
=\sum_{\sigma\in S_N}\operatorname{sgn}(\sigma)\,
\chi_{M^{\alpha(\sigma)}},
\qquad
\alpha_i(\sigma)=\lambda_i-i+\sigma(i).
\]
含负部分的项规定为零；其他项将 \(\alpha(\sigma)\) 重排成分拆 \(\mu(\sigma)\)。

对任意 \(k\)，有
\[
\sum_{i=1}^k\alpha_i(\sigma)
=\sum_{i=1}^k\lambda_i+
\sum_{i=1}^k\sigma(i)-\frac{k(k+1)}2
\ge\sum_{i=1}^k\lambda_i.
\]
重排后的前 \(k\) 项和不会减小，故 \(\mu(\sigma)\) 支配 \(\lambda\)。若 \(\sigma\ne1\)，存在一个 \(k\) 使 \(\bigl\{\sigma(1),\ldots,\sigma(k)\bigr\}\ne\{1,\ldots,k\}\)，此时上式严格，故 \(\mu(\sigma)\ne\lambda\)。

现在应用\cref{b5:prop:specht-dominance}：\(M^\mu\) 只含支配 \(\mu\) 的 Specht 类型，并且 \(S^\mu\) 的重数是一。因此
\[
\theta_\lambda=\chi_{S^\lambda}
+\sum_{\eta\triangleright\lambda}c_\eta\chi_{S^\eta},
\qquad c_\eta\in\mathbb Z,
\]
其中 \(\triangleright\) 表示严格支配。另一方面，特征映射等距且 Schur 函数范数为一，故
\[
1=\langle s_\lambda,s_\lambda\rangle
=\langle\theta_\lambda,\theta_\lambda\rangle
=1+\sum_{\eta\triangleright\lambda}c_\eta^2.
\]
最后一式使用的是已经在\cref{b5:thm:specht-classification}建立的不可约分类与特征标正交性。于是全部 \(c_\eta=0\)，得到 \(\theta_\lambda=\chi_{S^\lambda}\)。证明没有用 Schur 函数反过来建立 Specht 分类。整数基及乘法结论随后由诱导乘积公式得到。
\end{proof}

\subsection{特征标计算的组合化}
\label{b5:subsec:0702:03}

\begin{corollary}[Frobenius 特征标公式]\label{b5:cor:frobenius-coefficient}
设 \(n\) 为非负整数，\(\lambda,\mu\vdash n\)，\(S^\lambda\) 为复 Specht 模，\(\chi_{S^\lambda}(\mu)\) 为其特征标在循环型 \(\mu\) 上的值，\(z_\mu\) 为相应中心化子阶。取整数 \(N\ge n\)，把 \(\lambda\) 补零至 \(N\) 项，令 \(x=(x_1,\ldots,x_N)\)、\(\delta_N=(N-1,\ldots,0)\)、\(\Delta_N(x)=\prod_{i<j}(x_i-x_j)\)。对这些变量中的 Schur 多项式和幂和多项式，有
\[
p_\mu(x)=\sum_{\lambda\vdash n}\chi_{S^\lambda}(\mu)s_\lambda(x),
\qquad
\chi_{S^\lambda}(\mu)
=[x^{\lambda+\delta_N}]\Delta_N(x)p_\mu(x).
\]
等价地，\(s_\lambda\) 中 \(p_\mu\) 的系数乘以 \(z_\mu\)，就是 \(\chi_{S^\lambda}(\mu)\)。
\end{corollary}
\begin{proof}
由\eqref{b5:eq:specht-schur}，\(\langle p_\mu,s_\lambda\rangle=\chi_{S^\lambda}(\mu)\)。这些特征标值为实数：一个置换与其逆具有相同循环型，而有限群特征标满足 \(\chi(g^{-1})=\overline{\chi(g)}\)。所以交换内积的两个参数不会改变该值。在 Schur 标准正交基下展开 \(p_\mu\)，得第一式。乘以 \(\Delta_N\)，不同严格递减指数列的交错式互不混淆，读取 \(x^{\lambda+\delta_N}\) 的系数便得第二式。
\end{proof}

例如
\[
h_2=\frac{p_1^2+p_2}{2},\qquad
h_3=\frac{p_1^3+3p_1p_2+2p_3}{6},
\qquad
s_{(2,1)}=\frac{p_1^3-p_3}{3}.
\]
由于 \(z_{(1^3)}=6\)、\(z_{(2,1)}=2\)、\(z_{(3)}=3\)，立即读出 \(S_3\) 的二维特征标为 \((2,0,-1)\)。

\ProseParagraphBreak
对四次情形，
\[
s_{(2,2)}=h_2^2-h_1h_3
=\frac{p_1^4}{12}+\frac{p_2^2}{4}-\frac{p_1p_3}{3}.
\]
按循环型 \((1^4),(2,1,1),(2,2),(3,1),(4)\) 排列，特征标值为 \((2,0,2,-1,0)\)。同一个行列式因而同时编码维数、换位上的迹以及较长轮换上的迹。

\ProseParagraphBreak
若实际计算 \(S^\lambda\) 在指定循环型上的特征标，可以先用 Jacobi--Trudi 求 \(s_\lambda\)，再用\eqref{b5:eq:newton-h}逐个把所需 \(h_r\) 化成幂和，最后读取一个系数。不需要预先求整张特征标表。Etingof 等人的叙述在\cite[Sections 5.14--5.15, pp. 114--118]{EtingofEtAl2011RepresentationTheory}中以 Young 置换模、三角性及 Cauchy 行列式建立相应公式；这里已把等距性与诱导乘法一并写出，以便继续处理张量分解。

% !TeX root = ../../Book5.tex
% 纲要标题：### 7.3 Schur–Weyl 对偶
% 纲要位置：工作记录/计划纲要.md，第 3965 行。
% * 张量幂上对称群与一般线性群的交换作用
% * 特征零下的双中心化子定理
% * 分解及维数条件
% * 采用特征零下的半单性、极化及中心化子计算建立双中心化子结论，不前置第18章最高权分类或第19章 Weyl 特征标公式
% * 在张量次数大于向量空间维数时明确可出现的分拆行数限制；交换作用不自动意味着双方表示均忠实
\section{Schur–Weyl 对偶}
\label{b5:sec:0703}

张量幂中有两种不同的对称性：可以同时改变每个因子中的向量，也可以重新排列因子。二者交换很容易验证；要证明它们互为中心化子，还须证明全部与一种作用交换的线性算子都能由另一种作用线性生成。本节用半单代数的矩阵块和极化证明这一结论，无须调用一般线性群表示的最高权分类。

\subsection{张量幂上的对称群与一般线性群作用}
\label{b5:subsec:0703:01}

设 \(k\) 为特征零域，\(V\) 为 \(d\ge1\) 维 \(k\) 向量空间，\(n\ge1\)，记 \(T=V^{\otimes n}\)。定义两个左作用：
\[
P_\sigma(v_1\otimes\cdots\otimes v_n)
=v_{\sigma^{-1}(1)}\otimes\cdots\otimes v_{\sigma^{-1}(n)},
\qquad
R_g=g^{\otimes n},
\]
其中 \(\sigma\in S_n\)、\(g\in\operatorname{GL}(V)\)。逆号保证 \(P_\sigma P_\tau=P_{\sigma\tau}\)；若省掉它，就会写成反向乘法。逐个因子计算给出 \(P_\sigma R_g=R_gP_\sigma\)。

\ProseParagraphBreak
记两种作用在线性自同态代数中张成的代数为
\[
A=\operatorname{im}\bigl(k[S_n]\longrightarrow\operatorname{End}_k(T)\bigr),
\qquad
B=\operatorname{span}_k\bigl\{\,g^{\otimes n}\mid g\in\operatorname{GL}(V)\,\bigr\}.
\]
\(B\) 确为含单位的子代数，因为 \(g^{\otimes n}h^{\otimes n}=(gh)^{\otimes n}\)。这里讨论的是像代数；两个原始作用有核，并不妨碍它们的像成为彼此的中心化子。

\subsection{特征零下的双中心化子定理}
\label{b5:subsec:0703:02}

\begin{lemma}[分裂半单代数的双中心化子]\label{b5:lem:split-double-centralizer}
设 \(k\) 为域，\(E\) 为有限维非零 \(k\) 向量空间，\(A\subseteq\operatorname{End}_k(E)\) 为含单位的分裂半单子代数，即 \(A\simeq\prod_iM_{r_i}(k)\)。令 \(C=\operatorname{End}_A(E)\)。则 \(C\) 也是分裂半单代数，且
\[
\operatorname{End}_C(E)=A.
\]
更具体地，若 \(S_i\) 遍历 \(A\) 的简单左模的同构类代表，并令 \(L_i=\operatorname{Hom}_A(S_i,E)\)，以 \(c\cdot f=c\circ f\) 定义左 \(C\) 作用，则
\[
E\simeq\bigoplus_i S_i\otimes_k L_i,\qquad
A\simeq\prod_i\operatorname{End}_k(S_i),\qquad
C\simeq\prod_i\operatorname{End}_k(L_i).
\]
所有 \(L_i\) 均非零，并且是 \(C\) 的全部互不同构简单模。
\end{lemma}
\begin{proof}
分裂半单性给出 \(E\) 的简单直和分解。\(A\) 在 \(E\) 上忠实，因为 \(A\) 已是自同态代数的子代数；若某个简单类型未出现，其矩阵块的单位元就会零化 \(E\)，与忠实性矛盾。因此各 \(L_i\ne0\)。评价映射 \(s\otimes f\mapsto f(s)\) 给出第一式。

与 \(A\) 交换的算子须保持其中心幂等元划出的各块。块内与所有矩阵单位 \(E_{ab}\otimes1\) 交换的算子恰为 \(1\otimes U\)，其中 \(U\in\operatorname{End}_k(L_i)\)：先与 \(E_{aa}\) 交换排除非对角块，再与 \(E_{ab}\) 交换使所有对角块相同。因此得到 \(C\) 的乘积描述。反过来，重复这个矩阵单位计算，与全部 \(1\otimes\operatorname{End}_k(L_i)\) 交换的算子正是 \(\operatorname{End}_k(S_i)\otimes1\)，即 \(A\)。

最后，每个 \(L_i\) 是第 \(i\) 个全矩阵代数的列向量简单模，不同块由中心幂等元区分。这些也是乘积矩阵代数的全部简单模。
\end{proof}

本引理说明，只要算出一侧的中心化子，另一侧的等式就可由半单性推出。但“两个作用交换”本身只给包含关系，不足以调用这个结论。

\subsection{分解及维数条件}
\label{b5:subsec:0703:03}

由\cref{b5:thm:specht-classification}，\(k[S_n]\) 在特征零时是分裂半单代数，简单模为 \(S^\lambda_k\)。因此先不判断一般线性群一侧的不可约性，也已有自然分解
\begin{equation}\label{b5:eq:tensor-isotypic}
T\simeq\bigoplus_{\lambda\vdash n}S^\lambda_k\otimes_k L_\lambda(V),
\qquad
L_\lambda(V)=\operatorname{Hom}_{k[S_n]}(S^\lambda_k,T).
\end{equation}
一般线性群在重数空间上作用为 \(g\cdot f=g^{\otimes n}\circ f\)。这个定义良定，正是因为 \(g^{\otimes n}\) 与 \(S_n\) 作用交换；评价映射与两个作用均相容。

\ProseParagraphBreak
例如 \(n=2\) 时，两个 Specht 模分别是平凡模与符号模。投影 \(\bigl(1+P_{(12)}\bigr)/2\)、\(\bigl(1-P_{(12)}\bigr)/2\) 把 \(V^{\otimes2}\) 分成对称和反对称张量。一般分拆对应的重数空间可以看作这两种构造的推广；接下来将证明它们除了为零的情形以外都是不可约表示。

\subsection{半单性、极化与中心化子计算}
\label{b5:subsec:0703:04}

\begin{lemma}[极化]\label{b5:lem:tensor-polarization}
设 \(k\) 为特征零域，\(U\) 为有限维 \(k\) 向量空间，\(n\ge1\)。置换张量因子作用下的不变子空间 \((U^{\otimes n})^{S_n}\)，由纯幂 \(u^{\otimes n}\)（\(u\in U\)）线性生成。
\end{lemma}
\begin{proof}
平均投影把纯张量的生成集送到对称化张量的生成集。对 \(u_1,\ldots,u_n\in U\)，有极化恒等式
\[
\sum_{J\subseteq\{1,\ldots,n\}}(-1)^{n-|J|}
\left(\sum_{j\in J}u_j\right)^{\otimes n}
=\sum_{\sigma\in S_n}u_{\sigma(1)}\otimes\cdots\otimes u_{\sigma(n)}.
\]
展开左端，一项若只使用真子集 \(I\) 中的指标，其系数为 \(\sum_{J\supseteq I}(-1)^{n-|J|}=0\)。若使用全部 \(n\) 个指标，在 \(n\) 个位置中每个指标必须恰出现一次，系数为一。这证明恒等式。除以 \(n!\)，每个平均纯张量都成为纯幂的线性组合；特征零保证可以作此除法。
\end{proof}

\begin{theorem}[Schur--Weyl 对偶]\label{b5:thm:schur-weyl}
设 \(k\) 为特征零域，\(V\) 为有限维非零 \(k\) 向量空间，\(n\ge1\) 为整数。在 \(T=V^{\otimes n}\) 中，\(S_n\) 按置换张量因子作左作用。令 \(A\) 为这个 \(k[S_n]\) 作用的像，\(B\) 为全部 \(g^{\otimes n}\)（\(g\in\operatorname{GL}_k(V)\)）的 \(k\)-线性生成空间，\(S_k^\lambda\) 为分拆 \(\lambda\vdash n\) 的 Specht 模，\(L_\lambda(V)=\operatorname{Hom}_{k[S_n]}(S_k^\lambda,T)\)，其一般线性群作用为 \(g\cdot f=g^{\otimes n}\circ f\)。则
\[
\operatorname{End}_A(T)=B,\qquad
\operatorname{End}_B(T)=A.
\]
两者均为分裂半单代数。在自然分解 \(T\cong\bigoplus_{\lambda\vdash n}S_k^\lambda\otimes_k L_\lambda(V)\) 中，非零 \(L_\lambda(V)\) 是互不同构的不可约有限维 \(\operatorname{GL}_k(V)\) 表示。
\end{theorem}
\begin{proof}
令 \(U=\operatorname{End}_k(V)\)。自然映射
\[
U^{\otimes n}\longrightarrow\operatorname{End}_k(V^{\otimes n}),
\quad
a_1\otimes\cdots\otimes a_n\longmapsto
\left(v_1\otimes\cdots\otimes v_n\mapsto
a_1v_1\otimes\cdots\otimes a_nv_n\right)
\]
是向量空间同构：选择 \(V\) 的基，两侧相应的张量矩阵单位一一对应。共轭 \(F\mapsto P_\sigma F P_\sigma^{-1}\) 在左侧对应置换 \(U\) 的张量因子。所以 \(\operatorname{End}_A(T)\) 对应 \((U^{\otimes n})^{S_n}\)。

由\cref{b5:lem:tensor-polarization}，后者由 \(a^{\otimes n}\)、\(a\in U\) 生成。还需把任意 \(a\) 换成可逆算子。多项式 \(\det(a+tI)\) 不是零多项式，所以只有有限多个 \(t\in k\) 使 \(a+tI\) 不可逆。域 \(k\) 无限，可取 \(n+1\) 个互异的可用值 \(t_0,\ldots,t_n\)。向量值多项式 \((a+tI)^{\otimes n}\) 次数至多 \(n\)，由 Lagrange 插值，其在 \(t=0\) 的值 \(a^{\otimes n}\) 是这 \(n+1\) 个可逆算子张量幂的线性组合。因此 \(\operatorname{End}_A(T)\subseteq B\)。反向包含已由两个作用交换得到，所以相等。

\(A\) 是分裂半单群代数 \(k[S_n]\) 的商，仍为若干全矩阵块的乘积。现在对 \(E=T\) 使用\cref{b5:lem:split-double-centralizer}，得到另一中心化子等式及 \(B\) 的分裂半单性。该引理还说明各非零重数空间是 \(B\) 的互不同构简单模。一个子空间对所有 \(g^{\otimes n}\) 不变，当且仅当对它们的线性生成代数 \(B\) 不变；群交织映射与 \(B\) 交织映射也相同。因此这恰是一般线性群表示的不可约性及非同构性。
\end{proof}

这一证明只用 Specht 分类、矩阵单位和极化，全部计算在有限群表示及向量空间的线性映射中完成。有关双中心化子及可逆矩阵张量幂的论证，也可参阅\cite[Theorem 5.18.1; Proposition 5.19.1 and Corollary 5.19.2, pp. 120--122]{EtingofEtAl2011RepresentationTheory}。

\subsection{分拆行数限制与作用的忠实性}
\label{b5:subsec:0703:05}

\begin{proposition}\label{b5:prop:schur-functor-character}
设 \(k\) 为特征零域，\(V\) 为 \(d\ge1\) 维 \(k\) 向量空间，\(n\ge1\) 为整数，\(\lambda\vdash n\)，\(S_k^\lambda\) 为 Specht 模，\(L_\lambda(V)=\operatorname{Hom}_{k[S_n]}(S_k^\lambda,V^{\otimes n})\)，其中 \(S_n\) 置换张量因子。一般线性群在该 Hom 空间上以 \(g\cdot f=g^{\otimes n}\circ f\) 作用。对 \(g\in\operatorname{GL}_k(V)\)，设 \(x_1,\ldots,x_d\) 为其在 \(k\) 的代数闭包中的特征值，按重数列出，\(s_\lambda\) 为 Schur 多项式。则
\[
\operatorname{tr}\bigl(g\mid L_\lambda(V)\bigr)=s_\lambda(x_1,\ldots,x_d).
\]
空间 \(L_\lambda(V)\) 非零当且仅当 \(\ell(\lambda)\le d\)。在此条件下，把 \(\lambda\) 补零到 \(d\) 项，有
\begin{equation}\label{b5:eq:schur-dimension}
\dim_k L_\lambda(V)
=s_\lambda(1,\ldots,1)
=\prod_{1\le i<j\le d}\frac{\lambda_i-\lambda_j+j-i}{j-i}.
\end{equation}
\end{proposition}
\begin{proof}
设 \(\sigma\in S_n\) 的循环型为 \(\mu\)。按 \(V\) 的一组基计算迹，每个轮换使相应矩阵指标首尾连接，一个 \(r\) 轮换贡献 \(\operatorname{tr}(g^r)\)。因此
\[
\operatorname{tr}(g^{\otimes n}P_\sigma)
=\prod_r\operatorname{tr}(g^r)^{m_r(\mu)}
=p_\mu(x_1,\ldots,x_d).
\]
另一方面，对\eqref{b5:eq:tensor-isotypic}取迹，得
\[
p_\mu(x)=\sum_{\lambda\vdash n}
\chi_{S^\lambda}(\mu)\operatorname{tr}\bigl(g\mid L_\lambda(V)\bigr).
\]
\cref{b5:cor:frobenius-coefficient}给出同样的展开，其系数为 \(s_\lambda(x)\)。\(S_n\) 的特征标表可逆，故比较得到所需迹公式。该比较在任意特征零域都合法：Specht 模来自有理系数模型，特征标表及其正交逆矩阵具有有理系数；所用对称函数恒等式也在 \(\mathbb Q\) 上成立。故无需把一般的 \(k\) 嵌入复数域。

令 \(g=I\)，得到维数等于 \(s_\lambda(1^d)\)。若 \(\ell(\lambda)>d\)，有限变量 Schur 多项式为零，所以该空间维数为零。若 \(\ell(\lambda)\le d\)，在交错式商中先代入 \(x_i=q^{i-1}\)。对指数 \(\lambda_j+d-j\) 使用 Vandermonde 公式，再约去分母，得到
\[
s_\lambda(1,q,\ldots,q^{d-1})
=q^{\sum_i(i-1)\lambda_i}
\prod_{i<j}\frac{1-q^{\lambda_i-\lambda_j+j-i}}{1-q^{j-i}}.
\]
每个商在 \(q=1\) 的值是相应指数之比，得到\eqref{b5:eq:schur-dimension}。所有分子、分母都是正整数，因此该有理数严格为正；它又是已经算出的维数，故 \(L_\lambda(V)\ne0\)。
\end{proof}

例如 \(\dim V=2\) 时，三次张量幂只出现 \((3)\) 与 \((2,1)\)，不出现三行的 \((1,1,1)\)。其重数空间维数分别为四与二，Specht 维数分别为一与二，所以
\[
2^3=1\cdot4+2\cdot2.
\]
这是两边维数的乘积之和，不能把重数空间的维数直接当成 Specht 模的维数。

\begin{proposition}\label{b5:prop:tensor-faithfulness}
设 \(k\) 为特征零域，\(V\) 为 \(d\ge1\) 维 \(k\) 向量空间，\(n\ge1\) 为整数。\(S_n\) 在 \(V^{\otimes n}\) 上置换张量因子，\(\operatorname{GL}_k(V)\) 以 \(g\mapsto g^{\otimes n}\) 作用。则 \(k[S_n]\) 的作用忠实，当且仅当 \(d\ge n\)。当 \(n\ge2\) 时，群 \(S_n\) 的作用忠实，当且仅当 \(d\ge2\)。一般线性群作用的核为
\[
\{\,cI_V\mid c\in k^\times,\ c^n=1\,\}.
\]
\end{proposition}
\begin{proof}
群代数的核是未出现的 Specht 类型所对应矩阵块之和。由\cref{b5:prop:schur-functor-character}，所有分拆都出现当且仅当 \(d\ge n\)；当 \(d<n\) 时至少遗漏 \((1^n)\)，所以得到第一条。

若 \(d=1\)，所有因子置换都作用为恒等。若 \(d\ge2\)，选独立向量 \(u,v\)。对非恒等 \(\sigma\)，取被它移动的位置 \(i\)，构造只在第 \(i\) 位放 \(v\)、其他位置放 \(u\) 的基张量。置换后 \(v\) 的位置改变，所得张量不同，故群作用忠实。

最后，若 \(g^{\otimes n}=I\)，则 \((gv)^{\otimes n}=v^{\otimes n}\) 对每个 \(v\ne0\) 成立，因而 \(gv\) 与 \(v\) 共线：若不共线，可取线性泛函在 \(v\) 上为零、在 \(gv\) 上非零，对张量的各因子应用该泛函便产生矛盾。一个线性算子若保持每条直线，必为标量；在基向量及两两之和上比较即可证明各标量相同。于是 \(g=cI\)，再代回得到 \(c^n=1\)。反向显然。
\end{proof}

特别地，在 \(d=2<n\) 时，对称群本身仍可忠实作用，而群代数作用已经不忠实。群元素之间互不相同，并不意味着这些元素的作用矩阵线性无关。

% !TeX root = ../../Book5.tex
% 纲要标题：### 7.4 张量分解
% 纲要位置：工作记录/计划纲要.md，第 3973 行。
% * Schur 函子
% * Pieri 规则及低阶 Littlewood–Richardson 例子
% * Littlewood–Richardson 诱导分解先在有理数上确定重数，再扩张到任意特征零域；一般线性群分解用迹多项式与 Schur 基的线性无关性比较
% * 与第二卷对称幂、外幂的统一
% ---
\section{张量分解}
\label{b5:sec:0704}

Schur--Weyl 对偶给每个分拆配置了一个重数空间。改变向量空间或施加线性映射时，这些空间也会相容地改变，因而形成函子。对称函数的乘法随后给出其张量积规则；需要注意，这种乘法对应对称群一侧的诱导，不能误认为同一对称群上的内部张量积。

\subsection{Schur 函子}
\label{b5:subsec:0704:01}

\begin{definition}\label{b5:def:schur-functor}
设 \(k\) 为特征零域，\(n\) 为非负整数，\(\lambda\vdash n\)，固定一个 Specht 左 \(k[S_n]\)-模 \(S^\lambda_k\)。对有限维 \(k\) 向量空间 \(V\)，在 \(V^{\otimes n}\) 上采用置换张量因子的左 \(S_n\) 作用，令
\[
\mathbb S_\lambda(V)
=\operatorname{Hom}_{k[S_n]}(S^\lambda_k,V^{\otimes n}).
\]
对有限维 \(k\) 向量空间之间的线性映射 \(f:V\to W\)，以及 \(u\in\mathbb S_\lambda(V)\)，令
\[
\mathbb S_\lambda(f)(u)=f^{\otimes n}\circ u.
\]
这称为分拆 \(\lambda\) 对应的\term[Schur-hanzi]{Schur 函子}[Schur functor]。空分拆对应常值空间 \(k\)；正次数时规定 \(\mathbb S_\lambda(0)=0\)。
\end{definition}

因为 \(f^{\otimes n}\) 与因子置换相容，上述复合仍为 \(S_n\) 同态。恒等映射和复合也被保持，所以确实得到函子。选定基时，其作用由 \(f\) 的矩阵条目的 \(n\) 次齐次多项式给出；通常它不是关于 \(f\) 的线性运算。例如 \(\operatorname{Sym}^2(f+g)\) 中有交叉项，不能写成两个平方的和。

\ProseParagraphBreak
这个定义中的空间就是上一节的 \(L_\lambda(V)\)。因此它在非零时是不可约的一般线性群表示，其迹函数由 Schur 多项式 \(s_\lambda\) 给出：把群元素在代数闭包中的特征值代入即可求迹。分拆的行数超过 \(\dim V\) 时，这个空间为零。在不同维数下使用同一分拆，不需要重新定义构造。

\begin{definition}\label{b5:def:lr-coefficient}
设 \(a,b\) 为非负整数，\(\lambda\vdash a\)、\(\mu\vdash b\)。在有理系数对称函数代数 \(\Lambda_{\mathbb Q}\) 中，用 Schur 基展开乘积：
\[
s_\lambda s_\mu=\sum_{\nu\vdash a+b}c_{\lambda\mu}^{\nu}s_\nu.
\]
系数 \(c_{\lambda\mu}^{\nu}\) 称为\term[Littlewood-Richardson-xishu]{Littlewood--Richardson 系数}[Littlewood--Richardson coefficient]。\footnote{利特尔伍德（Dudley Ernest Littlewood，1903-1979）与理查森（Archibald Read Richardson，1881-1954）均为英国数学家，合作研究对称群表示与对称函数的乘法系数。}
\end{definition}

\begin{proposition}\label{b5:prop:schur-tensor-coefficients}
设 \(k\) 为特征零域，\(a,b\) 为非负整数，\(\lambda\vdash a\)、\(\mu\vdash b\)，\(c_{\lambda\mu}^{\nu}\) 为 Schur 函数乘积的 Littlewood--Richardson 系数，\(S^\eta\) 表示 \(k\) 上相应的 Specht 左模，\(\boxtimes\) 表示外张量积，\(S_a\times S_b\le S_{a+b}\) 按两个连续字母块嵌入。则所有 \(c_{\lambda\mu}^{\nu}\) 是非负整数，且
\[
\operatorname{Ind}_{S_a\times S_b}^{S_{a+b}}
(S^\lambda\boxtimes S^\mu)
\simeq\bigoplus_{\nu\vdash a+b}(S^\nu)^{\oplus c_{\lambda\mu}^{\nu}}.
\]
对任意有限维 \(k\) 向量空间 \(V\)，以 \(\mathbb S_\eta\) 记 Schur 函子，还有 \(\operatorname{GL}_k(V)\)-表示的同构
\begin{equation}\label{b5:eq:schur-tensor-decomposition}
\mathbb S_\lambda(V)\otimes\mathbb S_\mu(V)
\simeq
\bigoplus_{\substack{\nu\vdash a+b\\\ell(\nu)\le\dim V}}
\mathbb S_\nu(V)^{\oplus c_{\lambda\mu}^{\nu}}.
\end{equation}
\end{proposition}
\begin{proof}
若 \(a+b=0\)，则 \(\lambda,\mu\) 都是空分拆，唯一的乘积系数为一，两式两边均为 \(k\)。以下设 \(a+b>0\)。先在复数域上，由诱导乘积公式及 Specht--Schur 对应得到第一式；其系数是实际分解重数，所以是非负整数。令
\[
I_{\mathbb Q}=\operatorname{Ind}_{S_a\times S_b}^{S_{a+b}}
\bigl(S_{\mathbb Q}^\lambda\boxtimes S_{\mathbb Q}^\mu\bigr).
\]
由\eqref{b5:eq:symmetric-split-algebra}，\(\mathbb Q[S_{a+b}]\) 是分裂半单代数，故
\[
I_{\mathbb Q}\cong\bigoplus_{\nu\vdash a+b}
(S_{\mathbb Q}^\nu)^{\oplus m_\nu},
\qquad m_\nu\in\mathbb Z_{\geq0}.
\]
整数标准基给出 \(k\otimes_{\mathbb Q}S_{\mathbb Q}^\eta\cong S_k^\eta\)。在群代数张量模型中结合张量积，又得到诱导与扩张标量相容的同构
\[
k\otimes_{\mathbb Q}I_{\mathbb Q}
\cong\operatorname{Ind}_{S_a\times S_b}^{S_{a+b}}
\bigl(S_k^\lambda\boxtimes S_k^\mu\bigr).
\]
先取 \(k=\mathbb C\)。各个 \(S_{\mathbb Q}^\nu\) 复化后仍是互不同构的简单模，比较复数域上的分解即得 \(m_\nu=c_{\lambda\mu}^{\nu}\)。再取任意特征零域 \(k\)，上述同构给出第一式。

对第二式，若 \(V=0\)，两边都是零。以下设 \(V\ne0\)。每个非零 \(\mathbb S_\lambda(V)\) 可选为 \(V^{\otimes a}\) 中一个一般线性群直和项，\(\mathbb S_\mu(V)\) 同样可选为 \(V^{\otimes b}\) 的直和项；次数零的因子为 \(k\)。张量所得是 \(V^{\otimes(a+b)}\) 的直和项，故由 Schur--Weyl 对偶完全可约，全部不可约项都形如右侧。取 \(V\) 的基，对角矩阵上的迹由多项式 \(s_\lambda(x)s_\mu(x)\) 给出。行数不超过 \(\dim V\) 的 Schur 多项式线性无关，而特征零域无限，对所有非零对角元素成立的迹等式就是多项式恒等式。因此各个不可约项的重数只能为 \(c_{\lambda\mu}^{\nu}\)，得到第二式。
\end{proof}

上述论证解释了分解系数的表示意义，一般 Littlewood--Richardson 表格规则尚需另作证明。下面证明最常用的单行、单列情形，再计算一个出现重数二的例子。

\subsection{Pieri 规则与 Littlewood–Richardson 低阶例子}
\label{b5:subsec:0704:02}

\begin{definition}\label{b5:def:horizontal-vertical-strip}
设分拆 \(\mu\) 的 Young 图包含分拆 \(\lambda\) 的 Young 图。差图 \(\mu/\lambda\) 若每列至多有一格，称为\term[hengtiao]{横条}[horizontal strip]；若每行至多有一格，称为\term[shutiao]{竖条}[vertical strip]。其大小是新增格数 \(|\mu|-|\lambda|\)。
\end{definition}

下面的 Pieri 规则\footnote{皮耶里（Mario Pieri，1860-1913），意大利数学家，研究几何基础，其名字也见于对称函数的乘法规则。}用横条、竖条描述 Schur 函数分别乘以单行、单列函数的结果。

\begin{theorem}[Pieri 规则]\label{b5:thm:pieri}
在有理系数对称函数代数 \(\Lambda_{\mathbb Q}\) 中，记 \(s_\lambda\) 为 Schur 函数、\(h_r,e_r\) 分别为完全、初等对称函数。对分拆 \(\lambda\) 和整数 \(r\ge0\)，有
\begin{align}
s_\lambda h_r&=\sum_{\mu/\lambda\text{ 为大小 }r\text{ 的横条}}s_\mu,
\label{b5:eq:pieri-horizontal}\\
s_\lambda e_r&=\sum_{\mu/\lambda\text{ 为大小 }r\text{ 的竖条}}s_\mu.
\label{b5:eq:pieri-vertical}
\end{align}
每项的系数均为一。
\end{theorem}
\begin{proof}
取 \(N\ge|\lambda|+r\)，令 \(\alpha=\lambda+\delta_N\)。先证横条公式。由\cref{b5:thm:jacobi-trudi}证明中按单项式轨道得到的交错多项式基，\(a_\alpha h_r\) 可唯一写成严格递减非负指数列 \(\beta\) 对应的 \(a_\beta\) 的线性组合。\(a_\beta\) 中单项式 \(x^\beta\) 的系数为一，而其他基元素中该系数为零，所以展开系数正是 \(a_\alpha h_r\) 中 \(x^\beta\) 的系数。由于 \(h_r\) 中每个次数 \(r\) 的单项式系数都是一，该系数为
\[
\sum_{\sigma\in S_N}\operatorname{sgn}(\sigma)
\prod_{i=1}^N\mathbf1_{\bigl\{\beta_i\ge\alpha_{\sigma(i)}\bigr\}}
=\det\bigl(\mathbf1_{\{\beta_i\ge\alpha_j\}}\bigr)_{i,j=1}^N,
\]
其中已限定 \(|\beta|=|\alpha|+r\)。

\(\alpha,\beta\) 都严格递减。因此矩阵第 \(i\) 行先为若干个零，随后全为一；其第一个一的位置 \(q_i\) 随 \(i\) 不减。如果有零行或重复行，行列式为零。其余情形的 \(N\) 个不同起始位置只能依次为 \(1,\ldots,N\)，即 \(q_i=i\)，矩阵上三角且对角元全为一。这个条件等价于
\[
\beta_1\ge\alpha_1,\qquad
\alpha_i\le\beta_i<\alpha_{i-1}\quad(2\le i\le N).
\]
写 \(\beta=\mu+\delta_N\)，它化为
\[
\mu_i\ge\lambda_i,\qquad
\mu_i\le\lambda_{i-1}\quad(i\ge2).
\]
恰好表示新增格子不能在同一列重复，即横条条件。除以 \(\Delta_N\)，得到第一式。

再证竖条公式。以 \(\operatorname{Alt}\) 表示对变量置换带符号求和；则 \(a_\alpha=\operatorname{Alt}(x^\alpha)\)。因为 \(e_r\) 对称，
\[
a_\alpha e_r
=\sum_{\substack{J\subseteq\{1,\ldots,N\}\\|J|=r}}
a_{\alpha+\mathbf1_J}.
\]
原指数相邻至少相差一，增加零或一后仍不增。若有相等指数，交错式为零；否则仍严格递减，不需要重新排列，也不产生负号。非零项正好满足 \(\mu_i-\lambda_i\in\{0,1\}\) 且 \(\mu\) 为分拆，也就是每行至多增加一格。每个 \(\mu\) 来自唯一的 \(J\)，系数为一。再除以 \(\Delta_N\) 即得第二式。稳定检测保证有限变量证明给出对称函数恒等式。
\end{proof}

例如
\[
s_{(2,1)}h_2=s_{(4,1)}+s_{(3,2)}+s_{(3,1,1)}+s_{(2,2,1)}.
\]
\((3,2)\) 的新增格位于第一行第三列和第二行第二列，属于横条；不能把“横条”误解为全部新格须在同一行。

\begin{example}\label{b5:exm:lr-square-21}
利用 \(s_{(2,1)}=h_2h_1-h_3\)，计算 \(s_{(2,1)}^2\)。先按横条规则乘 \(h_2\)，得到上面的四项，再逐项添一格：
\begin{align*}
s_{(2,1)}h_2h_1
={}&s_{(5,1)}+2s_{(4,2)}+2s_{(4,1,1)}+s_{(3,3)}
+3s_{(3,2,1)}\\
&+s_{(3,1,1,1)}+s_{(2,2,2)}+s_{(2,2,1,1)}.
\end{align*}
另一方面，
\[
s_{(2,1)}h_3
=s_{(5,1)}+s_{(4,2)}+s_{(4,1,1)}+s_{(3,2,1)}.
\]
相减得到
\begin{align}
s_{(2,1)}^2
={}&s_{(4,2)}+s_{(4,1,1)}+s_{(3,3)}+2s_{(3,2,1)}
\notag\\
&+s_{(3,1,1,1)}+s_{(2,2,2)}+s_{(2,2,1,1)}.
\label{b5:eq:lr-square-21}
\end{align}
因此一般 Littlewood--Richardson 系数不一定是零或一。对三维空间，只保留至多三行的项；它们的维数依次为 \(27,10,10,8,1\)，其中八维项出现两次，合计 \(64\)，与 \(\dim\mathbb S_{(2,1)}(k^3)^2=8^2\) 相符。
\end{example}

\subsection{Schur 函子与对称幂、外幂}
\label{b5:subsec:0704:03}

\begin{proposition}\label{b5:prop:schur-symmetric-exterior}
设 \(k\) 为特征零域，\(V\) 为有限维 \(k\) 向量空间，\(n\ge1\) 为整数，\(\mathbb S_\lambda\) 为分拆 \(\lambda\) 对应的 Schur 函子。则自然地有
\[
\mathbb S_{(n)}(V)\simeq\operatorname{Sym}^nV,\qquad
\mathbb S_{(1^n)}(V)\simeq\bigwedge^nV.
\]
因而对任意分拆 \(\lambda\) 与非负整数 \(r\)，张量积
\(\mathbb S_\lambda(V)\otimes\operatorname{Sym}^rV\) 按横条分解，
\(\mathbb S_\lambda(V)\otimes\bigwedge^rV\) 按竖条分解；行数超过 \(\dim V\) 的项删去。
\end{proposition}
\begin{proof}
\(S^{(n)}\) 为平凡模，故 \(\mathbb S_{(n)}(V)\) 是对称张量子空间。对称幂的商映射 \(q:V^{\otimes n}\to\operatorname{Sym}^nV\)，在该子空间上是同构，其逆为
\[
v_1\cdots v_n\longmapsto
\frac1{n!}\sum_{\sigma\in S_n}
v_{\sigma(1)}\otimes\cdots\otimes v_{\sigma(n)}.
\]
右式对因子排列不变，所以良定；与商映射复合为恒等，对已对称张量再平均也为恒等。故两个复合都正确，且与线性映射相容。

\(S^{(1^n)}\) 为符号模，故相应 Hom 空间识别为反对称张量。外幂的商映射在此子空间上的逆为
\[
v_1\wedge\cdots\wedge v_n\longmapsto
\frac1{n!}\sum_{\sigma\in S_n}\operatorname{sgn}(\sigma)
v_{\sigma(1)}\otimes\cdots\otimes v_{\sigma(n)}.
\]
有两个相等向量时，右端按交换它们的置换配对相消，所以通过外幂的定义关系；两次复合同样为恒等。最后把这两种识别代入\eqref{b5:eq:schur-tensor-decomposition}，并使用\cref{b5:thm:pieri}，即得所述分解。
\end{proof}

例如
\[
\operatorname{Sym}^2V\otimes V
\simeq\operatorname{Sym}^3V\oplus\mathbb S_{(2,1)}(V),
\qquad
\bigwedge^2V\otimes V
\simeq\mathbb S_{(2,1)}(V)\oplus\bigwedge^3V.
\]
因此同一个 \(\mathbb S_{(2,1)}\) 同时出现在“先对称两个因子”和“先反对称两个因子”这两种构造中。对称、反对称只是一般 Young 对称性的两端情形。

\ProseParagraphBreak
这些分解由有限群与张量计算得到。研究一般线性 Lie 代数时，分拆还会作为最高权出现，从而给出另一种解释。Schur 多项式的迹与维数计算可对照\cite[Sections 5.21--5.22, pp. 128--130]{EtingofEtAl2011RepresentationTheory}。


\begin{chaptersummary}
循环型决定幂和基，不可约特征标决定 Schur 基；Frobenius 特征映射把这两组信息放在同一个分次代数中。其等距性来自中心化子阶 \(z_\mu\)，乘法相容性来自把轮换分配到两个不变子集的计数。Jacobi--Trudi 与 Young 置换模的三角性进一步确定了不可约类型的标签，而不只是证明两边维数相等。

\ProseParagraphBreak
张量幂上的因子置换与同时线性变换彼此交换。极化证明它们已经给出全部中心化子；半单矩阵块再产生双向对应。分拆超过空间维数的行数时，相应重数空间为零，群代数作用因此可能不忠实，即使对称群本身仍忠实作用。

\ProseParagraphBreak
Schur 函子的张量积由 Schur 函数乘法控制。单行与单列乘法分别按横条、竖条展开，Pieri 规则给出计算低阶 Littlewood--Richardson 系数的方法。一般系数可解释为诱导重数；本章只证明单行和单列的组合规则。
\end{chaptersummary}

\BookExercises

\begin{exercise}\label{b5:ex:07-stable-specialization}
求 \(\dim_{\mathbb Q}\Lambda_{\mathbb Q}^5\)。再求将对称函数代入两个变量的映射在次数五上的核，并说明取多少个变量可以保证对所有次数五的对称函数单射。
\end{exercise}
\begin{solution}
五的分拆为 \((5),(4,1),(3,2),(3,1,1),(2,2,1),(2,1,1,1),(1^5)\)，故维数为七。两个变量中，前三个单项式对称函数仍互相独立，其余四个为零。因此核以这四个 \(m_\lambda\) 为基，维数四。取五个变量时全部七个基元素均存在且独立，所以单射；少于五个变量时 \(e_5=m_{(1^5)}\) 被送到零，因此五也是最少的数目。
\end{solution}

\begin{exercise}\label{b5:ex:07-degree-four-generators}
将 \(h_4\) 与 \(e_4\) 写成幂和函数的多项式，并由此计算 \(\langle h_4,e_4\rangle\)。
\end{exercise}
\begin{solution}
由指数生成函数取四次项，
\begin{align*}
h_4&=\frac{p_1^4+6p_1^2p_2+3p_2^2+8p_1p_3+6p_4}{24},\\
e_4&=\frac{p_1^4-6p_1^2p_2+3p_2^2+8p_1p_3-6p_4}{24}.
\end{align*}
五种循环型的 \(z_\mu\) 依次为 \(24,4,8,3,4\)。Hall 内积因而为
\[
\frac{24-36\cdot4+9\cdot8+64\cdot3-36\cdot4}{24^2}=0.
\]
这也与 \(h_4=s_{(4)}\)、\(e_4=s_{(1^4)}\) 正交一致。
\end{solution}

\begin{exercise}\label{b5:ex:07-schur-two-variables}
用行列式求 \(s_{(3,1)}\)，再把它代入两个变量 \(x,y\)。求 \(s_{(3,1)}(1,1)\) 与 \(s_{(3,1)}(1)\)。
\end{exercise}
\begin{solution}
Jacobi--Trudi 给出 \(s_{(3,1)}=h_3h_1-h_4\)。两变量中 \(h_3h_1-h_4=x^3y+x^2y^2+xy^3=xy(x^2+xy+y^2)\)。所以两项取值分别为三与零；后一结论也由分拆有两行、而变量只有一个得到。
\end{solution}

\begin{exercise}\label{b5:ex:07-h-m-duality}
证明完全对称函数基与单项式对称函数基在 Hall 配对下互为对偶：
\[
\langle h_\lambda,m_\mu\rangle=\delta_{\lambda\mu}.
\]
\end{exercise}
\begin{solution}
在 Cauchy 乘积中先按 \(y\) 的单项式展开：
\[
\prod_{i,j}(1-x_iy_j)^{-1}
=\prod_j\left(\sum_{r\ge0}h_r(x)y_j^r\right)
=\sum_\lambda h_\lambda(x)m_\lambda(y).
\]
最后一式将同一变量置换轨道的单项式合并，每个轨道的系数为 \(h_\lambda(x)\)。同一个 Cauchy 核又为 \(\sum_\nu p_\nu(x)p_\nu(y)/z_\nu\)，是 Hall 配对的逆矩阵张量。因此在第二因子的 \(m_\lambda\) 基下，第一因子的系数恰为对偶基，即 \(h_\lambda\)，得到结论。
\end{solution}

\begin{exercise}\label{b5:ex:07-omega}
令代数自同构 \(\omega:\Lambda_{\mathbb Q}\to\Lambda_{\mathbb Q}\) 由 \(\omega(p_r)=(-1)^{r-1}p_r\) 确定。证明它为保次数的等距对合，并交换 \(h_r\) 与 \(e_r\)。
\end{exercise}
\begin{solution}
幂和自由生成，所以该赋值唯一延拓为代数自同构；再施一次各生成元都不变，因此 \(\omega^2=1\)。它将 \(p_\lambda\) 乘以符号 \((-1)^{|\lambda|-\ell(\lambda)}\)，故保持幂和基的正交范数，因而等距。对 \(H(t)\) 的指数表达式施 \(\omega\)，得到 \(E(t)\)，所以 \(\omega h_r=e_r\)；再施一次得到反向等式。
\end{solution}

\begin{exercise}\label{b5:ex:07-class-indicator}
设 \(1_\mu\) 为 \(S_n\) 中循环型 \(\mu\) 的示性函数。求 \(\operatorname{ch}(1_\mu)\) 与 \(p_\mu\) 对应的类函数，并说明后一类函数一般不是实际表示的特征标。
\end{exercise}
\begin{solution}
由定义，\(\operatorname{ch}(1_\mu)=p_\mu/z_\mu\)，故 \(p_\mu\) 对应 \(z_\mu1_\mu\)。若 \(\mu\ne(1^n)\)，该类函数在单位元处为零，却在指定共轭类上非零。实际表示若在单位元处迹为零，则维数为零，特征标处处为零，所以这里不是实际特征标。它仍是虚特征标，因为由 Frobenius 公式，其不可约展开系数为整数 \(\chi_{S^\lambda}(\mu)\)。
\end{solution}

\begin{exercise}\label{b5:ex:07-sign-twist}
设 \(U\) 为 \(S_n\) 的复表示。用\cref{b5:ex:07-omega}中的 \(\omega\)证明
\[
\operatorname{ch}(U\otimes\mathrm{sgn})=\omega\operatorname{ch}(U).
\]
计算 \(S_3\) 的 \(S^{(2,1)}\) 张量符号表示后是否改变。
\end{exercise}
\begin{solution}
循环型 \(\mu\) 上，符号为 \((-1)^{n-\ell(\mu)}\)，而 \(\omega p_\mu\) 恰乘同一个符号。逐项代入特征映射定义得到公式。\(s_{(2,1)}=(p_1^3-p_3)/3\) 在 \(\omega\) 下不变，所以该二维表示张量符号后同构于自身。
\end{solution}

\begin{exercise}\label{b5:ex:07-induced-sign-trivial}
分解 \(\operatorname{Ind}_{S_2\times S_2}^{S_4}(\mathrm{sgn}\boxtimes1)\)，并核对维数。
\end{exercise}
\begin{solution}
特征映射给出 \(e_2h_2=s_{(1,1)}h_2\)。从两格竖列添大小二的横条，得到 \((3,1)\) 与 \((2,1,1)\)，所以分解为 \(S^{(3,1)}\oplus S^{(2,1,1)}\)。两者维数均为三，和为六，恰为指数 \([S_4:S_2\times S_2]\)。
\end{solution}

\begin{exercise}\label{b5:ex:07-subset-permutation}
设 \(0\le r\le n/2\)。分解 \(S_n\) 在全部 \(r\) 元子集上的复置换表示，并说明它是否有重复的不可约类型。
\end{exercise}
\begin{solution}
一个 \(r\) 元子集的稳定子群为 \(S_r\times S_{n-r}\)，故特征映射为 \(h_rh_{n-r}\)。在长度 \(n-r\) 的单行上添 \(r\) 格横条，只能得到 \((n-i,i)\)，其中 \(0\le i\le r\)。这些形状都确实可以得到，而且系数为一。因此置换表示为 \(\bigoplus_{i=0}^rS^{(n-i,i)}\)，没有重复类型。
\end{solution}

\begin{exercise}\label{b5:ex:07-two-products}
以 \(S_2\) 的平凡表示为例，说明 Frobenius 特征映射不把同一个对称群上的张量积送到普通的对称函数乘法。
\end{exercise}
\begin{solution}
\(1_{S_2}\otimes1_{S_2}\simeq1_{S_2}\)，其像为次数二的 \(h_2\)。普通乘积 \(h_2^2\) 却是次数四的 \(s_{(4)}+s_{(3,1)}+s_{(2,2)}\)，对应从 \(S_2\times S_2\) 到 \(S_4\) 的诱导。次数已不同，因而两种乘法不可能混同。
\end{solution}

\begin{exercise}\label{b5:ex:07-standard-character}
对 \(n\ge2\)，用 \(s_{(n-1,1)}=h_{n-1}h_1-h_n\) 证明
\[
\chi_{S^{(n-1,1)}}(\sigma)=m_1(\sigma)-1,
\]
其中 \(m_1(\sigma)\) 为 \(\sigma\) 的不动点数。
\end{exercise}
\begin{solution}
\(h_{n-1}h_1\) 对应从 \(S_{n-1}\times S_1\) 诱导的平凡表示，即 \(n\) 个字母的置换表示。其特征标等于不动字母数 \(m_1\)。减去 \(h_n\) 对应的平凡特征标一，得到所求式。这同时说明该 Specht 模就是去掉常数直线后的标准表示。
\end{solution}

\begin{exercise}\label{b5:ex:07-two-row-character}
对 \(n\ge4\)，证明
\[
\chi_{S^{(n-2,2)}}(\sigma)
=\binom{m_1(\sigma)}2+m_2(\sigma)-m_1(\sigma),
\]
其中 \(m_2(\sigma)\) 为二轮换数。
\end{exercise}
\begin{solution}
\(\sigma\) 固定的二元子集或者由两个不动点组成，或者是某个二轮换的支集，所以二元子集置换特征标为 \(\binom{m_1}2+m_2\)。由横条规则，该置换表示为 \(S^{(n)}\oplus S^{(n-1,1)}\oplus S^{(n-2,2)}\)。前两项特征标之和为 \(1+(m_1-1)=m_1\)，减去即得公式。
\end{solution}

\begin{exercise}\label{b5:ex:07-regular-schur}
证明 \(p_1^n=\sum_{\lambda\vdash n}f^\lambda s_\lambda\)，其中 \(f^\lambda=\dim S^\lambda\)。据此解释 \(\langle p_1^n,p_1^n\rangle=n!\)。
\end{exercise}
\begin{solution}
\(p_1=h_1\) 对应 \(S_1\) 的平凡表示，\(n\) 次诱导乘积就是从平凡子群到 \(S_n\) 诱导，即正则表示。正则表示的 Specht 重数为其维数，给出所求展开。Schur 正交性随后给出范数平方为 \(\sum_\lambda(f^\lambda)^2=n!\)。也可直接由幂和范数 \(z_{(1^n)}=n!\) 得到同样数值。
\end{solution}

\begin{exercise}\label{b5:ex:07-centralizer-dimensions}
设 \(V=\mathbb C^2\)，\(T=V^{\otimes3}\)。求对称群作用像 \(A\) 及一般线性群作用生成代数 \(B\) 的维数，并说明为什么不能把两个数相乘来计算 \(\dim\operatorname{End}(T)\)。
\end{exercise}
\begin{solution}
只出现 \((3),(2,1)\)。相应 Specht 维数为一、二，重数空间维数为四、二。所以 \(\dim A=1^2+2^2=5\)，\(\dim B=4^2+2^2=20\)。而 \(\dim\operatorname{End}(T)=8^2=64\)，不是一百。两个代数在各等型块中分别作用于不同因子，但在不同块上不是任意组合；乘法映射 \(A\otimes B\to\operatorname{End}(T)\) 还有交叉块被送到零，不能据此相乘维数。
\end{solution}

\begin{exercise}\label{b5:ex:07-commuting-insufficient}
在二维复空间中，取 \(A=B=\mathbb C I\)。检验它们相互交换，但不满足 Schur--Weyl 型的双中心化子结论，并指出缺少什么。
\end{exercise}
\begin{solution}
标量矩阵与一切矩阵交换，所以两者当然相互交换。然而 \(\operatorname{End}_A(\mathbb C^2)=M_2(\mathbb C)\ne B\)。\(A\) 虽然半单，但 \(B\) 并未取为它的全部中心化子。因此半单性与交换关系合在一起仍不够；还须实际证明其中一个中心化子等式。
\end{solution}

\begin{exercise}\label{b5:ex:07-polarization-cubic}
在特征零域上，将
\[
u\otimes u\otimes v+u\otimes v\otimes u+v\otimes u\otimes u
\]
写成三个向量的三次纯幂的线性组合。
\end{exercise}
\begin{solution}
展开 \((u+v)^{\otimes3}-(u-v)^{\otimes3}\)，含偶数个 \(v\) 的项相消，含一个 \(v\) 的三项各出现两次，另有 \(2v^{\otimes3}\)。因此所求表达式为
\[
\frac12\bigl((u+v)^{\otimes3}-(u-v)^{\otimes3}-2v^{\otimes3}\bigr).
\]
这里张量乘法不交换，所以展开时须保留三个不同的位置。
\end{solution}

\begin{exercise}\label{b5:ex:07-alternating-kernel}
设 \(\dim_{\mathbb C}V=2\)。证明 \(\mathbb C[S_3]\to\operatorname{End}(V^{\otimes3})\) 的核恰由
\[
a=\sum_{\sigma\in S_3}\operatorname{sgn}(\sigma)\sigma
\]
线性生成，而 \(S_3\) 本身仍忠实作用。
\end{exercise}
\begin{solution}
唯一未出现的分拆是 \((1,1,1)\)，它的矩阵块一维；因而核一维。\(a/6\) 是符号类型的中心投影，作用像为三阶反对称张量，而 \(\bigwedge^3V=0\)，所以 \(a\) 在核中。它的单位元系数为一，故非零，必生成整个核。另一方面，两个独立向量已可构造只有一个特殊位置的基张量来检测任意非恒等置换，故群作用忠实。
\end{solution}

\begin{exercise}\label{b5:ex:07-general-linear-kernel}
对 \(d\ge1\)，分别求 \(\operatorname{GL}_d(\mathbb R)\) 在二次与三次张量幂上的核，再与复数域比较。
\end{exercise}
\begin{solution}
核由满足 \(c^n=1\) 的标量 \(cI\) 组成。实数域上二次幂的核为 \(\{I,-I\}\)，三次幂的核只有 \(I\)。复数域上分别有两个和三个标量根，所以三次张量表示在实数域上忠实，在复数域上不忠实。这里讨论的是有理点群的作用，不能忽略底域。
\end{solution}

\begin{exercise}\label{b5:ex:07-schur-dimensions}
对 \(d\ge1\)，求 \(\dim\mathbb S_{(2,1)}(k^d)\) 与 \(\dim\mathbb S_{(2,2)}(k^d)\)，其中 \(k\) 为特征零域。
\end{exercise}
\begin{solution}
在 \(d\) 个一处，\(h_r(1^d)=\binom{d+r-1}{r}\)，因为它数次数 \(r\) 的单项式。代入两个 Jacobi--Trudi 公式，
\begin{align*}
\dim\mathbb S_{(2,1)}(k^d)
&=\binom{d+1}{2}d-\binom{d+2}{3}
=\frac{d(d^2-1)}3,\\
\dim\mathbb S_{(2,2)}(k^d)
&=\binom{d+1}{2}^2-d\binom{d+2}{3}
=\frac{d^2(d^2-1)}{12}.
\end{align*}
\(d=1\) 时均为零，与两行限制相符；\(d=3\) 时分别为八与六。
\end{solution}

\begin{exercise}\label{b5:ex:07-determinant-shift}
设 \(k\) 为特征零域，\(V\) 为 \(d\ge1\) 维 \(k\) 向量空间，\(\lambda\) 为至多有 \(d\) 行的分拆，\(r\) 为非负整数，\(\det V=\bigwedge^dV\)，\(\lambda+(r^d)\) 表示补零后各部分加 \(r\)。证明 \(\operatorname{GL}_k(V)\)-表示同构
\[
\mathbb S_{\lambda+(r^d)}(V)
\simeq\mathbb S_\lambda(V)\otimes(\det V)^{\otimes r}.
\]
\end{exercise}
\begin{solution}
在交错式商中，各指数同时增加 \(r\)，分子提出 \((x_1\cdots x_d)^r\)，故
\[
s_{\lambda+(r^d)}(x)=(x_1\cdots x_d)^r s_\lambda(x).
\]
\(\det V=\bigwedge^dV\) 为一维表示。右侧是 \(V^{\otimes(|\lambda|+rd)}\) 的直和项，因而完全可约，其特征标恰为左侧 Schur 多项式。该张量次数的 Schur 特征标线性无关，所以表示同构。
\end{solution}

\begin{exercise}\label{b5:ex:07-gl2}
设 \(k\) 为特征零域，\(V\) 为二维 \(k\) 向量空间，\(a,b\) 为满足 \(a\ge b\ge0\) 的整数，\(\det V=\bigwedge^2V\)。证明
\[
\mathbb S_{(a,b)}(V)\simeq
\operatorname{Sym}^{a-b}V\otimes(\det V)^{\otimes b},
\]
并求其维数。
\end{exercise}
\begin{solution}
从单行分拆 \((a-b,0)\) 开始，向两行同时添 \(b\) 格，即得 \((a,b)\)。用\cref{b5:ex:07-determinant-shift}中的行列式计算，特征标为 \((xy)^b h_{a-b}(x,y)\)，所以得到所述同构。行列式因子一维，而二维空间的 \(r\) 次对称幂以 \(x^{r-i}y^i\)、\(0\le i\le r\) 为基，所以维数为 \(a-b+1\)。
\end{solution}

\begin{exercise}\label{b5:ex:07-symmetric-exterior-product}
设 \(k\) 为特征零域，\(V\) 为有限维 \(k\) 向量空间，\(a,b\ge1\) 为整数。分解 \(\operatorname{GL}_k(V)\)-表示 \(\operatorname{Sym}^aV\otimes_k\bigwedge^bV\)，并在 \(\dim_kV=b\) 时简化答案。
\end{exercise}
\begin{solution}
在单行 \((a)\) 上添 \(b\) 格竖条，第一行或者添一格，或者不添，得到
\[
\operatorname{Sym}^aV\otimes\bigwedge^bV
\simeq\mathbb S_{(a+1,1^{b-1})}(V)
\oplus\mathbb S_{(a,1^b)}(V).
\]
当 \(\dim V=b\) 时，第二项有 \(b+1\) 行而为零。第一种形状 \((a+1,1^{b-1})\) 对每行同时减一后成为 \((a,0,\ldots,0)\)，故第一项为 \(\operatorname{Sym}^{a}V\otimes\det V\)。这与左侧 \(\bigwedge^bV=\det V\) 的直接计算一致。
\end{solution}

\begin{exercise}\label{b5:ex:07-square-symmetric-exterior}
设 \(k\) 为特征零域，\(V\) 为有限维 \(k\) 向量空间。分别分解 \(\operatorname{GL}_k(V)\)-表示 \(\operatorname{Sym}^2V\otimes_k\operatorname{Sym}^2V\) 与 \(\bigwedge^2V\otimes_k\bigwedge^2V\)，并说明二维空间中哪些项消失。
\end{exercise}
\begin{solution}
横条与竖条规则给出
\begin{align*}
h_2^2&=s_{(4)}+s_{(3,1)}+s_{(2,2)},\\
e_2^2&=s_{(2,2)}+s_{(2,1,1)}+s_{(1,1,1,1)}.
\end{align*}
分别把 \(s_\lambda\) 换成 \(\mathbb S_\lambda(V)\) 得到分解。二维时第一式三项均存在，维数为五、三、一，总计九；第二式只剩 \(\mathbb S_{(2,2)}(V)=(\det V)^{\otimes2}\)，维数为一。
\end{solution}

\begin{exercise}\label{b5:ex:07-lr-dimension-two}
令 \(V\) 为二维特征零空间。从\eqref{b5:eq:lr-square-21}求 \(\mathbb S_{(2,1)}(V)^{\otimes2}\)，再用一般线性群的二维表达式核对。
\end{exercise}
\begin{solution}
删去三行及更多行的分拆，得到
\[
\mathbb S_{(2,1)}(V)^{\otimes2}
\simeq\mathbb S_{(4,2)}(V)\oplus\mathbb S_{(3,3)}(V).
\]
另一方面，\(\mathbb S_{(2,1)}(V)\simeq V\otimes\det V\)，其平方为 \(V^{\otimes2}\otimes(\det V)^{\otimes2}\)。将 \(V^{\otimes2}\) 分成 \(\operatorname{Sym}^2V\oplus\det V\)，得到 \(\operatorname{Sym}^2V\otimes(\det V)^{\otimes2}\) 与 \((\det V)^{\otimes3}\)，正对应上述两项。维数为 \(3+1=4\)。
\end{solution}

\begin{exercise}\label{b5:ex:07-cauchy-degree-two}
设 \(V,W\) 为特征零域上的有限维空间。证明自然同构
\[
\operatorname{Sym}^2(V\otimes W)
\simeq
(\operatorname{Sym}^2V\otimes\operatorname{Sym}^2W)
\oplus(\bigwedge^2V\otimes\bigwedge^2W),
\]
并写出 \(\bigwedge^2(V\otimes W)\) 的对应分解。
\end{exercise}
\begin{solution}
把 \((V\otimes W)^{\otimes2}\) 重排成 \(V^{\otimes2}\otimes W^{\otimes2}\)。交换两个 \(V\otimes W\) 因子，对应 \(P_V\otimes P_W\)。两个交换算子的特征值都是 \(\pm1\)，相同符号相乘为正，异号相乘为负。因此正特征空间给出题中分解，负特征空间给出
\[
\bigwedge^2(V\otimes W)
\simeq(\operatorname{Sym}^2V\otimes\bigwedge^2W)
\oplus(\bigwedge^2V\otimes\operatorname{Sym}^2W).
\]
所有重排与投影均与线性映射相容，所以同构自然。
\end{solution}

\begin{exercise}\label{b5:ex:07-schur-functor-nonlinearity}
设 \(V=k\)、\(\lambda=(2)\)，其中 \(k\) 为特征零域。直接说明 \(\mathbb S_\lambda\) 保持映射复合，却不把映射之和变成映射之和。再求 \(\mathbb S_\lambda(cI_V)\) 对一般 \(\lambda\vdash n\) 的作用。
\end{exercise}
\begin{solution}
一维空间上的线性映射由标量 \(a\) 表示，对称平方将它送到 \(a^2\)。所以 \((ab)^2=a^2b^2\)，但取 \(a=b=1\) 时 \((a+b)^2=4\ne2=a^2+b^2\)。一般地，\((cI)^{\otimes n}=c^nI\)，故对整个 Hom 重数空间也作用为 \(c^nI\)。函子性要求保持复合，不要求关于态射加法线性。
\end{solution}

\begin{exercise}\label{b5:ex:07-trace-cubic}
设 \(g\) 为特征零域上有限维空间的可逆算子，记 \(t_r=\operatorname{tr}(g^r)\)。证明
\begin{align*}
\operatorname{tr}(\operatorname{Sym}^3g)
&=\frac{t_1^3+3t_1t_2+2t_3}{6},\\
\operatorname{tr}(\bigwedge^3g)
&=\frac{t_1^3-3t_1t_2+2t_3}{6}.
\end{align*}
对三维单个幺幂 Jordan 块计算这两个迹。
\end{exercise}
\begin{solution}
对称幂与外幂的特征标为 \(h_3,e_3\)。将幂和 \(p_r\) 代入 \(t_r\)，使用其指数生成函数的三次项，得到两个公式；它们对不可对角化的 \(g\) 同样成立，因为三角化后的对角特征值仍计算全部幂的迹。三维幺幂 Jordan 块满足 \(t_1=t_2=t_3=3\)，故两个值为十与一。
\end{solution}

\begin{exercise}\label{b5:ex:07-modular-obstruction}
设 \(k\) 的特征为二，\(V=k^2\)，取基 \(u,v\)。证明非零张量 \(u\otimes v+v\otimes u\) 在对称平方商映射 \(V^{\otimes2}\to\operatorname{Sym}^2V\) 下为零。由此说明正文把对称张量与对称幂相互识别的证明为何不能直接搬到此处。
\end{exercise}
\begin{solution}
两个基张量互不相同，所以其和非零；交换它们后和不变，故它是对称张量。在对称平方中二者都映到 \(uv\)，和映到 \(2uv=0\)。因此商映射在对称张量子空间上不单射。正文的逆需要除以 \(2!\)，在特征二中不存在；不变子空间与对称商空间不能再由那组平均映射识别。
\end{solution}

% !TeX root = ../../Book5.tex
% 纲要标题：## 第8章*　模表示与 Brauer 特征标
% 纲要位置：工作记录/计划纲要.md，第 3981 行。
\OptionalChapter{模表示与 Brauer 特征标}{b5:ch:08}
\BookChapterNumber{b5:ch:08}

\begin{chapterabstract}
本章研究底域特征整除群阶时的有限群表示，借投射覆盖与块分析不可分解模，再以 Brauer 特征标比较模特征与特征零的表示信息。
\end{chapterabstract}

\begin{learninggoals}
\begin{itemize}
\item 能构造半单性失效的例子，区分简单模、不可分解模及其根层。
\item 能计算循环 \(p\) 群的模、投射覆盖、基座和 Cartan 矩阵。
\item 从中心本原幂等元得到块，并在小群的群代数中验证块分解。
\item 能在 \(p\)-正则共轭类上计算 Brauer 特征标，并由格约化确定分解数。
\item 能按定义计算 Brauer 映射及具体块的缺陷群，明确特征标记录合成因子的范围。
\end{itemize}
预备知识为有限群表示、群代数、合成列及有限维代数的根；投射覆盖的构造使用第二卷的幂等元提升。Brauer 特征标部分另固定一个适当的分裂 \(p\)-模系统，届时会说明特征零域、离散赋值环及剩余域各自的作用。
\end{learninggoals}

设 \(G=C_p=\langle g\rangle\)，而域 \(k\) 的特征为素数 \(p\)。在群代数中令 \(u=g-1\)，则 \(u^p=g^p-1=0\)，并有 \(kG\cong k[u]/(u^p)\)。这个代数的逐次商 \(u^j kG/u^{j+1}kG\) 都是一维平凡表示，但乘 \(u\) 的作用把它们接成一个不能直接拆开的模。简单成分因而不再决定整个表示，特征零下用平均投影取得不变补空间的办法也失效。投射覆盖、块与 Brauer 特征标将分别记录这里留下的结构和数值信息。投射覆盖及有限维代数 Cartan 矩阵的基础构造，可对照\cite[Sections 9.1--9.3, pp. 209--212]{EtingofEtAl2011RepresentationTheory}。

\ProseParagraphBreak
引用 Webb 著作时，本章采用作者2016年2月23日公开稿的节号与结果编号。作者说明该稿与正式刊本的数学内容相同；稿件页码与刊本页码另有区别，以下引用使用节号。格约化与 Brauer 特征标见\cite[Sections 9.4, 10.1; Section 12.5]{Webb2016FiniteGroupRepresentations}。

% !TeX root = ../../Book5.tex
% 纲要标题：### 8.1 半单性的失效
% 纲要位置：工作记录/计划纲要.md，第 3983 行。

\section{半单性的失效}
\label{b5:sec:0801}


% 纲要范围：
%
% * 特征 p 整除群阶
% * 群代数的根及不可分解模
% * 循环 p 群的完整基本例子
%

当底域的特征整除群阶时，平均投影中的除法不再有意义。表示仍然可以有合成列，但这些简单商之间的扩张未必分裂；因此研究简单模还不够，还须研究不可分解模及其根层。本章以循环 \(p\) 群为反复计算的例子，始终区分直和分解与合成因子。

\subsection{\texorpdfstring{特征 \(p\) 整除群阶的情形}{特征 p 整除群阶的情形}}
\label{b5:subsec:0801:01}

\begin{definition}\label{b5:def:modular-representation}
设 \(G\) 为有限群，\(k\) 为特征 \(p>0\) 的域。\(G\) 的有限维 \(k\) 表示称为特征 \(p\) 的\term[mo-biaoshi]{模表示}[modular representation]。本章重点讨论 \(p\,\bigm\vert\,|G|\) 的情形；若 \(p\mathrel{\vcenter{\hbox{\scalebox{1.25}{\(\nmid\)}}}}|G|\)，\cref{b5:thm:maschke}仍保证完全可约。
\end{definition}
由\cref{b5:prop:maschke-converse}，当 \(p\,\bigm\vert\,|G|\) 时，增广映射
\[
\varepsilon:kG\longrightarrow k,\qquad \sum_g a_gg\longmapsto\sum_g a_g
\]
没有 \(kG\) 线性截面。具体地，一个不动元必须为 \(c\sum_g g\)，它的增广为 \(c|G|=0\)，所以不能提升目标中的一。这既给出不完全可约的正则模，也给出取不动空间不保持满射的例子。
\ProseParagraphBreak
最小的矩阵例子是 \(C_p=\langle s\rangle\) 在 \(k^2\) 上作用为
\[
s\longmapsto\begin{pmatrix}1&1\\0&1\end{pmatrix}.
\]
第一坐标轴和商都为平凡表示，但没有不变补。这个表示与平凡二维表示有相同的两个简单合成因子，却不同构：前者的生成元不是恒等。后面定义的 Brauer 特征标会记录合成因子的重数，不会额外记录这条扩张是否分裂。

\subsection{群代数的根与不可分解模}
\label{b5:subsec:0801:02}

设 \(A\) 为有限维含幺 \(k\) 代数。其 Jacobson 根可写为全部简单左模的零化子的交，记为 \(J=J(A)\)。若 \(M\) 为有限维左模，则 \(JM\) 是其\term[mo-gen]{根}[radical of a module]，即全部极大子模的交；商 \(M/JM\) 称为\term[mo-ding]{顶}[top of a module]。根层 \(J^iM/J^{i+1}M\) 记录从顶逐步向下的简单成分。它们通常不能在原模中同时选成直和项。
\begin{proposition}\label{b5:prop:finite-radical-layers}
设 \(k\) 为域，\(A\) 为有限维含幺 \(k\) 代数，\(J=J(A)\)。则 \(J\) 幂零，\(A/J\) 半单。对任意有限维左 \(A\) 模 \(M\)，全部根层 \(J^iM/J^{i+1}M\) 都是半单模，且 \(JM=M\) 蕴涵 \(M=0\)。
\end{proposition}
\begin{proof}
% 跨卷目标：b2:prop:finite-algebra-semisimple-quotient（第二卷命题15.1.3）。
取左正则模 \(A\) 的有限合成列，设长度为 \(\ell\)。\(J\) 零化每个简单商，故乘以 \(J\) 将每一步送入前一步，迭代得到 \(J^\ell A=0\)，即 \(J^\ell=0\)。有限维代数的半单商结构由第二卷命题15.1.3及该节建立的半单商分解给出：\(A/J\) 为有限个除环上矩阵代数的直积。每个根层被 \(J\) 零化，是 \(A/J\) 模，因而半单。最后若 \(JM=M\)，则 \(M=J^\ell M=0\)。
\end{proof}
上述最后一步是这里需要的 Nakayama 引理\footnote{中山正（假名：なかやま ただし，罗马音：Tadashi Nakayama，1912-1964），日本数学家，研究环论和群表示论。}。它也说明为什么顶不能为零：任何非零有限维模都有非零的顶。有限维模按不可分解项的直和分解及其唯一性可用第二卷定理12.4.2；该定理的有限长度条件在此由有限维性满足。不可分解模的各根层则可能包含多种简单模。
\begin{proposition}\label{b5:prop:p-group-local}
设 \(p\) 为素数，\(P\) 为有限 \(p\)-群，\(k\) 为特征 \(p\) 的域，\(\varepsilon:kP\to k\) 为增广映射 \(\sum_g a_gg\mapsto\sum_g a_g\)。唯一的简单左 \(kP\)-模是平凡一维模 \(k\)，且
\[
J(kP)=\ker\varepsilon,\qquad kP/J(kP)\cong k.
\]
因此 \(kP\) 是局部代数。
\end{proposition}
\begin{proof}
按 \(|P|\) 归纳。若 \(P\ne1\)，第一卷定理4.2.5给出 \(p\,\bigm\vert\,|Z(P)|\)：类方程中各非中心类的大小都被 \(p\) 整除。取中心中一个非单位元，再取其适当的 \(p\) 次幂，便得阶恰为 \(p\) 的元 \(z\)。在简单模 \(S\) 上，\(z-1\) 是模自同态，且 \((z-1)^p=z^p-1=0\)。Schur 引理使非零自同态可逆，幂零自同态不可能可逆，所以 \(z-1\) 作用为零。于是 \(S\) 通过较小的 \(P/\langle z\rangle\) 作用，归纳使它平凡。平凡作用下简单性又迫使维数为一。根是这个唯一简单模的零化子，恰为增广理想。模掉根是域，所以非单位元恰为根，代数局部。
\end{proof}
例如一个非循环 \(p\) 群仍然只有一个简单模，但其不可分解模远不止一种。简单类型的数目很少，并不意味着扩张和直和分解的分类也很短。

\subsection{\texorpdfstring{循环 \(p\) 群的基本例子}{循环 p 群的基本例子}}
\label{b5:subsec:0801:03}

\begin{theorem}\label{b5:thm:cyclic-p-modules}
设 \(p\) 为素数，\(k\) 为特征 \(p\) 的域，\(r\ge1\) 为整数，\(G=C_{p^r}=\langle s\rangle\)，并记 \(m=p^r\)。则
\[
kG\cong k[t]/(t^m),\qquad t\longmapsto s-1.
\]
全部有限维不可分解左 \(kG\)-模的同构类代表恰为
\[
M_d=k[t]/(t^d),\qquad 1\le d\le m.
\]
每个有限维左 \(kG\)-模唯一分解成这些模的直和，唯一性指各项同构类的重数唯一。记 \(J=J(kG)\)，则对 \(i\ge0\) 及 \(1\le a,b,d\le m\) 有
\[
J^iM_d=t^iM_d,\qquad
\operatorname{soc}(M_d)=kt^{d-1},\qquad
\dim_k\operatorname{Hom}_{kG}(M_a,M_b)=\min(a,b).
\]
其中\term[jizuo]{基座}[socle] \(\operatorname{soc}(M_d)\) 是全部简单子模之和。
\end{theorem}
\begin{proof}
特征 \(p\) 中重复使用 \((x-1)^p=x^p-1\)，得到 \((s-1)^m=0\)。由 \(s=1+t\) 可知所给代数同态满射，两侧维数都是 \(m\)，故同构。模结构因此由满足 \(T^m=0\) 的线性算子 \(T\) 决定。幂零算子的 Jordan 链分解，或第二卷定理4.3.2在 \(k[t]\) 上的应用，给出模为 \(k[t]/(t^d)\) 的直和，且 \(d\le m\)。这里多项式只有不可约因子 \(t\)，所以不要求 \(k\) 代数闭。
\(M_d\) 的自同态由一的像决定，环为 \(k[t]/(t^d)\)。这个环的常数项非零元素可用有限几何级数求逆，常数项零元素幂零，故它局部，没有非平凡幂等元。因此 \(M_d\) 不可分解；不同 \(d\) 的维数不同。任一分解中，大小至少为 \(j\) 的块数等于
\[
\dim\ker T^j-\dim\ker T^{j-1},
\]
所以块的重数唯一。根层由 \(J=(t)\) 直接给出；简单子模被 \(t\) 零化，而 \(\ker(t:M_d\to M_d)=kt^{d-1}\)，得到基座。最后，\(M_a\to M_b\) 的同态由一个满足 \(t^a v=0\) 的 \(v\in M_b\) 决定，这个核的维数为 \(\min(a,b)\)。
\end{proof}
这些模的全部子模也成一条链。因为 \(M_d\) 是 \(k[t]\) 的循环商模，子模对应包含 \((t^d)\) 的理想，恰为 \((t^j)\)、\(0\le j\le d\)。例如 \(M_3\) 的根层各为一个平凡模，却不是三个平凡模的直和；其不动空间只有一维，后者有三维。

% !TeX root = ../../Book5.tex
% 纲要标题：### 8.2 投射与块
% 纲要位置：工作记录/计划纲要.md，第 3989 行。

\section{投射与块}
\label{b5:sec:0802}


% 纲要范围：
%
% * 不可分解投射模
% * 投射覆盖与 Cartan 矩阵
% * 块分解的基本概念
%

非半单模往往不能拆成简单模，但仍可用投射模覆盖它。不可分解投射模由简单顶来区分，Cartan 矩阵记录它们内部的合成因子；块则把整个代数分成互不联系的部分。本节先把这三种层次分别说清，再在小群中计算。

\subsection{不可分解投射模}
\label{b5:subsec:0802:01}

若 \(P\) 为有限维左 \(A\) 模，且每个到 \(P\) 的满射都分裂，便称 \(P\) 投射；等价地，从 \(P\) 出发的同态可沿任意满射提升。有限生成情形中，这又等价于 \(P\) 是有限自由模的直和因子。将满射 \(A^n\to P\) 分裂即可得到后一条件，反向的提升可先在自由模上逐基选取。
\begin{definition}\label{b5:def:modular-projective-cover}
设 \(k\) 为域，\(A\) 为有限维含幺 \(k\) 代数，\(P,M\) 为有限维左 \(A\)-模。若 \(\pi:P\to M\) 为满射模同态，\(P\) 投射，并且 \(\ker\pi\) 是多余子模，即
\(U+\ker\pi=P\) 对子模 \(U\subseteq P\) 只能在 \(U=P\) 时成立，就称 \(\pi\) 为 \(M\) 的\term[toushe-fugai]{投射覆盖}[projective cover]。
\end{definition}
\begin{proposition}\label{b5:prop:projective-cover-top}
设 \(k\) 为域，\(A\) 为有限维含幺 \(k\) 代数，\(J=J(A)\)，\(P,M\) 为有限维左 \(A\)-模，且 \(P\) 投射、\(\pi:P\to M\) 为满射模同态。\(\pi\) 是投射覆盖，当且仅当 \(\ker\pi\subseteq JP\)，也当且仅当它诱导的映射 \(P/JP\to M/JM\) 为同构。
\end{proposition}
\begin{proof}
若 \(N=\ker\pi\) 多余，则任何极大子模都包含 \(N\)，否则与 \(N\) 的和为 \(P\)，矛盾。故 \(N\subseteq JP\)。反过来，若 \(N\subseteq JP\) 且 \(U+N=P\)，则 \(P/U=J(P/U)\)，由\cref{b5:prop:finite-radical-layers}为零。最后 \(\pi(JP)=JM\)，商映射的核为 \((N+JP)/JP\)，得到等价性。
\end{proof}
对 \(A=k[t]/(t^m)\)，自然满射 \(A\to M_d\) 是投射覆盖，因为核 \((t^d)\) 包含在根 \((t)\) 中。当 \(d<m\) 时它不分裂：若分裂，循环模 \(A\) 就有非平凡直和分解，与其局部自同态环矛盾。所以这些较短的不可分解模都不是投射模。

\subsection{投射覆盖与 Cartan 矩阵}
\label{b5:subsec:0802:02}

\begin{theorem}\label{b5:thm:split-projective-covers}
设 \(k\) 为域，\(A\) 为有限维含幺 \(k\) 代数，且对正整数 \(s,n_1,\ldots,n_s\) 有
\[
A/J(A)\cong\prod_{i=1}^sM_{n_i}(k).
\]
取简单左 \(A\)-模的同构类代表 \(S_1,\ldots,S_s\)，其中 \(\dim_kS_i=n_i\)，记 \([M:S_i]\) 为 \(S_i\) 在有限维左 \(A\)-模 \(M\) 的合成列中的重数。则每个 \(S_i\) 有不可分解投射覆盖 \(P_i\to S_i\)，且作为左 \(A\)-模有
\[
A\cong\bigoplus_iP_i^{\oplus n_i},\qquad
\dim_k\operatorname{Hom}_A(P_i,M)=[M:S_i]
\]
对任意有限维左 \(A\) 模 \(M\) 成立。全部有限维投射模均为这些 \(P_i\) 的直和，各 \(P_i\) 的类型由顶 \(S_i\) 唯一确定。
\end{theorem}
\begin{proof}
记 \(J=J(A)\)。利用根的幂零性及第二卷定理15.3.4，把商中各矩阵块的对角矩阵单位提升为一组完整正交幂等元 \(e_{ij}\)、\(1\le j\le n_i\)。令 \(P_{ij}=Ae_{ij}\)。它们是正则模的直和因子，顶为简单列模 \(S_i\)，且自然商的核为 \(Je_{ij}\)，由\cref{b5:prop:projective-cover-top}为覆盖。顶简单也使 \(P_{ij}\) 不可分解：若它分成两个非零项，两项的顶都非零，会分解这个简单顶。
\ProseParagraphBreak
同一个简单模 \(S\) 的任意两个投射覆盖 \(\pi:P\to S\)、\(\pi':P'\to S\)，存在与覆盖映射相容的同构 \(P\to P'\)。先由\cref{b5:prop:projective-cover-top}得到 \(\ker\pi\subseteq JP\)；另一方面，\(J\) 零化 \(S\)，故 \(JP\subseteq\ker\pi\)。因此 \(\ker\pi=JP\)，同样有 \(\ker\pi'=JP'\)，覆盖映射分别给出 \(P/JP\cong S\)、\(P'/JP'\cong S\)。投射性给出相互提升 \(u:P\to P'\)、\(v:P'\to P\)，满足 \(\pi'u=\pi\)、\(\pi v=\pi'\)。由 \(\pi vu=\pi\)，复合 \(vu\) 在 \(P/JP\) 上诱导恒等。于是 \(P=vu(P)+JP\)，商 \(P/vu(P)\) 等于其根，由\cref{b5:prop:finite-radical-layers}中的 Nakayama 结论为零。因此 \(vu\) 满射，有限维性使它可逆。同理 \(uv\) 可逆，故 \(u\) 为同构，且已有 \(\pi'u=\pi\)。这使 \(P_{ij}\) 的同构类型只依赖 \(i\)，记为 \(P_i\)，正交分解给出正则模公式。
\ProseParagraphBreak
一般 \(M\) 的顶可写成 \(M/JM\cong\bigoplus_iS_i^{\oplus m_i}\)。令 \(Q=\bigoplus_iP_i^{\oplus m_i}\)，选取同构 \(\bar\pi:Q/JQ\xrightarrow{\sim}M/JM\)。\(Q\) 的投射性使复合 \(Q\to Q/JQ\xrightarrow{\bar\pi}M/JM\) 沿商映射 \(M\to M/JM\) 提升为 \(\pi:Q\to M\)。于是 \(M=\pi(Q)+JM\)，余核 \(C=M/\pi(Q)\) 满足 \(JC=C\)，由同一 Nakayama 结论得到 \(C=0\)。诱导的顶映射为 \(\bar\pi\)，故 \(\ker\pi\subseteq JQ\)，由\cref{b5:prop:projective-cover-top}得到 \(M\) 的投射覆盖。若 \(M\) 投射，这个覆盖分裂，而多余的核在分裂中只能为零，所以 \(M\) 就是这些 \(P_i\) 的直和。
\ProseParagraphBreak
最后 \(\operatorname{Hom}_A(P_i,S_j)\) 的维数为 \(\delta_{ij}\)，因为所有映射都通过简单顶，再用分裂假设下的 Schur 引理。函子 \(\operatorname{Hom}_A(P_i,-)\) 正合，沿 \(M\) 的合成列逐步取维数，便得到重数公式。
\end{proof}
\begin{definition}\label{b5:def:modular-cartan-matrix}
设 \(k\) 为域，\(A\) 为有限维含幺 \(k\) 代数，且 \(A/J(A)\cong\prod_{i=1}^sM_{n_i}(k)\)，其中 \(s,n_i\) 为正整数。取全部简单左 \(A\)-模的不同构代表 \(S_i\)，以及它们的有限维投射覆盖 \(P_i\to S_i\)，以 \([P_j:S_i]\) 表示合成因子的重数，记
\[
c_{ij}=[P_j:S_i]=\dim_k\operatorname{Hom}_A(P_i,P_j).
\]
矩阵 \(C_A=(c_{ij})\) 称为 \(A\) 的\term[Cartan-juzhen-daishu]{Cartan 矩阵}[Cartan matrix of an algebra]。本章按“行对应简单模，列对应投射覆盖”排列；它不是后来由根系定义的另一种 Cartan 矩阵。
\end{definition}
在循环 \(p^r\) 群的例子中，唯一的 \(P_1\) 是 \(kG=M_{p^r}\)，故 Cartan 矩阵为一阶矩阵 \((p^r)\)。这个数字记录合成长度，不是不可分解投射模的种类数。

\subsection{块分解的基本概念}
\label{b5:subsec:0802:03}

\begin{definition}\label{b5:def:algebra-block}
设 \(k\) 为域，\(A\) 为有限维含幺 \(k\) 代数，\(b\in Z(A)\) 为非零幂等元。若 \(b\) 不能再写成两个非零正交中心幂等元之和，就称 \(b\) 为块幂等元，称以 \(b\) 为单位的代数 \(Ab\) 为 \(A\) 的一个\term[kuai]{块}[block]。
\end{definition}
反复分解中心幂等元，维数下降使过程终止，得到
\[
1=b_1+\cdots+b_\ell,\qquad A\cong\prod_{\nu=1}^{\ell}Ab_\nu.
\]
这个分解唯一到重排：若另有 \(1=\sum c_\mu\)，则
\(b_\nu=\sum_\mu b_\nu c_\mu\) 是正交中心幂等分解，本原性使恰有一项非零，再对 \(c_\mu\) 同样分析，就使两组块对应相等。任何模都有 \(M=\bigoplus_\nu b_\nu M\)，不可分解模只属于其中一个块。不同块之间的模同态为零，扩张也分裂，因为对扩张中项分别乘各 \(b_\nu\) 即给出分解。
\begin{example}\label{b5:exm:s3-two-blocks}
设 \(k\) 为特征二的域，\(G=S_3=\langle r,s\mid r^3=s^2=1,\ srs=r^{-1}\rangle\)。令
\[
e=1+r+r^2,\qquad f=1-e.
\]
有 \(e^2=e\)，且 \(e\) 中心，故 \(f\) 也是中心幂等元。理想 \(kGe\) 以 \(e,se\) 为基，乘法由 \((se)^2=e\) 给出，所以它同构于 \(kC_2\)，是二维局部块。
\end{example}
另一块是矩阵代数。取二维表示
\[
R=\begin{pmatrix}0&1\\1&1\end{pmatrix},\qquad
S=\begin{pmatrix}0&1\\1&0\end{pmatrix}.
\]
有 \(R^3=S^2=I\)、\(SRS=R^{-1}\)，故它们定义表示映射 \(\rho:kG\to M_2(k)\)。又有 \(I+R+R^2=0\)，所以 \(\rho(e)=0\)、\(\rho(f)=I\)，而 \(f,rf,sf,rsf\) 的像分别为 \(I,R,S,RS\)。这四个矩阵线性无关：若
\[
aI+bR+cS+dRS
=\begin{pmatrix}a+d&b+c\\b+c+d&a+b+d\end{pmatrix}=0,
\]
比较四个条目便得 \(d=a=b=c=0\)。因此 \(\rho(kGf)\) 的维数为四，等于 \(\dim_kM_2(k)\)，限制映射 \(\rho:kGf\to M_2(k)\) 满射。中心幂等分解给出 \(kG=kGe\oplus kGf\)，前面已由基 \(e,se\) 得到 \(\dim_k(kGe)=2\)，故
\[
\dim_k(kGf)=6-\dim_k(kGe)=4.
\]
满射两侧维数相等，所以 \(kGf\cong M_2(k)\)。于是恰有两个简单模，分别为平凡模 \(S_1\) 和上述二维模 \(S_2\)。后者已经投射，前者的覆盖为 \(kGe\)，因此
\[
C_{kS_3}=\begin{pmatrix}2&0\\0&1\end{pmatrix}.
\]
这里 \(p\,\bigm\vert\,|S_3|\)，却仍有一个半单矩阵块。模特征使整个群代数不半单，并不要求每个块都非半单。

% !TeX root = ../../Book5.tex
% 纲要标题：### 8.3 Brauer 特征标
% 纲要位置：工作记录/计划纲要.md，第 3995 行。

\section{Brauer 特征标}
\label{b5:sec:0803}


% 纲要范围：
%
% * p-正则元素
% * 普通特征标与 Brauer 特征标的区别
% * 分解映射与分解数
%

正特征中的矩阵迹会把整数重数约化模 \(p\)，所以丢掉一部分信息。Brauer 特征标改为先提升适当的单位根，再在特征零中相加。为了让这种提升与普通表示的约化对应，必须同时说明基域、赋值环和整数格；只说“把矩阵模 \(p\)”并不足以定义分解数。

\subsection{\texorpdfstring{\(p\)-正则元素}{p-正则元素}}
\label{b5:subsec:0803:01}

\begin{definition}\label{b5:def:p-regular-element}
设 \(G\) 为有限群，\(p\) 为素数。若 \(g\in G\) 的阶与 \(p\) 互素，就称 \(g\) 为\term[p-zhengze-yuansu]{\(p\)-正则元素}[\(p\)-regular element]。全体此类元素记为 \(G_{p'}\)；其余元素称为 \(p\)-奇异元素。这个集合在共轭下稳定，但通常不是子群。
\end{definition}
设 \(k\) 的特征为 \(p\)，\(g\) 的阶为 \(m\) 且 \(p\nmid m\)。在任一表示上，\(g\) 的最小多项式整除 \(X^m-1\)；后者与导数 \(mX^{m-1}\) 互素，所以在分裂域上 \(g\) 可对角化，特征值都是阶与 \(p\) 互素的单位根。这正是定义 Brauer 特征标时选择这些元素的原因。
\begin{lemma}\label{b5:lem:modular-trace-regular-part}
设 \(G\) 为有限群，\(p\) 为素数，\(k\) 为特征 \(p\) 的域，\(V\) 为有限维左 \(kG\)-模。每个 \(g\in G\) 唯一写成 \(g=g_pg_{p'}=g_{p'}g_p\)，其中两因子都属于 \(\langle g\rangle\)，前者的阶为 \(p\) 的幂，后者的阶与 \(p\) 互素。并且
\[
\operatorname{tr}_k(g|V)=\operatorname{tr}_k(g_{p'}|V).
\]
\end{lemma}
\begin{proof}
在阶为 \(p^a m\)、\(p\nmid m\) 的循环群中，中国剩余定理将指数分成模 \(p^a\) 与模 \(m\) 的两部分，给出存在与唯一性。扩到代数闭包不改变迹。此时 \(g_{p'}\) 对角化，而 \(g_p\) 与它交换，故保持其每个特征空间。在这些空间上，\((g_p-I)^{p^a}=0\)，所以 \(g_p\) 的特征值全为一。若某特征空间上 \(g_{p'}\) 的值为 \(\lambda\)，则 \(g\) 在该空间的迹为 \(\lambda\) 乘该空间维数，与 \(g_{p'}\) 相同。相加即得。
\end{proof}
这里的等式只比较 \(k\) 中的迹，并没有恢复被模 \(p\) 丢掉的整数。比如 \(p\) 个平凡模的直和，在所有群元素上的 \(k\) 值迹都是零，却当然不是零模。

\subsection{普通特征标与 Brauer 特征标}
\label{b5:subsec:0803:02}

\begin{definition}\label{b5:def:p-modular-system}
设 \(p\) 为素数，\(\mathcal O\) 为离散赋值环，其分式域 \(K\) 的特征为零，唯一非零极大理想为 \((\pi)\)，剩余域 \(k=\mathcal O/(\pi)\) 的特征为 \(p\)。三元组 \((K,\mathcal O,k)\) 称为一个\term[p-mo-xitong]{\(p\)-模系统}[\(p\)-modular system]。对有限群 \(G\)，若 \(KG\) 与 \(kG\) 的半单商均为相应基域上全矩阵代数的直积，就称它为 \(G\) 的分裂 \(p\)-模系统。
\end{definition}
这里用到的环性质是：\(\mathcal O\) 是主理想整环，有限生成无挠模自由，每个非零元素等于单位乘以某个 \(\pi\) 的幂。本节把这样的系统作为给定数据。
\ProseParagraphBreak
为讨论分解矩阵，以下使用分裂 \(p\)-模系统，并进一步假设 \(K\) 含有所有 \(|G|\) 次单位根。分裂域与模系统的准备可参阅\cite[Section 9.2, Theorem 9.2.6; Section 9.4]{Webb2016FiniteGroupRepresentations}。分裂条件是这些证明的前提，任意有限域未必满足它。
\begin{lemma}\label{b5:lem:prime-to-p-root-lift}
设 \(G\) 为有限群，\((K,\mathcal O,k)\) 为 \(p\)-模系统，且 \(K\) 含有全部 \(|G|\) 次单位根。若正整数 \(m\,\bigm\vert\,|G|\) 且 \(p\nmid m\)，则约化映射给出 \(K\) 中 \(m\) 次单位根与 \(k\) 中 \(m\) 次单位根的群同构。记相应唯一提升为 \(\lambda\mapsto\widehat\lambda\)。
\end{lemma}
\begin{proof}
单位根的赋值为零，所以都在 \(\mathcal O^\times\) 中。多项式 \(X^m-1\) 在 \(K\) 中分成 \(m\) 个一次因子，模 \(\pi\) 后仍是这些因子的乘积；而在特征 \(p\) 中它仍可分，所以这些根的约化两两不同，并给出全部根。约化保持乘法，故得到群同构。
\end{proof}
这些提升落在特征零的分圆域中。固定一个将该分圆域嵌入 \(\mathbb C\) 的映射，以下把提升后的单位根看成复单位根。选择一旦固定，所有表示和张量运算都使用同一个提升。
\begin{definition}\label{b5:def:brauer-character}
设 \(G\) 为有限群，\(p\) 为素数，\((K,\mathcal O,k)\) 为分裂 \(p\)-模系统，\(K\) 含有全部 \(|G|\) 次单位根，且已固定这些单位根到复数的对应。记 \(G_{p'}\) 为 \(p\)-正则元素集合，\(\lambda\mapsto\widehat\lambda\) 为约化映射 \(\mathcal O\to k\) 所确定的阶与 \(p\) 互素单位根的唯一提升，再按固定对应看成复单位根。对有限维左 \(kG\)-模 \(M\)，令 \(d=\dim_kM\)。对 \(g\in G_{p'}\)，将 \(g\) 在 \(M\) 上的特征值按重数写为 \(\lambda_1,\ldots,\lambda_d\)。函数 \(\varphi_M:G_{p'}\to\mathbb C\)，由
\[
\varphi_M(g)=\widehat\lambda_1+\cdots+\widehat\lambda_d
\]
确定，称为 \(M\) 的\term[Brauer-tezhengbiao]{Brauer 特征标}[Brauer character]。其定义域只取 \(G_{p'}\)，值在特征零中；简单左 \(kG\)-模的 Brauer 特征标称为不可约 Brauer 特征标。
\end{definition}
\begin{proposition}\label{b5:prop:brauer-character-properties}
设 \(G\) 为有限群，\(p\) 为素数，\((K,\mathcal O,k)\) 为分裂 \(p\)-模系统，\(K\) 含有全部 \(|G|\) 次单位根，并固定这些根到复数的对应。设 \(L,M,N\) 为有限维左 \(kG\)-模，\(S_i\) 为全部简单左 \(kG\)-模的不同构代表，\(\varphi_X\) 表示模 \(X\) 的 Brauer 特征标，\([M:S_i]\) 表示合成因子的重数。这些特征标在 \(p\)-正则共轭类上为常数，且对 \(p\)-正则元素 \(g\in G\) 有
\[
\varphi_M(1)=\dim_kM,\qquad
\varphi_{M\otimes N}=\varphi_M\varphi_N,\qquad
\varphi_{M^*}(g)=\varphi_M(g^{-1})=\overline{\varphi_M(g)}.
\]
对短正合列 \(0\to L\to M\to N\to0\)，有 \(\varphi_M=\varphi_L+\varphi_N\)。因而
\[
\varphi_M=\sum_i[M:S_i]\varphi_{S_i}.
\]
\end{proposition}
\begin{proof}
共轭矩阵有相同的特征值，单位元的全部特征值为一，得到前两个直接性质。张量积的特征值为所有两两乘积；单位根提升是群同态，故其提升也两两相乘，给出乘法公式。对偶把特征值取逆，而复单位根的逆为其共轭，得到对偶公式。
将一个 \(p\)-正则元素生成的循环群限制到给定正合列。因为该循环群的阶可逆，Maschke 定理使列在这个子群上分裂，所以中项的特征值多重集为两端的并。逐个提升后相加便得加法性；沿合成列归纳得到最后一式。
\end{proof}
在特征二的二维 \(S_3\) 模中，三循环的矩阵为
\(\begin{psmallmatrix}0&1\\1&1\end{psmallmatrix}\)，其 \(k\) 值迹为一。然而其两个特征值是非平凡三次单位根，提升到复数后之和为 \(-1\)。因此 Brauer 特征标的值是 \(-1\)，不能把域中的迹“一”随意提升成复数一。这些定义和性质可对照\cite[Section 10.1, Lemma 10.1.1; Proposition 10.1.3]{Webb2016FiniteGroupRepresentations}。

\subsection{分解映射与分解数}
\label{b5:subsec:0803:03}

\begin{definition}\label{b5:def:stable-integral-lattice}
设 \(G\) 为有限群，\(p\) 为素数，\((K,\mathcal O,k)\) 为 \(p\)-模系统，\((\pi)\) 为 \(\mathcal O\) 的极大理想，\(V\) 为有限维左 \(KG\)-模。若 \(L\subseteq V\) 是被 \(G\) 保持的有限自由 \(\mathcal O\) 子模，且其一组 \(\mathcal O\) 基也是 \(V\) 的 \(K\) 基，就称 \(L\) 为 \(V\) 中的一个满的 \(G\)-稳定\term[zhenggeshu-ge]{格}[lattice]。商 \(L/\pi L\) 是一个有限维左 \(kG\)-模，称为该格的约化。
\end{definition}
\begin{theorem}[Brauer--Nesbitt]\label{b5:thm:lattice-reduction-independence}
设 \(G\) 为有限群，\(p\) 为素数，\((K,\mathcal O,k)\) 为 \(p\)-模系统，\((\pi)\) 为 \(\mathcal O\) 的极大理想，\(V\) 为有限维左 \(KG\)-模。则 \(V\) 有满的 \(G\)-稳定格；任意两个这种格 \(L_1,L_2\) 的约化 \(L_1/\pi L_1\)、\(L_2/\pi L_2\) 具有相同的简单左 \(kG\)-模合成因子及重数，但不必同构。
\end{theorem}
\begin{proof}
先任取 \(K\) 基 \(v_1,\ldots,v_d\)，令 \(L\) 为所有 \(gv_i\) 的 \(\mathcal O\) 线性生成空间。它有限生成、无挠且被 \(G\) 保持；由主理想整环上的模结构，它自由，清分母可知其秩恰为 \(d\)，所以是满格。
两满格可通约：分别把一组格基用另一组表示，给有限多个系数清去 \(\pi\) 的负幂，就能把某个 \(\pi^aL_1\) 包含于 \(L_2\)，再夹在 \(L_2\) 的某个 \(\pi\) 次幂之上。乘以 \(\pi^a\) 给出原格与缩放格的 \(\mathcal O G\) 模同构，所以只需比较 \(L_1\subseteq L_2\)。
商 \(L_2/L_1\) 有有限长度，可以沿合成列在两格间插入有限多个稳定格，归约到 \(L_2/L_1\) 简单的情形。因为这个商被某个 \(\pi\) 次幂零化，而 \(\pi\) 在简单商上的乘法若非零就满射，故只能被 \(\pi\) 本身零化。因此 \(\pi L_2\subseteq L_1\)，得到链
\[
L_2\supseteq L_1\supseteq\pi L_2\supseteq\pi L_1.
\]
两个约化的合成因子共有 \(L_1/\pi L_2\) 这一部分，剩余的商分别为 \(L_2/L_1\) 和 \(\pi L_2/\pi L_1\)。乘以 \(\pi\) 给出这两个商的模同构：满射显然，核等式由格无挠得到。所以两约化的合成因子多重集相同；沿插入的有限链传递即得一般结论。
\end{proof}
这条格无关性称为 Brauer--Nesbitt 定理\footnote{内斯比特（Cecil James Nesbitt，1912-2001），美国数学家，与布劳尔合作研究模表示论；布劳尔已见本卷前注。}的格约化形式，见\cite[Theorem 9.4.8]{Webb2016FiniteGroupRepresentations}。必须保留“合成因子”这一表述：取 \(G=C_2\)、\(\mathcal O=\mathbb Z_{(2)}\)、\(V=\mathbb QG\)，格
\(\mathcal O\{1,s\}\) 的约化是不可分解的正则模，格
\(\mathcal O(1+s)\oplus\mathcal O(1-s)\) 的约化却是两个平凡模的直和。
\begin{definition}\label{b5:def:decomposition-map}
设 \(G\) 为有限群，\(p\) 为素数，\((K,\mathcal O,k)\) 为分裂 \(p\)-模系统，\((\pi)\) 为 \(\mathcal O\) 的极大理想。令 \(G_0(kG)\) 为以有限维左 \(kG\)-模的同构类生成、并对每个短正合列 \(0\to L\to M\to N\to0\) 加入 \([M]=[L]+[N]\) 的交换群。Jordan--Hölder 定理\footnote{若尔当（Camille Jordan，1838-1922），法国数学家，研究群论与代数方程；赫尔德（Otto Ludwig Hölder，1859-1937），德国数学家，研究群论与分析。}使它是以全部简单左 \(kG\)-模的不同构代表类 \([S_i]\) 为基的自由交换群。令 \(R_K(G)\) 为有限维左 \(KG\)-模的表示群。对每个这样的模 \(V\)，选择其中的满的 \(G\)-稳定格 \(L\)，则对应
\[
d:R_K(G)\longrightarrow G_0(kG),\qquad
[V]\longmapsto[L/\pi L]
\]
称为\term[fenjie-yingshe]{分解映射}[decomposition map]。取全部简单左 \(KG\)-模的不同构代表 \(V_\alpha\)，以及其中的满的 \(G\)-稳定格 \(L_\alpha\)，定义
\[
d_{\alpha i}=[L_\alpha/\pi L_\alpha:S_i].
\]
这些非负整数称为\term[fenjieshu]{分解数}[decomposition numbers]，矩阵 \(D=(d_{\alpha i})\) 称为分解矩阵，行对应特征零简单模，列对应特征 \(p\) 简单模。
\end{definition}
前一定理保证它与格的选择无关。直和格给出加法性；张量格给出乘法性，所以若两群都配张量乘积，\(d\) 还是环同态。这里 \(G_0\) 使用的是短正合关系，不是把非分裂扩张误认成实际直和。
\begin{proposition}\label{b5:prop:ordinary-brauer-restriction}
设 \(G\) 为有限群，\(p\) 为素数，\((K,\mathcal O,k)\) 为分裂 \(p\)-模系统，\((\pi)\) 为 \(\mathcal O\) 的极大理想，\(K\) 含有全部 \(|G|\) 次单位根，并固定这些根到复数的对应。设 \(V\) 为有限维左 \(KG\)-模，\(L\) 为其中的满的 \(G\)-稳定格，\(\chi_V\) 为按此对应得到的普通特征标，\(G_{p'}\) 为 \(p\)-正则元素集合，\(\varphi_M\) 表示左 \(kG\)-模 \(M\) 的 Brauer 特征标。取简单左 \(KG\)-模代表 \(V_\alpha\) 及其中的稳定满格 \(L_\alpha\)，记其普通特征标为 \(\chi_\alpha\)，取简单左 \(kG\)-模代表 \(S_i\)，记 \(d_{\alpha i}=[L_\alpha/\pi L_\alpha:S_i]\)。则
\[
\chi_V|_{G_{p'}}=\varphi_{L/\pi L},\qquad
\chi_\alpha|_{G_{p'}}=\sum_i d_{\alpha i}\varphi_{S_i}.
\]
\end{proposition}
\begin{proof}
设 \(g\in G_{p'}\) 的阶为 \(m\)。它在 \(L\) 上有系数位于 \(\mathcal O\) 的特征多项式，约化后就是在 \(L/\pi L\) 上的特征多项式。其根均为 \(m\) 次单位根，且这些单位根的约化互异；因而每个根及其重数在约化后对应到唯一的根及相同重数。把根再提升回来，就恢复 \(\chi_V(g)\)。第二式由 Brauer 特征标对合成因子的加法性得到。
\end{proof}
\begin{theorem}\label{b5:thm:brauer-character-basis}
设 \(G\) 为有限群，\(p\) 为素数，\((K,\mathcal O,k)\) 为分裂 \(p\)-模系统，\(K\) 含有全部 \(|G|\) 次单位根，并固定这些根到复数的对应，\(G_{p'}\) 为 \(p\)-正则元素集合。不可约 Brauer 特征标构成 \(G_{p'}\) 上复类函数空间的一组基。因此简单左 \(kG\)-模的同构类型数等于 \(p\)-正则共轭类数；两个有限维左 \(kG\)-模有相同 Brauer 特征标，当且仅当它们有相同的简单合成因子及重数。
\end{theorem}
\begin{proof}
先证线性无关。简单模 \(S_i\) 的 \(k\) 值迹作为 \(kG\) 上的线性函数两两独立：在
\(kG/J\cong\prod_i\operatorname{End}_k(S_i)\)
中，取第 \(i\) 块的一个对角矩阵单位，其各个简单迹为 \(\delta_{ij}\)，再提升到群代数即可检验独立性。群元素是 \(kG\) 的基，因此这些迹限制到 \(G\) 上仍独立。由\cref{b5:lem:modular-trace-regular-part}，所有值由 \(p\)-正则元素上的值确定，故相应矩阵在 \(k\) 中有一个非零的满行秩子式。
把这个子式中的各迹换成在 \(\mathcal O\) 中取值的 Brauer 特征标，约化后恢复原子式。因此提升后的行列式非零。在固定的分圆域嵌入下它仍非零，于是复值 Brauer 特征标线性无关。
再证生成性。先在分裂的特征零群代数 \(KG\) 上对 \(\operatorname{Hom}_K(V_\alpha,V_\beta)\) 的共轭作用取平均。平均算子的像是交织空间，分裂假设与 Schur 引理使其维数为 \(\delta_{\alpha\beta}\)，故取迹得到
\[
\frac1{|G|}\sum_{g\in G}\chi_\alpha(g^{-1})\chi_\beta(g)
=\delta_{\alpha\beta}.
\]
这些等式在固定的分圆域嵌入下仍成立，所以复值普通特征标线性无关。另一方面，\(Z(KG)\) 的共轭类和基说明其维数为共轭类数；半单分裂的矩阵块分解又说明该维数等于简单模数。因此这些特征标恰构成全部复类函数的一组基。任一 \(G_{p'}\) 上的类函数可以在其余共轭类补零，故普通不可约特征标的限制生成所需空间。\cref{b5:prop:ordinary-brauer-restriction}又把每一个这样的限制写成不可约 Brauer 特征标的线性组合，得到生成性。最后用\cref{b5:prop:brauer-character-properties}和刚证的独立性比较系数，得到合成因子的判别。
\end{proof}
对于特征二的 \(S_3\)，\cref{b5:exm:s3-two-blocks}已经给出两个简单模。两类 \(2\)-正则元素由 \(1,r\) 代表，故
\[
\begin{array}{c|rr}
&1&r\\ \hline
\varphi_{S_1}&1&1\\
\varphi_{S_2}&2&-1
\end{array}
\qquad
D=
\begin{pmatrix}
1&0\\1&0\\0&1
\end{pmatrix}.
\]
三行普通表示依次为平凡、符号和标准二维表示。前两者约化后相同；二维整系数标准表示约化后是前述简单模。由此直接计算
\(D^{\mathsf T}D=\operatorname{diag}(2,1)=C_{kS_3}\)。
这是本例的计算；一般的 Cartan--分解矩阵关系及投射模提升需要另作证明，可参阅\cite[Section 9.5, Corollary 9.5.6]{Webb2016FiniteGroupRepresentations}。

% !TeX root = ../../Book5.tex
% 纲要标题：### 8.4 模表示理论补充
% 纲要位置：工作记录/计划纲要.md，第 4001 行。

\section{模表示理论补充}
\label{b5:sec:0804}


% 纲要范围：
%
% * 本章不同于第5章 Brauer 诱导理论
% * 缺陷群和块论的学习地图
% * 不把特征零分类原样搬到正特征
%
% ---
%

本章的几个计算已经说明，特征零表示、模表示和投射模之间有联系，却不能彼此替换。本节用约化与诱导的相容性说明第一种联系，再用一个能够直接计算的截断映射说明块的缺陷群。更深的块分类还需要其他工具，这里先把问题和所需理论说明清楚。

\subsection{模表示与 Brauer 诱导理论的区别}
\label{b5:subsec:0804:01}

\chapref{b5:ch:05}的 Brauer 诱导定理在特征零表示环中，把普通特征标写成由某些子群的一维特征标诱导而来的整数线性组合。本章的 Brauer 特征标则在特征 \(p\) 中先取表示，再把 \(p\)-正则特征值提升到特征零；分解矩阵记录整格约化的合成因子。两种理论处理不同的构造，不能因为都冠以 Brauer 之名就将其结论互换。
\begin{proposition}\label{b5:prop:reduction-induction}
设 \(G\) 为有限群，\(H\le G\)，\(p\) 为素数，并给定对 \(G,H\) 都分裂的 \(p\)-模系统 \((K,\mathcal O,k)\)。记 \(d_G:R_K(G)\to G_0(kG)\)、\(d_H:R_K(H)\to G_0(kH)\) 为分解映射，其中表示群及 Grothendieck 群\footnote{格罗滕迪克（Alexandre Grothendieck，1928-2014），法籍数学家，重塑代数几何的概形与上同调方法。}都取有限维左模。它们与限制、诱导相容：
\[
d_G\operatorname{Ind}_H^G
=\operatorname{Ind}_H^Gd_H,\qquad
d_H\operatorname{Res}_H^G
=\operatorname{Res}_H^Gd_G.
\]
右侧的诱导、限制作用于由短正合列定义的 \(G_0\) 群。
\end{proposition}
\begin{proof}
限制一个稳定满格仍是稳定满格，其约化显然与限制交换。若 \(L\) 是 \(KH\)-模 \(V\) 中的满格，则
\(\mathcal OG\otimes_{\mathcal OH}L\)
是 \(\operatorname{Ind}_H^G V\) 中的稳定满格：按左陪集代表，\(\mathcal OG\) 是有限自由右 \(\mathcal OH\) 模，这个张量是若干份 \(L\) 的直和，且扩到 \(K\) 后得到所需诱导模。模 \(\pi\) 后自然同构于
\(kG\otimes_{kH}(L/\pi L)\)。
诱导因右侧自由性保持正合列，限制也保持正合列，所以二者确实诱导 \(G_0\) 上的映射。结合格选择无关性得到公式。
\end{proof}
因此可以先诱导再约化，也可以先约化再诱导，得到相同的合成因子信息。但这并不把任意模表示恢复为完全可约表示；\(G_0\) 本来就已经忘掉了非分裂扩张的区别。

\subsection{缺陷群与块论简介}
\label{b5:subsec:0804:02}

设 \(G\) 为有限群，\(p\) 为素数，\(k\) 为特征 \(p\) 的域，\(P\le G\) 为 \(p\)-子群。先把 \(P\) 的共轭作用当作置换作用来看。固定空间 \((kG)^P\) 以各个共轭轨道的轨道和为基，其中单点轨道恰好是 \(C_G(P)\) 中的元素。若一个轨道含 \(g\)，令 \(Q=C_P(g)\) 为它在 \(P\) 中的稳定子群，则其轨道和为
\[
\sum_{x\in P/Q}xgx^{-1}.
\]
这里 \(g\in(kG)^Q\)。对任意子群 \(Q\le P\)，相对迹
\[
\operatorname{Tr}_Q^P:(kG)^Q\longrightarrow(kG)^P,\qquad
z\longmapsto\sum_{x\in P/Q}xzx^{-1}
\]
与代表元选择无关，因为 \(z\) 被 \(Q\) 固定；所得和被 \(P\) 固定。非单点轨道的稳定子群 \(Q\) 是真子群，所以该轨道和正是从较小子群来的相对迹。

\ProseParagraphBreak

反过来，若 \(Q<P\)，则 \(\operatorname{Tr}_Q^P(z)\) 在每个 \(h\in C_G(P)\) 上的系数为 \([P:Q]\) 乘以 \(z\) 在 \(h\) 上的系数；这个指数是正的 \(p\) 次幂，在 \(k\) 中为零。因此保留单点轨道、删去其余轨道的操作，其核恰为全部真子群相对迹的像之和。它保留的是固定空间中不能由这些较小子群的相对迹表示的部分：
\[
(kG)^P\big/\sum_{Q<P}
\operatorname{Tr}_Q^P\bigl((kG)^Q\bigr)
\cong kC_G(P)
\]
作为 \(k\) 向量空间同构。下面还要验证，这个系数截断与乘法相容。

\begin{proposition}\label{b5:prop:brauer-truncation}
设 \(G\) 为有限群，\(p\) 为素数，\(k\) 为特征 \(p\) 的域，\(P\le G\) 为 \(p\)-子群，\(C_G(P)\) 为其中心化子。令 \(P\) 以共轭作用于 \(kG\)，并记固定子代数为 \((kG)^P\)。映射
\[
\operatorname{Br}_P^G:(kG)^P\longrightarrow kC_G(P),\qquad
\sum_{g\in G}a_gg\longmapsto\sum_{g\in C_G(P)}a_gg
\]
是保持单位的 \(k\) 代数同态，称为此处的\term[Brauer-yingshe]{Brauer 映射}[Brauer map]。
\end{proposition}
\begin{proof}
加法、标量与单位的保持直接来自定义。取两个 \(P\) 共轭不变的元素 \(x=\sum a_gg\)、\(y=\sum b_hh\)，固定 \(z\in C_G(P)\)。乘积中 \(z\) 的系数为
\[
\sum_{gh=z}a_gb_h.
\]
\(P\) 同时共轭二元组 \((g,h)\)，保持条件 \(gh=z\)，且每个轨道内的系数乘积相同。非平凡轨道的大小是正的 \(p\) 次幂，在 \(k\) 中贡献为零。剩余的固定二元组恰好满足 \(g,h\in C_G(P)\)，其总贡献就是 \(\operatorname{Br}_P^G(x)\operatorname{Br}_P^G(y)\) 的 \(z\) 系数。逐个 \(z\) 比较即得乘法性。
\end{proof}
\begin{definition}\label{b5:def:block-defect-group}
设 \(G\) 为有限群，\(p\) 为素数，\(k\) 为特征 \(p\) 的域，\(b\) 为 \(kG\) 的块幂等元，对 \(p\)-子群 \(D\le G\)，记 \(\operatorname{Br}_D^G:(kG)^D\to kC_G(D)\) 为保留中心化子上系数的 Brauer 映射。若 \(D\) 满足 \(\operatorname{Br}_D^G(b)\ne0\)，并且在满足这一条件的 \(p\)-子群中按包含关系极大，就称 \(D\) 为块 \(kGb\) 的一个\term[quexian-qun]{缺陷群}[defect group]。
\end{definition}
\symidx{\(\operatorname{Br}_D^G\)：群代数的 Brauer 映射}
由于 \(\operatorname{Br}_{1}^G(b)=b\ne0\)，且 \(G\) 有限，这样的极大子群存在。其进一步的共轭与结构性质属于块论；这里采用支集截断的描述以便计算。这个描述及其与一般定义的关系，可参阅\cite[Corollaries 12.5.4--12.5.5]{Webb2016FiniteGroupRepresentations}。
\ProseParagraphBreak
对\cref{b5:exm:s3-two-blocks}中的 \(S_3\) 与特征二，取 \(P=\langle s\rangle\)。中心化子 \(C_G(P)=P\)，故
\[
\operatorname{Br}_P^G(1+r+r^2)=1,\qquad
\operatorname{Br}_P^G(r+r^2)=0.
\]
三个非平凡 \(2\) 子群均为这些共轭子群，因此二维局部块的缺陷群为一个 Sylow \(2\) 子群，四维矩阵块的缺陷群为平凡群。后者称为缺陷零块。在任意有限 \(p\) 群上，唯一块的幂等元为一，而对整个群的截断仍为一，所以其缺陷群就是整个群。

\subsection{特征零与正特征下的表示分类}
\label{b5:subsec:0804:03}

不可约 Brauer 特征标分类的是简单模，并由加法性确定任意模的合成因子。对于扩张信息，它有明确局限：循环 \(p\) 群上，\(M_d\) 与 \(k^{\oplus d}\) 的 Brauer 特征标都只在单位元处取整数 \(d\)，但前者不可分解，后者在 \(d>1\) 时可分解。分解矩阵同样保留约化后的合成因子，不指定某个整格约化的全部模结构。
\ProseParagraphBreak
有限群在正特征中甚至可以有无限多个不可分解表示。这里给一个不用分类定理的例子。设 \(k\) 是特征二的无限域，\(G=C_2\times C_2=\langle x,y\rangle\)，并在 \(k^2\) 上取
\[
N=\begin{pmatrix}0&1\\0&0\end{pmatrix},\qquad
\rho_\lambda(x)=I+N,\qquad \rho_\lambda(y)=I+\lambda N
\quad(\lambda\in k).
\]
因为 \(N^2=0\)，两个矩阵交换且平方都为一，故给出表示。其自同态必须与 \(N\) 交换，恰为 \(aI+bN\)，组成局部环，所以每个表示不可分解。若两个参数的表示同构，同构矩阵 \(T\) 与 \(N\) 交换，再由
\(T(\lambda N)=\mu NT\)
得到 \((\lambda-\mu)NT=0\)。\(T\) 可逆且 \(N\ne0\)，故 \(\lambda=\mu\)。于是得到一个参数化的无限族。
\ProseParagraphBreak
这个例子解释了为何后续会研究投射覆盖、块、相对投射性和稳定范畴，而不只寻找另一张有限的特征标表。一般缺陷群如何控制块、何时只有有限多种不可分解模，以及不同块之间的等价，都是进一步的问题。本编接下来改从箭图研究有限维代数；那里的表示可以在任意特征上定义，所用的无环性、关系和代数闭性条件仍须逐项写明。


\begin{chaptersummary}
在模特征下，群代数的根和模的根层保留了非半单部分的信息。循环 \(p^r\) 群的群代数是 \(k[t]/(t^{p^r})\)，其不可分解模由幂零链的长度区分；它们虽然只有一种简单合成因子，直和结构仍不相同。投射覆盖按简单顶排列，Cartan 矩阵记录投射模中的简单重数，块则把整个模范畴分成彼此没有同态和扩张的部分。

\ProseParagraphBreak

Brauer 特征标只在 \(p\)-正则元素上定义，把剩余域中的特征值提升到特征零后求和。这样避免了 \(k\) 值迹把整数重数约化模 \(p\) 的损失；在本章的分裂条件下，整个 Brauer 特征标函数确定简单合成因子及其整数重数，但不能区分具有相同合成因子的不同扩张。分解数描述普通不可约表示约化后的简单成分；块的缺陷群则借助 Brauer 映射把块与 \(p\) 子群联系起来。
\end{chaptersummary}

\BookExercises

\begin{optionalexercise}\label{b5:ex:08-jordan-data}
设 \(\operatorname{char}k=2\)，\(C_8=\langle s\rangle\) 在十维空间 \(V\) 上作用。令 \(T=s-1\)，已知
\[
\dim\ker T^j=3,5,7,8,9,10\qquad(j=1,\ldots,6).
\]
求不可分解直和分解，以及顶、基座、不动空间和 \(T^5\) 的秩。
\end{optionalexercise}
\begin{solution}
相邻核维数之差为 \(3,2,2,1,1,1\)，分别给出长度至少为 \(1,\ldots,6\) 的幂零块数。故块长为 \(6,3,1\)，
\(V\cong M_6\oplus M_3\oplus M_1\)。
三个直和项各有一维顶和一维基座，故两者维数均为 \(3\)。不动空间为 \(\ker(s-1)=\ker T\)，也为三维。最后
\(\operatorname{rank}T^5=10-9=1\)。只有简单合成因子及其总数不能恢复这些块长，而各次核维数可以。
\end{solution}

\begin{optionalexercise}\label{b5:ex:08-norm-invariants}
设 \(G=C_m\)，其中 \(m=p^r\)，\(\operatorname{char}k=p\)，\(t=s-1\)。对 \(M_d=k[t]/(t^d)\)、\(1\leq d\leq m\)，计算不动空间、余不动空间 \(M_d/tM_d\) 及两者之间的自然映射。再证明群的求和算子 \(\sum_{j=0}^{m-1}s^j\) 等于乘 \(t^{m-1}\)，并计算它在这些模上的像。
\end{optionalexercise}
\begin{solution}
不动空间是 \(kt^{d-1}\)，余不动空间由 \(1\) 的类生成。自然映射由包含不动空间再取商给出：\(d=1\) 时为同构，\(d>1\) 时因 \(t^{d-1}\in tM_d\) 而为零。
在多项式环中，
\[
\sum_{j=0}^{m-1}(1+t)^j
=\frac{(1+t)^m-1}{t}=t^{m-1},
\]
其中商式是几何级数恒等式，而特征 \(p\) 及 \(m=p^r\) 给出 \((1+t)^m=1+t^m\)。故在 \(d<m\) 时求和算子为零；在 \(d=m\) 时，其像为 \(kt^{m-1}\)，并诱导从一维余不动空间到一维不动空间的同构。
\end{solution}

\begin{optionalexercise}\label{b5:ex:08-elementary-abelian-radical}
设 \(P=C_p\times C_p=\langle s,u\rangle\)，\(\operatorname{char}k=p\)。证明
\[
kP\cong k[x,y]/(x^p,y^p),\qquad x=s-1,\ y=u-1.
\]
求根 \(J\)、每层 \(J^i/J^{i+1}\) 的维数、最小的 \(N\) 使 \(J^N=0\)，以及正则模的基座。
\end{optionalexercise}
\begin{solution}
两生成元交换，且 \(x^p=y^p=0\)。所给商以 \(x^a y^b\)、\(0\leq a,b<p\) 为基，共 \(p^2\) 个；到 \(kP\) 的映射满射且维数相等，所以同构。
根为 \(J=(x,y)\)，第 \(i\) 层以全次数为 \(i\) 的上述单项式为基。因此层维数依次为
\[
\dim J^i/J^{i+1}=
\begin{cases}
i+1,&0\leq i\leq p-1,\\
2p-1-i,&p\leq i\leq2p-2,\\
0,&i\geq2p-1.
\end{cases}
\]
最高非零单项式为 \(x^{p-1}y^{p-1}\)，故 \(J^{2p-2}\neq0\)、\(J^{2p-1}=0\)。一个元素同时被 \(x,y\) 零化，当且仅当只含这个最高单项式；因此基座为 \(kx^{p-1}y^{p-1}\)。
\end{solution}

\begin{optionalexercise}\label{b5:ex:08-elementary-abelian-family}
设 \(k\) 为特征 \(p\) 的无限域。记 \(V_\lambda\) 为 \(C_p\times C_p\) 的二维表示，两生成元分别作用为 \(I+N\)、\(I+\lambda N\)，其中 \(N^2=0,N\neq0\)。求任意参数 \(\lambda,\mu\) 之间的同态空间及非零同态的秩。由此说明不同构的不可分解模之间也可以有非零同态。
\end{optionalexercise}
\begin{solution}
两生成元的矩阵交换且 \(p\) 次幂为恒等，故确为表示。取基使
\(N=\left(\begin{smallmatrix}0&1\\0&0\end{smallmatrix}\right)\)。
同态矩阵先须与 \(N\) 交换，因而为 \(f=aI+bN\)；第二生成元的条件为
\((\lambda-\mu)fN=0\)，而 \(fN=aN\)。
所以当 \(\lambda=\mu\) 时，同态空间为 \(kI+kN\)，其中 \(a\neq0\) 的映射可逆、\(a=0,b\neq0\) 的映射秩为一。当 \(\lambda\neq\mu\) 时，同态空间仅为 \(kN\)，所有非零映射秩为一，且两个这样的跨参数映射的复合为零。
每个自同态环同构于局部环 \(k[t]/(t^2)\)，所以模不可分解；不同参数之间没有可逆同态，却仍有一维同态空间。非半单情形下，不能把简单模的 Schur 引理直接改成不可分解模的断言。
\end{solution}

\begin{optionalexercise}\label{b5:ex:08-hom-through-projectives}
设 \(p\) 为素数，\(k\) 为特征 \(p\) 的域，\(r\ge1\) 为整数，\(m=p^r\)、\(A=kC_{p^r}\cong k[t]/(t^m)\)、\(M_d=A/(t^d)\)，其中 \(1\le d\le m\)。给定 \(1\le a,b\le m\)，
在 \(\operatorname{Hom}_A(M_a,M_b)\) 中，求可经过有限维投射模分解的映射组成的子空间及其维数，再求模掉该子空间后的维数。
\end{optionalexercise}
\begin{solution}
由\cref{b5:thm:split-projective-covers}及唯一简单模的投射覆盖为 \(A\)，每个有限维投射模都是有限个 \(A\) 的直和。因此只须求经过 \(A\) 的映射之和。
映射 \(M_a\to A\) 把 \(1\) 送到被 \(t^a\) 零化的元素，即 \(t^{m-a}A\) 中的元素；随后映到 \(M_b\)，所得 \(1\) 的像恰好遍历 \(t^{m-a}M_b\)。这空间的维数为 \(\max(0,a+b-m)\)，且已对加法封闭，所以就是全部可经投射模分解的映射。全部 Hom 的维数为 \(\min(a,b)\)，故商空间维数为
\[
\min(a,b)-\max(0,a+b-m)=\min(a,b,m-a,m-b).
\]
这也直接显示，当任一端为 \(M_m=A\) 时，商空间为零。
\end{solution}

\begin{optionalexercise}\label{b5:ex:08-cyclic-extension-parameters}
设 \(p\) 为素数，\(k\) 为特征 \(p\) 的域，\(r\ge1\) 为整数，\(m=p^r\)、\(A=k[t]/(t^m)\)，并对 \(1\le d\le m\) 令左模 \(M_d=A/(t^d)\)。给定 \(1\le a,b\le m\)，按固定子模与商模的等价方式分类左 \(A\)-模的短正合列
\(0\to M_b\to E\to M_a\to0\)。
证明扩张类可由向量空间
\[
\ker(t^{m-a}:M_b\to M_b)/t^aM_b
\]
参数化，并求这个空间的维数。证明中直接使用生成元，不要求先调用消解计算。
\end{optionalexercise}
\begin{solution}
把子模生成元记为 \(u\)，取商模生成元 \(1\in M_a\) 的提升 \(v\in E\)。因为 \(t^a\) 在商上为零，有 \(t^av=w\in M_b\)。又因 \(t^mE=0\)，必有 \(t^{m-a}w=0\)。
反过来，给定满足此条件的 \(w\)，在向量空间
\(M_b\oplus\operatorname{span}_k(v,tv,\ldots,t^{a-1}v)\)
上定义 \(t\) 的作用：在 \(M_b\) 上用原作用，沿第二组基逐次后移，并把最后一项送到 \(w\)。条件保证 \(t^m=0\)，所以得到所需扩张，且子模的嵌入不发生额外退化。
改变提升为 \(v+z\)、\(z\in M_b\)，把 \(w\) 改为 \(w+t^az\)。固定两端的任何扩张同构都只能这样改变商生成元的提升，因此等价类恰为所列商空间；参数零恰好对应可以取 \(t^av=0\) 的分裂扩张。
核的维数为 \(\min(m-a,b)\)，而 \(t^aM_b\) 的维数为 \(\max(0,b-a)\)。相减得到
\[
\min(m-a,b)-\max(0,b-a)=\min(a,b,m-a,m-b).
\]
\end{solution}

\begin{optionalexercise}\label{b5:ex:08-c4-extension-middle}
设 \(k\) 为特征 \(2\) 的域，\(A=kC_4\cong k[t]/(t^4)\)，\(M_d=A/(t^d)\)、\(1\le d\le4\)。
对左 \(A\)-模短正合列 \(0\to M_2\to E\to M_2\to0\)，记指定子模的生成元为 \(u\)、商模生成元的提升为 \(v\)，把扩张参数写成
\(t^2v=\alpha u+\beta tu\)。分别讨论 \(\alpha\neq0\)、\(\alpha=0,\beta\neq0\)、\(\alpha=\beta=0\) 时中间项的不可分解直和类型。
\end{optionalexercise}
\begin{solution}
上一题中 \(m=4,a=b=2\)，整个 \(M_2\) 都可作为参数，且 \(t^2M_2=0\)，所以参数确为 \(\alpha+\beta t\)，不同参数给出不同扩张类。
若 \(\alpha\neq0\)，有 \(t^3v=\alpha tu\neq0\)。向量
\(v,tv,t^2v,t^3v\) 线性无关，故 \(E\cong M_4\)。
若 \(\alpha=0,\beta\neq0\)，有 \(t^2v=\beta tu\neq0\)、\(t^3v=0\)，故四维模含有一个长度三的幂零块，剩余一维给出 \(E\cong M_3\oplus M_1\)。
若参数为零，\(v\) 生成一个与指定子模互补的 \(M_2\)，故 \(E\cong M_2\oplus M_2\)。
扩张类的两个标量参数与中间项的三种同构类型记录的是不同的分类要求。
\end{solution}

\begin{optionalexercise}\label{b5:ex:08-projective-sylow-dimension}
设 \(G\) 为有限群，\(k\) 的特征为 \(p\)，\(P\) 为一个 Sylow \(p\) 子群。证明每个有限维投射 \(kG\)-模的维数都被 \(|P|\) 整除，并给出维数可被 \(|P|\) 整除却不投射的模。
\end{optionalexercise}
\begin{solution}
将 \(G\) 分为互不相交的集合 \(Pg\)，可知左 \(kP\)-模 \(kG\) 自由。因此投射 \(kG\)-模限制到 \(P\) 后，是有限自由 \(kP\)-模的直和因子，仍投射。
\(kP\) 是局部代数，唯一简单模为 \(k\)，其投射覆盖为 \(kP\)。由\cref{b5:thm:split-projective-covers}，有限维投射 \(kP\)-模都是有限自由模，维数为 \(|P|\) 的倍数。
反例取 \(G=C_p\) 及 \(p\) 维平凡模 \(k^{\oplus p}\)。如果它投射，便应为秩一自由 \(kC_p\)-模；但自由模的生成元 \(s\) 作用并非恒等，而平凡模中作用恒等，矛盾。
\end{solution}

\begin{optionalexercise}\label{b5:ex:08-projective-tensor}
设 \(G\) 为有限群，\(k\) 为任意域，\(P\) 为有限维投射 \(kG\)-模，\(V\) 为任意有限维 \(kG\)-模。证明带对角群作用的 \(P\otimes_k V\) 仍投射。
\end{optionalexercise}
\begin{solution}
先看 \(kG\otimes_kV\)。从对角作用的版本到只在第一因子作用的版本，定义
\[
g\otimes v\longmapsto g\otimes g^{-1}v.
\]
它的逆为 \(g\otimes v\mapsto g\otimes gv\)。对 \(h\in G\)，对角作用后的像为
\(hg\otimes(hg)^{-1}(hv)=hg\otimes g^{-1}v\)，恰为目标上的左作用，故这是模同构。目标按 \(V\) 的一组基为 \(\dim V\) 份正则模，因而自由。
取 \(P\) 在某个 \((kG)^r\) 中的直和补，张量保持此直和分解，所以 \(P\otimes V\) 为自由模 \((kG\otimes V)^r\) 的直和因子，因而投射。
\end{solution}

\begin{optionalexercise}\label{b5:ex:08-two-block-tensor}
设 \(G=C_p\)，\(\operatorname{char}k=p\)，\(M_2\) 中生成元的矩阵为 \(I+N\)，其中 \(N^2=0,N\neq0\)。分别在 \(p=2\) 和 \(p>2\) 时求对角作用的 \(M_2\otimes_kM_2\) 的不可分解分解。
\end{optionalexercise}
\begin{solution}
取基 \(u,v\) 满足 \(Nu=0,Nv=u\)。记张量基为
\(a=u\otimes u,b=u\otimes v,c=v\otimes u,d=v\otimes v\)。
若 \(D=s-1\)，则
\[
Da=0,\qquad Db=Dc=a,\qquad Dd=b+c+a.
\]
因此 \(\operatorname{rank}D=2\)，且 \(D^2d=2a\)，其余基向量的二次像为零，\(D^3=0\)。
若 \(p=2\)，\(D^2=0\)，四维空间上的秩二幂零算子由两个长度二块组成，故张量积为 \(M_2^{\oplus2}\)。
若 \(p>2\)，二次幂秩为一，块长为 \(3,1\)，故张量积为 \(M_3\oplus M_1\)。\(p=2\) 时 \(M_2=kC_2\) 投射，其张量积的结果也符合上一题。
\end{solution}

\begin{optionalexercise}\label{b5:ex:08-nonsymmetric-cartan}
设 \(A=T_2(k)\)，取顶点幂等元 \(e_i=E_{ii}\) 和简单模
\(S_i=Ae_i/J(A)e_i\)。按本章“行是简单模、列是投射覆盖”的约定计算 Cartan 矩阵，说明一般有限维代数的 Cartan 矩阵不必对称。
\end{optionalexercise}
\begin{solution}
根为 \(kE_{12}\)。\(Ae_1\) 只有基向量 \(E_{11}\)，所以 \(P_1=S_1\)。\(Ae_2\) 的基为 \(E_{12},E_{22}\)，根子模 \(kE_{12}\) 同构于 \(S_1\)，商为 \(S_2\)。因此
\[
C_A=\begin{pmatrix}1&1\\0&1\end{pmatrix}.
\]
它不对称。这个计算也提示 Cartan 矩阵的指标顺序必须固定，不能根据名称把它当成根系中的对称化数据。
\end{solution}

\begin{optionalexercise}\label{b5:ex:08-c6-blocks}
设 \(\operatorname{char}k=3\)，\(G=C_3\times C_2=\langle r\rangle\times\langle s\rangle\)。
求 \(kG\) 的块、简单模、不可分解投射模及 Cartan 矩阵。再证明与 \(s\mapsto-1,r\mapsto1\) 的一维表示作张量会交换两个块中的模。
\end{optionalexercise}
\begin{solution}
因为 \(2\) 可逆，中心幂等元
\(e_+=(1+s)/2\)、\(e_-=(1-s)/2\)
正交且和为一。每个角理想 \(kGe_\pm\) 都同构于局部代数 \(kC_3\cong k[t]/(t^3)\)，故它们恰为两个块。
简单模 \(S_+,S_-\) 均为一维，\(r\) 平凡作用，\(s\) 分别按 \(1,-1\) 作用。投射覆盖为 \(P_\pm=kGe_\pm\)，每个有三个相同的简单合成因子，所以 \(C=\operatorname{diag}(3,3)\)。
与指定一维表示作张量后，\(r\) 的矩阵不变而 \(s\) 的矩阵乘 \(-1\)，因此 \(e_+\)、\(e_-\) 的作用互换。块分解使同态和扩张分开，却不表示张量任意表示都保持原块。
\end{solution}

\begin{optionalexercise}\label{b5:ex:08-s3-augmentation}
设 \(\operatorname{char}k=2\)，采用\cref{b5:exm:s3-two-blocks}的 \(S_3\) 块幂等元 \(e,f\)。令 \(I=\ker(\varepsilon:kS_3\to k)\)。将 \(I\) 分解为不可分解模，并求 \(I\) 的不动空间及它与 \(J(kS_3)\) 的关系。
\end{optionalexercise}
\begin{solution}
局部块 \(kGe\) 以 \(e,se\) 为基，二者的增广都为 \(3=1\)，故其增广核是
\(k(e+se)=k\sum_{g\in G}g\)，为平凡一维模。
又因 \(\varepsilon(f)=1-\varepsilon(e)=0\)，矩阵块 \(kGf\) 整体包含在 \(I\) 中。该块同构于 \(M_2(k)\)，作为左模是二维简单模 \(S_2\) 的两份直和，因此
\[
I\cong S_1\oplus S_2^{\oplus2}.
\]
矩阵块没有根，局部二维块的根就是 \(k(e+se)\)，所以
\(J(kS_3)=k\sum_g g\subsetneq I\)。
不动空间只来自 \(S_1\)；非平凡简单模 \(S_2\) 不含平凡子模。因此 \(I^G=J(kS_3)\) 为一维。增广理想与群代数的根在一般有限群中可以不同，即使特征整除群阶。
\end{solution}

\begin{optionalexercise}\label{b5:ex:08-regular-set-and-parts}
证明 \(S_3\) 的 \(3\)-正则元素不构成子群。再设 \(g\) 的阶为 \(12\)，按 \(p=2\) 写出其交换的 \(p\) 部分与 \(p'\) 部分，并核验两部分的阶。
\end{optionalexercise}
\begin{solution}
恒等元及三个换位的阶与 \(3\) 互素，但两个不同换位的乘积是三循环，所以该集合不封闭。
对阶为 \(12\) 的 \(g\)，取
\[
g_2=g^9,\qquad g_{2'}=g^4.
\]
前者的阶为 \(12/\gcd(12,9)=4\)，后者的阶为 \(3\)，且二者都在 \(\langle g\rangle\) 中，所以交换；乘积为 \(g^{13}=g\)。指数 \(9\) 在模 \(4\) 时为 \(1\)、模 \(3\) 时为 \(0\)，指数 \(4\) 则相反，这就是相应的中国剩余分解。
\end{solution}

\begin{optionalexercise}\label{b5:ex:08-trace-brauer-distinction}
设 \(k\) 的特征为 \(p\)，\(G=C_p\)。比较正则模 \(M_p=kG\)、\(p\) 份平凡模的直和和零模的 \(k\) 值迹函数、Brauer 特征标及不动空间。哪些资料能够区分这三个对象？
\end{optionalexercise}
\begin{solution}
\(G\) 的唯一 \(p\)-正则元素是恒等元。前两个模的 Brauer 特征标在该元素上都取整数 \(p\)，零模取 \(0\)，所以 Brauer 特征标区分它们与零模，却不能区分这两个非零模。
三者的 \(k\) 值迹函数都为零：在两个 \(p\) 维模上，每个群元素都幂幺，迹为域中的 \(p=0\)。不动空间则分别为一维、\(p\) 维及零维，因而能够区分三者。普通的模特征矩阵迹会进一步丢掉整数重数，Brauer 特征标保留重数，但仍不记录扩张是否分裂。
\end{solution}

\begin{optionalexercise}\label{b5:ex:08-cyclic-mixed-brauer}
设 \(G=C_{p^r}\times C_\ell=\langle s\rangle\times\langle h\rangle\)，其中 \(p\nmid\ell\)，并固定本章的分裂 \(p\)-模系统。令 \(\alpha\in k\) 为本原 \(\ell\) 次单位根，\(\zeta=\widehat\alpha\)。求全部简单模、不可分解模及它们的 Brauer 特征标，再求分解矩阵和 Cartan 矩阵。
\end{optionalexercise}
\begin{solution}
中心幂等元
\[
e_j=\frac1\ell\sum_{v=0}^{\ell-1}\alpha^{-jv}h^v
\qquad(0\leq j<\ell)
\]
将群代数分成 \(\ell\) 个块，每个块同构于 \(kC_{p^r}\)。第 \(j\) 块的唯一简单模 \(S_j\) 由 \(s\mapsto1,h\mapsto\alpha^j\) 给出，全部不可分解模为
\(M_d\otimes S_j\)，其中 \(1\leq d\leq p^r\)。
\(p\)-正则元素恰为 \(h^v\)，故
\[
\varphi_{M_d\otimes S_j}(h^v)=d\,\zeta^{jv}.
\]
在特征零中，所有简单表示均一维，由 \(s\) 的一个 \(p^r\) 次单位根及 \(h\mapsto\zeta^j\) 确定。前者约化后必须为 \(1\)，因为特征 \(p\) 中 \(X^{p^r}-1=(X-1)^{p^r}\)。所以每个这样的普通表示约化为 \(S_j\)，分解矩阵的相应行就是第 \(j\) 个标准行向量，且每一种行重复 \(p^r\) 次。
每个简单模的投射覆盖是该块的正则模，具有 \(p^r\) 个相同简单因子。因此
\(C=p^rI_\ell\)，也在本例直接满足 \(D^{\mathsf T}D=C\)。
\end{solution}

\begin{optionalexercise}\label{b5:ex:08-s3-characteristic-three}
设 \(k\) 为特征 \(3\) 的分裂域。求 \(S_3\) 的全部不可约 Brauer 特征标及分解矩阵，普通不可约表示依次按平凡、符号和标准二维表示排列。判断二维增广子模
\[
U=\bigl\{(x,y,z)\in k^3:x+y+z=0\bigr\}
\]
是否半单，并给出它的简单合成因子。
\end{optionalexercise}
\begin{solution}
记 \(r=(123)\)。在任意非零模上，\((r-1)^3=0\)，所以有非零的 \(C_3\)-不动向量。因 \(C_3\) 正规，全部 \(C_3\)-不动向量形成 \(S_3\)-子模；在简单模上便是全空间。因此所有简单模都来自 \(S_3/C_3\cong C_2\)，恰为平凡模 \(S_+\) 和符号模 \(S_-\)。
\(3\)-正则类由恒等元及换位 \(s\) 代表，所以表及分解矩阵为
\[
\begin{array}{c|rr}
&1&s\\ \hline
\varphi_+&1&1\\
\varphi_-&1&-1
\end{array}
\qquad
D=\begin{pmatrix}1&0\\0&1\\1&1\end{pmatrix}.
\]
最后一行来自标准普通特征标在两类上的值 \((2,0)=\varphi_++\varphi_-\)。
在 \(U\) 中，向量 \((1,1,1)\) 非零且坐标和为 \(0\)，给出平凡子模；一维商为符号模，例如一个换位在 \(U\) 上的迹为 \(0\)，减去子模上的迹 \(1\)，商的迹为 \(-1\)。三循环在 \(U\) 上并非恒等，例如它把 \((1,-1,0)\) 送到 \((0,1,-1)\)。若 \(U\cong S_+\oplus S_-\)，三循环却应在两项上都平凡，矛盾。所以它不是半单模。
\end{solution}

\begin{optionalexercise}\label{b5:ex:08-s3-three-cartan}
沿用上一题，令 \(t=r-1\)、\(e_\pm=(1\pm s)/2\)。
求 \(kS_3\) 的根以及投射模 \(P_\pm=kS_3e_\pm\) 的根层，进而计算 Cartan 矩阵，并证明此时群代数只有一个块。
\end{optionalexercise}
\begin{solution}
群关系给出 \(t^3=0\) 及
\[
sts=r^{-1}-1=(1+t)^2-1=-t+t^2.
\]
因此由 \(t\) 生成的双边理想以 \(t,t^2,ts,t^2s\) 为基，立方为零，商为半单代数 \(kC_2\)，故它就是 \(J(kS_3)\)。
\(P_\pm\) 的基为 \(e_\pm,te_\pm,t^2e_\pm\)，顶分别为 \(S_\pm\)。在第 \(i\) 个根层上，\(r\) 平凡作用，而
\[
s\,t^ie_\pm\equiv(\pm1)(-1)^i t^ie_\pm
\pmod{J^{i+1}P_\pm}.
\]
所以 \(P_+\) 从顶到底的因子为 \(S_+,S_-,S_+\)，\(P_-\) 的因子为 \(S_-,S_+,S_-\)。两模投射且顶简单，故为相应不可分解覆盖，得到
\[
C=\begin{pmatrix}2&1\\1&2\end{pmatrix}
=D^{\mathsf T}D.
\]
一个不可分解模只属于一个块；\(P_+\) 的合成因子同时包含两种简单模，所以它们必属于同一块。若有其他非零块，它必有简单模，与简单类型已经穷尽矛盾。因此群代数只有一个块。
\end{solution}

\begin{optionalexercise}\label{b5:ex:08-naive-brauer-inner-product}
对上一题的两个 Brauer 特征标，定义
\[
(f,g)_0=\frac1{|S_3|}\sum_{x\in(S_3)_{3'}}f(x)\overline{g(x)}.
\]
计算两者的平方长度及内积，说明为何不能直接把普通不可约特征标的正交关系套到这里。
\end{optionalexercise}
\begin{solution}
\(3\)-正则元素包括恒等元及三个换位，因此
\[
(\varphi_+,\varphi_+)_0=(\varphi_-,\varphi_-)_0=\frac46=\frac23,
\qquad
(\varphi_+,\varphi_-)_0=\frac{1-3}{6}=-\frac13.
\]
它们既未归一化，也不正交。不可约 Brauer 特征标形成类函数的一组基，并不表示在任意仿照普通公式定义的内积下都成为标准正交基；本章的基定理只给出线性独立与生成性。
\end{solution}

\begin{optionalexercise}\label{b5:ex:08-s3-two-tensor}
在特征 \(2\) 的 \(S_3\) 模中，记 \(S_1\) 为平凡模，\(S_2\) 为二维简单投射模，\(P_1\) 为 \(S_1\) 的投射覆盖。先由 Brauer 特征标求 \(S_2\otimes S_2\) 的简单重数，再证明实际模同构
\[
S_2\otimes S_2\cong P_1\oplus S_2.
\]
\end{optionalexercise}
\begin{solution}
两类 \(2\)-正则元素由 \(1,r\) 代表，\(\varphi_{S_1}=(1,1)\)、\(\varphi_{S_2}=(2,-1)\)。张量特征标为 \((4,1)\)，解
\((4,1)=a(1,1)+b(2,-1)\)
得到 \(a=2,b=1\)。这一步只确定两个平凡因子与一个二维简单因子。
由于 \(S_2\) 投射，\cref{b5:ex:08-projective-tensor}保证张量积也投射。全部不可分解投射类型为 \(P_1,S_2\)，它们的简单重数分别是 \((2,0)\)、\((0,1)\)。因此投射直和分解只能各取一份，得到所求同构。这里从合成重数走到实际直和，额外使用了投射性，不能仅凭特征标相同就作这一步推断。
\end{solution}

\begin{optionalexercise}\label{b5:ex:08-s3-two-lattices}
设 \((K,\mathcal O,k)\) 为本章特征 \(3\) 的分裂模系统，\(\pi\) 为一致化元，并让 \(S_3\) 置换三个坐标。在
\(V=K^3/K(1,1,1)\) 中取满格
\[
\Lambda=\mathcal O^3/\mathcal O(1,1,1),\qquad
L=\text{\(\bigl\{(x,y,z)\in\mathcal O^3:x+y+z=0\bigr\}\) 在 \(V\) 中的像}.
\]
证明 \(\Lambda/L\cong\mathcal O/3\mathcal O\)，并证明两个约化不同构，却有相同的简单合成因子及 Brauer 特征标。
\end{optionalexercise}
\begin{solution}
对角向量有一个单位坐标，故商 \(\Lambda\) 为秩二自由模；和为零的子模也自由秩二，在 \(V\) 中的映射单射，因为特征零中常向量的坐标和为 \(3c\)，只有 \(c=0\) 时为零。它们都是被群保持的满格。
求和映射识别 \(\mathcal O^3/\ker\Sigma\cong\mathcal O\)，而对角子模在其下的像为 \(3\mathcal O\)，所以 \(\Lambda/L\cong\mathcal O/3\mathcal O\)。
约化后，\(L/\pi L\) 是 \(k^3\) 的二维增广子模，含一维平凡子模，商为符号模。另一方面，\(\Lambda/\pi\Lambda\cong k^3/k(1,1,1)\)。它的增广子模像是一维符号模，商为平凡模，因此简单重数相同。
但第二个约化没有非零群不动元：若向量 \(v\in k^3\) 的类不动，则每个换位使 \(v\) 的变化都在对角直线上。这个变化有一个坐标为零，故只能是零；所有换位因此固定 \(v\)，使 \(v\) 本身为常向量，其商类为零。第一个约化的不动空间却由 \((1,1,1)\) 生成，维数为一，故两个约化不同构。它们的 Brauer 特征标都为 \(\varphi_++\varphi_-\)。
\end{solution}

\begin{optionalexercise}\label{b5:ex:08-decomposition-column-rank}
设 \(G\) 为有限群，\(p\) 为素数，\((K,\mathcal O,k)\) 为 \(G\) 的分裂 \(p\)-模系统，\(K\) 含有全部 \(|G|\) 次单位根，并固定这些根到复数的对应。记 \(d:R_K(G)\to G_0(kG)\) 为分解映射，\(D\) 为分解矩阵。证明 \(D\) 的列在 \(\mathbb C\) 上线性无关，并据此计算
\(\mathbb Q\otimes\ker d\) 的维数：用普通共轭类数与 \(p\)-正则共轭类数表示。说明这一秩结论本身为何不能证明分解映射在整数上满射。
\end{optionalexercise}
\begin{solution}
普通不可约特征标限制到 \(G_{p'}\) 后生成全部复类函数，而不可约 Brauer 特征标是一组基。恒等式
\(\chi_\alpha|_{G_{p'}}=\sum_i d_{\alpha i}\varphi_i\)
因此说明 \(D\) 的行向量生成整个 \(\mathbb C^s\)，其中 \(s\) 为简单模数，即 \(p\)-正则类数。所以 \(D\) 的秩为 \(s\)，各列线性无关。
普通表示群的秩为全部共轭类数 \(c\)，分解映射的矩阵是 \(D^{\mathsf T}\)。该整数矩阵在 \(\mathbb Q\) 与 \(\mathbb C\) 上秩相同，故有理化后的核维数为 \(c-s\)。但满秩的整数矩阵未必在整数格上满射，例如一阶矩阵 \((2)\) 的有理秩为一，整数像却为 \(2\mathbb Z\)。因此若要证明整数上的更强结论，还需额外论证。
\end{solution}

\begin{optionalexercise}\label{b5:ex:08-brauer-domain}
设 \(\operatorname{char}k=2\)，\(G=S_3=\langle r,s\rangle\)，\(P=\langle s\rangle\)。
把支集截断到 \(C_G(P)\) 的线性映射暂时定义在整个 \(kG\) 上。证明它一般不保持乘法，并说明这与本章的 Brauer 映射命题不矛盾。
\end{optionalexercise}
\begin{solution}
\(C_G(P)=P\)，两个三循环 \(r,r^{-1}\) 都不属于 \(P\)。因此全空间上的截断分别把它们送到零，却把其乘积 \(rr^{-1}=1\) 送到 \(1\)，不保持乘法。
命题\cref{b5:prop:brauer-truncation}的定义域是 \((kG)^P\)，其中元素须对 \(P\) 的共轭作用不变。这里 \(srs=r^{-1}\neq r\)，所以两个被选的元素都不在该定义域。轨道系数相同正是乘法性证明中消去非平凡 \(p\) 轨道的条件，不能删去。
\end{solution}

\begin{optionalexercise}\label{b5:ex:08-defect-change-prime}
在特征 \(3\) 中，求\cref{b5:ex:08-c6-blocks}的 \(C_6\) 两个块的缺陷群，以及\cref{b5:ex:08-s3-three-cartan}的 \(S_3\) 唯一块的缺陷群。将后者与正文中特征 \(2\) 的 \(S_3\) 比较。
\end{optionalexercise}
\begin{solution}
\(C_6\) 交换，所以任一子群的中心化子都是全群，Brauer 截断在群代数上就是恒等。因此两个非零块幂等元对 Sylow \(3\) 子群 \(C_3\) 的 Brauer 像都非零；没有更大的 \(3\) 子群，两个块的缺陷群都是 \(C_3\)。
特征 \(3\) 的 \(S_3\) 只有一个块，幂等元为 \(1\)。对 \(P=\langle r\rangle\) 有 \(\operatorname{Br}_P^G(1)=1\neq0\)，且 \(P\) 是 Sylow \(3\) 子群，所以它是缺陷群。
特征 \(2\) 时，\(S_3\) 则有一个缺陷群为 \(C_2\) 的二维局部块和一个缺陷群平凡的四维矩阵块。改变模特征会同时改变块的数目、块的结构以及相应的缺陷群。
\end{solution}

\begin{optionalexercise}\label{b5:ex:08-induction-reduction-example}
设 \(\operatorname{char}k=3\)，\(G=S_3\)，\(H=\langle s\rangle\cong C_2\)。
分别诱导 \(H\) 的平凡与符号表示，确定所得模的不可分解类型及简单重数。再从特征零的诱导分解出发约化，核验与直接计算一致。
\end{optionalexercise}
\begin{solution}
令 \(e_\pm=(1\pm s)/2\)。一维 \(kH\)-模 \(k_\pm\) 同构于 \(kHe_\pm\)，乘法诱导
\[
kG\otimes_{kH}k_\pm\cong kGe_\pm=P_\pm.
\]
由\cref{b5:ex:08-s3-three-cartan}，两者都是不可分解投射模，简单重数分别为
\(2S_++S_-\) 与 \(S_++2S_-\)。
特征零中，Frobenius 互反性给出
\[
\operatorname{Ind}_{C_2}^{S_3}\mathbf1=\mathbf1+\chi_{\mathrm{std}},
\qquad
\operatorname{Ind}_{C_2}^{S_3}\operatorname{sgn}
=\operatorname{sgn}+\chi_{\mathrm{std}}.
\]
例如标准二维表示限制到一个换位时有特征值 \(1,-1\)，两个一维类型各出现一次，而群的平凡、符号表示各自只贡献对应的一个类型。按\cref{b5:ex:08-s3-characteristic-three}的分解矩阵约化，标准表示给出 \([S_+]+[S_-]\)，所以两式分别变为 \(2[S_+]+[S_-]\)、\([S_+]+2[S_-]\)，与直接计算相同。特征零的直和经约化后只保证合成类相同，此处所得的两个模本身都不可分解。
\end{solution}

% !TeX root = ../../Book5.tex
% 纲要标题：## 第9章　箭图与有限维代数的表示
% 纲要位置：工作记录/计划纲要.md，第 4009 行。
\chapter{箭图与有限维代数的表示}
\label{b5:ch:09}
\BookChapterNumber{b5:ch:09}

\begin{chapterabstract}
本章用箭图组织向量空间及线性映射，并把箭的复合写成路径代数的乘法。先建立箭图表示与左模的等价，再用选基和复合映射的秩分类线性定向的区间表示。无环箭图的短投射消解给出可直接计算的 Hom--Ext 正合列。最后从有限维代数的根商提取顶点、箭及关系，说明代数闭域上的有限维代数经过 Morita 约化后都可由容许关系箭图描述。
\end{chapterabstract}

\begin{learninggoals}
\begin{itemize}
\item 能按既定路径乘法构造路径代数，判断有限维性及关系理想的容许性。
\item 能在左路径代数模与箭图表示之间转换，并用逐顶点计算处理同态、核、商及正合列。
\item 能分解线性定向的区间表示，从各段复合的秩恢复区间重数。
\item 能写出顶点投射表示和短投射消解，计算态射及扩张，并核验扩张的方向。
\item 能从分裂基本代数的 \(J/J^2\) 读出 Ext 箭图，区分箭图第一层资料与更高乘法关系。
\end{itemize}
前四节需要向量空间、线性映射、模、直和及投射模的基本知识，所用 \(\operatorname{Ext}^1\) 将通过短正合扩张定义。\secref{b5:sec:0905}还会用到第二卷关于有限维代数根的幂零性、幂等元提升及满幂等元的 Morita 等价。
\end{learninggoals}

一个算子的标准形已经是表示分类的例子；若同时出现多个空间和映射，分别化简每个矩阵通常会破坏其他矩阵已经取得的形式。箭图把允许的换基精确地固定在各个顶点上，使“同时化简”成为一个统一问题。没有关系时，路径记录一切可作的复合；有关系时，路径之间的线性恒等式则记录额外限制。路径代数及最初的区间例子可对照\cite[Section 2.8, pp. 19--22; Section 6.2, pp. 150--153]{EtingofEtAl2011RepresentationTheory}。

% !TeX root = ../../Book5.tex
% 纲要标题：### 9.1 箭图与路径代数
% 纲要位置：工作记录/计划纲要.md，第 4011 行。

\section{箭图与路径代数}
\label{b5:sec:0901}


% 纲要范围：
%
% * 顶点、箭和有向路径
% * 路径代数及幂等元
% * 有关系箭图与有限维性
%

一个线性算子可以看成一个向量空间上的一支箭；两个线性映射若首尾相接，还可以讨论它们的复合。箭图把这些信息分开记录：顶点说明有哪些空间，箭说明哪些空间之间给了映射。先只保留箭与路径，就能构造一个代数，随后再把线性映射赋给它。

\subsection{顶点、箭和有向路径}
\label{b5:subsec:0901:01}

\begin{definition}\label{b5:def:quiver}
设 \(Q_0,Q_1\) 为有限集合，\(s,t:Q_1\to Q_0\) 为两个映射。由这些数据组成的有向图 \(Q=(Q_0,Q_1,s,t)\) 称为一个有限\term[jiantu]{箭图}[quiver]；\(Q_0\) 的元素称为顶点，\(Q_1\) 的元素称为箭。箭 \(a\) 从 \(s(a)\) 指向 \(t(a)\)，记为 \(a:s(a)\to t(a)\)。允许重箭以及起点等于终点的圈。
\end{definition}
\symidx{\(Q_0,Q_1,s,t\)：箭图的顶点集、箭集及起点、终点映射}
箭图不要求同一对顶点间至多一支箭。两支不同的箭以后可以表示两个完全无关的线性映射；把它们合并，会改变所研究的对象。这里允许圈，但“无环箭图”将专指没有正长度有向闭路径的箭图，并非仅仅没有一支箭构成的圈。
\begin{definition}\label{b5:def:quiver-path}
设 \(Q=(Q_0,Q_1,s,t)\) 为有限箭图，\(s,t\) 分别给出箭的起点、终点，\(r\ge1\) 为整数。若箭 \(a_1,\ldots,a_r\in Q_1\) 满足 \(t(a_\nu)=s(a_{\nu+1})\) 对 \(1\le\nu<r\) 成立，则依次经过这些箭的序列称为长度 \(r\) 的\term[youxiang-lujing]{有向路径}[oriented path]，记作 \(a_r\cdots a_1\)。它从 \(s(a_1)\) 指向 \(t(a_r)\)。每个顶点 \(i\in Q_0\) 另有一条长度零的路径 \(e_i\)，其起点与终点都是 \(i\)。
\end{definition}
例如在 \(1\xrightarrow{a}2\xrightarrow{b}3\) 中，长度二的路径记为 \(ba\)。这个记法与线性映射的复合一致：先作 \(a\)，再作 \(b\)。相反，\(ab\) 在这个箭图中不能首尾相接。无向图是一条线并不规定唯一的箭图；\(1\to2\leftarrow3\) 的两支箭就不能组成长度二的路径。

\subsection{路径代数及幂等元}
\label{b5:subsec:0901:02}

\begin{definition}\label{b5:def:path-algebra}
设 \(k\) 为域，\(Q\) 为非空有限箭图。以 \(Q\) 的全部有向路径为 \(k\) 基，取只有有限多个非零系数的线性组合。对基路径 \(p:i\to j\)、\(q:\ell\to m\)，规定 \(qp\) 在 \(j=\ell\) 时为先走 \(p\) 再走 \(q\) 的拼接路径，在 \(j\ne\ell\) 时为零；再按 \(k\) 双线性延拓乘法。所得代数记为 \(kQ\)，称为 \(Q\) 的\term[lujing-daishu]{路径代数}[path algebra]。
\end{definition}
\symidx{\(kQ\)：箭图的路径代数}
拼接满足结合律。对三条路径，如果相邻的两处都能接合，两种括号给出同一条箭序列；如果某处不能接合，两种乘法都为零。逐项展开便得到任意元素的结合律。长度零路径满足
\[
1=\sum_{i\in Q_0}e_i,\qquad e_ie_j=\delta_{ij}e_i,\qquad
a=e_{t(a)}a e_{s(a)}.
\]
尤其 \(e_j(kQ)e_i\) 以全部从 \(i\) 到 \(j\) 的路径为基。任意正长度路径是箭的乘积，所以这些幂等元和箭生成 \(kQ\)。给它们指定到另一个代数中的像，只要满足上式的顶点关系，就唯一决定一个代数同态：每条路径的像只能是各箭像的乘积，而上述关系保证不可拼接的乘积被送到零。
\begin{proposition}\label{b5:prop:finite-path-algebra}
设 \(k\) 为域，\(Q\) 为非空有限箭图。代数 \(kQ\) 有限维，当且仅当 \(Q\) 没有正长度有向闭路径。
\end{proposition}
\begin{proof}
若有正长度闭路径 \(c\)，则 \(e_i,c,c^2,\ldots\) 是长度不同的基路径，线性无关，所以代数无限维。反过来，若没有闭路径，任何路径都不能重复经过一个顶点，否则两次经过之间就是闭路径。因此路径长度至多为 \(|Q_0|-1\)。每个长度的箭序列只有有限多个，所以全部路径有限。
\end{proof}
在 \(1\to2\to3\) 中，基为 \(e_1,e_2,e_3,a,b,ba\)，维数六；在 \(1\to2\leftarrow3\) 中，基只含三个顶点和两支箭，维数五。两个箭图的无向图相同，路径代数却不相同。若只有一个顶点和一个圈 \(a\)，则 \(kQ\cong k[t]\)；有两个圈时，则得到两个不定元的自由结合代数，两个圈的乘法一般不交换。

\subsection{有关系箭图与有限维性}
\label{b5:subsec:0901:03}

\begin{definition}\label{b5:def:admissible-quiver-relations}
设 \(k\) 为域，\(Q\) 为非空有限箭图。由全部箭生成的双边理想记为 \(\mathfrak a\)，称为\term[jian-lixiang]{箭理想}[arrow ideal]。它以正长度路径为基，\(\mathfrak a^r\) 以长度至少 \(r\) 的路径为基。若双边理想 \(I\) 满足
\[
\mathfrak a^N\subseteq I\subseteq\mathfrak a^2
\quad\text{对某个整数}\;N\ge2,
\]
就称 \(I\) 为\term[rongxu-guanxi-lixiang]{容许关系理想}[admissible ideal]，称 \((Q,I)\) 为有关系箭图，并研究商代数 \(kQ/I\)。
\end{definition}
\symidx{\(\mathfrak a\)：路径代数的箭理想}
一个关系可以是具有相同起点和终点的若干路径的线性组合。一般元素也可按 \(e_j I e_i\) 拆成这种形式。商中要求每个关系为零；例如在 \(1\xrightarrow{a}2\xrightarrow{b}3\) 中加入 \(ba=0\)，以后对应的两个线性映射就须满足复合为零。
\begin{proposition}\label{b5:prop:bound-quiver-radical}
设 \(k\) 为域，\(Q\) 为非空有限箭图，\(\mathfrak a\subseteq kQ\) 为由全部箭生成的双边理想，\(I\subseteq kQ\) 为双边理想，且对某个整数 \(N\ge2\) 有 \(\mathfrak a^N\subseteq I\subseteq\mathfrak a^2\)。令 \(A=kQ/I\)。则 \(A\) 有限维，其 Jacobson 根为
\[
J(A)=\mathfrak a/I,\qquad A/J(A)\cong\prod_{i\in Q_0}k.
\]
\end{proposition}
\begin{proof}
所有长度至少 \(N\) 的路径在商中为零，所以有限多个较短路径的像生成 \(A\)。理想 \(\mathfrak a/I\) 幂零；幂零双边理想包含在 Jacobson 根中，因为对其元素 \(x\) 和任意 \(a\in A\)，\(1-ax\) 有有限几何级数给出的逆。模掉这个理想，只剩互相正交的顶点幂等元，商为 \(\prod_i k\)，其根为零。因此 \(J(A)\) 也包含于 \(\mathfrak a/I\)，得到等式。
\end{proof}
容许条件既使长路径消失，又不对顶点与箭的第一层施加线性关系。有限维性本身不保证一个给定关系理想容许。例如一个圈的关系 \(a^2-a=0\) 给出二维商，但关系含有一次项；这个商同构于 \(k\times k\)，它的根为零，圈的像却不是零。本章最后将从有限维代数的根重新选择箭图与容许关系，而不是把任意展示都当成这种形式。

% !TeX root = ../../Book5.tex
% 纲要标题：### 9.2 表示与模
% 纲要位置：工作记录/计划纲要.md，第 4017 行。

\section{表示与模}
\label{b5:sec:0902}


% 纲要范围：
%
% * 向量空间、线性箭及交换关系
% * 箭图表示与路径代数模的等价
% * 左模约定与路径复合方向
%

路径代数把箭的复合变成乘法，箭图表示则把它变成线性映射的复合。本节不仅比较两者的对象，还比较保持作用的映射；这使直和、核、商和扩张都能在两种描述之间转换。

\subsection{向量空间、线性箭及交换关系}
\label{b5:subsec:0902:01}

\begin{definition}\label{b5:def:quiver-representation}
设 \(k\) 为域，\(Q\) 为有限箭图。给每个顶点 \(i\) 一个有限维 \(k\) 向量空间 \(V_i\)，给每支箭 \(a:i\to j\) 一个线性映射 \(V_a:V_i\to V_j\)，所得数据 \(V=(V_i,V_a)\) 称为 \(Q\) 的有限维 \(k\)\term[jiantu-biaoshi]{表示}[representation of a quiver]。对路径 \(p=a_r\cdots a_1\)，记 \(V_p=V_{a_r}\cdots V_{a_1}\)，并令 \(V_{e_i}=1_{V_i}\)。
对两个表示 \(V,W\)，一族线性映射 \(f_i:V_i\to W_i\) 若满足
\[
W_a f_i=f_j V_a\qquad(a:i\to j),
\]
就称为表示同态 \(f:V\to W\)。若所有 \(f_i\) 都可逆，则称 \(f\) 为表示同构。
\end{definition}
对于有关系箭图 \((Q,I)\)，再要求每个关系 \(\sum_p c_p p\in I\) 所给的映射 \(\sum_p c_pV_p\) 为零。只需检查一组双边理想生成关系：左右乘路径就是在两侧复合映射，加法则是相加。记这些表示及其同态组成的范畴为 \(\operatorname{rep}_k(Q,I)\)；\(I=0\) 时写为 \(\operatorname{rep}_k Q\)。
\symidx{\(\operatorname{rep}_k(Q,I)\)：有关系箭图的有限维表示范畴}
\begin{definition}\label{b5:def:quiver-subrepresentation}
设 \(k\) 为域，\(Q\) 为有限箭图，\(V=(V_i,V_a)\)、\(W=(W_i,W_a)\) 为 \(Q\) 的有限维 \(k\) 表示。若一族子空间 \(U_i\subseteq V_i\)、\(i\in Q_0\)，满足 \(V_a(U_i)\subseteq U_j\) 对每支箭 \(a:i\to j\) 成立，则这些子空间及箭映射的限制组成 \(V\) 的子表示 \(U\)。商表示在各顶点取 \(V_i/U_i\)，箭映射由 \(V_a\) 诱导。直和 \(V\oplus W\) 则在各顶点及各支箭分别取直和。向量 \(\mathbf d(V)=(\dim_kV_i)_{i\in Q_0}\) 称为 \(V\) 的\term[weishu-xiangliang]{维数向量}[dimension vector]。
\end{definition}
\symidx{\(\mathbf d(V)\)：箭图表示的维数向量}
表示同态的核与余核也逐顶点计算。例如 \(W_a f_i=f_jV_a\) 使 \(V_a\) 将 \(\ker f_i\) 映入 \(\ker f_j\)，而同一个等式保证商上的箭映射良定义。因此一列表示正合，当且仅当它在每个顶点的向量空间列都正合。由此，\(\operatorname{rep}_k(Q,I)\) 是交换范畴。

\subsection{箭图表示与路径代数模的等价}
\label{b5:subsec:0902:02}

\begin{theorem}\label{b5:thm:quiver-module-equivalence}
设 \(k\) 为域，\(Q\) 为非空有限箭图，\(I\) 为 \(kQ\) 的双边理想，\(e_i\) 为顶点 \(i\) 的长度零路径。有限维左 \(kQ/I\) 模范畴与满足 \(I\) 中全部关系的有限维 \(k\) 箭图表示范畴 \(\operatorname{rep}_k(Q,I)\) 等价。对应在对象上为
\[
M\longmapsto(e_iM,\ a|_{e_iM}),\qquad
(V_i,V_a)\longmapsto\bigoplus_{i\in Q_0}V_i.
\]
第二个对象上，顶点幂等元作坐标投影，路径按相应线性映射的复合作用；两种对应保持直和与正合列。
\end{theorem}
\begin{proof}
若 \(M\) 是左模，单位分解为 \(\sum_i e_i\)，故 \(M=\bigoplus_i e_iM\)。由 \(a=e_jae_i\)，左乘 \(a\) 将 \(e_iM\) 映入 \(e_jM\)，而 \(I\) 在 \(M\) 上作用为零，得到满足关系的表示。模同态把每个 \(e_iM\) 送入目标模的相应分量，并与每条箭对应的线性映射相容，因而给出表示同态。
反过来，给定表示，在直和上令一条从 \(i\) 到 \(j\) 的路径只作用于第 \(i\) 个坐标，并用 \(V_p\) 将其送到第 \(j\) 个坐标。如果两条路径不接合，两个坐标算子的复合为零；如果接合，复合正是拼接路径的算子。因此这些算子延拓为 \(kQ\) 的左作用。关系条件使作用通过 \(kQ/I\)。一族表示同态的直和保持顶点分量，并与两端的箭映射相容，所以也与全部路径作用相容，成为模同态。
从模出发再取直和，由 \((m_i)_i\mapsto\sum_i m_i\) 得到与作用相容的自然同构；从表示出发则在各坐标恢复原空间和原箭映射。以上构造对同态也互逆，故为范畴等价。直和逐坐标保持，正合性由各顶点的核与余核计算保持。
\end{proof}
例如单圈表示就是一个向量空间及其一个线性算子，因而等同于 \(k[t]\) 模；两个圈的表示是两个任意线性算子，不要求它们交换。若另加关系 \(ab-ba=0\)，才变为两个可交换算子。箭图与关系分别记录“有哪些映射”和“映射应满足什么恒等式”。

\subsection{左模约定与路径复合方向}
\label{b5:subsec:0902:03}

本章始终用左模。若 \(a:i\to j\)，左乘 \(a=e_jae_i\) 从 \(e_iM\) 到 \(e_jM\)，与箭向相同。若改用右模，则右乘 \(a\) 从 \(Me_j\) 到 \(Me_i\)，箭向反转。因此把一个右模版本直接抄成左模，往往会同时弄反路径乘积、投射模和扩张箭的方向。
\ProseParagraphBreak
一个能随时检查约定的例子是上三角矩阵代数 \(T_2(k)\)。令 \(e_1=E_{11}\)、\(e_2=E_{22}\)，则 \(E_{12}=e_1E_{12}e_2\)，对应的箭为 \(2\to1\)。把箭送到 \(E_{12}\)，顶点送到两个对角矩阵，便给出此箭图的路径代数与 \(T_2(k)\) 的同构：两侧均以这三个元素为基，全部乘法关系相同。
\begin{proposition}\label{b5:prop:quiver-change-of-basis}
设 \(k\) 为域，\(Q\) 为有限箭图，\(\mathbf d=(d_i)_{i\in Q_0}\) 为非负整数向量。固定每个顶点空间为 \(k^{d_i}\)，则所有 \(k\) 表示由矩阵组
\[
\operatorname{Rep}_k(Q,\mathbf d)
=\prod_{a:i\to j}\operatorname{Mat}_{d_j\times d_i}(k)
\]
给出。换基群 \(\prod_i\operatorname{GL}_{d_i}(k)\) 作用为
\[
(g_i)_i\cdot(X_a)_a=(g_jX_ag_i^{-1})_{a:i\to j}.
\]
两个矩阵组在同一轨道中，当且仅当相应表示同构。加关系时，取满足关系矩阵方程的稳定子集即可。
\end{proposition}
\begin{proof}
每支箭的线性映射由任意相应尺寸的矩阵给定，故得到所列直积。表示同构的方程为 \(Y_ag_i=g_jX_a\)，每个 \(g_i\) 可逆时恰好等价于 \(Y_a=g_jX_ag_i^{-1}\)。关系的矩阵值从起点到终点作同样的左右换基，因而其为零的条件在作用下保持。
\end{proof}
维数向量只规定矩阵的尺寸，并不确定表示。例如一个圈上的一维表示由标量 \(\lambda\in k\) 给出，一维换基不改变这个标量，所以不同的 \(\lambda\) 已经给出不同构的表示。下一节先讨论没有圈、又只有少数箭的情形。

% !TeX root = ../../Book5.tex
% 纲要标题：### 9.3 区间表示与线性定向箭图
% 纲要位置：工作记录/计划纲要.md，第 4023 行。

\section{区间表示与线性定向箭图}
\label{b5:sec:0903}


% 纲要范围：
%
% * A_2、A_3 和线性定向 A_n
% * 直和与不可分解表示
% * 用线性代数分解区间表示
%

一个箭头的分类只需矩阵的秩。沿一条有向直线增加箭头以后，需要同时记录每段复合的秩；这些数仍然足以恢复分解。本节给出实际选基的办法，为后面更一般的箭图分类提供可计算的比较。

\subsection{\texorpdfstring{\(A_2\)、\(A_3\) 与线性定向 \(A_n\)}{A2、A3 与线性定向 An}}
\label{b5:subsec:0903:01}

\begin{definition}\label{b5:def:interval-quiver-representation}
设 \(k\) 为域，\(n\ge1\) 为整数，\(Q\) 为线性定向箭图 \(1\to2\to\cdots\to n\)。对满足 \(1\le a\le b\le n\) 的整数 \(a,b\)，在顶点 \(a,a+1,\ldots,b\) 放置 \(k\)，在其余顶点放置零空间；区间内部的箭映射为恒等，其他箭映射为零。所得表示记为 \(I[a,b]\)，称为\term[qujian-biaoshi]{区间表示}[interval representation]。
\end{definition}
\symidx{\(I[a,b]\)：线性定向箭图的区间表示}
对 \(A_2=(1\to2)\)，给定映射 \(f:V_1\to V_2\)，令 \(r=\operatorname{rank}f\)、\(d_i=\dim V_i\)。在 \(V_1\) 中补全 \(\ker f\) 的基，并用其余基向量的像作为 \(\operatorname{im}f\) 的基，再补全到 \(V_2\)，得到
\[
V\cong I[1,2]^{\oplus r}\oplus I[1,1]^{\oplus(d_1-r)}
\oplus I[2,2]^{\oplus(d_2-r)}.
\]
因此两端空间的维数与箭的秩共同决定同构类；只给维数则不够。
\ProseParagraphBreak
对于 \(1\xrightarrow{f}2\xrightarrow{g}3\)，还须比较 \(\operatorname{im}f\) 与 \(\ker g\)。先在它们的交中选基，分别补到两个子空间，再补到 \(V_2\)；把 \(\operatorname{im}f\) 的这些基向量提升到 \(V_1\)，补上 \(\ker f\) 的基，并在 \(V_3\) 中补全相应的像。这就把表示分成六种区间。记 \(r_f,r_g,r_{gf}\) 为三个映射的秩，则相应重数为
\[
\begin{aligned}
m_{13}&=r_{gf},&m_{12}&=r_f-r_{gf},&m_{23}&=r_g-r_{gf},\\
m_{11}&=d_1-r_f,&m_{22}&=d_2-r_f-r_g+r_{gf},&
m_{33}&=d_3-r_g.
\end{aligned}
\]
重数 \(m_{22}\ge0\) 来自
\[
\dim(\operatorname{im}f\cap\ker g)=r_f-r_{gf}
\le\dim\ker g=d_2-r_g.
\]
这个不等式正好给出所需的非负性；对应的向量是 \(\ker g\) 中不来自前一顶点的部分。
\ProseParagraphBreak
另一种定向 \(1\to2\leftarrow3\) 也可直接分解。对两支箭的像 \(U,W\subseteq V_2\)，取适应 \(U\cap W\subseteq U,W\subseteq U+W\) 的基，再分别提升到两端。交中的每个基方向给一个在三个顶点都为 \(k\)、两箭均为恒等的表示；\(U,W\) 各自的补给相邻两顶点的表示；剩余基及两端的核给单点表示。因此仍有六种不可分解模型，但区间中的箭须按此定向解释。

\subsection{直和与不可分解表示}
\label{b5:subsec:0903:02}

\begin{definition}\label{b5:def:indecomposable-quiver-representation}
设 \(k\) 为域，\(Q\) 为有限箭图，\(V\) 为 \(Q\) 的非零有限维 \(k\) 表示。若 \(V\) 不能写成两个非零表示的直和，就称其为\term[jiantu-bukefenjie]{不可分解表示}[indecomposable representation]。若 \(V\) 没有非零真子表示，就称其为简单表示。两者不同：在箭图 \(1\to2\) 中，区间表示 \(I[1,2]\) 不可分解，却有支集只在顶点2的子表示。
\end{definition}
一个区间表示的自同态在每个非零顶点都是标量，恒等箭迫使这些标量相同。因此
\[
\operatorname{End}_Q\bigl(I[a,b]\bigr)\cong k.
\]
若它有非平凡直和分解，相应投影会给出不等于零、一的幂等自同态，与域中无此幂等元矛盾。所以每个区间不可分解。不同区间的维数向量不同，故不同构。
下面使用 Krull--Schmidt 分解\footnote{克鲁尔（Wolfgang Krull，1899-1971），德国数学家，研究可换代数与模的分解；施密特（\OriginalName{Отто Юльевич Шмидт}，1891-1956），苏联数学家，也研究地球物理与天文学。}。
\begin{proposition}[Krull--Schmidt]\label{b5:thm:quiver-krull-schmidt}
设 \(k\) 为域，\(Q\) 为有限箭图，\(I\) 为 \(kQ\) 的双边理想。每个满足关系 \(I\) 的有限维 \(k\) 表示都能分解为有限个不可分解表示的直和；各直和项的同构类型及重数唯一到重排。无关系情形取 \(I=0\)。
\end{proposition}
\begin{proof}
对给定表示 \(V\)，每次把一个非零可分解直和项 \(X\) 拆成 \(Y\oplus Z\)，其中 \(Y,Z\) 都非零，非零直和项的个数便增加一。各项的总维数之和始终为 \(\sum_{i\in Q_0}\dim_kV_i\)，而每个非零项至少占一维，所以项数不超过这个数。因此拆分过程有限步终止，得到存在性；零表示对应空直和。
为证唯一性，先看不可分解表示 \(X\) 的自同态 \(u\)。核与像的升降链在有限维下稳定，取充分大的 \(N\)，便有
\[
X=\ker u^N\oplus\operatorname{im}u^N.
\]
确实，稳定性使 \(u\) 在稳定像上可逆，而该像与稳定核相交为零，维数相加为 \(\dim X\)；二者都是子表示。于是 \(u\) 或可逆，或幂零。其自同态环的非单位元构成理想：左右乘任意自同态不可能把不可逆线性算子变可逆；若两个非单位元 \(u,v\) 的和可逆，则
\(a=(u+v)^{-1}u\)、\(1-a=(u+v)^{-1}v\) 都非单位，但 \(a\) 幂零使 \(1-a\) 可逆，矛盾。
现在设 \(V=\bigoplus_iX_i=\bigoplus_jY_j\)。把 \(X_1\) 包含到 \(V\)，再分别投到各 \(Y_j\) 并投回 \(X_1\)，所得复合之和为 \(1_{X_1}\)。非单位元构成理想，所以至少一个复合可逆。相应映射 \(X_1\to Y_j\) 为分裂单射；\(Y_j\) 不可分解，故它是同构。更具体地，\(V\to X_1\) 在这个 \(Y_j\) 上为同构，所以
\(V=Y_j\oplus\bigoplus_{i>1}X_i\)：任意元可先减去唯一具有相同 \(X_1\) 投影的 \(Y_j\) 元。与原来的 \(Y_j\) 分解取商，得到
\(\bigoplus_{i>1}X_i\cong\bigoplus_{\ell\ne j}Y_\ell\)。
按直和项数归纳即得唯一性。
\end{proof}
此论证只使用有限维性，没有使用路径代数本身有限维。因此带圈的有限箭图，只要所讨论的表示有限维，也满足这个分解唯一性结论。

\subsection{区间表示的线性代数分解}
\label{b5:subsec:0903:03}

\begin{theorem}\label{b5:thm:linear-quiver-intervals}
设 \(k\) 为域，\(n\ge1\) 为整数，\(Q=(1\to2\to\cdots\to n)\)，\(V=(V_i,V_a)\) 为有限维 \(k\) 表示。则 \(V\) 是区间表示 \(I[a,b]\) 的直和，其中 \(1\le a\le b\le n\)。记
\[
r_{ij}=\operatorname{rank}_k(V_i\to V_j)\quad(1\le i\le j\le n),
\qquad r_{ii}=\dim_k V_i,
\]
其中 \(V_i\to V_j\) 为沿唯一有向路径的箭映射复合，并约定 \(r_{0j}=r_{i,n+1}=0\)，则 \(I[a,b]\) 的重数为
\[
m_{ab}=r_{ab}-r_{a-1,b}-r_{a,b+1}+r_{a-1,b+1}.
\]
特别地，恰有 \(n(n+1)/2\) 个不可分解同构类。
\end{theorem}
\begin{proof}
对 \(n\) 归纳。\(n=1\) 是选基。假定已把前 \(n-1\) 个顶点分成区间。在顶点 \(n-1\) 已为零的区间不参与最后一支箭，可直接保留。其余各区间都结束于 \(n-1\)，把其末端空间记为 \(E=V_{n-1}\)。按这些区间的起点定义旗
\[
0=F_0\subseteq F_1\subseteq\cdots\subseteq F_{n-1}=E,
\]
其中 \(F_i\) 由起点不超过 \(i\) 的末端基向量生成，亦即 \(V_i\to E\) 的像。令 \(f:E\to V_n\) 为最后的箭，\(K=\ker f\)。可以选一组同时适应旗 \(F_i\) 与子空间 \(K\) 的基：逐步补全 \(K\cap F_i\) 的基，再补全商旗 \((F_i+K)/K\) 的基，并把后者的新增向量提升到 \(F_i\)。
任意保持这面旗的可逆换基都延伸为这些末端区间之和的自同构。为看清这一点，将起点为 \(a\) 的末端向量送到起点 \(b\le a\) 的末端向量，给出 \(I[a,n-1]\to I[b,n-1]\) 的同态：在重合区间取同一个标量，在其余顶点取零。旗保持的矩阵恰由这些允许的项组成；其逆也保持旗，故同样延拓，且两复合在每个顶点都是恒等。
因此可以用刚才选出的适应基代替原区间末端基。属于 \(K\) 的向量使相应区间仍结束于 \(n-1\)；其余向量的像线性无关，将它们作为 \(\operatorname{im}f\) 的基，相应区间便延长到 \(n\)。最后补全到 \(V_n\) 的基，新增的向量各给一个 \(I[n,n]\)。这完成存在性的归纳。
对任一区间和，复合 \(V_i\to V_j\) 在一个直和项上秩为一，当且仅当该区间同时包含 \(i,j\)。故
\[
r_{ij}=\sum_{\substack{a\le i\\b\ge j}}m_{ab}.
\]
先对起点作相邻差，再对终点作相邻差，就得到所列重数公式。每个区间不可分解且不同构已经证明；反过来，一个不可分解表示的区间分解只能有一项，因此类数为区间数。
\end{proof}
这个公式也给出判别同构的有限算法：计算各段复合的秩，再取四项差即可。仅计算相邻箭的秩仍不够。例如两个表示都可有维数向量 \((1,2,1)\) 及相邻秩 \(1,1\)，但复合秩分别为一和零；它们的区间分解因而不同。

% !TeX root = ../../Book5.tex
% 纲要标题：### 9.4 箭图表示的同调初步
% 纲要位置：工作记录/计划纲要.md，第 4029 行。

\section{箭图表示的同调初步}
\label{b5:sec:0904}


% 纲要范围：
%
% * 无环箭图的投射表示
% * 短投射消解与 Ext^1
% * 第10章根系分类的准备
% * 本节的 \(\operatorname{Ext}^1\) 可先按短正合扩张定义，并就地说明与短投射消解计算的一致性；不要求完整同调代数课程
%

区间分解说明线性定向箭图很容易分类，但一般箭图不能同时把所有箭化为这种形式。另一种有效做法是把一个表示写成若干简单的投射表示之间的商。无环箭图只需一层关系，因而可以用一条短消解同时计算同态与扩张。

\subsection{无环箭图的投射表示}
\label{b5:subsec:0904:01}

设 \(k\) 为域，\(Q\) 为有限箭图，\(P\) 为 \(Q\) 的有限维 \(k\) 表示。若对有限维 \(k\) 表示之间的每个满射同态 \(E\to V\)，任意表示同态 \(P\to V\) 都能提升到 \(P\to E\)，则称 \(P\) 为投射表示。这里满射按顶点理解；这与对应左模的投射性一致。
\begin{proposition}\label{b5:prop:vertex-projective}
设 \(k\) 为域，\(Q\) 为有限无环箭图，\(A=kQ\)，\(i\in Q_0\)，\(e_i\) 为顶点 \(i\) 的长度零路径。左 \(A\)-模 \(P_i=Ae_i\) 在顶点 \(j\) 以全部路径 \(i\to j\) 为基，箭通过在路径末端继续拼接而作用。对 \(Q\) 的任意有限维 \(k\) 表示 \(W\)，在 \(e_i\) 处求值给出自然同构
\[
\operatorname{Hom}_Q(P_i,W)\cong W_i,\qquad f\longmapsto f(e_i).
\]
因此 \(P_i\) 投射，且 \(\operatorname{End}_Q(P_i)\cong k\)，所以不可分解。
\end{proposition}
\begin{proof}
模同态满足 \(f(e_i)=e_if(e_i)\in W_i\)，又满足 \(f(p)=W_p f(e_i)\)。反过来，任给 \(w\in W_i\)，这个公式在路径基上线性延拓，便给出唯一模同态。该求值映射对 \(W\) 的同态相容，因而自然。顶点取值函子保持正合列，特别是满射，所以该 Hom 函子也保持满射，得到投射性。无环性使 \(i\to i\) 只有长度零路径；取 \(W=P_i\)，自同态由 \(e_i\) 的标量倍决定。其自同态环是域，因而无非平凡直和投影。
\end{proof}
对 \(1\to2\to\cdots\to n\)，有 \(P_i=I[i,n]\)。箭向改变时投射表示也改变：在 \(1\to2\leftarrow3\) 中，\(P_2\) 只支集在顶点2，\(P_1,P_3\) 则各沿一支箭延伸。判断投射性应看“从这个顶点出发的路径”，不是只凭支集大小。

\subsection{\texorpdfstring{短投射消解与 \(\operatorname{Ext}^1\)}{短投射消解与 Ext1}}
\label{b5:subsec:0904:02}

\begin{theorem}\label{b5:thm:acyclic-quiver-resolution}
设 \(k\) 为域，\(Q\) 为有限无环箭图，\(A=kQ\)，\(V=(V_i,V_a)\) 为 \(Q\) 的有限维 \(k\) 表示。记 \(e_i\) 为顶点 \(i\) 的长度零路径、\(P_i=Ae_i\)，\(V_p\) 为沿路径 \(p\) 的箭映射复合。存在左 \(A\)-模的短正合列
\[
0\longrightarrow
\bigoplus_{a:i\to j}P_j\otimes_k V_i
\xrightarrow{d}
\bigoplus_{i\in Q_0}P_i\otimes_k V_i
\xrightarrow{\epsilon}V\longrightarrow0.
\]
其中左 \(A\) 作用在每项的第一个因子上，\(\epsilon(p\otimes v)=V_pv\)。对于属于箭 \(a:i\to j\) 的左端分量，微分为
\[
d(p\otimes v)=pa\otimes v-p\otimes V_av,
\]
第一项放在中项的 \(i\) 分量，第二项放在 \(j\) 分量。两项直和都是有限维投射表示。
\end{theorem}
\begin{proof}
\(\epsilon\) 与 \(d\) 都为 \(A\) 线性，且 \(\epsilon d=0\)，因为路径 \(pa\) 的作用是先 \(V_a\) 再 \(V_p\)。每个顶点的向量由 \(e_i\otimes v\) 提升，所以 \(\epsilon\) 满。
证明 \(d\) 单射时，对左端的一个非零元素按路径 \(p\) 收集其向量系数，选出现的最大路径长度 \(L\)。正项 \(pa\otimes v\) 长度为 \(L+1\)，负项长度至多 \(L\)。一条正长度路径唯一决定其第一支箭 \(a\) 及剩余路径 \(p\)，所以这些最高长度正项之间也不能相消。因此像非零。
在中项模掉 \(\operatorname{im}d\) 后，关系
\(pa\otimes v=p\otimes V_av\)
可逐次删去路径的第一支箭，把 \(p\otimes v\) 化为 \(e_{t(p)}\otimes V_pv\)。于是每个商类都是长度零项的和；若其 \(\epsilon\) 像为零，各顶点的向量系数都为零，商类也为零。故 \(\ker\epsilon=\operatorname{im}d\)。最后由\cref{b5:prop:vertex-projective}，各项是有限个 \(P_i\) 的直和，因而投射。
\end{proof}
为了说明这条消解怎样计算扩张，先把本节所用的 \(\operatorname{Ext}^1\) 定义清楚。对有限维 \(A\) 模 \(V,W\)，考虑短正合列
\[
0\longrightarrow W\longrightarrow E\longrightarrow V\longrightarrow0.
\]
两列若有一个中项同构使两端恒等及相应图表交换，就看成同一个扩张类，记全体类为 \(\operatorname{Ext}_A^1(V,W)\)。零类是分裂列。同样的定义也适用于一般左模；本章的主要应用仍取有限维对象。两个扩张先沿 \(V\to V\oplus V\) 拉回，再沿 \(W\oplus W\to W\) 的相加映射推出，定义它们的和；沿标量映射 \(W\to W\) 推出定义标量乘法。下面的具体描述同时证明这些运算给出 \(k\) 向量空间。
\begin{proposition}\label{b5:prop:projective-presentation-extensions}
设 \(k\) 为域，\(A\) 为含幺 \(k\) 代数，\(V,W\) 为左 \(A\)-模，并给定左 \(A\)-模的短正合列
\(0\to K\xrightarrow{\iota}P\to V\to0\)，其中 \(P\) 投射。则有自然同构
\[
\operatorname{Ext}_A^1(V,W)\cong
\operatorname{coker}\left(
\operatorname{Hom}_A(P,W)\xrightarrow{\iota^*}
\operatorname{Hom}_A(K,W)\right).
\]
这里 \(\iota^*\) 把同态 \(f:P\to W\) 送到 \(f\circ\iota:K\to W\)，\(\operatorname{Ext}_A^1(V,W)\) 取短正合扩张类及其 Baer 和。
\end{proposition}
\begin{proof}
给定 \(f:K\to W\)，取
\[
E_f=(W\oplus P)/\Bigl\{\,\bigl(f(x),-\iota(x)\bigr):x\in K\,\Bigr\}.
\]
映射 \(w\mapsto\bigl[(w,0)\bigr]\) 单射，商到 \(V\) 的映射为 \(\bigl[(w,p)\bigr]\mapsto\bar p\)，其核正是这份 \(W\)，所以得到扩张。反过来，对任一扩张，利用 \(P\) 的投射性将 \(P\to V\) 提升为 \(u:P\to E\)。其在 \(K\) 上的限制落入 \(W\)，给出 \(f\)；映射 \((w,p)\mapsto w+u(p)\) 诱导 \(E_f\cong E\)。
两个提升的差是 \(P\to W\)，故限制到 \(K\) 后只改变一个来自 \(\operatorname{Hom}_A(P,W)\) 的项。反过来，若两个 \(f\) 的差可延伸到 \(P\)，用这个延伸修正 \(W\oplus P\) 的第一坐标，便得到固定两端的扩张同构。因此商空间恰好记录扩张类。按定义作拉回和推出时，两个限制映射相加，标量推出则把限制映射乘相应标量；所以对应保持上述运算，自然性也由同样的构造得到。
\end{proof}

\subsection{同调方法与箭图表示的分类}
\label{b5:subsec:0904:03}

\begin{theorem}\label{b5:thm:quiver-hom-ext-euler}
设 \(k\) 为域，\(Q\) 为有限无环箭图，\(V=(V_i,V_a)\)、\(W=(W_i,W_a)\) 为有限维 \(k\) 表示。记 \(\operatorname{Ext}_Q^1(V,W)=\operatorname{Ext}_{kQ}^1(V,W)\)，其中两表示视作左 \(kQ\)-模。存在自然正合列
\[
\begin{split}
0\longrightarrow\operatorname{Hom}_Q(V,W)
&\longrightarrow\bigoplus_{i\in Q_0}\operatorname{Hom}_k(V_i,W_i)\\
&\xrightarrow{\delta}
\bigoplus_{a:i\to j}\operatorname{Hom}_k(V_i,W_j)
\longrightarrow\operatorname{Ext}_Q^1(V,W)\longrightarrow0,
\end{split}
\]
其中 \(\delta\bigl((f_i)_i\bigr)=(W_af_i-f_jV_a)_{a:i\to j}\)。若 \(d_i=\dim_k V_i\)、\(e_i=\dim_k W_i\)，则
\[
\dim_k\operatorname{Hom}_Q(V,W)-\dim_k\operatorname{Ext}_Q^1(V,W)
=\sum_i d_i e_i-\sum_{a:i\to j}d_i e_j.
\]
\end{theorem}
\begin{proof}
将\cref{b5:prop:projective-presentation-extensions}用于\cref{b5:thm:acyclic-quiver-resolution}。由\cref{b5:prop:vertex-projective}及张量的线性泛性质，
\[
\operatorname{Hom}_Q(P_i\otimes_k V_i,W)
\cong\operatorname{Hom}_k(V_i,W_i).
\]
将同态应用于 \(d(p\otimes v)\)，再令 \(p=e_j\)，便得到 \(W_af_i-f_jV_a\)。这个映射的核正是表示同态的交换条件，而余核由前一命题计算扩张，所以所列序列正合。所有空间有限维，维数的交错和为零，得到公式。
\end{proof}
若 \(S_i\) 表示只在顶点 \(i\) 放置 \(k\) 的表示，则对于不同顶点 \(i,j\)，有
\[
\dim\operatorname{Ext}_Q^1(S_i,S_j)
=\#\{\,a:i\to j\,\}.
\]
这里右侧的方向与左模约定一致：在 \(1\to2\) 中，非分裂列是
\(0\to S_2\to I[1,2]\to S_1\to0\)。
\ProseParagraphBreak
令 \(V=W\)，公式左侧比较自同态与自扩张，右侧只含维数向量。下一章将把这个整数二次式与反射联系起来。若采用第四卷的高次 Ext 定义，上面的长度一投射消解还直接给出所有 \(\operatorname{Ext}_{kQ}^r(V,W)=0\)、\(r\ge2\)。加上关系以后，长度一消解一般不再成立。

\subsection{短正合扩张与消解计算的一致性}
\label{b5:subsec:0904:04}

上述余核还能从扩张中的箭矩阵直接读出。给定
\(0\to W\to E\to V\to0\)，在每个顶点选择向量空间截面，识别
\(E_i=W_i\oplus V_i\)。箭矩阵必为
\[
E_a=
\begin{pmatrix}
W_a&\xi_a\\0&V_a
\end{pmatrix},
\qquad \xi_a:V_i\to W_j\quad(a:i\to j).
\]
没有关系时，各个 \(\xi_a\) 可以任意选择。将截面 \(s_i\) 换成
\(s_i+f_i\)、\(f_i:V_i\to W_i\)，右上角变为
\[
\xi'_a=\xi_a+W_af_i-f_jV_a.
\]
这是将 \(E_a\bigl(s_i(v)+f_i(v)\bigr)\) 减去 \(s_j(V_av)+f_j(V_av)\) 的直接计算。因此两组矩阵表示相同的扩张类，当且仅当相差 \(\delta\bigl((f_i)_i\bigr)\)；中项同构固定两端时也恰有这种分块上三角形式。
\ProseParagraphBreak
这与消解所得的对应一致，而不只是维数相同。由各顶点的截面构造中项到 \(E\) 的提升
\[
P_i\otimes V_i\longrightarrow E,\qquad p\otimes v\longmapsto E_p s_i(v).
\]
在箭 \(a:i\to j\) 所对应的关系 \(a\otimes v-e_j\otimes V_av\) 上，它的值为
\(E_as_i(v)-s_j(V_av)=\xi_a(v)\)，正是前一命题取到的限制映射。
\begin{example}\label{b5:exm:arrow-extension-parameter}
设 \(k\) 为域，\(Q=(1\to2)\)。扩张 \(0\to S_2\to E\to S_1\to0\) 由一个标量 \(c\) 给出，箭映射为 \(k\xrightarrow{c}k\)。此时两个顶点的 \(\operatorname{Hom}_k(V_i,W_i)\) 都为零，所以不同 \(c\) 给出不同扩张类，\(\operatorname{Ext}_Q^1(S_1,S_2)\cong k\)。然而所有 \(c\ne0\) 的中项都同构于 \(I[1,2]\)。扩张等价要求固定两端的指定映射，因此比只比较中项同构更严格。
\end{example}
对于有关系箭图，还须把分块矩阵代入关系。比如 \(1\xrightarrow{a}2\xrightarrow{b}3\) 加上 \(ba=0\) 时，右上角应满足
\(W_b\xi_a+\xi_bV_a=0\)，并非每组 \(\xi_a,\xi_b\) 都可用。这里已经能看见关系怎样限制扩张；不能继续把无关系箭图的 Euler 公式\footnote{欧拉（Leonhard Euler，1707-1783），瑞士数学家，研究数论、分析与图论，奠定多项基础。}原样照搬。

% !TeX root = ../../Book5.tex
% 纲要标题：### 9.5 基本代数、Ext 箭图与关系
% 纲要位置：工作记录/计划纲要.md，第 4036 行。

\section{基本代数、Ext 箭图与关系}
\label{b5:sec:0905}


% 纲要范围：
%
% * 对代数闭域上的有限维代数，调用第二卷的投射覆盖和 Morita 理论构造基本代数，并说明分裂基域假设的用途
% * 从本原幂等元及 \(J/J^2\) 构造箭图，按第9.2节既定的路径和左模约定核验箭的方向及 \(\operatorname{Ext}^1\) 维数
% * 证明基本代数可写成路径代数模容许关系理想，并用箭理想的幂界解释容许性和有限维性
% * 以三角矩阵代数、双数代数及一个带关系例子连接第二卷的根滤过；一般有关系代数不直接适用第10章无关系箭图的分类
% * 从第9章可直接进入第10章；所需 ADE 二次型和反射将在第10.1节就地建立，不要求先读 Lie 代数章节
%
% ---
%

前面从箭图造出代数。现在反过来：一个有限维代数的简单模之间有哪些最初的联系，能否由箭表示？答案要去掉半单商中重复的矩阵大小，以根的第一层作箭，以更高层的乘法作关系。代数闭基域在这一步有明确作用。本节把讨论扩大到带关系的有限维代数，需要第二卷的根、幂等元提升与 Morita 等价；第10.1—10.3节关于 Euler 型、反射函子与 Gabriel 定理的证明仍以第9.1—9.4节的无关系箭图为基础，不使用这里的约化与展示。

\subsection{基本代数与分裂基域假设}
\label{b5:subsec:0905:01}

设 \(A\) 为代数闭域 \(k\) 上的有限维含幺代数，记 \(J=J(A)\)。第二卷定理15.1.4证明了 \(J\) 幂零；半单结构定理把商写为
\[
A/J\cong\prod_{i=1}^r M_{n_i}(k).
\]
每个简单模的自同态除环在代数闭域上为 \(k\)，所以这里没有额外的除代数系数。第二卷定理15.3.4又允许将各矩阵块的对角矩阵单位提升为一组完整正交幂等元。
\begin{definition}\label{b5:def:split-basic-algebra}
设 \(k\) 为域，\(B\) 为非零有限维含幺 \(k\) 代数。若
\(B/J(B)\cong k^r\) 对某个 \(r\ge1\) 成立，就称 \(B\) 为一个分裂\term[jiben-daishu]{基本代数}[basic algebra]。在代数闭域上，“分裂”自动满足；一般基本代数的半单商可以是若干不同除代数的直积，本节不把它们都换成 \(k\)。
\end{definition}
\begin{proposition}\label{b5:prop:basic-morita-reduction}
设 \(k\) 为代数闭域，\(A\) 为非零有限维含幺 \(k\) 代数。可取一个幂等元 \(e\in A\)，使 \(AeA=A\)，且 \(B=eAe\) 是分裂基本代数。将 \(Ae\) 按乘法看成 \((A,B)\)-双模，则函子
\[
M\longmapsto eM,\qquad N\longmapsto Ae\otimes_B N
\]
给出左 \(A\) 模与左 \(B\) 模的等价，并保持有限维性。
\end{proposition}
\begin{proof}
在每个 \(M_{n_i}(k)\) 块中选一个秩一对角幂等元，取相容提升 \(e_i\)，并令 \(e=\sum_i e_i\)。第二卷命题15.4.5关于基本角代数的结论在这里给出
\[
J(eAe)=eJe,\qquad eAe/eJe\cong k^r.
\]
这些条件也可直接检查：\(eJe\) 幂零，商角环恰是在每个矩阵块保留一个对角位置，因此半单且为 \(k^r\)。
商代数中每个所选秩一幂等元生成整个对应矩阵块，故 \(AeA+J=A\)。令 \(M=A/AeA\)，则 \(JM=M\)；迭代得到 \(M=J^NM=0\)，所以 \(AeA=A\)。这验证了第二卷命题16.3.4的满幂等元条件，该命题用乘法映射
\(Ae\otimes_B eM\to M\)、\(ae\otimes m\mapsto am\)
和 \(eAe\otimes_BN\to N\) 证明所列函子互为拟逆。这里 \(Ae\) 有限维，故两个函子也保持有限维对象。
\end{proof}
非代数闭域上的限制不能忽略。例如实代数 \(\mathbb C\) 已是除代数，只有一个简单模类型；一个普通单顶点、无箭的实箭图只给出 \(\mathbb R\)，不能用它记录 \(\mathbb C\) 的除代数结构。本节的普通箭图结论以分裂条件为基础。

\subsection{根商构造的箭图、箭向与扩张维数}
\label{b5:subsec:0905:02}

以下设 \(B/J\cong k^r\)，并选取提升标准坐标元的一组完整正交幂等元 \(e_1,\ldots,e_r\)，其中 \(J=J(B)\)。简单模为 \(S_i=Be_i/Je_i\)，每个都一维。选取每个空间
\[
e_j(J/J^2)e_i
\]
的一组 \(k\) 基，对每个基元画一支箭 \(i\to j\)。这里包括 \(i=j\) 的圈。根只保留第一层，是因为 \(J^2\) 中的元素已经可能由两支箭的复合产生，不应当再次作为独立箭。
\begin{proposition}\label{b5:prop:ext-quiver-direction}
设 \(k\) 为域，\(B\) 为非零有限维含幺分裂基本 \(k\) 代数，\(J=J(B)\)，\(B/J\cong k^r\)。取一组完整正交幂等元 \(e_1,\ldots,e_r\in B\)，其模 \(J\) 的像分别为 \(k^r\) 的标准坐标幂等元，令简单左模 \(S_i=Be_i/Je_i\)。则对 \(1\le i,j\le r\) 有
\[
\operatorname{Ext}_B^1(S_i,S_j)
\cong \operatorname{Hom}_k\bigl(e_j(J/J^2)e_i,k\bigr).
\]
以 \(1,\ldots,r\) 为顶点，对每个 \(e_j(J/J^2)e_i\) 的基元画一支箭 \(i\to j\)，则所得箭图中从 \(i\) 到 \(j\) 的箭数为 \(\dim_k\operatorname{Ext}_B^1(S_i,S_j)\)。
\end{proposition}
\begin{proof}
\(Be_i\) 是正则模的直和因子，因而投射。将\cref{b5:prop:projective-presentation-extensions}用于
\(0\to Je_i\to Be_i\to S_i\to0\)。
从 \(Be_i\) 到 \(S_j\) 的任何模同态都把 \(Je_i\) 送到 \(JS_j=0\)，故该 Hom 映射的限制为零。于是
\[
\operatorname{Ext}_B^1(S_i,S_j)\cong\operatorname{Hom}_B(Je_i,S_j).
\]
右侧的同态又零化 \(J^2e_i\)，所以可在 \(Je_i/J^2e_i\) 上计算。这个模被 \(J\) 零化，是 \(k^r\) 模；映到 \(S_j\) 的同态只保留第 \(j\) 个坐标，并在该坐标上为任意 \(k\) 线性泛函。选 \(e_j\) 在 \(S_j\) 中的像为基，便得到所列同构。
\end{proof}
这个箭图称为 \(B\) 的\term[Ext-jiantu]{Ext 箭图}[Ext quiver]。箭的具体基依赖选择，但每一对顶点间的箭数由扩张空间决定。注意顺序是 \(\operatorname{Ext}^1(S_i,S_j)\) 对应 \(i\to j\)：其扩张以 \(S_j\) 为子模、\(S_i\) 为商模。

\subsection{容许关系、路径代数与有限维性}
\label{b5:subsec:0905:03}

\begin{theorem}\label{b5:thm:basic-algebra-bound-quiver}
设 \(k\) 为域，\(B\) 为非零有限维含幺分裂基本 \(k\) 代数，\(J=J(B)\)，\(B/J\cong k^r\)。取一组完整正交幂等元 \(e_1,\ldots,e_r\)，其模 \(J\) 的像为各标准坐标幂等元。以 \(1,\ldots,r\) 为顶点，对每个 \(e_j(J/J^2)e_i\) 的一组 \(k\) 基中的元素画一支箭 \(i\to j\)，构成有限箭图 \(Q\)。则存在 \(kQ\) 的容许关系理想 \(I\)，使
\[
B\cong kQ/I.
\]
特别地，每个代数闭域上的非零有限维代数都与某个这种商代数 Morita 等价。
\end{theorem}
\begin{proof}
记 \(J=J(B)\)。将顶点送到选定的 \(e_i\)，将每支箭 \(i\to j\) 送到其基元的某个提升 \(x_a\in e_jJe_i\)。顶点关系成立，由路径代数的构造得到同态 \(\phi:kQ\to B\)。令 \(C\) 为这些像生成的子代数，\(L\) 为箭像的线性生成空间。基的选择给出 \(J=L+J^2\)，而顶点像生成 \(B/J\)，故 \(B=C+J\)。
将 \(J=L+J^2\) 的乘积展开，可得
\(J^m\subseteq L^m+J^{m+1}\subseteq C+J^{m+1}\)。
逐层代入 \(B=C+J\)，最终用 \(J^N=0\) 得到 \(B=C\)。所以 \(\phi\) 满射。
令 \(I=\ker\phi\)，并令 \(\mathfrak a\) 为箭理想。因每支箭映入 \(J\)，有 \(\mathfrak a^N\subseteq I\)，必要时增大 \(N\) 使 \(N\ge2\)。为证 \(I\subseteq\mathfrak a^2\)，把核中一个元素拆成长度零项、长度一项及较长路径。其像模 \(J\) 为零，迫使全部顶点系数为零；再模 \(J^2\)，箭像的基独立性迫使全部长度一系数为零。故它只含长度至少二的路径。所求理想容许，第一同构定理给出商表示。最后结合\cref{b5:prop:basic-morita-reduction}，得到一般代数闭基域情形。
\end{proof}
关系理想通常不唯一，因为箭的提升与基可以改变。定理给出的是一种适合模范畴的展示，并不宣称所有代数都有唯一的一组路径关系。这里的生成性也用到了根的幂零性；若只知道根的第一层而没有这种有限性，就不能用同样的有限步推导结束。

\subsection{三角矩阵、双数与带关系代数}
\label{b5:subsec:0905:04}

上三角矩阵代数 \(T_3(k)\) 的根由 \(E_{12},E_{23},E_{13}\) 生成，根平方由 \(E_{13}=E_{12}E_{23}\) 生成。因此箭图只有 \(3\to2\to1\) 的两支箭；其六条路径对应六个上三角矩阵单位，不必再加入关系。这也复核了左模约定下的箭向。
\ProseParagraphBreak
双数代数 \(k[t]/(t^2)\) 则只有一个顶点和一支圈，关系为圈的平方为零。它的简单模 \(S=k\) 满足
\[
\operatorname{Ext}^1(S,S)\cong k,
\]
因为 \(J/J^2=kt\)。同样的单圈若改加 \(t^3=0\)，Ext 箭图仍相同，代数维数却从二变成三，最长不可分解的幂零链也从二变成三。箭图本身没有记录全部乘法，关系不可省去。
\begin{example}\label{b5:exm:zero-composition-bound-quiver}
设 \(k\) 为域，\(Q=(1\xrightarrow{a}2\xrightarrow{b}3)\)，\(B=kQ/(ba)\)。其基为 \(e_1,e_2,e_3,a,b\)，根平方为零，Ext 箭图仍是原来的三顶点直线。然而维数向量 \((1,1,1)\) 不再有不可分解表示：两支箭由标量 \(\alpha,\beta\) 给出，关系要求 \(\beta\alpha=0\)，故至少一支箭为零，表示沿该处拆成直和。在无关系的 \(kQ\) 上，\(I[1,3]\) 则是这个维数向量的不可分解表示。
\end{example}
因此下一章关于无关系路径代数的根分类，不能直接用于具有同一箭图的任意商代数。

\subsection{ADE 二次型与反射简介}
\label{b5:subsec:0905:05}

对有限无环箭图，\cref{b5:thm:quiver-hom-ext-euler}已经给出
\[
q_Q(\mathbf d)=\sum_i d_i^2-\sum_{a:i\to j}d_i d_j
=\dim\operatorname{End}_Q(V)-\dim\operatorname{Ext}_Q^1(V,V),
\quad \mathbf d=\mathbf d(V).
\]
右边来自表示的线性方程，左边却只用图和整数向量。在线性定向 \(A_n\) 中，一个区间向量含有 \(b-a+1\) 个一，区间内有 \(b-a\) 支箭，所以其值为一。我们已经用选基法证明了这些区间给出全部不可分解表示；这个数值因而有实际分类意义。
\ProseParagraphBreak
其他无向图是否也有类似现象，需要继续研究 \(q_Q\) 的正定性以及改变箭向的办法。\chapref{b5:ch:10}将定义整数格上的反射，证明所需的 ADE 图判别，再构造反射函子；这些论证不依赖 Lie 代数的根系分类。对于带关系代数，\cref{b5:exm:zero-composition-bound-quiver}已经说明，仅仅保留同一个图上的二次式还不足以恢复不可分解表示。


\begin{chaptersummary}
路径乘法采用“先右后左”的复合次序，因此箭 \(a:i\to j\) 满足 \(a=e_jae_i\)，在左模上从 \(e_iM\) 映到 \(e_jM\)。顶点分解把左模与箭图表示联系起来，箭上的交换条件则把模同态化为一组线性方程。这个对应同时保留直和、核、余核和扩张。

\ProseParagraphBreak

在线性定向的 \(A_n\) 上，区间表示给出全部不可分解对象，各段复合的秩决定区间重数。一般无环箭图虽然未必有这种简单分类，仍有长度一的投射消解。它把 \(\operatorname{Hom}\) 与 \(\operatorname{Ext}^1\) 分别表示为同一个线性映射的核与余核；扩张中的分块箭矩阵又解释了这一余核的含义。扩张以 \(S_j\) 为子模、\(S_i\) 为商模时，对应的是从 \(i\) 指向 \(j\) 的箭。

\ProseParagraphBreak

反过来，有限维分裂基本代数的根第一层 \(J/J^2\) 决定箭，而根的更高乘法决定关系。代数闭域上的一般有限维代数先经满幂等元约化为基本代数，再得到路径代数模容许关系的展示。相同的 Ext 箭图可以承载不同的关系和不同的表示分类；下一章的 Gabriel 定理只处理没有附加关系的无环箭图。
\end{chaptersummary}

\BookExercises

\begin{exercise}\label{b5:ex:09-diamond-paths}
设箭图有四个顶点及四支箭
\(a:1\to2\)、\(b:1\to3\)、\(c:2\to4\)、\(d:3\to4\)。
列出路径代数的基与维数。再取 \(B=kQ/(ca-db)\)，证明关系理想容许，并求 \(B\)、\(J(B)\)、\(J(B)^2\) 的维数。
\end{exercise}
\begin{solution}
全部路径为
\[
e_1,e_2,e_3,e_4,\ a,b,c,d,\ ca,db,
\]
所以 \(\dim kQ=10\)。没有长度三路径，故箭理想 \(\mathfrak a\) 满足 \(\mathfrak a^3=0\)。关系 \(ca-db\) 位于 \(\mathfrak a^2\)，且不能再与正长度路径非零拼接，所以它生成的一维理想正是 \(k(ca-db)\)，容许条件 \(\mathfrak a^3\subseteq I\subseteq\mathfrak a^2\) 成立。商中两条长度二路径合为一个非零基向量，故 \(\dim B=9\)。根由四支箭及这个长度二路径生成，维数为 \(5\)；根平方由后者生成，维数为 \(1\)，根立方为零。
\end{solution}

\begin{exercise}\label{b5:ex:09-triangular-all-sizes}
设 \(Q\) 是箭图 \(n\to n-1\to\cdots\to1\)。证明 \(kQ\cong T_n(k)\)，其中 \(T_n(k)\) 为上三角矩阵代数，并说明为什么在本章左模约定下应取这个箭向。
\end{exercise}
\begin{solution}
把 \(e_i\) 送到 \(E_{ii}\)，把箭 \(i\to i-1\) 送到 \(E_{i-1,i}\)。从 \(i\) 到 \(j\leq i\) 的唯一路径被送到 \(E_{ji}\)，包括 \(j=i\) 的长度零路径。矩阵单位满足
\(E_{\ell j}E_{ji}=E_{\ell i}\)，不可拼接时乘积为零，故这定义代数同态。全部路径与全部上三角矩阵单位一一对应，所以同态为同构。
左乘 \(E_{ji}=E_{jj}E_{ji}E_{ii}\) 将左模的第 \(i\) 个顶点空间送到第 \(j\) 个顶点空间；因此矩阵下标与箭向正好相符。若不同时改变顶点编号和约定，不能把箭向直接反过来。
\end{solution}

\begin{exercise}\label{b5:ex:09-truncated-cycle}
设 \(Q\) 有两顶点及箭 \(a:1\to2\)、\(b:2\to1\)，令 \(\mathfrak a\) 为箭理想，\(B=kQ/\mathfrak a^3\)。求 \(B\) 的路径基、根的各次幂，以及左投射模 \(Be_1,Be_2\) 的维数向量。
\end{exercise}
\begin{solution}
商中保留长度零、一、二的路径，基为
\(e_1,e_2,a,b,ba,ab\)，所以维数为 \(6\)。
\(J(B)\) 由 \(a,b,ba,ab\) 张成，\(J(B)^2\) 由 \(ba,ab\) 张成，\(J(B)^3=0\)。
\(Be_1\) 以从顶点 \(1\) 出发的路径 \(e_1,a,ba\) 为基，分别终止于 \(1,2,1\)，故维数向量为 \((2,1)\)；同理 \(Be_2\) 的基为 \(e_2,b,ab\)，向量为 \((1,2)\)。原箭图有有向圈，路径代数无限维，但容许关系使商代数有限维。
\end{solution}

\begin{exercise}\label{b5:ex:09-nonadmissible-square}
单圈箭图给出 \(k[t]\)。讨论二维商 \(A=k[t]/(t^2-1)\)：分别在 \(\operatorname{char}k\neq2\) 与 \(\operatorname{char}k=2\) 时求其根和一个具有容许关系的箭图展示，并解释原关系为何不容许。
\end{exercise}
\begin{solution}
原箭理想为 \((t)\)，而 \(t^2-1\notin(t^2)\)，所以原理想不容许。特征不为 \(2\) 时，互素因子 \(t-1,t+1\) 给出
\(A\cong k\times k\)，根为零，改用两个孤立顶点、零关系理想即可。
特征为 \(2\) 时，令 \(u=t-1\)，有 \(t^2-1=u^2\)，所以
\(A\cong k[u]/(u^2)\)，根为 \(ku\)，根平方为零。此时用一个顶点、一支新圈 \(u\) 及关系 \(u^2=0\)，便得到容许展示。代数的有限维性不保证任意一组给定生成元都适合作为根第一层的箭。
\end{solution}

\begin{exercise}\label{b5:ex:09-square-change-basis}
沿\cref{b5:ex:09-diamond-paths}的四支箭分别放置非零标量 \(a,b,c,d\)，各顶点空间均为 \(k\)。证明无关系时，同构类由 \(\lambda=ca/(db)\) 决定；加入关系 \(ca-db=0\) 后，这些表示只剩一个同构类。
\end{exercise}
\begin{solution}
顶点换基为 \(g_1,g_2,g_3,g_4\in k^\times\)。两条复合分别变为
\(g_4(ca)g_1^{-1}\) 和 \(g_4(db)g_1^{-1}\)，所以其比值不变。取
\(g_1=1,g_2=a^{-1},g_3=b^{-1},g_4=(ca)^{-1}\)，前三支箭 \(a,b,c\) 全变为 \(1\)，第四支箭变为 \(db/(ca)=\lambda^{-1}\)。因此相同比值足以保证同构，不同比值不可能同构。关系恰好要求 \(\lambda=1\)，所以有关系时得到唯一的全非零类型。
\end{solution}

\begin{exercise}\label{b5:ex:09-dual-opposite-quiver}
设 \(V\) 为有限箭图 \(Q\) 的有限维表示。证明取各顶点对偶空间和各箭的转置映射给出反向箭图 \(Q^{\mathrm{op}}\) 的表示，并说明它与把有限维左 \(kQ\)-模取 \(k\)-对偶所得右模相符。
\end{exercise}
\begin{solution}
箭 \(a:i\to j\) 的对偶为 \(V_a^*:V_j^*\to V_i^*\)，方向反转。复合满足 \((V_bV_a)^*=V_a^*V_b^*\)，正好对应反向路径。若 \(M\) 为左模，在 \(M^*=\operatorname{Hom}_k(M,k)\) 上定义
\[
(\ell\cdot x)(m)=\ell(xm)\qquad(x\in kQ).
\]
则 \((\ell\cdot x)\cdot y=\ell\cdot(xy)\)，所以得到右模。其第 \(i\) 个顶点空间为 \(M^*e_i\cong(e_iM)^*\)，右乘 \(a=e_jae_i\) 从第 \(j\) 个空间到第 \(i\) 个空间，正是上述转置箭。态射对偶反转方向；有限维向量空间的对偶保持短正合列并反转箭头，故两种描述也在态射及正合列上相符。
\end{solution}

\begin{exercise}\label{b5:ex:09-kernel-cokernel}
在 \(Q:1\to2\) 上，取
\[
V=\left(k^2\xrightarrow{(x,y)\mapsto x}k\right),\qquad
W=\left(k\xrightarrow{z\mapsto(z,0)}k^2\right).
\]
令 \(f_1(x,y)=x\)、\(f_2(z)=(z,0)\)。验证 \(f=(f_1,f_2)\) 为表示同态，求其核、像和余核，并判断它是否为单射或满射。
\end{exercise}
\begin{solution}
两条复合都把 \((x,y)\) 送到 \((x,0)\)，故交换条件成立。核在顶点 \(1\) 为 \(k(0,1)\)、顶点 \(2\) 为零，所以是 \(S_1\)。像在两顶点分别为 \(k\) 和 \(k(1,0)\)，箭为同构，所以同构于 \(I[1,2]\)。余核只在顶点 \(2\) 有一维空间，因而为 \(S_2\)。核与余核都非零，故 \(f\) 既不是单射也不是满射；某一顶点上的满射或单射不足以替代所有顶点上的条件。
\end{solution}

\begin{exercise}\label{b5:ex:09-interval-hom}
设 \(k\) 为域，\(n\ge1\) 为整数，\(Q=(1\to2\to\cdots\to n)\)，\(1\le a\le b\le n\)、\(1\le c\le d\le n\)，\(I[a,b],I[c,d]\) 为 \(k\) 上的区间表示。证明
\[
\dim_k\operatorname{Hom}_Q\bigl(I[a,b],I[c,d]\bigr)
=\begin{cases}
1,&c\leq a\leq d\leq b,\\
0,&\text{其他情形}.
\end{cases}
\]
在非零情形写出一个基态射。
\end{exercise}
\begin{solution}
两区间不相交时，逐顶点的映射全为零。若相交，重合区间内部的恒等箭迫使全部分量为同一个标量 \(\lambda\)。若 \(a<c\)，在目标首次出现的箭 \(c-1\to c\) 上，源端的映射为恒等而目标从零空间进入，交换条件迫使 \(\lambda=0\)。若 \(b<d\)，在源最后出现后的箭 \(b\to b+1\) 上，同样迫使 \(\lambda=0\)。非零态射因此要求 \(c\leq a\)、\(d\leq b\)，连同重合条件得到 \(c\leq a\leq d\leq b\)。
这些条件成立时，在重合顶点 \(a,\ldots,d\) 取恒等映射，其余分量取零。内部箭交换；两端边界处要么两边复合都为零，要么仍是同一恒等复合，故得到态射。任意态射都是它的标量倍，所以空间恰为一维。
\end{solution}

\begin{exercise}\label{b5:ex:09-interval-endomorphism-algebra}
对 \(Q:1\to2\to3\)，令
\(U=I[1,3]\)、\(V=I[2,3]\)、\(W=I[2,2]\)。
求 \(\operatorname{End}_Q(U\oplus V\oplus W)\) 的矩阵形式、维数及 Jacobson 根。
\end{exercise}
\begin{solution}
由上一题，除了三个恒等类型的 Hom 外，只有
\(\operatorname{Hom}(V,U)\cong k\) 和
\(\operatorname{Hom}(V,W)\cong k\) 非零。
分别选择重合区间上的恒等映射为基，以 \(U,V,W\) 排列直和项，自同态代数可写为
\[
\left\{
\begin{pmatrix}
a&b&0\\0&c&0\\0&d&e
\end{pmatrix}:a,b,c,d,e\in k
\right\}.
\]
维数为 \(5\)。令三个对角元全为零所得二维理想的平方为零，模掉它得到 \(k^3\)。幂零理想包含在根中，半单商又使根不可能更大，所以该理想正是 Jacobson 根。这也说明各直和项的自同态虽只有标量，整个直和的自同态代数仍含有非对角的态射信息。
\end{solution}

\begin{exercise}\label{b5:ex:09-a3-matrix-decomposition}
在 \(1\to2\to3\) 上取维数向量 \((3,4,2)\)，箭矩阵为
\[
A=\begin{pmatrix}1&0&0\\0&1&0\\0&0&0\\0&0&0\end{pmatrix},
\qquad
B=\begin{pmatrix}1&0&0&0\\0&0&1&0\end{pmatrix}.
\]
求区间分解，并说明为何只知道两支箭各自的秩仍不够。
\end{exercise}
\begin{solution}
有 \(\operatorname{rank}A=2\)、\(\operatorname{rank}B=2\)、\(\operatorname{rank}(BA)=1\)。代入区间重数公式，得到
\[
I[1,3]\oplus I[1,2]\oplus I[2,3]\oplus I[1,1]\oplus I[2,2].
\]
若保持 \(A\) 不变，将 \(B\) 换成选取第三、第四坐标的映射，则 \(B\) 的秩仍为 \(2\)，但 \(BA=0\)；所得分解为
\(I[1,2]^{\oplus2}\oplus I[2,3]^{\oplus2}\oplus I[1,1]\)。
两个表示的维数向量及相邻箭秩完全相同，区间分解却不同。
\end{solution}

\begin{exercise}\label{b5:ex:09-a4-rank-table}
对线性定向 \(A_4\)，设各段复合的秩为
\[
(r_{ij})_{i\leq j}=
\begin{pmatrix}
2&2&1&0\\
 &3&2&1\\
 & &3&2\\
 & & &2
\end{pmatrix}.
\]
求全部非零区间重数，并据此说明这些秩确实能够由一个表示实现。
\end{exercise}
\begin{solution}
用\cref{b5:thm:linear-quiver-intervals}的四项差公式，非零重数恰为
\[
m_{12}=m_{13}=m_{24}=m_{34}=1.
\]
因此取
\(I[1,2]\oplus I[1,3]\oplus I[2,4]\oplus I[3,4]\)
即可。其顶点维数为 \((2,3,3,2)\)；每段复合的秩等于同时覆盖该段首尾顶点的区间数，逐项计数恢复所列全部数值。这既给出分解，也提供了可实现性的构造。
\end{solution}

\begin{exercise}\label{b5:ex:09-impossible-ranks}
有人给 \(1\to2\to3\) 提出维数向量 \((2,2,2)\)，两支箭的秩都为 \(2\)，但复合秩为 \(1\)。说明这组资料为何不可能，并从区间重数公式找到同一个矛盾。
\end{exercise}
\begin{solution}
两个二维空间间的满秩映射都可逆，故复合也可逆，秩应为 \(2\)。从重数公式看，单点区间 \(I[2,2]\) 的重数应为
\(2-2-2+1=-1\)，不可能为一个直和重数。分别检查每个秩不超过定义域和值域维数，只给出了必要条件，尚未保证不同复合之间相容。
\end{solution}

\begin{exercise}\label{b5:ex:09-interval-submodules}
在 \(1\to2\to\cdots\to n\) 上求区间表示 \(I[a,b]\) 的全部子表示。写出一个简单组成列，并求相应商表示。
\end{exercise}
\begin{solution}
各个非零顶点空间为一维，子空间只能为零或整个 \(k\)。若某顶点上的子空间为 \(k\)，其后区间内部的恒等箭迫使所有后续顶点也取 \(k\)。因此除零子表示外，子表示恰为 \(I[c,b]\)，其中 \(a\leq c\leq b\)。相应商为 \(I[a,c-1]\)，在 \(c=a\) 时理解为零。
逐步加入右端顶点，得到
\[
0\subset I[b,b]\subset I[b-1,b]\subset\cdots\subset I[a,b].
\]
各相邻商依次为 \(S_b,S_{b-1},\ldots,S_a\)，所以组成长度为 \(b-a+1\)。这些组成因子并不构成直和分解；当 \(a<b\) 时，原区间仍不可分解。
\end{solution}

\begin{exercise}\label{b5:ex:09-diamond-projective}
对\cref{b5:ex:09-diamond-paths}的箭图，写出 \(P_1=kQe_1\) 的各顶点空间与箭映射。再对 \(B=kQ/(ca-db)\) 求 \(Be_1\)，解释一个关系怎样改变这个投射表示。
\end{exercise}
\begin{solution}
\(P_1\) 在顶点 \(1,2,3,4\) 的基分别为
\[
(e_1),\quad(a),\quad(b),\quad(ca,db).
\]
前两支箭把 \(e_1\) 分别送到 \(a,b\)，后两支箭把 \(a,b\) 分别送到 \(ca,db\)。因此维数向量为 \((1,1,1,2)\)，顶点 \(4\) 的两条到达路径彼此独立。
在 \(Be_1\) 中，\(ca=db\)；它们的共同像仍非零，所以顶点 \(4\) 变为一维，其他顶点不变，各箭在所选基下均为恒等。维数向量变为 \((1,1,1,1)\)。\(Be_1\) 仍投射，因为 \(B=Be_1\oplus B(1-e_1)\) 是左模直和。关系改变路径之间的线性独立性，而不是取消顶点幂等元给出的投射直和项。
\end{solution}

\begin{exercise}\label{b5:ex:09-projective-simple-sinks}
设 \(k\) 为域，\(Q\) 为有限无有向环箭图，\(i\in Q_0\)，\(S_i\) 为只在顶点 \(i\) 取 \(k\) 的简单表示。证明 \(S_i\) 投射，当且仅当没有从 \(i\) 出发的箭。
\end{exercise}
\begin{solution}
没有出箭时，从 \(i\) 出发的路径只有 \(e_i\)，所以 \(P_i=kQe_i=S_i\)，投射性由\cref{b5:prop:vertex-projective}给出。若有出箭，则 \(P_i\) 的正长度路径张成非零真子表示，商为 \(S_i\)。若 \(S_i\) 投射，满射 \(P_i\to S_i\) 就会分裂，使 \(P_i\) 成为两个非零子表示的直和。可是同一命题已经证明 \(\operatorname{End}_Q(P_i)=k\)，所以 \(P_i\) 不可分解，矛盾。因此有出箭时 \(S_i\) 不投射。
\end{solution}

\begin{exercise}\label{b5:ex:09-ext-arrows-not-paths}
设 \(Q\) 有箭 \(a:1\to2\)、\(b:2\to3\)、\(c:1\to3\)。
求 \(\operatorname{Ext}^1_Q(S_1,S_3)\) 与
\(\operatorname{Ext}^1_Q(S_3,S_1)\) 的维数。再写出 \(S_1\) 的短投射消解，说明箭 \(c\) 与路径 \(ba\) 在其中的不同作用。
\end{exercise}
\begin{solution}
两简单表示之间的 Hom 为零。两种方向的 Euler 数分别为 \(-1\) 与 \(0\)，所以扩张维数分别为 \(1\) 与 \(0\)。扩张箭数只计直接箭 \(c\)，不能把长度二的路径 \(ba\) 再数一次。
投射消解为
\[
0\longrightarrow P_2\oplus P_3
\xrightarrow{u}P_1\longrightarrow S_1\longrightarrow0.
\]
\(P_2\) 以 \(e_2,b\) 为基，\(u\) 把它们分别送到 \(a,ba\)；\(P_3\) 只有基向量 \(e_3\)，它被送到 \(c\)。这三条正长度路径线性无关，恰好生成 \(P_1\to S_1\) 的核。路径 \(ba\) 是箭 \(a\) 所生成的投射子表示内部继续沿 \(b\) 得到的向量；箭 \(c\) 则需要另一个独立投射生成项。
\end{solution}

\begin{exercise}\label{b5:ex:09-interval-ext}
设 \(k\) 为域，\(n\ge1\) 为整数，\(Q=(1\to2\to\cdots\to n)\)，\(1\le a\le b\le n\)、\(1\le c\le d\le n\)，\(I[a,b],I[c,d]\) 为 \(k\) 上的区间表示。证明
\[
\dim_k\operatorname{Ext}^1_Q\bigl(I[a,b],I[c,d]\bigr)
=\begin{cases}
1,&a<c\leq b+1\leq d,\\
0,&\text{其他情形}.
\end{cases}
\]
提示：先把 \(I[a,b]\) 写成两个顶点投射表示之间的商。
\end{exercise}
\begin{solution}
若 \(b=n\)，则 \(I[a,b]=P_a\) 投射，Ext 为零，所列非零条件也因 \(d\leq n\) 而不可能成立。若 \(b<n\)，有
\[
0\longrightarrow P_{b+1}\longrightarrow P_a
\longrightarrow I[a,b]\longrightarrow0,
\]
左映射在两区间重合的顶点取恒等。对任意表示 \(W\)，由\cref{b5:prop:projective-presentation-extensions}及顶点取值同构，得到
\[
\operatorname{Ext}^1_Q\bigl(I[a,b],W\bigr)
\cong W_{b+1}/\operatorname{im}(W_a\to W_{b+1}),
\]
其中映射是沿唯一有向路径的复合。
取 \(W=I[c,d]\)。分子非零恰好要求 \(c\leq b+1\leq d\)，此时它为 \(k\)。在这一条件下，从 \(W_a\) 到它的复合非零恰好当 \(c\leq a\)，因为 \(a\leq b<b+1\leq d\)。故余核一维恰好要求再加 \(a<c\)；其余情形余核为零。
\end{solution}

\begin{exercise}\label{b5:ex:09-overlapping-extension}
在 \(1\to2\to3\to4\) 上构造非分裂短正合列
\[
0\longrightarrow I[2,4]\longrightarrow
I[1,4]\oplus I[2,2]\longrightarrow I[1,2]\longrightarrow0.
\]
写出全部非零顶点分量，并核验它代表
\(\operatorname{Ext}^1_Q\bigl(I[1,2],I[2,4]\bigr)\) 的非零类。
\end{exercise}
\begin{solution}
左映射的第一分量为自然包含 \(I[2,4]\to I[1,4]\)，第二分量为到顶点 \(2\) 的自然商映射。右映射的第一分量是截去顶点 \(3,4\) 的商映射 \(I[1,4]\to I[1,2]\)，第二分量取自然包含 \(I[2,2]\to I[1,2]\) 的负号。顶点 \(2\) 上的两个映射因此是
\[
k\xrightarrow{x\mapsto(x,x)}k^2
\xrightarrow{(u,v)\mapsto u-v}k,
\]
它们正合；在顶点 \(1\) 上右映射为恒等，在顶点 \(3,4\) 上左映射为恒等。各箭交换条件由上述区间包含及商映射满足，所以全列正合。
如果分裂，中间项应同构于 \(I[2,4]\oplus I[1,2]\)，但它的实际不可分解分解为 \(I[1,4]\oplus I[2,2]\)，与\cref{b5:thm:quiver-krull-schmidt}的唯一性矛盾。故扩张非零。上一题的条件为 \(1<2\leq3\leq4\)，因此相应 Ext 空间恰为一维。
\end{solution}

\begin{exercise}\label{b5:ex:09-submodule-projective}
设 \(k\) 为域，\(Q\) 为有限无有向环箭图，\(P\) 为有限维投射 \(k\) 表示，\(U\subseteq P\) 为子表示。利用本章的长度一投射消解证明 \(U\) 也投射。
\end{exercise}
\begin{solution}
令 \(M=P/U\)。由\cref{b5:thm:acyclic-quiver-resolution}，取
\(0\to Q_1\to Q_0\to M\to0\)，其中 \(Q_0,Q_1\) 都投射。记两个满射为 \(u:P\to M\)、\(v:Q_0\to M\)，构造拉回
\[
E=\bigl\{\,(p,q)\in P\oplus Q_0\mid u(p)=v(q)\,\bigr\},
\]
投影到 \(Q_0\) 给出 \(0\to U\to E\to Q_0\to0\)，因 \(Q_0\) 投射而分裂，故 \(E\cong U\oplus Q_0\)。投影到 \(P\) 又给出 \(0\to Q_1\to E\to P\to0\)，因 \(P\) 投射而分裂，故 \(E\cong Q_1\oplus P\) 投射。投射对象的直和因子仍投射：把一个待提升映射先接直和投影，再对整个投射对象提升，最后限制到该因子即可。因此 \(U\) 投射。
\end{solution}

\begin{exercise}\label{b5:ex:09-relation-kills-extension}
设 \(Q:1\xrightarrow{a}2\xrightarrow{b}3\)，\(B=kQ/(ba)\)。
令 \(V=S_1\)，\(W=I[2,3]\)，两者都满足关系。
分别计算 \(\operatorname{Ext}_{kQ}^1(V,W)\) 与
\(\operatorname{Ext}_{B}^1(V,W)\)，并说明无关系的 Euler 公式为何不能直接用于后者。
\end{exercise}
\begin{solution}
在无关系箭图上，上一节的区间 Ext 公式给出一维扩张空间；非零类的中间项为 \(I[1,3]\)。也可直接看箭矩阵：扩张只可能有一个参数 \(\xi_a:V_1\to W_2\)，而 \(\xi_b=0\)，且所有顶点上的 \(V_i\to W_i\) 都为零，所以无截面变化能消去该参数。
在 \(B\)-模扩张中，关系要求
\[
W_b\xi_a+\xi_bV_a=0.
\]
这里 \(W_b=1\)，所以 \(\xi_a=0\)，每个扩张都分裂，Ext 为零。用维数向量计算的原 Euler 数仍为
\(\bigl\langle(1,0,0),(0,1,1)\bigr\rangle_Q=-1\)，但在 \(B\) 上 Hom 与 Ext 都为零。关系排除了那个非分裂中间项，因为 \(I[1,3]\) 的两箭复合为恒等，不满足 \(ba=0\)。
\end{solution}

\begin{exercise}\label{b5:ex:09-truncated-polynomial-types}
设 \(m=2\) 或 \(3\)，\(A_m=k[t]/(t^m)\)。分类全部有限维不可分解左 \(A_m\)-模，求唯一简单模的自扩张维数，并比较两个代数的 Ext 箭图。
\end{exercise}
\begin{solution}
一个有限维左模等同于一个满足 \(T^m=0\) 的幂零算子。其最小多项式整除 \(t^m\)，在 \(k\) 上分裂，所以第二卷定理4.5.3的 Jordan 标准形适用；各 Jordan 块大小至多为 \(m\)，每块对应 \(k[t]/(t^r)\)，其中 \(1\leq r\leq m\)。其自同态环为 \(k[t]/(t^r)\)：同态由 \(1\) 的像唯一决定。该环只有平凡幂等元，因此每一块不可分解；不同块的维数不同，故不同构。这给出 \(m\) 个不可分解类型。
两个代数都只有简单模 \(S=k\)，根为 \((t)\)，\(J/J^2\) 一维。由\cref{b5:prop:ext-quiver-direction}，
\(\dim\operatorname{Ext}^1_{A_m}(S,S)=1\)，Ext 箭图都是单顶点单圈。关系分别为 \(t^2=0\) 与 \(t^3=0\)，所以相同的箭图并不能区分二者，也不能决定不可分解类型数。
\end{solution}

\begin{exercise}\label{b5:ex:09-matrix-dual-basic}
令 \(D=k[t]/(t^2)\)、\(A=M_2(D)\)、\(e=E_{11}\)。
求 \(J(A)\)、\(A/J(A)\)、\(eAe\)，并说明 \(A\) 的基本代数与容许关系箭图。再求 \(A\) 的有限维不可分解左模的可能维数。
\end{exercise}
\begin{solution}
理想 \(M_2((t))\) 的平方为零，商为 \(M_2(k)\)，所以
\(J(A)=M_2((t))\)、\(A/J(A)\cong M_2(k)\)。
角代数 \(eAe\cong D\)，且矩阵单位使 \(AeA=A\)，故它给出本章的满幂等元 Morita 约化。基本代数是 \(D\)，对应一个顶点、一支圈及圈平方为零的关系；\(A\) 自己的半单商是二阶矩阵代数，因而不是分裂基本代数。
上一题给出 \(D\) 的两个不可分解类型 \(k,D\)。右 \(D\)-模 \(Ae\) 自由秩为 \(2\)，所以拟逆 \(Ae\otimes_D-\) 把它们分别送到维数为 \(2\) 与 \(4\) 的 \(A\)-模。范畴等价保持并反映直和分解，故这就是全部不可分解类型及其维数。
\end{solution}

\begin{exercise}\label{b5:ex:09-radical-square-zero-family}
设 \(k\) 为无限域，
\[
A=k\langle x,y\rangle/(x^2,xy,yx,y^2).
\]
求其根与 Ext 箭图。再在 \(k^2\) 上令 \(x\) 作用为
\(N=\left(\begin{smallmatrix}0&1\\0&0\end{smallmatrix}\right)\)，
令 \(y\) 作用为 \(\lambda N\)，证明 \(\lambda\in k\) 给出无穷多个两两不同构的不可分解左 \(A\)-模。
\end{exercise}
\begin{solution}
\(A\) 的基为 \(1,x,y\)，根为 \(kx\oplus ky\)，根平方为零，半单商为 \(k\)。因此它分裂基本，Ext 箭图有一个顶点、两支圈，关系使任意两圈的乘积为零。
给定的两个算子满足全部关系，因为 \(N^2=0\)。与 \(N\) 对易的矩阵恰为 \(aI+bN\)，所以同时与 \(N,\lambda N\) 对易的代数同构于 \(k[t]/(t^2)\)，没有非平凡幂等元。每个模因而不可分解。
若可逆矩阵 \(g\) 给出参数 \(\lambda,\mu\) 两模的同构，则
\(gN=Ng\) 且 \(\lambda gN=\mu Ng\)。由 \(gN\neq0\)，得到 \(\lambda=\mu\)。无限域因此给出无穷多个同构类型。三维代数的有限维性并不保证有限表示型。
\end{solution}

\begin{exercise}\label{b5:ex:09-nonsplit-field-obstruction}
把 \(\mathbb C\) 看作实代数。证明它不可能与任何具有容许关系的普通实箭图代数 \(\mathbb RQ/I\) Morita 等价。
\end{exercise}
\begin{solution}
\(\mathbb C\) 作为除代数只有一个简单左模的同构类型，其自同态实代数为 \(\mathbb C^{\mathrm{op}}=\mathbb C\)。另一方面，容许关系给出
\[
(\mathbb RQ/I)/J(\mathbb RQ/I)\cong\mathbb R^{Q_0},
\]
所以每个简单模都是某个顶点上的一维实空间，自同态实代数为 \(\mathbb R\)。Morita 等价把简单对象送到简单对象，并由全忠实性保持它们的自同态环；这会要求 \(\mathbb C\cong\mathbb R\)，不可能，例如前者含有平方为 \(-1\) 的元素而后者没有。因此代数闭或分裂基域的假设不能从普通箭图展示定理中直接删掉。
\end{solution}

% !TeX root = ../../Book5.tex
% 纲要标题：## 第10章*　箭图、反射函子与 Gabriel 定理
% 纲要位置：工作记录/计划纲要.md，第 4046 行。
% > 本章紧接箭图基础，独立介绍箭图表示的分类理论。所需 ADE 图、整数格二次型和简单反射在本章建立，不要求预先完成 Lie 代数根系理论；后来学习 Lie 代数时再比较两种来源。
\OptionalChapter{箭图、反射函子与 Gabriel 定理}{b5:ch:10}
\BookChapterNumber{b5:ch:10}

\begin{chapterabstract}
本章分类有限无有向环箭图的有限维表示。Euler 型把顶点维数与态射、扩张的维数联系起来；对其二次型的直接分析得到 ADE 图。随后在顶点整数格上建立根与反射，用核、余核构造反射函子，并证明 Gabriel 定理：有限表示型恰好对应 ADE 图的不交并，不可分解表示由正根唯一确定。最后讨论 Kronecker 参数族、三箭图的野型性质，以及一个几乎分裂序列。
\end{chapterabstract}

\begin{learninggoals}
\begin{itemize}
\item 能计算 Euler 型，利用 Hom--Ext 公式检验具体表示的自扩张。
\item 能从图的分叉与三臂配方判断正定性，并在整数格上计算根和简单反射。
\item 能实施汇点的核反射及源点的余核反射，说明函子的互逆范围。
\item 分别论证 Gabriel 定理的必要性、存在性与唯一性，并由正根构造表示。
\item 区分有限型、单参数分类与野型嵌入，检验几乎分裂序列所要求的因子分解性质。
\end{itemize}
预备知识为\secref{b5:sec:0901}至\secref{b5:sec:0904}的箭图表示、路径代数、不可分解对象及短投射消解，以及有限维线性代数。本章独立建立所需的 ADE 二次型和反射，不依赖 Lie 代数根系或带关系代数的构造。
\end{learninggoals}

箭图把许多线性映射放在一起，因而很适合观察“具体矩阵问题怎样产生离散结构”。有限型分类中，图的形状、一个正定整数二次型和不可分解表示之间会形成精确对应。不过，维数向量并不保存任意表示的全部信息；能够由它唯一恢复的是正定情形中的不可分解类型。反射证明将逐步解释这个限制为什么必要。

\ProseParagraphBreak
这里的分类以无关系的路径代数为对象。第9.5节还讨论带非零关系的代数；对这类代数，不能只凭相同的 Ext 箭图套用这里的正根参数。

% !TeX root = ../../Book5.tex
% 纲要标题：### 10.1 Euler 型与维数向量
% 纲要位置：工作记录/计划纲要.md，第 4050 行。
% * 箭图的 Euler 双线性型
% * Hom、Ext 与维数向量
% * 二次型的正定性
% * 从有限无环箭图定义 \(q_Q(\mathbf d)=\sum_i d_i^2-\sum_{a:i\to j}d_id_j\)，独立处理其正定性与 ADE 图的关系
% * 在顶点整数格上定义简单反射和正根，证明本章需要的基本性质；与第17章根系的对应在构造后说明，不预设全部半单 Lie 代数理论
\section{Euler 型与维数向量}
\label{b5:sec:1001}

同一组顶点空间可以承载许多不同的箭映射。维数向量遗忘了这些映射，却仍能限制自同态和扩张的大小。本节从这个限制出发，把有限分类的问题化为一个实二次型的问题。这里的正定性始终指实数上的正定性，与表示的基域是否有序无关。

\subsection{箭图的 Euler 双线性型}
\label{b5:subsec:1001:01}

\begin{definition}\label{b5:def:quiver-euler-form}
设 \(Q=(Q_0,Q_1,s,t)\) 为有限无有向环箭图，\(Q_0=\{1,\ldots,n\}\)，其中 \(s,t\) 为箭的起点、终点映射。对 \(\mathbf x=(x_i)\)、\(\mathbf y=(y_i)\in\mathbb Z^{Q_0}\)，定义
\[
\langle\mathbf x,\mathbf y\rangle_Q
=\sum_{i\in Q_0}x_i y_i-\sum_{a:i\to j}x_i y_j.
\]
这一双线性型称为 \(Q\) 的\term{Euler 型}[Euler form]。它的二次型及对称化分别记为
\begin{equation}\label{b5:eq:quiver-tits-form}
\begin{aligned}
q_Q(\mathbf x)&=\sum_i x_i^2-\sum_{a:i\to j}x_i x_j,\\
B_Q(\mathbf x,\mathbf y)&=\langle\mathbf x,\mathbf y\rangle_Q+
\langle\mathbf y,\mathbf x\rangle_Q.
\end{aligned}
\end{equation}
二次型 \(q_Q\) 也称为\term{Tits 型}[Tits form]，以蒂茨\footnote{蒂茨（Jacques Tits，1930-2021），比利时裔法国数学家，研究群论与几何，提出楼理论。}命名。上述公式同时把它们延拓到实向量空间 \(\mathbb R^{Q_0}\)。
\end{definition}

箭的方向影响 Euler 型，但不影响 \(q_Q\) 或 \(B_Q\)。把方向忘掉而保留每条边的重数，所得无向图记为 \(\Gamma\)。若 \(m_{ij}\) 为顶点 \(i,j\) 间的边数，则 \(B_Q\) 的矩阵为 \(C=(c_{ij})\)，其中
\[
c_{ii}=2,\qquad c_{ij}=-m_{ij}\quad(i\neq j),\qquad
B_Q(\mathbf x,\mathbf x)=2q_Q(\mathbf x).
\]
本章不允许环，但允许平行箭；无向圈与有向圈是两种不同条件。

\subsection{Hom、Ext 与维数向量}
\label{b5:subsec:1001:02}

\begin{proposition}\label{b5:prop:quiver-euler-hom-ext}
设 \(k\) 为域，\(Q\) 为有限无有向环箭图，\(V,W\) 为 \(Q\) 的有限维 \(k\)-表示，维数向量分别为 \(\mathbf d,\mathbf e\)。记 \(\langle\ ,\ \rangle_Q\) 为箭图的 Euler 型，\(q_Q(\mathbf x)=\langle\mathbf x,\mathbf x\rangle_Q\)，将两表示视作左 \(kQ\)-模，则
\begin{equation}\label{b5:eq:quiver-euler-hom-ext}
\langle\mathbf d,\mathbf e\rangle_Q
=\dim_k\operatorname{Hom}_Q(V,W)
-\dim_k\operatorname{Ext}^1_{kQ}(V,W).
\end{equation}
特别地，\(q_Q(\mathbf d)=\dim_k\operatorname{End}_Q(V)-
\dim_k\operatorname{Ext}^1_{kQ}(V,V)\)。
\end{proposition}
\begin{proof}
\cref{b5:thm:quiver-hom-ext-euler}给出的正合列是
\[
\begin{aligned}
0\longrightarrow\operatorname{Hom}_Q(V,W)
&\longrightarrow\bigoplus_i\operatorname{Hom}_k(V_i,W_i)\\
&\longrightarrow\bigoplus_{a:i\to j}\operatorname{Hom}_k(V_i,W_j)
\longrightarrow\operatorname{Ext}^1_{kQ}(V,W)\longrightarrow0.
\end{aligned}
\]
它适用于当前的有限无环箭图和有限维表示。对有限维正合列取维数交错和，两个中间项的维数分别为 \(\sum_i d_i e_i\) 与 \(\sum_{a:i\to j}d_i e_j\)，即得公式。再令 \(W=V\)。
\end{proof}

例如，对 \(1\to2\)，令 \(V\) 为 \(k\xrightarrow{1}k\)，则自同态为两端相同的标量，故 \(\dim\operatorname{End}_Q(V)=1\)。因为 \(q_Q(1,1)=1\)，公式给出 \(\operatorname{Ext}^1_{kQ}(V,V)=0\)。若把箭映射换成零，同一维数向量上的表示成为 \(S_1\oplus S_2\)，自同态空间维数增为 \(2\)，自扩张空间的维数也增为 \(1\)。因此 \(q_Q(\mathbf d)=1\) 并不表示维数为 \(\mathbf d\) 的每个表示都不可分解。

\subsection{二次型的正定性}
\label{b5:subsec:1001:03}

\begin{definition}\label{b5:def:positive-quiver-form}
设 \(Q\) 为有限无有向环箭图，\(q_Q(\mathbf x)=\sum_{i\in Q_0}x_i^2-\sum_{a:i\to j}x_ix_j\) 为其实二次型。若对每个非零 \(\mathbf x\in\mathbb R^{Q_0}\) 都有 \(q_Q(\mathbf x)>0\)，则称 \(q_Q\) 为\term{正定二次型}[positive definite quadratic form]。
\end{definition}

对一条有 \(n\) 个顶点的无向链，二次型为
\begin{equation}\label{b5:eq:chain-positive-form}
q(\mathbf x)=\frac12\left(x_1^2+x_n^2+
\sum_{i=1}^{n-1}(x_i-x_{i+1})^2\right).
\end{equation}
若右端为零，则首项和全部相邻差都为零，故 \(\mathbf x=0\)。所以任何定向的链都给出正定型。这个配方也适用于 \(n=1\)，此时首尾两项都是 \(x_1^2\)。

\ProseParagraphBreak

相反，若无向图含有一个圈，在圈的顶点上赋值 \(1\)，其余顶点上赋值 \(0\)，便有 \(q(\mathbf x)\leq0\)：圈上的顶点数和边数相等，可能存在的弦只会继续减小二次型。这个例子提示我们，有限分类所容许的图必须相当稀疏。

\subsection{有限无环箭图的二次型与 ADE 图}
\label{b5:subsec:1001:04}

本章用图形本身定义 ADE 名称。\(A_n\) 是有 \(n\geq1\) 个顶点的链；\(D_n\) 是有一个三叉顶点、三个臂长分别为 \(1,1,n-3\) 的树，其中 \(n\geq4\)；\(E_6,E_7,E_8\) 的三个臂长依次为 \((1,2,2)\)、\((1,2,3)\)、\((1,2,4)\)。臂长计中心以外的顶点数。它们统称本章的\term{Dynkin 图}[Dynkin diagram]，也称 ADE 图，其形状见\cref{b5:fig:ade-diagrams}。

% !TeX root = ../Book5.tex
\begin{figure}[htbp]
\centering
\begin{tikzpicture}[x=0.85cm,y=0.85cm,every node/.style={font=\small}]
  % A_n: the dashed segment is a path, with no intermediate vertex when n=3.
  \node at (2,1.45) {\(A_n\quad(n\geq3)\)};
  \draw[thick] (0,0)--(1,0);
  \draw[thick,densely dashed] (1,0)--(4,0);
  \foreach \x in {0,1,4} \fill (\x,0) circle (1.6pt);
  \node at (2,-0.55) {\(n\) 个顶点};

  % D_4 is shown separately: all three arms have exactly one vertex.
  \node at (9.5,1.45) {\(D_4\)};
  \draw[thick] (8.5,0)--(10.5,0) (9.5,0)--(9.5,0.85);
  \foreach \x/\y in {8.5/0,10.5/0,9.5/0.85}
    \fill (\x,\y) circle (1.6pt);
  \fill[AlgebraBlue] (9.5,0) circle (2.2pt);
  \node at (9.5,-0.55) {臂长 \((1,1,1)\)};

  % D_n: the long arm has n-3 vertices; at n=5 its dashed segment is one edge.
  \node at (2,-1.25) {\(D_n\quad(n\geq5)\)};
  \draw[thick] (0,-2.7)--(2,-2.7) (1,-2.7)--(1,-1.85);
  \draw[thick,densely dashed] (2,-2.7)--(4,-2.7);
  \foreach \x/\y in {0/-2.7,2/-2.7,4/-2.7,1/-1.85}
    \fill (\x,\y) circle (1.6pt);
  \fill[AlgebraBlue] (1,-2.7) circle (2.2pt);
  \node at (2,-3.25) {臂长 \((1,1,n-3)\)};

  \node at (9.5,-1.25) {\(E_6\)};
  \draw[thick] (7.5,-2.7)--(11.5,-2.7) (9.5,-2.7)--(9.5,-1.85);
  \foreach \x in {7.5,8.5,10.5,11.5}
    \fill (\x,-2.7) circle (1.6pt);
  \fill (9.5,-1.85) circle (1.6pt);
  \fill[AlgebraBlue] (9.5,-2.7) circle (2.2pt);
  \node at (9.5,-3.25) {臂长 \((1,2,2)\)};

  \node at (2.5,-3.95) {\(E_7\)};
  \draw[thick] (0,-5.4)--(5,-5.4) (2,-5.4)--(2,-4.55);
  \foreach \x in {0,1,3,4,5} \fill (\x,-5.4) circle (1.6pt);
  \fill (2,-4.55) circle (1.6pt);
  \fill[AlgebraBlue] (2,-5.4) circle (2.2pt);
  \node at (2.5,-5.95) {臂长 \((1,2,3)\)};

  \node at (9.5,-3.95) {\(E_8\)};
  \draw[thick] (6.5,-5.4)--(12.5,-5.4) (8.5,-5.4)--(8.5,-4.55);
  \foreach \x in {6.5,7.5,9.5,10.5,11.5,12.5}
    \fill (\x,-5.4) circle (1.6pt);
  \fill (8.5,-4.55) circle (1.6pt);
  \fill[AlgebraBlue] (8.5,-5.4) circle (2.2pt);
  \node at (9.5,-5.95) {臂长 \((1,2,4)\)};
\end{tikzpicture}
\caption[ADE 图及臂长]{ADE 图及臂长。蓝色顶点为三叉点，臂长计中心之外的顶点数。
虚线表示链段，其中可以没有额外顶点：\(A_3\) 与 \(D_5\) 的虚线各是一条边。
\(A_1\) 是单顶点，\(A_2\) 是单边；\(D_4\) 的三条臂均长 \(1\)。}
\label{b5:fig:ade-diagrams}
\end{figure}


\begin{theorem}[正定型的图判别]\label{b5:thm:ade-positive-form}
设 \(Q\) 为有限无有向环箭图，\(\Gamma\) 为保留边重数的无向基础图，\(q_Q(\mathbf x)=\sum_{i\in Q_0}x_i^2-\sum_{a:i\to j}x_ix_j\)。则实二次型 \(q_Q\) 正定，当且仅当 \(\Gamma\) 的每个连通分支都是 ADE 图。若这一条件不成立，则存在非零 \(\mathbf d\in\mathbb Z_{\geq0}^{Q_0}\)，使 \(q_Q(\mathbf d)\leq0\)。
\end{theorem}
\begin{proof}
不同连通分支的变量不相交，二次型是各分支二次型之和，因而只需处理连通图。我们先排除四类子图。所取向量在子图之外均赋值零；子图顶点间若还有额外边，其贡献非正，不会破坏所得不等式。

若两顶点间有 \(m\geq2\) 条边，在这两点赋值 \(1\)，得到 \(q=2-m\leq0\)。若有无向圈，取上一小节的全 \(1\) 向量。所以正定图必须是没有重边的树。若一顶点至少连接四个邻点，在该顶点赋值 \(2\)，在四个邻点赋值 \(1\)，则 \(q=4+4-8=0\)。最后，若树有两个三叉顶点，取连接它们的简单路径；路径上各点赋值 \(2\)，再在两端各选两个不在路径上的邻点赋值 \(1\)。树的性质使这四个邻点互不相同，且没有额外相连的边。设路径上有 \(l\) 个顶点，则
\[
q=4l+4-4(l-1)-8=0.
\]
因此尚可能正定的图只有链和恰有一个三叉点的树。

链的正定性由\eqref{b5:eq:chain-positive-form}成立。设三叉树三个臂长为 \(p-1,q-1,r-1\)，按 \(2\leq p\leq q\leq r\) 排列。中心值记为 \(t\)。对长度为 \(p-1\) 的臂，补记 \(x_0=t,x_p=0\)，并令
\[
z_j=x_j-\left(1-\frac jp\right)t\qquad(0\leq j\leq p).
\]
两端有 \(z_0=z_p=0\)。对三条臂分别作此代换，展开相邻差的平方并用 \(\sum_{j=0}^{p-1}(z_j-z_{j+1})=0\)，得到
\begin{equation}\label{b5:eq:three-arm-square-completion}
q_Q(\mathbf x)=\frac12\sum_{\text{三个臂}}\sum_{j=0}^{p-1}
(z_j-z_{j+1})^2
+\frac12\left(\frac1p+\frac1q+\frac1r-1\right)t^2.
\end{equation}
第一个和式中每条臂的上界采用该臂自己的参数。当 \(1/p+1/q+1/r>1\) 时，等号为零强迫 \(t=0\)，继而每条臂的全部 \(z_j=0\)，所以二次型正定。当这个和不大于 \(1\) 时，取 \(t\) 为 \(p,q,r\) 的正公倍数，并令每条臂上 \(x_j=(1-j/p)t\)，便得到一个非负整数向量，其二次型不大于零。

还须解出严格不等式。若 \(p\geq3\)，三个倒数之和至多为 \(1\)，故 \(p=2\)。若 \(q=2\)，所有 \(r\geq2\) 都可行，给出 \(D_{r+2}\)。若 \(q\geq4\)，则和至多为 \(1/2+1/4+1/4=1\)；剩下 \(q=3\)，而 \(1/r>1/6\) 恰好给出 \(r=3,4,5\)，分别为 \(E_6,E_7,E_8\)。这证明了正定判别，也在所有被排除的情形构造了所需整数向量。
\end{proof}

\subsection{整数格上的简单反射与正根}
\label{b5:subsec:1001:05}

\begin{definition}\label{b5:def:quiver-roots-reflections}
设 \(\Gamma\) 为有限 ADE 图的不交并，顶点集为 \(\Gamma_0\)，\(m_{ij}\) 为顶点 \(i,j\) 间的边数。在 \(\mathbb Z^{\Gamma_0}\) 上令 \(q(\mathbf x)=\sum_i x_i^2-\sum_{\{i,j\}\text{ 为边}}x_ix_j\)，\(B(\mathbf x,\mathbf y)=q(\mathbf x+\mathbf y)-q(\mathbf x)-q(\mathbf y)\)。整数向量 \(\symbfit\alpha\in\mathbb Z^{\Gamma_0}\) 满足 \(q(\symbfit\alpha)=1\) 时，称为一个\term{根}[root]；所有坐标非负的根称为\term{正根}[positive root]。标准基向量 \(\mathbf e_i\) 称为\term{简单根}[simple root]。对顶点 \(i\in\Gamma_0\)，定义
\begin{equation}\label{b5:eq:quiver-simple-reflection}
s_i(\mathbf x)=\mathbf x-B(\mathbf x,\mathbf e_i)\mathbf e_i,
\qquad
(s_i\mathbf x)_i=-x_i+\sum_j m_{ij}x_j,
\end{equation}
其余坐标保持不变；线性变换 \(s_i\) 称为顶点 \(i\) 的\term{简单反射}[simple reflection]。
\end{definition}

\begin{lemma}\label{b5:lem:positive-root-reflection}
设 \(\Gamma\) 为有限 ADE 图的不交并，在顶点整数格上令 \(q(\mathbf x)=\sum_i x_i^2-\sum_{\{i,j\}\text{ 为边}}x_ix_j\)，\(B(\mathbf x,\mathbf y)=q(\mathbf x+\mathbf y)-q(\mathbf x)-q(\mathbf y)\)，\(\mathbf e_i\) 为标准基向量，\(s_i(\mathbf x)=\mathbf x-B(\mathbf x,\mathbf e_i)\mathbf e_i\)。以 \(q=1\) 的整数向量为根、非负根为正根，则有下列性质。
\begin{enumerate}
\item 根的集合有限；每个根或者是正根，或者是一个正根的负数。
\item \(s_i^2=1\)，且 \(s_i\) 保持 \(q\) 及整数格。若 \(\symbfit\alpha\) 为正根且不等于 \(\mathbf e_i\)，则 \(s_i\symbfit\alpha\) 仍为正根；\(s_i\mathbf e_i=-\mathbf e_i\)。
\item 每个非简单正根都可经某个简单反射变成坐标和较小的正根。因此每个根都可由简单根经有限次简单反射得到。
\end{enumerate}
\end{lemma}
\begin{proof}
正定性给出常数 \(c>0\)，使 \(q(\mathbf x)\geq c\sum_i x_i^2\)：可在欧氏单位球面上取 \(q\) 的正最小值。因此 \(q=1\) 的整数向量只有有限多个。将任意整数向量写为 \(\mathbf u-\mathbf v\)，其中 \(\mathbf u,\mathbf v\) 非负且支集不交。由于 \(B\) 的非对角元非正，有 \(B(\mathbf u,\mathbf v)\leq0\)。若两者都非零，正定性和整数性给出
\[
q(\mathbf u-\mathbf v)=q(\mathbf u)+q(\mathbf v)-B(\mathbf u,\mathbf v)\geq2.
\]
所以根不能同时有正、负坐标。

由 \(B(\mathbf e_i,\mathbf e_i)=2\)，代入反射公式即得 \(s_i^2=1\) 和 \(q(s_i\mathbf x)=q(\mathbf x)\)；系数均为整数。反射只改变第 \(i\) 个坐标。若一个正根反射后为负根，它原先在其他坐标必全为零，故由 \(q=1\) 得原根为 \(\mathbf e_i\)。这也证明了第二项中的正性断言。

最后，对正根 \(\symbfit\alpha\)，有
\[
\sum_i\alpha_i\cdot B(\symbfit\alpha,\mathbf e_i)
=B(\symbfit\alpha,\symbfit\alpha)=2.
\]
故某个 \(\alpha_i>0\) 的位置满足 \(B(\symbfit\alpha,\mathbf e_i)>0\)。若根非简单，第二项保证 \(s_i\symbfit\alpha\) 为正根，而其坐标和减少这个正整数。反复进行，正整数的坐标和最终不能再减少，终点只能是简单根。反向使用同一组反射即得最后的断言；负根先由 \(s_i\mathbf e_i=-\mathbf e_i\) 处理。
\end{proof}

这些根暂时只是整数格中的向量。它们与不可分解表示的对应将在\secref{b5:sec:1003}证明；以后 Lie 代数中也会出现相同的 ADE 图。根与反射的这套构造可对照\cite[Section 6.4, pp. 160--163]{EtingofEtAl2011RepresentationTheory}。

% !TeX root = ../../Book5.tex
% 纲要标题：### 10.2 反射函子
% 纲要位置：工作记录/计划纲要.md，第 4058 行。
% * 汇点和源点反射
% * 对不可分解表示的作用
% * 与简单反射的对应
% * 明确汇点反射采用核、源点反射采用余核；给出排除相应简单表示后互为逆的范围，不把反射函子当作整个模范畴上的无条件等价
\section{反射函子}
\label{b5:sec:1002}

简单反射只改动一个整数坐标。在表示中实现这一运算，需要构造新的顶点空间，并相应改变邻接的箭映射。汇点上的核、源点上的余核给出所需构造。反射可能把一个简单表示消去，因此它不会给出整个表示范畴之间的等价。

\subsection{汇点反射与源点反射}
\label{b5:subsec:1002:01}

\begin{definition}\label{b5:def:reflection-functors}
设 \(k\) 为域，\(Q\) 为有限无有向环箭图，\(i\) 为没有出箭的汇点，\(Q'\) 由反转所有以 \(i\) 为终点的箭得到。对 \(Q\) 的有限维 \(k\) 表示 \(V=(V_j,V_a)\)，令
\[
E_i(V)=\bigoplus_{a:j\to i}V_j,\qquad
\phi_V:E_i(V)\longrightarrow V_i,\qquad
(v_a)_a\longmapsto\sum_{a:j\to i}V_a(v_a).
\]
保留其他顶点空间，在 \(i\) 处取 \(\ker\phi_V\)；每条反转后的箭 \(a^*:i\to j\) 取包含映射 \(\ker\phi_V\hookrightarrow E_i(V)\) 与第 \(a\) 个投影的复合。所得 \(Q'\)-表示记为 \(F_i^+V\)，构造 \(F_i^+\) 称为\term{汇点反射函子}[sink reflection functor]。

若 \(i\) 是没有入箭的源点，另令 \(Q'\) 为反转所有以 \(i\) 为起点的箭所得箭图，并令
\[
E_i(V)=\bigoplus_{a:i\to j}V_j,\qquad
\psi_V:V_i\longrightarrow E_i(V),\qquad
v\longmapsto\bigl(V_a(v)\bigr)_a.
\]
在 \(i\) 处取 \(\operatorname{coker}\psi_V\)，其他顶点不变；每条反转后的箭取相应直和包含与商映射的复合。所得表示记为 \(F_i^-V\)，构造 \(F_i^-\) 称为\term{源点反射函子}[source reflection functor]。
\end{definition}

这里的直和以箭为指标。两条平行箭即使出自同一顶点，也对应两个独立的直和分量；漏掉这个重数会使维数公式错误。

\ProseParagraphBreak

还须说明这两个构造怎样作用于态射。若 \(f:V\to W\) 是表示态射，在汇点有
\[
f_i\phi_V=\phi_W\left(\bigoplus_a f_{s(a)}\right).
\]
因此右侧直和映射把 \(\ker\phi_V\) 送入 \(\ker\phi_W\)，它的限制就是新顶点上的映射。各直和投影与此映射相容，保证反转后的箭图仍交换。在源点同样有
\[
\left(\bigoplus_a f_{t(a)}\right)\psi_V=\psi_Wf_i,
\]
故直和映射诱导余核之间的映射；商映射与各直和包含相容。限制或取商都保持恒等态射、复合及态射的加法，所以确实得到两个加性 \(k\)-线性函子。汇点或源点反射不会制造有向圈，因为新箭图中的 \(i\) 分别成为源点或汇点，任何经过它的有向圈都不可能存在。

\subsection{反射函子对不可分解表示的作用}
\label{b5:subsec:1002:02}

记 \(S_i\) 为仅在顶点 \(i\) 处有一个 \(k\)、其余空间全为零的简单表示。以下分裂现象说明，反射唯一需要排除的不可分解对象正是 \(S_i\)。

\begin{lemma}\label{b5:lem:reflection-split-simple}
设 \(k\) 为域，\(Q\) 为有限无有向环箭图，\(V=(V_j,V_a)\) 为有限维 \(k\) 表示，\(i\in Q_0\)。记 \(S_i\) 为只在顶点 \(i\) 取 \(k\) 的简单表示，\(\phi_V:\bigoplus_{a:j\to i}V_j\to V_i\) 为各入箭映射之和，\(\psi_V:V_i\to\bigoplus_{a:i\to j}V_j\) 为各出箭映射组成的映射，\(F_i^+,F_i^-\) 分别为汇点、源点反射函子。
\begin{enumerate}
\item 若 \(i\) 是汇点，则 \(\phi_V\) 满射当且仅当 \(V\) 没有非零的 \(S_i\) 直和项。
\item 若 \(i\) 是源点，则 \(\psi_V\) 单射当且仅当 \(V\) 没有非零的 \(S_i\) 直和项。
\end{enumerate}
特别地，若 \(V\) 不可分解且 \(V\not\cong S_i\)，相应的满射或单射条件必成立；而 \(F_i^+S_i=0\) 或 \(F_i^-S_i=0\)。
\end{lemma}
\begin{proof}
汇点情形取向量空间补空间 \(V_i=\operatorname{im}\phi_V\oplus U\)。所有进入 \(i\) 的箭都落在第一项，而 \(i\) 没有出箭。因此 \(V\) 分解为两个子表示：一个在 \(i\) 处取 \(\operatorname{im}\phi_V\)、其余顶点取原空间，另一个只在 \(i\) 处取 \(U\)。后一项同构于 \(S_i^{\oplus\dim U}\)。反过来，若存在这样的非零直和项，则入箭在该项上全为零，\(\phi_V\) 不可能满射。

源点情形取 \(V_i=\ker\psi_V\oplus U\)。所有出箭在第一项上为零，且没有入箭，故只支撑于 \(\ker\psi_V\) 的部分是直和项；剩下的表示在 \(i\) 处取 \(U\)。存在非零 \(S_i\) 直和项也必使 \(\psi_V\) 有核。若 \(V\) 不可分解，分裂出的项只有在 \(V\cong S_i\) 时才可能非零。最后直接代入核或余核定义，\(S_i\) 均被送到零。
\end{proof}

这与普通线性代数中“把映射的核与像分开”很接近，但源、汇条件不可删去。若一个顶点既有入箭又有出箭，任选的补空间通常不能同时保持两侧的映射。

\subsection{反射函子与简单反射}
\label{b5:subsec:1002:03}

\begin{proposition}\label{b5:prop:reflection-dimension}
设 \(k\) 为域，\(Q\) 为有限无有向环箭图，\(i\) 为汇点或源点，\(V=(V_j,V_a)\) 为有限维 \(k\) 表示。记 \(m_{ij}\) 为顶点 \(i,j\) 间的箭数，定义整数格变换 \((s_i\mathbf x)_i=-x_i+\sum_jm_{ij}x_j\)，其余坐标不变，此处不要求二次型正定。记 \(\phi_V:\bigoplus_{a:j\to i}V_j\to V_i\) 为各入箭映射之和、\(\psi_V:V_i\to\bigoplus_{a:i\to j}V_j\) 为各出箭映射组成的映射，\(\mathbf d\) 为维数向量。若 \(i\) 为汇点且 \(\phi_V\) 满射，则对汇点反射函子有
\[
\mathbf d(F_i^+V)=s_i\bigl(\mathbf d(V)\bigr).
\]
若 \(i\) 为源点且 \(\psi_V\) 单射，则对源点反射函子有
\[
\mathbf d(F_i^-V)=s_i\bigl(\mathbf d(V)\bigr).
\]
\end{proposition}
\begin{proof}
汇点的短正合列为
\[
0\longrightarrow\ker\phi_V\longrightarrow
\bigoplus_{a:j\to i}V_j\longrightarrow V_i\longrightarrow0.
\]
故新顶点维数为 \(\sum_{a:j\to i}\dim V_j-\dim V_i\)，其他维数不变。这正是简单反射的坐标公式。源点则对
\(0\to V_i\to\bigoplus_{a:i\to j}V_j\to\operatorname{coker}\psi_V\to0\)
取维数，得到相同公式。
\end{proof}

满射或单射条件决定了公式能否适用。特别是 \(s_i\mathbf e_i=-\mathbf e_i\)，但反射函子把 \(S_i\) 送到零；负向量不是零表示的维数向量，不能把这两种操作混同。

\begin{example}\label{b5:exm:a3-sink-reflection}
对 \(Q:1\to2\leftarrow3\)，取三个顶点空间均为 \(k\)，两条箭都为恒等映射。汇点 \(2\) 上的总映射是
\[
k\oplus k\longrightarrow k,\qquad (u,v)\longmapsto u+v.
\]
其核由 \((1,-1)\) 生成。反射后得到 \(1\leftarrow2\to3\)，两条箭分别由 \(1\) 与 \(-1\) 给出。新维数仍为 \((1,1,1)\)，与 \(s_2(1,1,1)=(1,1,1)\) 相符。改变顶点 \(3\) 的基可把第二条箭也变成 \(1\)；特征为 \(2\) 时它原先就等于 \(1\)。
\end{example}

\subsection{核、余核与反射函子的互逆范围}
\label{b5:subsec:1002:04}

\begin{theorem}[反射等价]\label{b5:thm:reflection-equivalence}
设 \(k\) 为域，\(Q\) 为有限无有向环箭图，\(i\) 为汇点，\(Q'\) 为反转 \(i\) 处各箭所得箭图。对有限维 \(k\) 表示 \(V\) 属于 \(Q\)、\(W\) 属于 \(Q'\)，分别记 \(\phi_V:\bigoplus_{a:j\to i}V_j\to V_i\) 为入箭映射之和，\(\psi_W:W_i\to\bigoplus_{a^*:i\to j}W_j\) 为出箭映射组成的映射。令 \(\mathcal C\) 为 \(\operatorname{rep}_k Q\) 中 \(\phi_V\) 满射的对象组成的全子范畴，\(\mathcal D\) 为 \(\operatorname{rep}_k Q'\) 中 \(\psi_W\) 单射的对象组成的全子范畴。则汇点反射函子和源点反射函子
\[
F_i^+:\mathcal C\longrightarrow\mathcal D,\qquad
F_i^-:\mathcal D\longrightarrow\mathcal C
\]
互为拟逆。两边的对象条件分别等价于不含只在顶点 \(i\) 取 \(k\) 的简单表示 \(S_i\) 的非零直和项。
\end{theorem}
\begin{proof}
对 \(V\in\mathcal C\)，写 \(E=E_i(V)\)、\(K=\ker\phi_V\)。在 \(F_i^+V\) 中，源点总映射就是包含 \(K\hookrightarrow E\)，故它属于 \(\mathcal D\)。再次反射后，顶点 \(i\) 上的空间为 \(E/K\)。映射
\[
E/K\longrightarrow V_i,\qquad [v]\longmapsto\phi_V(v)
\]
因满射条件而为同构。每条入箭由 \(V_j\hookrightarrow E\to E/K\) 给出，经过该同构正好恢复原来的 \(V_a\)。其他顶点取恒等，因而得到表示同构 \(F_i^-F_i^+V\cong V\)。

对 \(W\in\mathcal D\)，记 \(L=W_i\)、\(E=E_i(W)\)，则 \(L\) 通过 \(\psi_W\) 单射嵌入 \(E\)。\(F_i^-W\) 的汇点总映射为商映射 \(E\to E/\psi_W(L)\)，所以它属于 \(\mathcal C\)。再取核便得到 \(\psi_W(L)\)，通过 \(\psi_W\) 与 \(L\) 同构。各投影限制到 \(\psi_W(L)\) 后恢复原来的出箭，故 \(F_i^+F_i^-W\cong W\)。

对于态射，上述两个同构分别由商空间的诱导映射和原单射 \(\psi_W\) 构成。表示同态与两端的 \(\phi,\psi\) 构成交换方块，故诱导的商映射和核映射也与上述同构构成交换方块，因此所得同构是自然的。两种对象条件的直和解释由\cref{b5:lem:reflection-split-simple}给出。
\end{proof}

\begin{corollary}\label{b5:cor:reflection-indecomposables}
设 \(k\) 为域，\(Q\) 为有限无有向环箭图，\(i\) 为汇点或源点，\(Q'\) 为反转 \(i\) 处各箭所得箭图。以 \(S_i\) 分别表示两箭图中只在顶点 \(i\) 取 \(k\) 的简单表示，\(m_{ij}\) 为顶点 \(i,j\) 间的箭数，\((s_i\mathbf x)_i=-x_i+\sum_jm_{ij}x_j\)，其余坐标不变。相应汇点或源点反射函子，在 \(Q\) 与 \(Q'\) 中不同于 \(S_i\) 的不可分解有限维 \(k\) 表示的同构类之间给出双射，并保持它们的自同态代数；维数向量按 \(s_i\) 变化。
\end{corollary}
\begin{proof}
这些不可分解对象都在\cref{b5:thm:reflection-equivalence}的相应全子范畴中。子范畴对直和项封闭：一个分块直和映射满射或单射，恰好当每个分块都如此。等价保持非零对象与直和分解，所以既保持也反映不可分解性；全忠实性给出自同态代数同构。维数变化由\cref{b5:prop:reflection-dimension}成立。
\end{proof}

反射等价不声称整个表示范畴等价，也不声称反射函子处处正合。例如对 \(1\to2\) 的短正合列
\[
0\longrightarrow S_2\longrightarrow
(k\xrightarrow{1}k)\longrightarrow S_1\longrightarrow0,
\]
取 \(F_2^+\) 后，中间项成为反向箭图上的 \(S_1\)，最后一项成为 \(k\xleftarrow{1}k\)；两者之间的映射不是满射。后文只在上述明确的子范畴中使用拟逆。构造及其基本性质可参看\cite[Section 6.6, pp. 164--169]{EtingofEtAl2011RepresentationTheory}。

% !TeX root = ../../Book5.tex
% 纲要标题：### 10.3 Gabriel 定理
% 纲要位置：工作记录/计划纲要.md，第 4065 行。
% * 有限表示型与 Dynkin 图
% * 不可分解表示和正根
% * 无环及基域假设的明确版本
% * 分别证明有限表示型的必要性、Dynkin 情形的存在与唯一性，以及维数向量对正根的对应；定理固定有限箭图、无有向环、代数闭基域和有限维表示
% * 非连通箭图逐个连通分支处理：有限表示型对应 ADE 图的不交并，而非默认箭图连通
\section{Gabriel 定理}
\label{b5:sec:1003}

反射函子把相邻定向的表示联系起来，但任意挑选反射并不能保证分类过程终止。我们固定一个不断重复的汇点次序，证明在正定情形下，每个不可分解表示终会变成顶点简单表示。反过来，从简单表示沿原次序倒行，就能由正根构造出表示。

\subsection{有限表示型与 Dynkin 图}
\label{b5:subsec:1003:01}

\begin{definition}\label{b5:def:finite-quiver-type}
设 \(k\) 为域，\(Q\) 为有限箭图。若有限维不可分解 \(k\)-表示只有有限多个同构类，则称 \(Q\) 在 \(k\) 上具有\term{有限表示型}[finite representation type]；否则称为无限表示型。
\end{definition}

有限表示型不表示全部表示只有有限多个同构类。例如只有一个顶点且没有箭时，不可分解表示只有 \(k\)，而全部表示有 \(k^m\)（\(m\geq0\)）这无穷多个类型。分类的对象是不能再作非平凡直和分解的部分。

\ProseParagraphBreak

接下来要证明两个不同方向的结论。正定型使反射过程终止，给出有限性及具体分类；非正定型则提供一个维数向量，使箭映射所含的自由参数多于换基能够消去的参数。后一论证不能由“二次型不是正定的”一句话代替。

\subsection{不可分解表示与正根}
\label{b5:subsec:1003:02}

对有限无有向环箭图，依次删去汇点，可以把顶点编号为 \(1,\ldots,n\)，使每条箭都从较大编号指向较小编号。先反射顶点 \(1\)，随后 \(2,\ldots,n\)：轮到 \(i\) 时，原先从 \(i\) 指向较小顶点的箭均已反转，而来自较大顶点的箭尚未反转，所以 \(i\) 确为当前箭图的汇点。一轮之后，每条箭的两端各反射一次，原方向恢复。

\begin{definition}\label{b5:def:quiver-coxeter-transformation}
设 \(Q\) 为有限无有向环箭图，将顶点编号为 \(1,\ldots,n\)，使每支箭都从较大编号指向较小编号。记 \(m_{ij}\) 为顶点 \(i,j\) 间的箭数，\(s_i\) 为整数格 \(\mathbb Z^{Q_0}\) 上的反射，其第 \(i\) 坐标为 \((s_i\mathbf x)_i=-x_i+\sum_jm_{ij}x_j\)，其余坐标不变。复合
\[
c=s_n s_{n-1}\cdots s_1
\]
称为这一编号下的\term{Coxeter 变换}[Coxeter transformation]。\footnote{考克斯特（Harold Scott MacDonald Coxeter，1907-2003），英裔加拿大数学家，研究反射群和多面体几何。}复合按从右到左作用，因此先施行 \(s_1\)。
\end{definition}
\symidx{\(c\)：编号确定的箭图 Coxeter 变换}

\begin{lemma}\label{b5:lem:coxeter-exits-positive-cone}
设 \(Q\) 为有限无有向环箭图，将顶点编号为 \(1,\ldots,n\)，使每支箭都从较大编号指向较小编号。令 \(q_Q(\mathbf x)=\sum_i x_i^2-\sum_{a:i\to j}x_ix_j\) 为正定实二次型，\(B_Q(\mathbf x,\mathbf y)=q_Q(\mathbf x+\mathbf y)-q_Q(\mathbf x)-q_Q(\mathbf y)\)，\(\mathbf e_i\) 为标准基向量，\(s_i(\mathbf x)=\mathbf x-B_Q(\mathbf x,\mathbf e_i)\mathbf e_i\)，\(c=s_n\cdots s_1\)。则 \(c\) 有有限阶，\(1\) 不是其在 \(\mathbb R^{Q_0}\) 上的特征值，并且对每个非零非负实向量 \(\mathbf x\in\mathbb R^{Q_0}\)，存在整数 \(r\ge1\)，使 \(c^r\mathbf x\) 至少有一个负坐标。
\end{lemma}
\begin{proof}
各 \(s_i\) 保持有限的根集合，且标准基向量都是根。由这些反射生成的线性变换群在根集合上的作用是忠实的：固定全部根便固定一组基。因此它是有限置换群的子群，\(c^h=1\) 对某个正整数 \(h\) 成立。

若 \(c\mathbf x=\mathbf x\)，依次施行 \(s_1,\ldots,s_n\)。只有 \(s_i\) 能改变第 \(i\) 个坐标，而且之后的反射不会再次改变它。最终向量等于原向量，所以第一步没有改变坐标 \(1\)，即 \(s_1\mathbf x=\mathbf x\)；随后第二步也不改变坐标 \(2\)。逐步下去，所有 \(s_i\mathbf x=\mathbf x\)。由反射公式，\(B_Q(\mathbf x,\mathbf e_i)=0\) 对每个 \(i\) 成立，故正定性强迫 \(\mathbf x=0\)。这排除了特征值 \(1\)。

令 \(T=1+c+\cdots+c^{h-1}\)，则 \((1-c)T=1-c^h=0\)。因为 \(1-c\) 可逆，得到 \(T=0\)。若 \(\mathbf x\neq0\) 非负且所有 \(c^r\mathbf x\) 都非负，则
\[
0=T\mathbf x=\mathbf x+c\mathbf x+\cdots+c^{h-1}\mathbf x
\]
是包含一个非零项的非负向量之和，矛盾。
\end{proof}

这个引理控制的是完整一轮反射。若从正根出发，在某一步第一次出现负坐标，则\cref{b5:lem:positive-root-reflection}说明前一步恰为当前顶点的简单根。因此，同一个终止论证既适用于表示的维数向量，也适用于尚未实现为表示的正根。此处采用的反射证明方法可对照\cite[Sections 6.7--6.8, pp. 169--172]{EtingofEtAl2011RepresentationTheory}。

\subsection{无环条件与基域假设}
\label{b5:subsec:1003:03}

本章分类定理固定代数闭域、有限无有向环箭图以及有限维表示，不附加箭之间的关系。无有向环条件保证汇点排序存在，也保证第九章使用的路径代数有限维。若加入关系，允许的箭矩阵不再遍历全部矩阵空间，下面的参数计数便不能原封不动地使用。

\ProseParagraphBreak

基域的作用也可以准确区分：核、余核、反射等价以及 Dynkin 情形的构造都只需要域公理；非 Dynkin 情形的参数论证需要域无限，代数闭域满足这一条件。有限域上一个固定维数向量只有有限多个矩阵组，所以不能照搬这个证明来推断存在无穷多个同维数类型。我们不把当前证明未经核对地推广到更一般的基域情形。

\subsection{有限型的必要性、存在性与唯一性}
\label{b5:subsec:1003:04}

先把必要性所需的参数计数写成可独立调用的引理。

\begin{lemma}\label{b5:lem:quiver-parameter-count}
设 \(k\) 为无限域，\(Q\) 为有限无有向环箭图，\(\mathbf d=(d_i)\in\mathbb Z_{\geq0}^{Q_0}\) 非零，\(q_Q(\mathbf d)=\sum_i d_i^2-\sum_{a:i\to j}d_id_j\)。若 \(q_Q(\mathbf d)\leq0\)，则维数向量为 \(\mathbf d\) 的有限维 \(k\)-表示有无穷多个同构类。
\end{lemma}
\begin{proof}
在每个顶点固定空间 \(k^{d_i}\)，所有箭映射组成仿射空间
\[
\operatorname{Rep}_k(Q,\mathbf d)
=\prod_{a:i\to j}M_{d_j\times d_i}(k)\cong k^N,\qquad
N=\sum_{a:i\to j}d_i d_j.
\]
换基群 \(G=\prod_{d_i>0}\operatorname{GL}_{d_i}(k)\) 的轨道恰好是同构类。记 \(m=\sum_i d_i^2\)，则假设给出 \(N\geq m\)。

固定一个表示矩阵组 \(A=(A_a)\)。轨道中的元素为
\[
\bigl(g_{t(a)}A_a g_{s(a)}^{-1}\bigr)_a.
\]
对所有 \(g_i\) 同时乘同一个非零标量不改变结果。选一个满足 \(d_{i_0}>0\) 的顶点；矩阵 \(g_{i_0}\) 的某个条目非零，可利用公共标量将它化为 \(1\)。非零条目的位置只有有限种选择。因此整个轨道是有限多个有理映射的像的并，每个映射至多使用 \(m-1\) 个独立标量参数，定义域只受相关行列式非零的限制。

说明为何这样的像包含在某个非零多项式的零点集中。若 \(N\) 个有理函数只依赖 \(l<N\) 个变量，取共同分母 \(D\)。对它们的所有全次数不超过 \(s\) 的单项式乘以 \(D^s\)，得到次数至多 \(Cs\) 的 \(l\) 元多项式，其中 \(C\) 与 \(s\) 无关。前者有 \(\binom{N+s}{N}\) 个，后者所在向量空间的维数至多为 \(\binom{l+Cs}{l}\)。当 \(s\) 足够大时，前一数量较大，故有非平凡线性关系。这就给出一个非零 \(N\) 元多项式，在有理映射的整个定义域上为零。当前 \(l=m-1<N\)，所以每个参数片都具有这种关系；将有限多个关系相乘，得到在整个轨道上为零的非零多项式。

如果表示只有有限多个同构类，再把各轨道所得的多项式相乘，便得到在整个 \(k^N\) 上恒为零的非零多项式。这不可能：对无限域，非零一元多项式只有有限个根；对变量数归纳，把一个多元多项式视为最后一个变量的多项式，先使某个非零系数不消失，再选择最后一个变量，即能找到非零值。因此轨道有无穷多个。
\end{proof}

\begin{theorem}[Gabriel 定理]\label{b5:thm:gabriel}
设 \(k\) 为代数闭域，\(Q\) 为有限无有向环箭图，允许平行箭，表示均为没有额外关系的有限维 \(k\)-表示，\(q_Q(\mathbf x)=\sum_{i\in Q_0}x_i^2-\sum_{a:i\to j}x_ix_j\)。则 \(Q\) 具有有限表示型，当且仅当其无向基础图是 ADE 图的不交并。此时，维数向量给出双射
\[
\left\{\text{不可分解表示的同构类}\right\}
\longleftrightarrow
\bigl\{\symbfit\alpha\in\mathbb Z_{\geq0}^{Q_0}:
q_Q(\symbfit\alpha)=1\bigr\}.
\]
每个这样的不可分解表示 \(V\) 还满足
\(\operatorname{End}_Q(V)\cong k\) 及
\(\operatorname{Ext}^1_{kQ}(V,V)=0\)。
\end{theorem}
\begin{proof}
先证必要性。
若基础图不是 ADE 图的不交并，\cref{b5:thm:ade-positive-form}给出非零非负整数向量 \(\mathbf d\)，使 \(q_Q(\mathbf d)\leq0\)。代数闭域无限，故\cref{b5:lem:quiver-parameter-count}给出这一维数向量上的无穷多个同构类。

假设非零不可分解表示只有有限多个同构类型 \(U_1,\ldots,U_t\)，并记 \(D=\sum_i d_i\)。对任一维数向量为 \(\mathbf d\) 的表示，每次把一个可分解的非零直和项拆成两个非零项，非零直和项的数目增加 \(1\)；而每个非零项的总维数至少为 \(1\)，所以项数始终至多为 \(D\)。因此分裂过程在有限步后终止，所得各项均不可分解，从而该表示同构于
\[
U_1^{\oplus m_1}\oplus\cdots\oplus U_t^{\oplus m_t}.
\]
其中每个 \(m_j\) 都是满足 \(0\leq m_j\leq D\) 的整数。这样的整数元组只有有限多个，每个元组给出的直和也只有一个同构类型，故维数向量为 \(\mathbf d\) 的表示至多有有限多个同构类型，与参数计数引理矛盾。

现在基础图为 ADE 图的不交并，按汇点次序 \(1,\ldots,n\) 反复施行反射。从不可分解 \(V\neq0\) 出发，只要当前对象不同于所反射顶点的简单表示，\cref{b5:cor:reflection-indecomposables}就保证下一对象仍非零不可分解，维数向量按同一个简单反射变化。若过程永远不遇到该简单表示，每轮之后的维数向量便依次为
\(\mathbf d,c\mathbf d,c^2\mathbf d,\ldots\)，全部非负，与\cref{b5:lem:coxeter-exits-positive-cone}矛盾。因此某一步得到简单表示 \(S_i\)。简单反射保持 \(q_Q\)，故原维数向量满足 \(q_Q(\mathbf d)=q_Q(\mathbf e_i)=1\)，即为正根。

下面证明同一正根至多对应一个不可分解表示。
给定两个维数向量相同的不可分解表示，按相同的汇点次序反射。只要当前维数向量不是该顶点的 \(\mathbf e_i\)，两者都在相应反射等价的范围内。第一次达到 \(\mathbf e_i\) 时，两者都只能是 \(S_i\)。沿已经施行的等价逐步取拟逆，即得原来的两个表示同构。这个论证同时表明，停止时刻由维数向量确定，而不取决于箭映射的具体矩阵。

还需证明每个正根都能实现。
给定正根 \(\symbfit\alpha\)，按相同次序只在整数格中反射。由 Coxeter 引理，有限步后必首次出现负坐标。设这一负向量出现前的正根为 \(\symbfit\beta\)，下一步反射顶点为 \(i\)。由\cref{b5:lem:positive-root-reflection}，必有 \(\symbfit\beta=\mathbf e_i\)。在当时的定向上取 \(S_i\)，并按原次序的反序使用源点反射函子。

需要核实倒行过程不会把对象消去。若某一步要倒行经过顶点 \(j\)，而当前维数向量是 \(\mathbf e_j\)，则原正向过程在这一步之前的向量应为 \(s_j\mathbf e_j=-\mathbf e_j\)，与首次出现负向量以前各项均为正根矛盾。因此每个倒行步骤都避开被消去的简单表示，\cref{b5:cor:reflection-indecomposables}保证对象非零不可分解，并按原坐标序列恢复维数。最终得到维数向量为 \(\symbfit\alpha\) 的不可分解 \(Q\)-表示。

正根有限，所以双射也给出有限表示型。最后，每一步反射等价都保持自同态代数，而终点 \(S_i\) 的自同态代数为 \(k\)，故 \(\operatorname{End}_Q(V)\cong k\)。在\eqref{b5:eq:quiver-euler-hom-ext}中令 \(W=V\)，代入 \(q_Q(\mathbf d)=1\)，得到自扩张空间的维数为零。
\end{proof}

证明还给出了一个构造方法：输入 ADE 箭图的定向及一个正根，先记录汇点反射使它首次变负之前的全部步骤，再从末端的简单表示逐步取余核。每一步只需有限维线性代数。选择核、余核的不同基可能产生不同矩阵，但反射等价保证所得表示的同构类相同。

\ProseParagraphBreak

例如，把 \(D_4\) 的中心记为 \(0\)，三片叶记为 \(1,2,3\)，定向均为 \(i\to0\)。正根 \((2;1,1,1)\) 先在中心反射，再依次在三片叶反射，其维数变化为
\[
(2;1,1,1)\xrightarrow{s_0}(1;1,1,1)
\xrightarrow{s_1}(1;0,1,1)
\xrightarrow{s_2}(1;0,0,1)
\xrightarrow{s_3}(1;0,0,0).
\]
每次所反射的顶点都是当时的汇点。最后的简单表示在中心取 \(k\)。倒行时，先对叶 \(3,2,1\) 作源点反射：各次余核都是 \(\operatorname{coker}(0\to k)=k\)，所以得到中心向三片叶的恒等映射。再在中心取余核，得到
\[
k\xrightarrow{z\mapsto(z,z,z)}k^3
\longrightarrow k^3/k(1,1,1).
\]
原定向的三支箭分别把 \(1\) 送到坐标向量的商类 \(\bar e_1,\bar e_2,\bar e_3\)。其中前两类构成一组基，而 \(\bar e_3=-\bar e_1-\bar e_2\)；把第三片叶的基乘以 \(-1\)，便得到三支箭的列向量 \((1,0)^{\mathsf T},(0,1)^{\mathsf T},(1,1)^{\mathsf T}\)。这就从简单表示实际构造出中心二维、三片叶一维的不可分解表示，特征为二时同样成立。

\subsection{非连通箭图与 ADE 图的不交并}
\label{b5:subsec:1003:05}

若 \(Q\) 的连通分支为 \(Q^{(1)},\ldots,Q^{(r)}\)，任何表示都按支集唯一分成各分支上的表示的直和。因而非零不可分解表示只支撑在一个分支上。二次型也按分支相加；在正定情形，整数向量的每个非零分支至少贡献 \(1\)，所以正根同样只支撑在一个分支上。Gabriel 对应在这里逐分支进行，不需要给非连通图强加一个连通的 Dynkin 名称。

\begin{example}\label{b5:exm:an-positive-roots}
对 \(A_n\)，设 \(x_0=x_{n+1}=0\)，则
\[
2q(\mathbf x)=\sum_{i=0}^{n}(x_{i+1}-x_i)^2.
\]
若 \(\mathbf x\) 是整数正根，右侧等于 \(2\)。全部整数差中只能有两个非零项，各自为 \(1\) 或 \(-1\)；它们的和为零，所以必为一个 \(1\) 和一个 \(-1\)。非负性又要求先出现 \(1\)。因此正根恰为某段连续顶点上全为 \(1\)、其余为零的向量，共有 \(n(n+1)/2\) 个。任意定向下，在这段顶点取 \(k\)，内部各箭取恒等映射，外部空间取零，就得到相应的不可分解表示：所有非零顶点上的自同态标量因箭交换条件而相等，所以自同态代数为 \(k\)，不可能存在非平凡直和分解所产生的幂等元。
\end{example}

例如 \(A_2\) 与 \(A_1\) 的不交并有 \(3+1=4\) 个不可分解类型，而不是 \(3\cdot1\) 个。直和组合描述全部表示，分支的并则描述不可分解类型；两种计数的对象不同。

% !TeX root = ../../Book5.tex
% 纲要标题：### 10.4 表示型与几乎分裂序列
% 纲要位置：工作记录/计划纲要.md，第 4073 行。
% * Kronecker 箭图
% * 有限、驯和野表示型的区别
% * Auslander–Reiten 理论的基本问题
% ---
\section{表示型与几乎分裂序列}
\label{b5:sec:1004}

ADE 图给出有限分类，图稍有扩大就可能出现带参数的不可分解表示。不过，“出现参数”并不等于“分类已无法描述”。本节用两条平行箭和三条平行箭比较两种不同的复杂程度，再用一个最小例子说明 Auslander--Reiten 理论\footnote{奥斯兰德（Maurice Auslander，1926-1994），美国数学家，研究同调代数与表示论；赖滕（Idun Reiten，1942-2025），挪威数学家，研究有限维代数的表示与倾斜理论。}研究什么。

\subsection{Kronecker 箭图}
\label{b5:subsec:1004:01}

\begin{definition}\label{b5:def:kronecker-quiver}
有两个顶点 \(1,2\)，以及两条从 \(1\) 指向 \(2\) 的平行箭的箭图，称为\term{Kronecker 箭图}[Kronecker quiver]。\footnote{克罗内克（Leopold Kronecker，1823-1891），德国数学家，研究代数数论与代数方程。}给定域 \(k\)，它的一个有限维 \(k\)-表示是一对有限维 \(k\) 向量空间之间具有相同定义域和值域的线性映射
\[
A,B:V_1\longrightarrow V_2.
\]
同构是同时作用于两条映射的定义域换基和值域换基。
\end{definition}

它的二次型为 \(q(d_1,d_2)=(d_1-d_2)^2\)，所以并不正定。对每个 \(\lambda\in k\)，取两端空间为 \(k\)，箭为 \(1,\lambda\)，记所得表示为 \(R_\lambda\)。若两端换基标量为 \(a,b\)，第一条箭的交换条件给出 \(a=b\)，第二条箭随即给出 \(\lambda=\mu\)。故 \(R_\lambda\cong R_\mu\) 当且仅当 \(\lambda=\mu\)。每个 \(R_\lambda\) 的自同态代数为 \(k\)，所以不可分解；另有两箭为 \(0,1\) 的 \(R_\infty\)，它与前述各项不同构。

\begin{proposition}\label{b5:prop:kronecker-jordan-family}
设 \(k\) 为代数闭域，\(Q\) 为两支箭均由顶点 \(1\) 指向顶点 \(2\) 的 Kronecker 箭图，\(r\geq1\) 为整数，\(\lambda\in k\)，\(J_r(\lambda)\) 为大小 \(r\) 的 Jordan 块。两顶点空间均取 \(k^r\)，则箭映射给出的表示
\[
R_{\lambda,r}=\bigl(I_r,J_r(\lambda)\bigr),\qquad
R_{\infty,r}=\bigl(J_r(0),I_r\bigr)
\]
均不可分解。对固定 \(r\)，它们由 \(\lambda\in k\cup\{\infty\}\) 给出两两不同的同构类，而且
\[
\dim_k\operatorname{End}_Q(R_{\lambda,r})
=\dim_k\operatorname{Ext}^1_{kQ}(R_{\lambda,r},R_{\lambda,r})=r,
\]
其中也允许 \(\lambda=\infty\)。
\end{proposition}
\begin{proof}
对有限 \(\lambda\)，第一条箭为恒等，故自同态的两端矩阵相同，并与 \(J_r(\lambda)\) 对易。令 \(N=J_r(\lambda)-\lambda I_r\)。Jordan 块有循环向量 \(v\)，使 \(v,Nv,\ldots,N^{r-1}v\) 为一组基。与 \(N\) 对易的算子由 \(v\) 的像唯一决定；若其像为 \(\sum_{j=0}^{r-1}a_jN^jv\)，该算子便是 \(\sum_j a_jN^j\)。因此
\[
\operatorname{End}_Q(R_{\lambda,r})\cong k[t]/(t^r).
\]
这个环的每个元素若常数项非零就可逆，若常数项为零则幂零，所以幂等元只有 \(0,1\)。非平凡直和分解会给出非平凡投影幂等元，故表示不可分解。无穷参数的情形只需用第二条恒等箭，计算相同。

两个有限参数表示的同构要求两个 Jordan 块相似，因而由唯一特征值分别为 \(\lambda,\mu\) 得 \(\lambda=\mu\)。无穷参数中第一条箭不可逆，而有限参数中第一条箭可逆，所以它们不同构。最后，自同态维数为 \(r\)，而 \(q(r,r)=0\)，由\eqref{b5:eq:quiver-euler-hom-ext}得到自扩张维数也为 \(r\)。
\end{proof}

上式给出的族并未穷尽全部不可分解表示。例如维数向量 \((1,2)\)、两箭分别取 \(1\mapsto(1,0)\) 与 \(1\mapsto(0,1)\) 的表示，其自同态也只有标量，却不属于维数为 \((r,r)\) 的那些族。完整分类还要处理维数不同的两端及一般矩阵铅笔；可把这一问题作为后续专题，参看\cite[Problem 6.9.1(d), p. 174]{EtingofEtAl2011RepresentationTheory}。

\subsection{有限、驯和野表示型的区别}
\label{b5:subsec:1004:02}

为了区分“由少量参数族描述”和“包含任意矩阵对的分类”，先给出常用的精确定义。以下只讨论代数闭域上的有限维代数；路径代数 \(kQ\) 在有限无有向环条件下属于这个范围。

\begin{definition}\label{b5:def:tame-wild-type}
设 \(k\) 为代数闭域，\(A\) 为有限维含幺 \(k\)-代数。
\begin{enumerate}
\item 若 \(A\) 有无限多个有限维不可分解左模的同构类，且对每个整数 \(d\geq1\)，存在有限多个 \(A\)-\(k[t]\) 双模 \(M_1,\ldots,M_s\)，每个作为右 \(k[t]\)-模自由且秩为 \(d\)，使除有限多个例外外，每个 \(d\) 维不可分解左 \(A\)-模都同构于某个
\[
M_j\otimes_{k[t]}k[t]/(t-\lambda)\qquad(\lambda\in k),
\]
则称 \(A\) 具有\term{驯表示型}[tame representation type]。
\item 令 \(R=k\langle x,y\rangle\) 为两个非交换变量上的自由结合代数。若存在 \(A\)-\(R\) 双模 \(M\)，作为右 \(R\)-模自由且有限秩，使函子
\[
M\otimes_R-:\{\text{有限维左 }R\text{-模}\}
\longrightarrow\{\text{有限维左 }A\text{-模}\}
\]
保持不可分解性，并且两个输入的像同构当且仅当输入本身同构，则称 \(A\) 具有\term{野表示型}[wild representation type]。
\end{enumerate}
\end{definition}

驯型定义要求在每个维数上穷尽几乎所有不可分解类型，而不是仅仅找到几个单参数族。上一小节的 Jordan 族本身还不足以证明 Kronecker 箭图为驯型。野型定义中的右自由条件保证张量函子正合；它所嵌入的 \(R\)-模分类恰好是任意两个方阵在同时相似下的分类，因为 \(x,y\) 可以独立作用为任意两个算子。

\begin{proposition}\label{b5:prop:three-kronecker-wild}
设 \(k\) 为代数闭域，\(Q\) 为具有两顶点及三条从 \(1\) 指向 \(2\) 的平行箭的箭图。则 \(kQ\) 为野表示型。
\end{proposition}
\begin{proof}
将一个有限维 \(R=k\langle x,y\rangle\)-模写为 \((V,X,Y)\)，其中 \(X,Y\) 是 \(x,y\) 的作用。给它对应 \(Q\)-表示
\[
F(V,X,Y)=(V,V;I_V,X,Y).
\]
两端映射 \((f_1,f_2)\) 与第一条箭交换，强迫 \(f_1=f_2=f\)；余下条件恰好为
\(fX=X'f\)、\(fY=Y'f\)。所以 \(F\) 在态射上给出双射，特别是保持自同态代数并反映同构。一个表示可分解当且仅当其自同态代数有非平凡幂等元：投影给出一个方向，另一个方向由幂等元的像与核给出。因此 \(F\) 保持不可分解性。

最后核实双模条件。取顶点 \(1,2\) 上都为右 \(R\)-模 \(R\)，三条箭分别为恒等映射、左乘 \(x\)、左乘 \(y\)。它们是右 \(R\)-线性映射，故组成一个 \(kQ\)-\(R\) 双模 \(M\)。其右 \(R\)-模为 \(R\oplus R\)，自由秩为 \(2\)，而 \(M\otimes_R V\) 正好是上述 \(F(V,X,Y)\)。于是满足野型定义。
\end{proof}

上述例子给出三种情形：ADE 箭图只有有限多个不可分解类型，Kronecker 箭图出现显式单参数族，三箭图包含任意矩阵对的同时相似问题。把一般有限维代数的无限型进一步分为驯型与野型，还需要更深入的分类理论。

\subsection{Auslander--Reiten 理论的基本问题}
\label{b5:subsec:1004:03}

Gabriel 定理列出了对象，却没有列出对象之间的全部态射。Auslander--Reiten 理论进一步研究：哪些非分裂态射能够控制其他态射的分解，哪些短正合列能够系统地组织不可分解对象之间的联系。

\begin{definition}\label{b5:def:almost-split-sequence}
设 \(k\) 为域，\(A\) 为有限维含幺 \(k\)-代数，\(X,E,Z\) 为有限维左 \(A\)-模，\(p:E\to Z\) 为模同态。若 \(p\) 不是分裂满射，且对每个有限维左 \(A\)-模 \(Y\) 及每个不是分裂满射的模同态 \(f:Y\to Z\)，存在模同态 \(h:Y\to E\) 使 \(f=ph\)，则称 \(p\) 为\term{右几乎分裂态射}[right almost split morphism]。对偶地，给定模同态 \(u:X\to E\)，若 \(u\) 不是分裂单射，且每个不是分裂单射的模同态 \(g:X\to Y\) 都经 \(u\) 分解，则称 \(u\) 为\term{左几乎分裂态射}[left almost split morphism]。一个非分裂短正合列
\[
0\longrightarrow X\xrightarrow{u}E\xrightarrow{p}Z\longrightarrow0
\]
若两端 \(X,Z\) 非零不可分解，且 \(u,p\) 分别具有上述左、右性质，则称为\term{几乎分裂序列}[almost split sequence]，也称 Auslander--Reiten 序列。\termidx{Auslander--Reiten 序列}
\end{definition}

这些定义中的“分裂”是范畴中的分裂，要求存在模态射的左逆或右逆；向量空间上的任意补空间不够。

\begin{example}\label{b5:exm:a2-almost-split}
对任意域 \(k\) 及箭图 \(Q:1\to2\)，有几乎分裂序列
\[
0\longrightarrow S_2\longrightarrow
P\longrightarrow S_1\longrightarrow0,\qquad
P=(k\xrightarrow{1}k).
\]
其中左映射在顶点 \(2\) 为恒等，右映射在顶点 \(1\) 为恒等。
\end{example}
\begin{proof}
序列逐顶点正合。若它分裂，则 \(P\cong S_1\oplus S_2\)，而后者箭映射为零，与 \(P\) 的箭映射秩为 \(1\) 矛盾。

设 \(Y=(Y_1\xrightarrow{T}Y_2)\)。态射 \(f:Y\to S_1\) 由任意线性泛函 \(f_1:Y_1\to k\) 给出。它是分裂满射，当且仅当有 \(v\in\ker T\) 使 \(f_1(v)=1\)：这样的 \(v\) 正好定义一个右逆 \(S_1\to Y\)。因此若 \(f\) 不是分裂满射，\(f_1\) 必在 \(\ker T\) 上为零。它于是经 \(T\) 的像分解为泛函，并可延拓为 \(h_2:Y_2\to k\)，满足 \(h_2T=f_1\)。两端取 \(h_1=f_1,h_2\)，得到态射 \(h:Y\to P\)，并有 \(f=ph\)。这证明右几乎分裂性质。

态射 \(g:S_2\to Y\) 由一个向量 \(w\in Y_2\) 给出。它有左逆，当且仅当存在泛函 \(\ell:Y_2\to k\)，满足 \(\ell T=0\)、\(\ell(w)=1\)，也就是 \(w\notin\operatorname{im}T\)。因此若 \(g\) 不是分裂单射，必有 \(w=T(v)\) 对某个 \(v\in Y_1\) 成立。把 \(1\in P_1\) 送到 \(v\)、把 \(1\in P_2\) 送到 \(w\)，就得到态射 \(P\to Y\)，使 \(g\) 经左映射分解。这证明左几乎分裂性质。两端简单表示当然不可分解。
\end{proof}

在一般理论中，人们把几乎分裂序列左端的同构类记作 \(\tau Z\)，并研究这些序列的存在、唯一性及中间项的不可分解成分；本章只讨论了上面的例子。当前例子说明 \(\tau S_1\cong S_2\)，但 \(\tau\) 并非所有模上的无条件自同构。特别是，若终点 \(Z\) 为投射模，任何以它为终点的短正合列都分裂，因此不存在这种非分裂序列。对象分类与态射结构在这里成为两个互相补充的问题。


\begin{chaptersummary}
Euler 型由顶点和箭的数值资料定义，却等于 Hom 与 Ext 的维数差。正定判别不依赖表示的箭矩阵：重边、圈、四叉及双三叉都会给出非正向量，剩下的三臂树由倒数和不等式筛选出 ADE 图。正定整数格中的根与简单反射因而可以在学习 Lie 代数以前建立。

\ProseParagraphBreak

反射函子把整数格的运算实现为核或余核。在排除相应顶点简单直和项的全子范畴上，它们互为拟逆；对不可分解对象，这正好排除唯一可能被消去的简单表示。Coxeter 变换控制重复反射的终止，由此证明每个 Dynkin 不可分解表示都来自一个简单表示，也证明每个正根能沿逆过程唯一实现。非正定情形的必要性则依赖另一个论证：换基至多消去有限数量的参数，无法把所构造的表示矩阵空间压缩为有限多个轨道。

\ProseParagraphBreak

Kronecker 箭图说明同一维数向量可以含有无穷多个不可分解类型，自扩张也可能非零；三箭图进一步包含任意矩阵对的同时相似问题。几乎分裂序列则把注意力转向对象之间的态射。有限分类、带参数分类和态射分解各自提出不同的问题，不能用其中一个问题的答案代替其他问题。
\end{chaptersummary}

\BookExercises

\begin{optionalexercise}\label{b5:ex:10-euler-orientation}
设 \(Q:1\to2\leftarrow3\)。写出满足 \(\langle\mathbf x,\mathbf y\rangle_Q=\mathbf x^{\mathsf T}E\mathbf y\) 的矩阵 \(E\)，并计算 \(q_Q(2,1,1)\)。反转箭 \(1\to2\) 后，指出 \(E\) 哪些条目改变，并说明二次型为何不变。
\end{optionalexercise}
\begin{solution}
由两条箭的贡献，
\[
E=\begin{pmatrix}1&-1&0\\0&1&0\\0&-1&1\end{pmatrix},
\qquad q_Q(2,1,1)=4+1+1-2-1=3.
\]
反转第一条箭使 \(E_{12}\) 由 \(-1\) 变成 \(0\)，\(E_{21}\) 由 \(0\) 变成 \(-1\)，其余不变。二次型原来的项 \(-x_1x_2\) 变成 \(-x_2x_1\)，故数值完全相同；非对称的 Euler 型则已经改变。
\end{solution}

\begin{optionalexercise}\label{b5:ex:10-hom-orientation}
对 \(Q:1\to2\)，记 \(P=(k\xrightarrow{1}k)\)。计算
\(\operatorname{Hom}_Q(S_1,P)\)、\(\operatorname{Hom}_Q(P,S_1)\)
及对应的两个 \(\operatorname{Ext}^1\)。将箭反转，对相同维数向量的表示重复 Hom 计算。
\end{optionalexercise}
\begin{solution}
态射 \(S_1\to P\) 的顶点 \(1\) 映射被箭交换条件迫为零，故 Hom 为零；态射 \(P\to S_1\) 的顶点 \(1\) 映射可以是任意标量，故 Hom 同构于 \(k\)。Euler 数分别为
\(\bigl\langle(1,0),(1,1)\bigr\rangle_Q=0\)、
\(\bigl\langle(1,1),(1,0)\bigr\rangle_Q=1\)，所以两个 Ext 都为零。
反向箭图的 \(P'=(k\xleftarrow{1}k)\) 满足
\(\operatorname{Hom}(S_1,P')\cong k\)、
\(\operatorname{Hom}(P',S_1)=0\)。
根集合不依赖定向，但根所标记对象之间的态射空间会随定向改变。
\end{solution}

\begin{optionalexercise}\label{b5:ex:10-affine-three-arm}
设树有一个三叉点，三个臂长为 \(1,2,5\)。构造一个所有坐标为正的整数向量，使二次型为零；再说明把最长臂末端删去后为什么变成正定型。
\end{optionalexercise}
\begin{solution}
三个参数为 \(p=2,q=3,r=6\)，其倒数之和为 \(1\)。中心取 \(6\)，三个臂从中心向外分别取
\[
(3),\qquad(4,2),\qquad(5,4,3,2,1).
\]
这些值沿每条臂作线性递减，故\eqref{b5:eq:three-arm-square-completion}中的全部 \(z_j\) 为零，最后的系数也为零，所以 \(q=0\)。删去最长臂末端后，臂长成为 \(1,2,4\)，参数为 \(2,3,5\)，倒数之和大于 \(1\)，这是正定的 \(E_8\) 图。
\end{solution}

\begin{optionalexercise}\label{b5:ex:10-d4-roots}
对 \(D_4\)，用 \(t\) 表示中心坐标，\(x,y,z\) 表示三个叶端坐标。求出全部正根，并数出它们的个数。
\end{optionalexercise}
\begin{solution}
配方得到
\[
q=t^2+x^2+y^2+z^2-t(x+y+z)
=\left(x-\frac t2\right)^2+
\left(y-\frac t2\right)^2+
\left(z-\frac t2\right)^2+\frac{t^2}{4}.
\]
对正根 \(q=1\)，故 \(t=0,1,2\)。若 \(t=0\)，恰有一个叶端坐标为 \(1\)，其余为零，共 \(3\) 个。若 \(t=1\)，条件为
\(x(x-1)+y(y-1)+z(z-1)=0\)；每个非负整数乘积均非负，故 \(x,y,z\) 各自为 \(0\) 或 \(1\)，共 \(8\) 个。若 \(t=2\)，三个平方均须为零，得 \((t;x,y,z)=(2;1,1,1)\)。总计 \(12\) 个。
\end{solution}

\begin{optionalexercise}\label{b5:ex:10-d4-three-lines}
令 \(D_4\) 的三个叶端都指向中心。叶端空间取 \(k\)，中心空间取 \(k^2\)，三条箭的像分别由 \((1,0)\)、\((0,1)\)、\((1,1)\) 生成。证明这个表示不可分解，并计算它的自扩张空间。
\end{optionalexercise}
\begin{solution}
中心自同态保持前两条坐标轴，故为对角矩阵 \(\operatorname{diag}(a,b)\)。它还须保持第三条直线，所以 \(a=b\)。叶端的标量由对应箭交换条件确定，均等于这个共同标量。因此自同态代数为 \(k\)，其幂等元只有 \(0,1\)，表示不可分解。维数向量为上一题的 \((2;1,1,1)\)，二次型为 \(1\)；Euler 公式给出 \(\dim\operatorname{Ext}^1=1-1=0\)。
\end{solution}

\begin{optionalexercise}\label{b5:ex:10-d4-reflect-center}
对上一题的表示施行中心汇点反射，写出新中心空间及三条出箭，再施行反向的源点反射恢复原表示。
\end{optionalexercise}
\begin{solution}
总入箭为
\[
k^3\longrightarrow k^2,\qquad(a,b,c)\longmapsto(a+c,b+c).
\]
其核由 \((-1,-1,1)\) 生成。取该向量为新中心的基，三条出箭的标量依次为 \(-1,-1,1\)。再次反射得到中心空间
\(k^3/k(-1,-1,1)\)；线性映射
\(\bigl[(a,b,c)\bigr]\mapsto(a+c,b+c)\)
是到原 \(k^2\) 的同构。三个标准基向量的类分别映到 \((1,0),(0,1),(1,1)\)，因而三条入箭也被恢复。
\end{solution}

\begin{optionalexercise}\label{b5:ex:10-reflect-rank}
对 \(Q:1\to2\) 的表示 \(V=(k^m\xrightarrow{A}k^n)\)，设 \(\operatorname{rank}A=r\)。求 \(F_2^+V\) 和 \(F_1^-V\) 的维数向量与唯一箭映射的秩，并写出 \(V\) 的不可分解直和分解。
\end{optionalexercise}
\begin{solution}
\(F_2^+V\) 在顶点 \(1,2\) 的空间为 \(k^m,\ker A\)，箭为核的包含，故维数向量为 \((m,m-r)\)、秩为 \(m-r\)。\(F_1^-V\) 的两端为 \(\operatorname{coker}A,k^n\)，箭从 \(2\) 指向 \(1\)，为商映射，故维数向量为 \((n-r,n)\)、秩为 \(n-r\)。
选择定义域的核补空间及值域的像补空间，把 \(A\) 化为一个 \(r\) 阶恒等块及零块，得到
\[
V\cong(k\xrightarrow{1}k)^{\oplus r}
\oplus S_1^{\oplus(m-r)}\oplus S_2^{\oplus(n-r)}.
\]
这也显示：核反射消去 \(S_2\)，余核反射消去 \(S_1\)。
\end{solution}

\begin{optionalexercise}\label{b5:ex:10-reflection-adjunction}
设 \(i\) 为 \(Q\) 的汇点，\(Q'\) 为反射后的箭图。对任意有限维表示 \(W\)（属于 \(Q'\)）与 \(V\)（属于 \(Q\)），证明存在自然同构
\[
\operatorname{Hom}_Q(F_i^-W,V)
\cong\operatorname{Hom}_{Q'}(W,F_i^+V).
\]
这里不要求排除任何简单直和项。
\end{optionalexercise}
\begin{solution}
给定 \(i\) 以外各顶点的映射 \(f_j:W_j\to V_j\)，先要求它们与不接触 \(i\) 的箭交换。记其直和为 \(E(f):E_i(W)\to E_i(V)\)。一个态射 \(F_i^-W\to V\) 在 \(i\) 处所需的映射为
\(E_i(W)/\operatorname{im}\psi_W\to V_i\)，并且与各入箭交换；它存在的必要充分条件为
\[
\phi_V E(f)\psi_W=0.
\]
满足条件时，该映射唯一由 \(\phi_VE(f)\) 取商诱导。另一方面，一个态射 \(W\to F_i^+V\) 在 \(i\) 处必须由 \(E(f)\psi_W\) 限制值域到 \(\ker\phi_V\) 得到，存在的条件也是上述等式，且同样唯一。因此两边由相同的外部映射族描述。复合态射与直和、取商及限制均相容，所以这一直对应自然。
\end{solution}

\begin{optionalexercise}\label{b5:ex:10-reflection-relations}
设 \(\Gamma\) 为 ADE 图。证明非相邻顶点的简单反射交换；相邻顶点满足
\(s_is_js_i=s_js_is_j\)。
\end{optionalexercise}
\begin{solution}
若 \(i,j\) 不相邻，则 \(B(\mathbf e_i,\mathbf e_j)=0\)，两个反射分别只改变沿互相正交的两个向量的分量，故交换。若相邻，在 \(\mathbf e_i,\mathbf e_j\) 张成的平面上，有
\[
s_i(\mathbf e_i)=-\mathbf e_i,\quad
s_i(\mathbf e_j)=\mathbf e_j+\mathbf e_i,\quad
s_j(\mathbf e_j)=-\mathbf e_j,\quad
s_j(\mathbf e_i)=\mathbf e_i+\mathbf e_j.
\]
直接复合可得两个三项乘积都把 \(\mathbf e_i\) 送到 \(-\mathbf e_j\)，把 \(\mathbf e_j\) 送到 \(-\mathbf e_i\)。正定性使这个平面与其 \(B\)-正交补直和为全空间，而两个反射都逐点固定正交补。因此两个乘积在全空间相同。
\end{solution}

\begin{optionalexercise}\label{b5:ex:10-a2-coxeter}
对 \(Q:2\to1\)，写出 \(s_1,s_2,c=s_2s_1\) 的矩阵，验证 \(c^3=1\) 与 \(1+c+c^2=0\)。从简单表示 \(S_2\) 出发，依次在 \(1,2,1,\ldots\) 处作汇点反射，列出遇到当前顶点简单表示以前的全部维数向量。
\end{optionalexercise}
\begin{solution}
矩阵为
\[
s_1=\begin{pmatrix}-1&1\\0&1\end{pmatrix},\quad
s_2=\begin{pmatrix}1&0\\1&-1\end{pmatrix},\quad
c=\begin{pmatrix}-1&1\\-1&0\end{pmatrix},\quad
c^2=\begin{pmatrix}0&-1\\1&-1\end{pmatrix}.
\]
相乘及相加即得所求等式。维数序列为
\((0,1)\mapsto(1,1)\mapsto(1,0)\)；
接下来要反射顶点 \(1\)，当前表示就是 \(S_1\)，于是停止。第一步的坐标和反而增大，说明按汇点次序进行的终止证明不能用“每一步高度都减小”代替 Coxeter 引理。
\end{solution}

\begin{optionalexercise}\label{b5:ex:10-a3-fixed-dimension}
在等向链 \(1\to2\to3\) 上固定维数向量 \((1,1,1)\)。分类这一维数向量上的全部表示，并指出哪个类型不可分解。
\end{optionalexercise}
\begin{solution}
两条箭由标量 \(a,b\) 给出。每个非零标量都可经相邻顶点的基变换化为 \(1\)，且是否为零在同构下不变。因此共有四个类型 \((a,b)=(0,0),(1,0),(0,1),(1,1)\)。它们依次分解为三个简单表示的直和、区间 \([1,2]\) 与 \(S_3\) 的直和、\(S_1\) 与区间 \([2,3]\) 的直和，以及整个区间 \([1,3]\)。仅最后一个不可分解。虽然该维数向量为正根，仍有三个可分解类型具有同一向量。
\end{solution}

\begin{optionalexercise}\label{b5:ex:10-disconnected-count}
设 \(Q\) 的无向基础图为 \(A_3\)、\(D_4\) 及两个孤立顶点的不交并，任取无环定向。求不可分解表示的同构类数，并解释为何横跨两个分支的非零表示必可分解。
\end{optionalexercise}
\begin{solution}
\(A_3\) 有 \(3\cdot4/2=6\) 个正根，\(D_4\) 由\cref{b5:ex:10-d4-roots}有 \(12\) 个，孤立顶点各有一个。因此共有 \(6+12+1+1=20\) 个不可分解类型。若表示在两个分支上都非零，把所有顶点空间按分支分组；由于分支之间没有箭，这两组各自形成子表示，原表示就是它们的非平凡直和。
\end{solution}

\begin{optionalexercise}\label{b5:ex:10-root-end-not-converse}
证明任意域上 \(A_2\) 的表示 \(k\xrightarrow{1}k\) 自同态代数为 \(k\)，却不是简单表示。结合 Kronecker 的 \(R_\lambda\)，解释“自同态只有标量”“简单”“不可分解”和“没有自扩张”不能全部混同。
\end{optionalexercise}
\begin{solution}
箭交换条件迫使两个顶点自同态的标量相同，故代数为 \(k\)。但仅在终点非零的 \(S_2\) 是其真非零子表示，所以它不简单。标量自同态代数没有非平凡幂等元，因而保证不可分解；它并不保证简单。Kronecker 的 \(R_\lambda\) 也具有标量自同态代数，但\cref{b5:prop:kronecker-jordan-family}在 \(r=1\) 时给出一维自扩张空间，说明这一条件也不保证没有自扩张。Dynkin 情形的自扩张消失还使用了 \(q=1\)。
\end{solution}

\begin{optionalexercise}\label{b5:ex:10-finite-field-orbits}
设 \(k\) 为有 \(q\) 个元素的有限域。计算 Kronecker 箭图维数向量 \((1,1)\) 的全部同构类数及不可分解类型数。这为何不与二次型 \(q_Q(1,1)=0\) 矛盾？
\end{optionalexercise}
\begin{solution}
两条箭是一对标量 \((a,b)\)，换基把它们同时乘一个非零标量。零对给出一个类型 \(S_1\oplus S_2\)；非零对的轨道是 \(\mathbb P^1(k)\)，有
\((q^2-1)/(q-1)=q+1\) 个。每个非零对的自同态只有共同标量，所以不可分解。总类型数为 \(q+2\)，不可分解类型数为 \(q+1\)。参数计数引理要求无限域，而有限域上的非零多项式可以在全部点上消失，例如 \(t^q-t\)，故它的最后一步在这里不成立。
\end{solution}

\begin{optionalexercise}\label{b5:ex:10-kronecker-hom-ext}
设 \(k\) 为代数闭域。对有限参数 \(\lambda,\mu\in k\) 与正整数 \(r,s\)，计算下列空间的维数：
\[
\operatorname{Hom}_Q(R_{\lambda,r},R_{\mu,s}),\qquad
\operatorname{Ext}_Q^1(R_{\lambda,r},R_{\mu,s}).
\]
\end{optionalexercise}
\begin{solution}
恒等箭使态射由一个矩阵 \(f:k^r\to k^s\) 给出，并满足
\(fJ_r(\lambda)=J_s(\mu)f\)。若 \(\lambda\neq\mu\)，源空间被某次 \(J_r(\lambda)-\lambda I\) 零化，而目标上的 \(J_s(\mu)-\lambda I\) 可逆；交换等式因此迫使 \(f=0\)。若 \(\lambda=\mu\)，态射等同于
\[
\operatorname{Hom}_{k[t]}\bigl(k[t]/(t^r),k[t]/(t^s)\bigr).
\]
它由 \(1\) 的像决定，该像必须被 \(t^r\) 零化。在目标中满足此条件的子空间维数为 \(\min(r,s)\)。两种情形的 Euler 数均为
\(rs+rs-2rs=0\)，所以 Ext 的维数等于 Hom 的维数：参数不同为 \(0\)，参数相同为 \(\min(r,s)\)。
\end{solution}

\begin{optionalexercise}\label{b5:ex:10-kronecker-unbalanced}
对 Kronecker 箭图，取 \(V_1=k^2,V_2=k^3\)，两箭为
\[
A=\begin{pmatrix}1&0\\0&1\\0&0\end{pmatrix},
\qquad
B=\begin{pmatrix}0&0\\1&0\\0&1\end{pmatrix}.
\]
证明自同态代数为 \(k\)，并计算自扩张维数。
\end{optionalexercise}
\begin{solution}
写源端自同态为 \(X=\left(\begin{smallmatrix}a&b\\c&d\end{smallmatrix}\right)\)，目标端为 \(Y\)。由 \(YA=AX\)，\(Y\) 的前两列为 \((a,c,0)^{\mathsf T}\)、\((b,d,0)^{\mathsf T}\)；由 \(YB=BX\)，后两列为 \((0,a,c)^{\mathsf T}\)、\((0,b,d)^{\mathsf T}\)。比较共同的第二列得 \(b=c=0,d=a\)，随后 \(Y=aI_3\)。故自同态代数为 \(k\)，表示不可分解。二次型值为 \((2-3)^2=1\)，Euler 公式给出自扩张维数为零。无限型箭图仍可有这样的无自扩张不可分解表示。
\end{solution}

\begin{optionalexercise}\label{b5:ex:10-kronecker-extension}
设 \(r\geq2\)。利用 \(k[t]/(t^r)\) 的子模 \(t\,k[t]/(t^r)\)，构造非分裂短正合列
\[
0\longrightarrow R_{\lambda,r-1}\longrightarrow
R_{\lambda,r}\longrightarrow R_{\lambda,1}\longrightarrow0.
\]
\end{optionalexercise}
\begin{solution}
把两端空间都识别为 \(k[t]/(t^r)\)，两箭取恒等与乘 \(\lambda+t\)。子空间 \(t\,k[t]/(t^r)\) 对这两个算子稳定，并通过 \(t^j\mapsto t^{j+1}\) 与 \(k[t]/(t^{r-1})\) 同构；商空间则为 \(k[t]/(t)\)。这给出所需逐顶点正合列。若它分裂，则乘 \(t\) 的算子将相似于大小 \(r-1\) 和大小 \(1\) 的两个幂零块，其幂零指数至多为 \(r-1\)。原模上乘 \(t\) 的幂零指数却为 \(r\)，矛盾。
\end{solution}

\begin{optionalexercise}\label{b5:ex:10-three-arrows-parameters}
在三箭图上取两端均为 \(k\)、三箭为 \(1,\lambda,\mu\) 的表示。证明不同参数对给出不同的不可分解类型，并计算它们的自扩张维数。说明仅仅存在这个双参数族为什么仍不是野型定义的全部要求。
\end{optionalexercise}
\begin{solution}
第一条恒等箭迫使同构两端的标量相同，另外两条箭遂迫使参数对相同。自同态也只有共同标量，故不可分解。二次型值为 \(1+1-3=-1\)，Euler 公式给出
\(\dim\operatorname{Ext}^1=1-(-1)=2\)。
野型定义要求通过一个固定双模，把所有维数上的任意矩阵对分类一起嵌入，并保持不可分解性、反映同构；一个固定维数上的双参数族并未提供这样的构造。\cref{b5:prop:three-kronecker-wild}中的函子才满足全部要求。
\end{solution}

\begin{optionalexercise}\label{b5:ex:10-free-pair-decomposition}
在 \(k^2\) 上取
\[
X=\begin{pmatrix}0&1\\0&0\end{pmatrix},\qquad
Y=\begin{pmatrix}1&0\\0&0\end{pmatrix}.
\]
证明由它们定义的 \(k\langle x,y\rangle\)-模可约但不可分解。将它送入三箭图构造 \(F(V,X,Y)=(V,V;I,X,Y)\)，指出一个真非零子表示，并解释它为何没有表示补空间。
\end{optionalexercise}
\begin{solution}
直线 \(ke_1\) 被两个算子保持，所以模可约。与 \(Y\) 对易的矩阵必为 \(\operatorname{diag}(a,b)\)，再与 \(X\) 对易便要求 \(a=b\)，故共同对易代数只有标量，没有非平凡幂等元，所以模不可分解。三箭图表示在两个顶点都取 \(ke_1\)，三条箭分别为 \(1,0,1\)，就给出一个真非零子表示。若它有表示补空间，投影到该子表示的两端映射会由恒等箭而相等，并同时与 \(X,Y\) 对易；这将给出共同对易代数中的非平凡幂等元，矛盾。共同不变旗与共同不变直和分解并不相同。
\end{solution}

\begin{optionalexercise}\label{b5:ex:10-almost-split-ext}
对 \(Q:1\to2\)，用标量 \(\alpha\in k\) 参数化全部扩张
\(0\to S_2\to E\to S_1\to0\) 的等价类。证明
\(\operatorname{Ext}^1_{kQ}(S_1,S_2)\cong k\)，并说明为何有无穷多个扩张类时仍只有两个可能的中间项同构类型。
\end{optionalexercise}
\begin{solution}
逐顶点正合使 \(E_1,E_2\) 都为 \(k\)，且由两端固定的包含和商映射确定其基。唯一箭便是某个标量 \(\alpha:k\to k\)。扩张等价要求在子模和商模上都是恒等，所以中间项两个顶点的映射都必须为恒等，迫使参数相等。由 Euler 数 \(\bigl\langle(1,0),(0,1)\bigr\rangle_Q=-1\) 及 \(\operatorname{Hom}_Q(S_1,S_2)=0\)，扩张空间维数为 \(1\)；也可以直接核验 Baer 和把两个箭参数相加。具体说，先在商 \(S_1\oplus S_1\) 上取对角拉回，箭变成 \(v\mapsto(\alpha v,\beta v)\)，再在子模 \(S_2\oplus S_2\) 上沿加法推出，箭即为 \(\alpha+\beta\)。标量作用同样对应参数乘标量，因此得到所述线性同构。参数 \(0\) 的中间项是 \(S_1\oplus S_2\)，所有非零参数的中间项都通过改变一个顶点的基同构于 \(k\xrightarrow{1}k\)。中间项的同构允许改变两端的固定识别，故比扩张等价粗。
\end{solution}

\begin{optionalexercise}\label{b5:ex:10-almost-split-factors}
令 \(Y=(k^2\xrightarrow{T}k^3)\)，其中 \(T(a,b)=(a,0,0)\)。分别判断由 \(f_1(a,b)=a\) 给出的 \(f:Y\to S_1\)，以及由 \(1\mapsto(1,0,0)\)、\(1\mapsto(0,1,0)\) 给出的两个 \(S_2\to Y\)，是否分裂。对不分裂的映射，显式写出经 \(P=(k\xrightarrow{1}k)\) 的分解。
\end{optionalexercise}
\begin{solution}
\(\ker T=k(0,1)\)，而 \(f_1\) 在它上面为零，所以 \(f\) 不分裂。取 \(h_1(a,b)=a\)、\(h_2(x,y,z)=x\)，有 \(h_2T=h_1\)，这给出 \(Y\to P\)，再接 \(P\to S_1\) 即为 \(f\)。
第一个 \(S_2\to Y\) 的像在 \(\operatorname{im}T\) 中，所以不分裂；取 \(P_1\to Y_1\) 将 \(1\) 送到 \((1,0)\)，\(P_2\to Y_2\) 将 \(1\) 送到 \((1,0,0)\)，就得到所需 \(P\to Y\)。
第二个映射有左逆：在顶点 \(2\) 取 \((x,y,z)\mapsto y\)，它在 \(\operatorname{im}T\) 上为零且把指定像送到 \(1\)。所以第二个映射是分裂单射，不在左几乎分裂性质所要求的因子分解范围内。
\end{solution}

